\documentclass[11pt]{article}
\usepackage{fontspec}
% Charter and Menlo (macOS); elsewhere, XCharter and DejaVu Sans Mono, from which they derive.
\IfFontExistsTF{Charter}{\setmainfont{Charter}}{\setmainfont{XCharter}}
\IfFontExistsTF{Menlo}{\setmonofont{Menlo}[Scale=0.82]}{\setmonofont{DejaVu Sans Mono}[Scale=0.82]}
\usepackage{polyglossia}\setmainlanguage{english}
\usepackage{microtype}
\usepackage[most]{tcolorbox}
\usepackage{marginnote}
% Each part (a theme) starts on a new page; part numbers in digits.
\usepackage{etoolbox}
\usepackage{titlesec}
\titleformat{\part}[display]{\normalfont\bfseries}{\large Part~\arabic{part}}{0.4em}{\huge}
\titlespacing*{\part}{0pt}{0pt}{2em}
\pretocmd{\part}{\clearpage}{}{}
% The main results are boxed.
\newtcolorbox{encadre}{enhanced, breakable, colback=white, colframe=black!55, boxrule=0.6pt, arc=0pt, left=6pt, right=6pt, top=2pt, bottom=2pt, before skip=8pt, after skip=8pt}
\usepackage[a4paper,margin=27mm]{geometry}
\usepackage{booktabs,longtable,array}
\usepackage{setspace}\setstretch{1.1}
\usepackage{xcolor}
\usepackage[hidelinks]{hyperref}
\usepackage{amsmath,amssymb,amsthm}
\usepackage{mathtools}% psmallmatrix
\usepackage{tikz-cd}% the figures
\usetikzlibrary{decorations.pathreplacing,calc}

\newcommand{\q}[1]{“#1”}
% Quillen's notebooks (Clay Mathematics Institute): one PDF per notebook.
\newcommand{\carnet}[2]{\href{https://www.claymath.org/library/Quillen/Working_papers/quillen\%20#1/#1-#2.pdf}{notebook~#1-#2}}
% The files of 1983 that carry a title instead of a number (notes of lectures, review notes).
\newcommand{\dossier}[2]{\href{https://www.claymath.org/library/Quillen/Working_papers/quillen\%201983/#1.pdf}{#2}}
% The file of the checked proof, once per section. The links point to main.
\newcommand{\leanblobmain}{https://github.com/Commutator-IO/grothendieck-project/blob/main/lean/Quillen}
\newcommand{\verifie}[1]{\footnote{Proof checked: \href{\leanblobmain/#1.lean}{\texttt{lean/Quillen/\detokenize{#1}.lean}}.}}
\definecolor{frise}{HTML}{2F4DA8}
\definecolor{piste}{HTML}{B53D1D}
% The state of the literature search, at the end of a section.
\newcommand{\piste}{\,\textcolor{piste}{\ensuremath{\blacklozenge}}}
\newcommand{\litterature}[1]{\par\smallskip\noindent\textit{Literature.} #1\par}
\theoremstyle{plain}
\newtheorem*{enonce}{Statement}
\newtheorem{lemme}{Lemma}[section]
\newtheorem{theoreme}[lemme]{Theorem}
\newtheorem{contrexemple}[lemme]{Counterexample}
\newtheorem{proposition}[lemme]{Proposition}
\newtheorem*{theoreme*}{Theorem}
\theoremstyle{definition}
\newtheorem{exemple}[lemme]{Example}
\theoremstyle{remark}
\newtheorem*{remarque}{Remark}
\raggedbottom\emergencystretch=1.5em
% Long Lean names (\texttt{A.b\_c\_d}) may break after an underscore.
\let\bkunderscore\_
\renewcommand{\_}{\bkunderscore\allowbreak}

\title{\Large Daniel Quillen's Notebooks: Statements Proved in Lean}
\author{Michel Hua\\ \small Independent researcher, Paris\\ \small \texttt{michel@commutator.io}}
\date{\small November 2026}

\begin{document}
\maketitle

\begin{abstract}
\noindent Daniel Quillen's research notebooks, from 1968 to 2003, are online at the Clay Mathematics Institute. This document picks out one hundred and twenty precise statements in them, in general topology, linear algebra, group theory, representation theory, category theory, combinatorics, order theory, commutative algebra, analysis and the analysis of theta functions, and proves them. Each proof follows a proof checked in the Lean~4 proof assistant \cite{lean4}, with its mathlib library \cite{mathlib2020}, and is written to be read without Lean. For each statement, the document says where it comes from, what it means, how it is proved, what is not formalized, and what is known of its presence in the literature. All one hundred and twenty statements hold in the form proved; where the page, as transcribed, needs a hypothesis, a factor, a surjection in place of an isomorphism or a corrected formula, the section says so. Four of them come from Quillen's notes of lectures by other mathematicians, and are attributed accordingly. Several hold under weaker hypotheses than those of the page or of the expected proof, one fails for $p = 2$, as the page suggests, one needs $2$ to be cancellable (with a counterexample in characteristic~$2$), and most are classical but absent from mathlib; the summary table gives the verdict for each.
\end{abstract}

\tableofcontents

\section*{Introduction}\addcontentsline{toc}{section}{Introduction}

\emph{The notebooks.} Daniel Quillen (1940--2011) kept working notebooks all his life, dated, in which he tries out ideas, redoes computations and notes statements. The Clay Mathematics Institute has digitized them and publishes them at \url{https://www.claymath.org/library/Quillen/}: several hundred notebooks, from 1968 to 2003, one PDF per notebook. This document cites each notebook by its year and number, in the form \q{notebook 1968-14, p.~27}, where the page is that of the PDF; each citation points to the Clay PDF. Some files of 1983 hold notes that Quillen took at lectures of other mathematicians, or review notes; they carry a title instead of a number and are cited by it, and a statement taken from notes of a lecture is attributed to the lecture.

\emph{The rule on rights.} The notebooks are protected by copyright. This document quotes almost nothing from them, a few words at most; it states each result in its own terms and refers to the page of the facsimile, which anyone can consult. The transcriptions made for this work remain private.

\emph{The method.} The notebooks were read in three stages. First, the precise statements are located: lemmas, propositions, identities, questions asked. Each statement is then the object of a literature search. Most of what one reads in the notebooks is found again in Quillen's published papers, on homotopical algebra, group cohomology and algebraic $K$-theory, cobordism, superconnections, determinants and cyclic cohomology; those statements are not taken up here. Finally, only the precise statements that the search did not find in print, or whose printed form does not match the page, are formalized. Four statements of this document are in that case. The others, in notebooks 1968-1, 1968-13 (restriction of representations), 1968-15 and~16, 1969-1, 1969-4, 1970-1, 1970-5, 1970-11, 1970-14, 1970-15, 1971-1, 1971-2, 1971-3, 1971-4, 1971-8, 1971-10, 1971-12 (cosets of ideals, marks, Kaloujnine's theorem and its corollary), 1971-14, 1971-15, 1972-1, 1972-6, 1972-8, 1972-9, 1972-9.5, 1972-11 (transverse subobjects), 1973-6, 1973-8, 1973-12, 1974-2, 1974-3, 1974-16, 1974-17, 1975-10, 1975-21, 1976-1, 1982-2, 1982-3, 1982-5, 1982-6, 1982-8, 1982-12, 1983-4, 1984-6, 1984-13, 1984-15, 1984-16, 1985-8, 1986-2, 1986-5, 1986-7, 1986-12, 1987-5, 1987-11, 1987-12, 1987-13, 1987-14, 1988-1 (the cocycle), 1988-2, 1988-3, 1988-4, 1988-5, 1988-7, 1988-8, 1988-9, 1988-10, 1989-1, 1989-2, 1989-3, 1989-4, 1989-5, 1990-2, 1990-3, 1990-4, 1990-5 and 1991-4, in the 1972 mixed papers, and in the lecture and review notes of 1983, are classical and found in the literature; they are formalized because the pages rely on them and mathlib does not have them. The red lozenge \textcolor{piste}{\ensuremath{\blacklozenge}}, in the title of a section, marks the four notebooks in which a proved statement is a candidate lead: the search did not find it in print in this form, which does not establish it as new.

\emph{What Lean guarantees.} Each section cites, in a footnote to its title, the Lean file that contains its proofs, in the \texttt{lean/Quillen/} library of the project, one file per notebook. Lean mechanically checks that a statement holds under the hypotheses as written. It does not guarantee that the statement faithfully renders the page, nor that it is new. A statement that the page formulates carries the name of Quillen, the notebook and the page; when the proof establishes more, we say \q{after}; a form that the edition adds (a generalization, an answer to a question, an intermediate lemma) is called the edition's statement.

\emph{Organization.} One section per statement. The sections are grouped by theme, in fifteen parts: general topology (\S\S\,\ref{sec:quillen-1968-13}--\ref{sec:quillen-1968-14}); algebraic $K$-theory and general linear groups (\S\S\,\ref{sec:quillen-1968-15}--\ref{sec:quillen-1995-2}); groups: commutators, perfect subgroups and permutations (\S\S\,\ref{sec:quillen-1971-2-normalizer}--\ref{sec:quillen-1990-3}); $p$-groups, nilpotent groups and their automorphisms (\S\S\,\ref{sec:quillen-1971-2}--\ref{sec:quillen-1971-12-nilpotent}); representations and group actions (\S\S\,\ref{sec:quillen-1968-1}--\ref{sec:quillen-1973-6-doset}); number theory and congruences (\S\S\,\ref{sec:quillen-1969-4}--\ref{sec:quillen-1987-5}); combinatorics, $q$-analogues and posets (\S\S\,\ref{sec:quillen-1970-14}--\ref{sec:quillen-1987-14-cyclic}); polynomials, power series and symmetric functions (\S\S\,\ref{sec:quillen-1969-1}--\ref{sec:quillen-1987-11}); rings, ideals and modules (\S\S\,\ref{sec:quillen-1971-12-cover}--\ref{sec:quillen-1991-4}); homological algebra and complexes (\S\S\,\ref{sec:quillen-1973-6-shift}--\ref{sec:quillen-1990-2-split}); Hochschild and cyclic cohomology (\S\S\,\ref{sec:quillen-1986-2}--\ref{sec:quillen-1990-2-green}); noncommutative algebras: traces, extensions and separability (\S\S\,\ref{sec:quillen-1988-1}--\ref{sec:quillen-1991-4-separable}); linear algebra, matrices and determinants (\S\S\,\ref{sec:quillen-1972-mixed-alt}--\ref{sec:quillen-1989-4-faithful}); Clifford algebras, fermions and geometry (\S\S\,\ref{sec:quillen-1982-12}--\ref{sec:quillen-1984-15}); analysis, integrals and special functions (\S\S\,\ref{sec:quillen-1972-4}--\ref{sec:quillen-1988-9-growth}). Inside a part, the sections follow the chronological order of the notebooks; a notebook that touches several themes appears in several parts. The summary table follows the same grouping. Each section presents the page in its context, gives the statement with its intuition and its proof, says what Lean checks and what is not proved, then the state of the literature search.

\section*{Summary table}\addcontentsline{toc}{section}{Summary table}

\begin{center}
{\small\setlength{\tabcolsep}{4pt}
\begin{longtable}{@{}>{\raggedright\arraybackslash}p{1.6cm}>{\raggedright\arraybackslash}p{5.6cm}>{\raggedright\arraybackslash}p{3.9cm}>{\raggedright\arraybackslash}p{2.6cm}r@{}}
\toprule
Notebook & Statement & Verdict & Lean file & \S\\
\midrule
\endhead
\multicolumn{5}{@{}l@{}}{\textbf{General topology}}\\*
1968-13, pp.~64--69 & The space $X \cup_\Phi \{\infty\}$ and the smash lemma & space found (Beer--Vipera); smash lemma more general (neither local compactness nor separation) & \href{\leanblobmain/Notebook1968_13.lean}{\texttt{Notebook1968\_13}} & \ref{sec:quillen-1968-13}\\
1968-14, pp.~27--28 & Compact preimages of closed sets away from $A$ & confirmed; more general (neither local compactness, nor $\sigma$-compactness, nor $A$ closed) & \href{\leanblobmain/Notebook1968_14.lean}{\texttt{Notebook1968\_14}} & \ref{sec:quillen-1968-14}\\
\addlinespace\multicolumn{5}{@{}l@{}}{\textbf{Algebraic $K$-theory and general linear groups}}\\*
1968-15, p.~28; 1968-16, pp.~1--2 & $E(n, A)$ perfect for $n \ge 3$; Whitehead's lemma & classical, found; over any unital ring & \href{\leanblobmain/Notebook1968_15.lean}{\texttt{Notebook1968\_15}}, \href{\leanblobmain/Notebook1968_16.lean}{\texttt{\_16}} & \ref{sec:quillen-1968-15}\\
1971-4, pp.~24--25 & A commutative monoid in its Grothendieck group: flat functor & classical, found; flat and final & \href{\leanblobmain/Notebook1971_4.lean}{\texttt{Notebook1971\_4}} & \ref{sec:quillen-1971-4}\\
1971-7, pp.~14, 23 & The monoid of injections of an infinite set has trivial group completion; so does that of monotone injections of $\mathbf Q$ & classical, found (Schwede, for $\mathbf N$); for any infinite set; ordered case not located, for any countable dense order & \href{\leanblobmain/Notebook1971_7.lean}{\texttt{Notebook1971\_7}} & \ref{sec:quillen-1971-7}\\
1973-6, pp.~30, 45 & Commutators of elementary matrices and Serre's hexagon; a Sylow reduction of an index & classical (Steinberg relations; Brown); second relation of p.~30 fails as transcribed, corrected & \href{\leanblobmain/Notebook1973_6.lean}{\texttt{Notebook1973\_6}} & \ref{sec:quillen-1973-6}\\
1973-8, p.~27 & $A \times M \cong A \times M'$ with $f(1, m)$ of unit first coordinate gives $M \cong M'$ & classical (Bass); explicit inverse & \href{\leanblobmain/Notebook1973_8Splitting.lean}{\texttt{Notebook1973\_8Splitting}} & \ref{sec:quillen-1973-8-splitting}\\
1973-12, p.~48 & For $I^2 = 0$, $\ker(GL_m(A) \to GL_m(A/I)) \cong (M_m(I), +)$, and units lift & classical; any ring, square-zero two-sided ideal & \href{\leanblobmain/Notebook1973_12.lean}{\texttt{Notebook1973\_12}} & \ref{sec:quillen-1973-12}\\
1974-17, p.~3 & $(\alpha, \beta) = (\alpha\sigma\alpha^{-1}\sigma^{-1}, \beta\tau\beta^{-1}\tau^{-1})$ in $GL_{3n}(A)$, $\sigma, \tau$ even block permutations & classical (Whitehead lemma); any ring & \href{\leanblobmain/Notebook1974_17.lean}{\texttt{Notebook1974\_17}} & \ref{sec:quillen-1974-17}\\
1982-5, pp.~15, 41, 55--60 & $\det M\det(M^{-1})_{22} = \det A$; $g \mapsto \tau(g^{-1}, g)$ a homomorphism; parametrices $P_n$ with errors $K_0^n$, $K_1^n$ & classical (Jacobi; Connes's pairing); any ring for the last two & \href{\leanblobmain/Notebook1982_5.lean}{\texttt{Notebook1982\_5}} & \ref{sec:quillen-1982-5}\\
1988-7, pp.~30, 33 & The index idempotent $(p; x^n)(\tilde q, x^n)$ for $qp = 1 - x$, $pq = 1 - y$; $(F - \varepsilon)^2$ central in $\langle F, \varepsilon\rangle$ & classical (Milnor; Connes); any ring & \href{\leanblobmain/Notebook1988_7.lean}{\texttt{Notebook1988\_7}} & \ref{sec:quillen-1988-7}\\
1989-1, pp.~9, 72 & The idempotent-lifting polynomial: $(x^2 - x)^{n+1} \mid f^2 - f$; $\sum_{j<n}x^{2j+1}$ a unit when $x^{2n-1} = 1$ & classical; algebraic proof of the second & \href{\leanblobmain/Notebook1989_1.lean}{\texttt{Notebook1989\_1}} & \ref{sec:quillen-1989-1}\\
1989-2, pp.~28, 30, 84 & $u_z(a)$ conjugate to $\operatorname{diag}(a, 0)$; Whitehead's identity; composition of triples $vu = 1 + [d, h]$ & classical; short identities & \href{\leanblobmain/Notebook1989_2Whitehead.lean}{\texttt{Notebook1989\_2Whitehead}} & \ref{sec:quillen-1989-2-whitehead}\\
1990-5, pp.~26--28, 67 & $ReR = R$ iff $(Re)^n$ maps onto $R$; the splitting lemma; $e = (1 + \varepsilon)/2$ with $FeF = 1 - e$ & classical (Morita) & \href{\leanblobmain/Notebook1990_5.lean}{\texttt{Notebook1990\_5}} & \ref{sec:quillen-1990-5}\\
1995-2, p.~65 & $za + (1 - a)$ is a unit of $R[z, z^{-1}]$ iff $a(1 - a)$ is nilpotent & not located in this form; elementary, commutative case & \href{\leanblobmain/Notebook1995_2.lean}{\texttt{Notebook1995\_2}} & \ref{sec:quillen-1995-2}\\
\addlinespace\multicolumn{5}{@{}l@{}}{\textbf{Groups: commutators, perfect subgroups and permutations}}\\*
1971-2, p.~14; 1975-10, p.~7 & $\Sigma_{n-1}$ is its own normalizer in $\Sigma_n$, $n \ge 3$ & classical; fails for $n = 2$ & \href{\leanblobmain/Notebook1971_2Normalizer.lean}{\texttt{Notebook1971\_2Normalizer}} & \ref{sec:quillen-1971-2-normalizer}\\
1971-3, pp.~32--33 & $T = z(gT)z^{-1}$ makes $g = (z^{-1}, T)$ a commutator & classical (Mather's swindle); the page's formula is the conjugate $T^{-1}gT$ & \href{\leanblobmain/Notebook1971_3.lean}{\texttt{Notebook1971\_3}} & \ref{sec:quillen-1971-3}\\
1971-8, p.~8 & $(G \times G)/\langle\!\langle \Delta G\rangle\!\rangle \cong G_{\mathrm{ab}}$ & classical; elementary & \href{\leanblobmain/Notebook1971_8.lean}{\texttt{Notebook1971\_8}} & \ref{sec:quillen-1971-8}\\
1971-10, pp.~1, 8 & $(aa', bb') = (a, b)$ for displaced elements, so $E = (G,G)$ is perfect; $x^q - x$ onto & classical (Epstein, Mather); any group, $E$ need not be normal & \href{\leanblobmain/Notebook1971_10.lean}{\texttt{Notebook1971\_10}} & \ref{sec:quillen-1971-10}\\
1972-1, p.~25 & A transposition acts on $A_n$ by an outer automorphism & classical; centralizer of $A_n$ in $\Sigma_n$, not its center; also $n = 3$ & \href{\leanblobmain/Notebook1972_1.lean}{\texttt{Notebook1972\_1}} & \ref{sec:quillen-1972-1}\\
1974-16, pp.~5--6 & The largest perfect subgroup; the unique perfect subgroup over a perfect $N$ in a central extension & classical; uniqueness without surjectivity & \href{\leanblobmain/Notebook1974_16.lean}{\texttt{Notebook1974\_16}} & \ref{sec:quillen-1974-16}\\
1990-3, pp.~2, 18; 1990-2, p.~11 & $\gamma\circ(x\ y)$ is an $(n-1)$-cycle iff $x$, $y$ are $\gamma$-adjacent; the Green operator $G$ with $G(1 - \kappa) = 1 - P$, $PG = 0$ & classical & \href{\leanblobmain/Notebook1990_3.lean}{\texttt{Notebook1990\_3}} & \ref{sec:quillen-1990-3}\\
\addlinespace\multicolumn{5}{@{}l@{}}{\textbf{$p$-groups, nilpotent groups and their automorphisms}}\\*
1971-2, p.~12 & An automorphism killing each element by some $(\psi - p)\cdots(\psi - p^n)$ makes the group uniquely $p$-divisible & folklore; algebraic step not located as a lemma; any $p$ & \href{\leanblobmain/Notebook1971_2.lean}{\texttt{Notebook1971\_2}} & \ref{sec:quillen-1971-2}\\
1971-11, pp.~68--71 & Automorphisms of order $p$ trivial on $\Omega_1$; quotient $P/\Omega_1(P)$ of a $p$-central $p$-group & confirmed; more general (neither finiteness nor bijectivity); false for $p = 2$; quotient: found (González-Sánchez--Weigel) & \href{\leanblobmain/Notebook1971_11.lean}{\texttt{Notebook1971\_11}}, \href{\leanblobmain/Notebook1971_11Central.lean}{\texttt{Central}} & \ref{sec:quillen-1971-11}\\
1971-12, pp.~23--24 & An abelian group with $FV = VF = p^d$, $F = p^{d+r}$ is $p^dM \oplus M[p^d]$ & not located as a stated lemma; elementary, $p$ need not be prime & \href{\leanblobmain/Notebook1971_12.lean}{\texttt{Notebook1971\_12}} & \ref{sec:quillen-1971-12}\\
1971-12, pp.~21, 34 & Kaloujnine: $[A_r, A_s] \le A_{r+s}$ for the automorphisms trivial on the $G_i/G_{i+r}$ & classical (Kaloujnine; P.~Hall); proof in the permutation group & \href{\leanblobmain/Notebook1971_12Kaloujnine.lean}{\texttt{Notebook1971\_12Kaloujnine}} & \ref{sec:quillen-1971-12-kaloujnine}\\
1971-12, p.~21 & Automorphisms of a nilpotent group trivial on $G^{\mathrm{ab}}$ form a nilpotent group; $[\Gamma_j, \Gamma_i] \le \Gamma_{i+j+1}$ & classical (Kaloujnine; P.~Hall); from Kaloujnine's theorem & \href{\leanblobmain/Notebook1971_12Nilpotent.lean}{\texttt{Notebook1971\_12Nilpotent}} & \ref{sec:quillen-1971-12-nilpotent}\\
\addlinespace\multicolumn{5}{@{}l@{}}{\textbf{Representations and group actions}}\\*
1968-1, p.~4 & A simple module is determined by its trace & classical, found; in every characteristic, over an algebraically closed field & \href{\leanblobmain/Notebook1968_1.lean}{\texttt{Notebook1968\_1}} & \ref{sec:quillen-1968-1}\\
1968-13, p.~4 & Restriction of representations is full if and only if $f$ is surjective & classical, found; equivalence if and only if $f$ is bijective & \href{\leanblobmain/Notebook1968_13Rep.lean}{\texttt{Notebook1968\_13Rep}} & \ref{sec:quillen-1968-13-rep}\\
1971-12, pp.~28, 30 & $\lvert X^{H'}\rvert \equiv \lvert X^H\rvert \bmod p$ for $H \lhd H'$, $H'/H$ a $p$-group & classical (Dress); $G$ need not be finite & \href{\leanblobmain/Notebook1971_12Marks.lean}{\texttt{Notebook1971\_12Marks}} & \ref{sec:quillen-1971-12-marks}\\
1973-6, p.~61 & $[G : K] = \sum_{HgK}\lvert H\rvert/\lvert H \cap gKg^{-1}\rvert$ & classical (orbit--stabiliser) & \href{\leanblobmain/Notebook1973_6Doset.lean}{\texttt{Notebook1973\_6Doset}} & \ref{sec:quillen-1973-6-doset}\\
\addlinespace\multicolumn{5}{@{}l@{}}{\textbf{Number theory and congruences}}\\*
1969-4, p.~45 & $(p^j)!/p^{(p^j-1)/(p-1)} \equiv (-1)^j$ modulo $p$ & classical, found (Legendre, Wilson); also for $p = 2$ & \href{\leanblobmain/Notebook1969_4.lean}{\texttt{Notebook1969\_4}} & \ref{sec:quillen-1969-4}\\
1970-5, p.~67 & $\prod_{j \le m}(q^{2j}-1)$ and $2^m m!(q-1)^m$ have the same $2$-adic valuation for $4 \mid q - 1$ (and the even case) & classical (lifting the exponent; Carter--Fong); fails for $q = 3$ & \href{\leanblobmain/Notebook1970_5.lean}{\texttt{Notebook1970\_5}} & \ref{sec:quillen-1970-5}\\
1970-11, p.~17 & $|GL_n(\mathbf F_q)|$ and $m!(2(q^2-1))^m(q-1)^e$ have the same $2$-part for $q \equiv 3 \bmod 4$ & classical (lifting the exponent) & \href{\leanblobmain/Notebook1970_11.lean}{\texttt{Notebook1970\_11}} & \ref{sec:quillen-1970-11}\\
1970-11, pp.~21--22; 1974-2, p.~26 & A nonzero multiple of $p^d - 1$ has digit sum $\ge d(p-1)$; first invariant in degree $r(2p - 3)$ & classical, elementary; pages crossed out in 1970-11, statement holds & \href{\leanblobmain/Notebook1970_11.lean}{\texttt{Notebook1970\_11}} & \ref{sec:quillen-1970-11-digits}\\
1971-1, pp.~3, 6--7 & $(1 + p + \beta p^2)^p \equiv 1 + p^2 \pmod{p^3}$ for $p$ odd; $\prod_a(1 + (1+ap)c) = 1 + c^p$ & elementary; fails for $p = 2$; $p$ need not be prime & \href{\leanblobmain/Notebook1971_1.lean}{\texttt{Notebook1971\_1}} & \ref{sec:quillen-1971-1}\\
1971-14, pp.~24--25 & $Z(z) = q^{g-1}z^{2g-2}Z(1/qz)$ for the zeta function of a curve, from Riemann--Roch & classical (Schmidt); finite sums replace the analytic continuation & \href{\leanblobmain/Notebook1971_14.lean}{\texttt{Notebook1971\_14}} & \ref{sec:quillen-1971-14}\\
1972-6, p.~40 & $v_2(3^t - 1)$, $v_\ell(p^t - 1)$ for a generator; denominators of $B_{2s}/4s$, $s \le 5$ & classical (Adams; lifting the exponent) & \href{\leanblobmain/Notebook1972_6.lean}{\texttt{Notebook1972\_6}} & \ref{sec:quillen-1972-6}\\
1972-8, p.~20 & $z \mapsto \sigma(z)/z$ on $k'^\times$ has kernel $k^\times$ and cokernel of order $q - 1$; weights $\lambda^q$ on $\Lambda^q$ & classical & \href{\leanblobmain/Notebook1972_8.lean}{\texttt{Notebook1972\_8}} & \ref{sec:quillen-1972-8}\\
1987-5, p.~64 & $\mathbf Z[1/p] + \mathbf Z_p = \mathbf Q_p$, $\mathbf Z[1/p] \cap \mathbf Z_p = \mathbf Z$ & classical, elementary & \href{\leanblobmain/Notebook1987_5.lean}{\texttt{Notebook1987\_5}} & \ref{sec:quillen-1987-5}\\
\addlinespace\multicolumn{5}{@{}l@{}}{\textbf{Combinatorics, $q$-analogues and posets}}\\*
1970-14, p.~4 & Euler characteristic of the Tits building: $\sum_{d}(-1)^{d+1}\sum[n]!/\prod[s_a]! = 1 + (-1)^n q^{n(n-1)/2}$ & classical (Solomon--Tits); identity of rational functions & \href{\leanblobmain/Notebook1970_14.lean}{\texttt{Notebook1970\_14}} & \ref{sec:quillen-1970-14}\\
1970-15, pp.~21--22 & Möbius function of the lattice of subspaces & classical, found; Gauss's binomial theorem & \href{\leanblobmain/Notebook1970_15.lean}{\texttt{Notebook1970\_15}} & \ref{sec:quillen-1970-15}\\
1971-15, pp.~3--4 & Every finite set of infinite, coinfinite subsets of $\mathbf N$ is coned off by $x \mapsto x \cup c$ & exact poset not located; elementary & \href{\leanblobmain/Notebook1971_15.lean}{\texttt{Notebook1971\_15}} & \ref{sec:quillen-1971-15}\\
1972-8, p.~47 & The $2 \times m$ chessboard graph: $m(m - 1)$ edges, disconnected for $m = 2$, connected for $m \ge 3$ & classical (chessboard complex); the page's $ee^4S^1$ does not fit & \href{\leanblobmain/Notebook1972_8Chess.lean}{\texttt{Notebook1972\_8Chess}} & \ref{sec:quillen-1972-8-chess}\\
1973-5, pp.~28--29 & A partial order is the intersection of its linear extensions; matrices supported on it & classical, found (Szpilrajn; Dushnik--Miller); with the subring of matrices supported on a preorder & \href{\leanblobmain/Notebook1973_5.lean}{\texttt{Notebook1973\_5}} & \ref{sec:quillen-1973-5}\\
1975-21, pp.~41--48 & A bounded lattice homomorphism of lower sets is the preimage along a unique monotone map & classical (Birkhoff duality on morphisms); mathlib TODO; only the source finite & \href{\leanblobmain/Notebook1975_21.lean}{\texttt{Notebook1975\_21}} & \ref{sec:quillen-1975-21}\\
1975-21, pp.~37--38 & Complements $V_j$ of $W(\{<j\})$ in $W(\{\le j\})$ split $W$: $W(Z) = \sum_{j\in Z}V_j$, independent & elementary; any modular lattice & \href{\leanblobmain/Notebook1975_21Filtration.lean}{\texttt{Notebook1975\_21Filtration}} & \ref{sec:quillen-1975-21-filtration}\\
1986-7, p.~56 & $\prod_{i=1}^s(1 - q^i)\sum_{0<i_1<\dots<i_s}q^{i_1+\dots+i_s} = q^{s(s+1)/2}$ & classical (Euler); any commutative ring & \href{\leanblobmain/Notebook1986_7Strict.lean}{\texttt{Notebook1986\_7Strict}} & \ref{sec:quillen-1986-7-strict}\\
1986-7, pp.~55--56 & $\sum_sq^{s^2}[m, s]_q[n, s]_q = [m + n, m]_q$; $\prod_{i>0}(1 - q^i)^{-1} = \sum_sq^{s^2}/(q)_s^2$ & classical (Durfee square); any commutative ring & \href{\leanblobmain/Notebook1986_7Durfee.lean}{\texttt{Notebook1986\_7Durfee}} & \ref{sec:quillen-1986-7-durfee}\\
1987-14, p.~45 & The cyclic embeddings $\mathbf Z/(k-1) \to \mathbf Z/k$ are $k(k-1)$, a $\mathbf Z/k \times \mathbf Z/(k-1)$-torsor & elementary; $k \ge 2$ & \href{\leanblobmain/Notebook1987_14Cyclic.lean}{\texttt{Notebook1987\_14Cyclic}} & \ref{sec:quillen-1987-14-cyclic}\\
\addlinespace\multicolumn{5}{@{}l@{}}{\textbf{Polynomials, power series and symmetric functions}}\\*
1969-1, p.~54 & $\sum_i X_i^q \prod_{j \ne i}(1 - aX_j)/(X_i - X_j) = a^{n-1-q}$ & classical, found (Lagrange interpolation); with the case $a = 0$ & \href{\leanblobmain/Notebook1969_1.lean}{\texttt{Notebook1969\_1}} & \ref{sec:quillen-1969-1}\\
1970-1, pp.~11, 15--16 & $\prod_{v \in V}(X - v)$ is additive for a finite additive subgroup $V$ of a field & classical (Ore; Dickson); any characteristic & \href{\leanblobmain/Notebook1970_1.lean}{\texttt{Notebook1970\_1}} & \ref{sec:quillen-1970-1}\\
1972-9, p.~66 & An additive polynomial in characteristic $p$ has only monomials $X^{p^i}$ & classical (Ore); one variable & \href{\leanblobmain/Notebook1972_9.lean}{\texttt{Notebook1972\_9}} & \ref{sec:quillen-1972-9}\\
1972, mixed papers, pp.~17, 33, 40 & $(\sum\varepsilon_i)^k = k!\,e_k$ when $\varepsilon_i^2 = 0$; $[d, hu' + vh'] = uu' - vv'$; $\dim\operatorname{Ker}(K^s \to K) = s - 1$ & elementary & \href{\leanblobmain/Notebook1972_Mixed.lean}{\texttt{Notebook1972\_Mixed}} & \ref{sec:quillen-1972-mixed}\\
1973-8, p.~62 & $\chi = \sum_k t^{2^k}$ satisfies $\chi + \chi^2 = t$ in characteristic $2$ & classical (Milnor's conjugation); with uniqueness & \href{\leanblobmain/Notebook1973_8.lean}{\texttt{Notebook1973\_8}} & \ref{sec:quillen-1973-8}\\
1974-3, p.~5 & A unique $(f', g')$ with $fg' - f'g = 1$ and $\deg f' < \deg f$ & elementary; struck clause $\deg f' = n - 1$ false & \href{\leanblobmain/Notebook1974_3.lean}{\texttt{Notebook1974\_3}} & \ref{sec:quillen-1974-3}\\
1982-8, p.~63 & The symmetrisation of $1/((z_1 - z_2)\cdots(z_n - z_1))$ vanishes for $n \ge 3$ & classical in substance ($U(1)$ decoupling); proof by partial fractions, not the page's & \href{\leanblobmain/Notebook1982_8.lean}{\texttt{Notebook1982\_8}} & \ref{sec:quillen-1982-8}\\
1985-8, p.~37 & $\sum_k\binom{2k}{k}x^k/4^k = (1-x)^{-1/2}$; $\sum_{i+j=n}\binom{2i}{i}\binom{2j}{j} = 4^n$ & classical; formal power series & \href{\leanblobmain/Notebook1985_8.lean}{\texttt{Notebook1985\_8}} & \ref{sec:quillen-1985-8}\\
1986-7, pp.~25, 50 & Idempotents with $\sum e_i = 1$, $e_te_s = 0$ for $t < s$ are orthogonal; antisymmetrisation $= \det(u_i^{a_j})$ & elementary; alternants & \href{\leanblobmain/Notebook1986_7.lean}{\texttt{Notebook1986\_7}} & \ref{sec:quillen-1986-7}\\
1986-7, pp.~28--29, 45 & $(-1)^nb_n$ is a Toeplitz--Hessenberg determinant for $(\sum c_iT^i)^{-1} = \sum b_iT^i$; $t\,c_t'/c_t = \sum(-1)^{d-1}p_dt^d$ & classical (Wronski; Newton); any commutative ring & \href{\leanblobmain/Notebook1986_7Series.lean}{\texttt{Notebook1986\_7Series}} & \ref{sec:quillen-1986-7-series}\\
1987-11, p.~77 & $s = \sum_{k<N}\binom{2k}k4^{-k}\alpha^k$ satisfies $s^2(1 - \alpha) = 1$ for $\alpha^N = 0$; $F = us$ an involution & classical (binomial series); any $\mathbf Q$-algebra & \href{\leanblobmain/Notebook1987_11.lean}{\texttt{Notebook1987\_11}} & \ref{sec:quillen-1987-11}\\
\addlinespace\multicolumn{5}{@{}l@{}}{\textbf{Rings, ideals and modules}}\\*
1971-12, p.~19 & A ring is not a finite union of cosets of ideals with infinite quotients & classical (Neumann's lemma); any ideals, any ring & \href{\leanblobmain/Notebook1971_12Cover.lean}{\texttt{Notebook1971\_12Cover}} & \ref{sec:quillen-1971-12-cover}\\
1971-13, p.~13 & Good submodules over $A[z]$ pass to quotients & not located; elementary, without finiteness & \href{\leanblobmain/Notebook1971_13.lean}{\texttt{Notebook1971\_13}} & \ref{sec:quillen-1971-13}\\
1972-2, pp.~20--21, 46 & Splitting $R^{(S)} = V \oplus Q$ adapted to an endomorphism; an integral over a triangle & splitting not located, elementary, without finite rank; integral elementary & \href{\leanblobmain/Notebook1972_2.lean}{\texttt{Notebook1972\_2}} & \ref{sec:quillen-1972-2}\\
1972-9, p.~71 & $k[t, x]$ is finite over $k[t, tx + x^q]$, $q \ge 2$ & elementary; any commutative $k$ & \href{\leanblobmain/Notebook1972_9Finite.lean}{\texttt{Notebook1972\_9Finite}} & \ref{sec:quillen-1972-9-finite}\\
1972-9.5, p.~3 & For $\pi$ injective on $E$, $\bigcap_n\pi^nE$ is $\pi$-divisible and contains every $\pi$-divisible part & standard & \href{\leanblobmain/Notebook1972_9_5.lean}{\texttt{Notebook1972\_9\_5}} & \ref{sec:quillen-1972-9-5}\\
1972-11, p.~36 & An extension of finitely presented modules is finitely presented & classical, found; over any ring; in mathlib only for split sequences & \href{\leanblobmain/Notebook1972_11.lean}{\texttt{Notebook1972\_11}} & \ref{sec:quillen-1972-11}\\
1972-11, p.~9 & Transverse subobjects: $A + B = X$ iff $X \to X/A \times X/B$ is onto & classical; onto, not an isomorphism as transcribed & \href{\leanblobmain/Notebook1972_11.lean}{\texttt{Notebook1972\_11}} & \ref{sec:quillen-1972-11-transverse}\\
1987-12, p.~81; 1990-5, p.~63 & $\operatorname{Ker}((A/I)\otimes J \to A/I) \cong (I \cap J)/IJ$; for $J = I$, $I/I^2$ & classical; the identification with $\operatorname{Tor}_1$ not formalised & \href{\leanblobmain/Notebook1987_12.lean}{\texttt{Notebook1987\_12}} & \ref{sec:quillen-1987-12}\\
1991-4, p.~20 & $J \cap K = JK + KJ$ for two-sided ideals with $J + K = R$ & classical, found; in mathlib only for commutative rings & \href{\leanblobmain/Notebook1991_4.lean}{\texttt{Notebook1991\_4}} & \ref{sec:quillen-1991-4}\\
\addlinespace\multicolumn{5}{@{}l@{}}{\textbf{Homological algebra and complexes}}\\*
1973-6, pp.~64, 67 & Inverses of finite filtration compose and are unique up to powers; stabilisation on finite sets and Noetherian modules & classical (Williams); in any category & \href{\leanblobmain/Notebook1973_6.lean}{\texttt{Notebook1973\_6}} & \ref{sec:quillen-1973-6-shift}\\
1984-16, pp.~19--20 & $\delta K + K\delta$ invertible and commuting with $\delta$ gives $\ker\delta = \operatorname{im}\delta$; Witten complex & classical (Atiyah--Bott) & \href{\leanblobmain/Notebook1984_16.lean}{\texttt{Notebook1984\_16}} & \ref{sec:quillen-1984-16}\\
1988-10, pp.~54--55 & $pb' = 0$, $ps = 1$, $b'h + sp = 1$ for $R = T(V)$: $R^{\otimes 4} \to R^{\otimes 3} \to R \otimes V \otimes R \to 0$ exact & classical; explicit homotopy not located & \href{\leanblobmain/Notebook1988_10Bar.lean}{\texttt{Notebook1988\_10Bar}} & \ref{sec:quillen-1988-10-bar}\\
1989-2, pp.~76--80 & The perturbation lemma: $d_M + pXi$, $X = (1 - \delta h)^{-1}\delta$, squares to zero; new chain maps and homotopy & classical (Brown; Gugenheim); side conditions only for $p'i' = 1$ & \href{\leanblobmain/Notebook1989_2.lean}{\texttt{Notebook1989\_2}} & \ref{sec:quillen-1989-2}\\
1989-2, pp.~94--95 & $\Phi(x, g) = ((g_1\cdots g_n)^{-1}x, g)$ is a bijection onto the cyclic bar construction commuting with the faces & classical (Waldhausen) & \href{\leanblobmain/Notebook1989_2Simplicial.lean}{\texttt{Notebook1989\_2Simplicial}} & \ref{sec:quillen-1989-2-simplicial}\\
1990-2, pp.~64--67, 71 & For a split extension, $S = kdh$: $iS = [d, h]$, $[L, S] = [d, kLh]$, and lifts $J = Lh - ig$ from $LS = [d, g]$ & classical; ungraded commutators & \href{\leanblobmain/Notebook1990_2Split.lean}{\texttt{Notebook1990\_2Split}} & \ref{sec:quillen-1990-2-split}\\
\addlinespace\multicolumn{5}{@{}l@{}}{\textbf{Hochschild and cyclic cohomology}}\\*
1986-2, pp.~65, 70 & $\operatorname{tr}(a_0[F,a_1])$ is a cyclic Hochschild cocycle; $\operatorname{tr}(F[F,a][F,b]) = 2\operatorname{tr}(a[F,b])$ for $F^2 = 1$ & classical (Connes); matrices & \href{\leanblobmain/Notebook1986_2.lean}{\texttt{Notebook1986\_2}} & \ref{sec:quillen-1986-2}\\
1987-13, pp.~41--42 & $\operatorname{Ker}N = \operatorname{Im}(1 - \sigma)$, $\operatorname{Im}N = \operatorname{Ker}(1 - \sigma)$ for $\sigma^n = 1$, $n$ invertible & classical (averaging) & \href{\leanblobmain/Notebook1987_13.lean}{\texttt{Notebook1987\_13}} & \ref{sec:quillen-1987-13}\\
1987-13, pp.~72--74 & $b(\varphi \cup \psi) = b\varphi \cup \psi + (-1)^p\varphi \cup b\psi$; $X \cup X = -\frac12 b(X^2)$; first-order deformations & classical (Gerstenhaber); no linearity needed & \href{\leanblobmain/Notebook1987_13Cup.lean}{\texttt{Notebook1987\_13Cup}} & \ref{sec:quillen-1987-13-cup}\\
1988-1, pp.~73--77 & $\operatorname{tr}(x[y,\rho])$ is a derivation in $y$ and depends only on $\rho$ modulo $A$-linear maps & elementary; matrices & \href{\leanblobmain/Notebook1988_1Cocycle.lean}{\texttt{Notebook1988\_1Cocycle}} & \ref{sec:quillen-1988-1-cocycle}\\
1988-5, pp.~87--88 & $\frac d{ds}\operatorname{tr}((sA + B)^{m+1}) = (m + 1)\operatorname{tr}(A(sA + B)^m)$ & classical (cyclicity of the trace) & \href{\leanblobmain/Notebook1988_5.lean}{\texttt{Notebook1988\_5}} & \ref{sec:quillen-1988-5}\\
1988-8, pp.~10, 16, 21 & $(k+1)\operatorname{tr}P_{k+1,l} = (k+l+1)\operatorname{tr}(AP_{k,l})$ for the word sums; $2\tau(t^n\varepsilon) = 0$; a gauge conjugation & classical (cyclicity); any commutative ring & \href{\leanblobmain/Notebook1988_8.lean}{\texttt{Notebook1988\_8}} & \ref{sec:quillen-1988-8}\\
1988-10, pp.~69--70 & $Nb = b'N$ and $b(1 - t) = (1 - t)b'$ & classical (Connes); transposed form on cochains & \href{\leanblobmain/Notebook1988_10.lean}{\texttt{Notebook1988\_10}} & \ref{sec:quillen-1988-10}\\
1989-4, pp.~70--71 & $(bs + sb)\psi = \psi - \psi^\lambda$ on Hochschild cochains & classical (Connes); signs of the box of p.~70 corrected & \href{\leanblobmain/Notebook1989_4Cochain.lean}{\texttt{Notebook1989\_4Cochain}} & \ref{sec:quillen-1989-4-cochain}\\
1989-5, pp.~95, 107, 124 & Cuntz's polynomial: $(1 - \lambda)W(\lambda) = 1 - N(\lambda)/(n + 1)$; $1 + \lambda$ a unit; $u - v$ for two homomorphisms; Cayley & classical or elementary; noncommutative rings & \href{\leanblobmain/Notebook1989_5.lean}{\texttt{Notebook1989\_5}} & \ref{sec:quillen-1989-5}\\
1990-2, pp.~2--3, 99 & $(\kappa - cN)^m = 1$ iff $mc = 1$ when $\kappa^m = 1 + N$; $(1 - \sigma)\sum j\sigma^{n-1-j} = n - \sum\sigma^i$ & elementary; any $\mathbf Q$-algebra & \href{\leanblobmain/Notebook1990_2.lean}{\texttt{Notebook1990\_2}} & \ref{sec:quillen-1990-2}\\
1990-2, p.~14 & $P_{km} = P_m$, $Q_{km} = Q_m$ for the Green operator of $1 - \kappa$, $\kappa^m = 1$ & elementary & \href{\leanblobmain/Notebook1990_2Green.lean}{\texttt{Notebook1990\_2Green}} & \ref{sec:quillen-1990-2-green}\\
\addlinespace\multicolumn{5}{@{}l@{}}{\textbf{Noncommutative algebras: traces, extensions and separability}}\\*
1988-1, pp.~71--72 & $\mathrm{Hom}_k(M, A)$ modulo commutators is the dual of $M$, through the trace & partly found (Lipman); for any $k$-algebra, not only commutative & \href{\leanblobmain/Notebook1988_1.lean}{\texttt{Notebook1988\_1}} & \ref{sec:quillen-1988-1}\\
1988-2, pp.~45, 62; 1988-4, p.~6 & The GNS product on $A \otimes B \otimes A$ is associative; a surjective $\partial$ with $\partial(x)y = x\partial(y)$ is injective & elementary; any linear $\rho$ & \href{\leanblobmain/Notebook1988_2.lean}{\texttt{Notebook1988\_2}} & \ref{sec:quillen-1988-2}\\
1988-4, pp.~59--60 & Cuntz--Krieger matrix units $s_iP_ks_j^*$; the $\mathcal O_2$ relations as a unitary $1 \times 2$ block & classical (Cuntz--Krieger); any ring with involution & \href{\leanblobmain/Notebook1988_4.lean}{\texttt{Notebook1988\_4}} & \ref{sec:quillen-1988-4}\\
1989-4, pp.~14, 18, 34, 38 & $(a + e)^{n+1} = a^{n+1} + \sum a^{n-k}ea^k$ when $exe = 0$; graded commutators; the modular operator & elementary & \href{\leanblobmain/Notebook1989_4.lean}{\texttt{Notebook1989\_4}} & \ref{sec:quillen-1989-4}\\
1989-4, pp.~53--54 & $a^{*(n+2)} = a^{n+2} + \sum_{i+j+k=n}a^i(Xa)a^j(Ya)a^k$ for $a_1 * a_2 = a_1a_2 + Xa_1\,Ya_2$ & elementary; the product is associative & \href{\leanblobmain/Notebook1989_4Square.lean}{\texttt{Notebook1989\_4Square}} & \ref{sec:quillen-1989-4-square}\\
1990-4, pp.~13--14 & A separability element makes every derivation inner, $D(a) = ac - ca$ & classical & \href{\leanblobmain/Notebook1990_4.lean}{\texttt{Notebook1990\_4}} & \ref{sec:quillen-1990-4}\\
1990-5, pp.~48, 68 & $\operatorname{tr}((g_1g_2)^{-1}d(g_1g_2)) \equiv \operatorname{tr}(g_1^{-1}dg_1) + \operatorname{tr}(g_2^{-1}dg_2)$ mod $[E, E]$; the derivation $x\xi - \xi x$ & classical; ungraded derivation & \href{\leanblobmain/Notebook1990_5Trace.lean}{\texttt{Notebook1990\_5Trace}} & \ref{sec:quillen-1990-5-trace}\\
1991-4, pp.~51--58 & $z = \frac1n\sum e_{ij}\otimes e_{ji}$: $(a\otimes1)z = z(1\otimes a)$, $m(z) = 1$, symmetric & classical (separable algebras); $s^2 = 1$ besides & \href{\leanblobmain/Notebook1991_4Separable.lean}{\texttt{Notebook1991\_4Separable}} & \ref{sec:quillen-1991-4-separable}\\
\addlinespace\multicolumn{5}{@{}l@{}}{\textbf{Linear algebra, matrices and determinants}}\\*
1972, mixed papers, p.~41 & A nondegenerate alternating form forces even dimension; odd alternating determinants vanish & classical; every characteristic, by universality & \href{\leanblobmain/Notebook1972_MixedAlt.lean}{\texttt{Notebook1972\_MixedAlt}} & \ref{sec:quillen-1972-mixed-alt}\\
1976-1, p.~49 & $\det(1 + t\nu) \in K^\times$ iff $\nu^2 = 0$ iff $\nu$ nilpotent, $\nu \in M_2(K)$ & elementary; every characteristic & \href{\leanblobmain/Notebook1976_1.lean}{\texttt{Notebook1976\_1}} & \ref{sec:quillen-1976-1}\\
1982-2, p.~22; 1982-6, p.~17 & Faddeev--LeVerrier: $C_{k+1} = KC_k + \Delta_{k+1}$, $(k+1)\Delta_{k+1} = -\operatorname{tr}(KC_k)$; Jacobi's formula & classical; not in mathlib (only the derivative at $0$); any commutative ring & \href{\leanblobmain/Notebook1982_6.lean}{\texttt{Notebook1982\_6}} & \ref{sec:quillen-1982-6}\\
1982-7, pp.~33--34 & A symmetric matrix with non-positive off-diagonal entries and a positive $v$ with $Av \ge 0$ is positive semidefinite & classical, found (Zariski's lemma; $M$-matrices); hypothesis $v^{t}Av \ge 0$ of the page refuted & \href{\leanblobmain/Notebook1982_7.lean}{\texttt{Notebook1982\_7}} & \ref{sec:quillen-1982-7}\\
1983-10, p.~31 & Traces of odd powers of $P - Q$ & question of the page; answer in finite dimension, over any commutative ring & \href{\leanblobmain/Notebook1983_10.lean}{\texttt{Notebook1983\_10}} & \ref{sec:quillen-1983-10}\\
1983, Lecture Notes 8 (Witten), pp.~4, 6 & $\operatorname{Tr}(S\Gamma S^*\Gamma) \le \operatorname{Tr}(SS^*)$; $\operatorname{Tr}(EMEM^*) \ge 0$ for $E \ge 0$ & classical (Cauchy--Schwarz); notes of a lecture; $\Gamma$ hermitian needed besides $\Gamma^2 = 1$ & \href{\leanblobmain/Notebook1983_LectureNotes8.lean}{\texttt{Notebook1983\_LectureNotes8}} & \ref{sec:quillen-1983-ln8}\\
1987-14, p.~49 & $\dim\operatorname{Ker}(\iota_v : \Lambda^qK^k \to \Lambda^{q-1}K^k) = \binom{k-1}{q}$ for $v = (1, \dots, 1)$ & classical (Koszul); any nonzero form & \href{\leanblobmain/Notebook1987_14.lean}{\texttt{Notebook1987\_14}} & \ref{sec:quillen-1987-14}\\
1989-3, pp.~99--100 & Powers of the block matrix $(A\ B;\ C\ 0)$ by $P_{n+2} = AP_{n+1} + BCP_n$; the involution $(z\ b;\ c\ {-z})$ & elementary; any ring & \href{\leanblobmain/Notebook1989_3.lean}{\texttt{Notebook1989\_3}} & \ref{sec:quillen-1989-3}\\
1989-4, p.~28 & $x \mapsto \operatorname{tr}(Tx)$ is faithful on $M_n(K)$ iff $T$ is invertible & classical & \href{\leanblobmain/Notebook1989_4Faithful.lean}{\texttt{Notebook1989\_4Faithful}} & \ref{sec:quillen-1989-4-faithful}\\
\addlinespace\multicolumn{5}{@{}l@{}}{\textbf{Clifford algebras, fermions and geometry}}\\*
1982-12, p.~88 & $(\gamma^1\cdots\gamma^m)^2 = (-1)^{m(m-1)/2}$; $\gamma^{2n+1}$ anticommutes with the $\gamma^i$ & standard; in any ring, from the relations & \href{\leanblobmain/Notebook1982_12.lean}{\texttt{Notebook1982\_12}} & \ref{sec:quillen-1982-12}\\
1982-13, p.~70 & The projections $1 - e(\eta_j)\,i(X_j)$ onto the joint kernel of the contractions & partly found (Mathai--Quillen; Kalkman); from the anticommutation relations alone, in any ring & \href{\leanblobmain/Notebook1982_13.lean}{\texttt{Notebook1982\_13}} & \ref{sec:quillen-1982-13}\\
1983, Lecture Notes 15 (Patterson), p.~1 & $1 + 2L(x, x') = q(p(x), p(x'))$ from the ball to the hyperboloid & classical; notes of a lecture; holds for twice the map as transcribed & \href{\leanblobmain/Notebook1983_LectureNotes15.lean}{\texttt{Notebook1983\_LectureNotes15}} & \ref{sec:quillen-1983-ln15}\\
1983, Review Getzler material, p.~7 & Composition of Cayley transforms: $Z = (1 - X)^{-1}(X + Y)(1 + XY)^{-1}(1 - X)$ & confirmed; elementary, in any ring, with $1 - Z$ invertible & \href{\leanblobmain/Notebook1983_ReviewGetzler.lean}{\texttt{Notebook1983\_ReviewGetzler}} & \ref{sec:quillen-1983-getzler}\\
1984-6, pp.~2, 48--49 & Clifford generators from a fermion; $B^{-1}A$, $DB^{-1}$ skew and $(B^{-1})^t = C - DB^{-1}A$ & standard (Berezin); three of the block relations suffice & \href{\leanblobmain/Notebook1984_6.lean}{\texttt{Notebook1984\_6}} & \ref{sec:quillen-1984-6}\\
1984-15, p.~55 & Pair symmetry $R_{ijkl} = R_{klij}$ from the skew symmetries and the cyclic identity & classical; needs $2$ cancellable, counterexample over $\mathbf F_2$ & \href{\leanblobmain/Notebook1984_15.lean}{\texttt{Notebook1984\_15}} & \ref{sec:quillen-1984-15}\\
\addlinespace\multicolumn{5}{@{}l@{}}{\textbf{Analysis, integrals and special functions}}\\*
1972-4, pp.~23--25 & Marginal of $N$ harmonic oscillators at total energy $\lambda N$: limit is the Gibbs density & classical, found (Khinchin); with the normalisation of the finite-$N$ law & \href{\leanblobmain/Notebook1972_4.lean}{\texttt{Notebook1972\_4}} & \ref{sec:quillen-1972-4}\\
1973-3, pp.~25, 79--81 & Gaussian sum over a translate of $\mathbf Z$ is maximal on $\mathbf Z$; fibres of the $\Sigma_3$-invariant of the cross-ratio & classical, found ($j/256$; Banaszczyk, Bost); dimension one, with the equality case & \href{\leanblobmain/Notebook1973_3.lean}{\texttt{Notebook1973\_3}} & \ref{sec:quillen-1973-3}\\
1982-3, p.~18 & $\prod(1 + a_n)$ converges iff $\sum a_n$ does, $a_n \ge 0$ & classical; mathlib has one direction & \href{\leanblobmain/Notebook1982_3.lean}{\texttt{Notebook1982\_3}} & \ref{sec:quillen-1982-3}\\
1983-4, pp.~71--72 & Cohomological equation of $\theta\partial_x - \partial_y$ on the torus & classical, found; both directions, for every irrational $\theta$ & \href{\leanblobmain/Notebook1983_4.lean}{\texttt{Notebook1983\_4}} & \ref{sec:quillen-1983-4}\\
1983, Lecture Notes 3 (Kirwan), p.~2 & Increasing ratios $v_j/v_{j+1}$ give $v_0^{n-1}v_n \le v_1^n$ & classical; notes of a lecture; Alexandrov--Fenchel not formalized & \href{\leanblobmain/Notebook1983_LectureNotes3.lean}{\texttt{Notebook1983\_LectureNotes3}} & \ref{sec:quillen-1983-ln3}\\
1986-5, pp.~17, 31 & $j^tj = 1$, $j^t\varepsilon j = A$; $\int_0^{2\pi}(2\sin\frac\theta2)^{2k}d\theta = 2\pi\binom{2k}k$ & elementary; Wallis & \href{\leanblobmain/Notebook1986_5.lean}{\texttt{Notebook1986\_5}} & \ref{sec:quillen-1986-5}\\
1986-12, pp.~23, 53 & $\|g^n - 1\| \le n\|g - 1\|$; $|t + 1| - |t - 1| \ge 2\operatorname{Re}t/(1 + |t|)$ and the Cayley bound & elementary & \href{\leanblobmain/Notebook1986_12.lean}{\texttt{Notebook1986\_12}} & \ref{sec:quillen-1986-12}\\
1988-3, pp.~46--47 & $(x^2 - 1)\partial_x + 4x$ is injective on $\mathbf R[x]$, with image the kernel of $q \mapsto \int_{-1}^1q(x)(x^2 - 1)dx$ & elementary & \href{\leanblobmain/Notebook1988_3.lean}{\texttt{Notebook1988\_3}} & \ref{sec:quillen-1988-3}\\
1988-9, p.~72; 1984-13, pp.~3--4; 1988-5, pp.~53, 85; 1988-7, p.~13; 1989-4, p.~56 & $\int_0^1t^a(1 - t)^bdt = a!\,b!/(a + b + 1)!$ and its consequences; $\int_0^\infty e^{-t^2}t^{2m+1}dt = m!/2$; the simplex integral & classical (Euler); not in mathlib in this real form & \href{\leanblobmain/Notebook1988_9.lean}{\texttt{Notebook1988\_9}} & \ref{sec:quillen-1988-9}\\
1988-9, pp.~95--97 & Chern--Simons values $(-1)^nn!/(2n)!$; $n!\,4^n \le (2n + 1)(2n)!/n!$, $(2n)!/n! \le n!\,4^n$ & classical; the page's $n!\,2^n$ corrected & \href{\leanblobmain/Notebook1988_9Growth.lean}{\texttt{Notebook1988\_9Growth}} & \ref{sec:quillen-1988-9-growth}\\
\bottomrule
\end{longtable}}
\end{center}

\part{General topology}
\section[Notebook 1968-13: the space $X \cup_\Phi \{\infty\}$ and the smash lemma]{Notebook 1968-13: the space $X \cup_\Phi \{\infty\}$ of a family of supports, and the smash lemma\protect\piste\protect\verifie{Notebook1968_13}}\label{sec:quillen-1968-13}

\emph{Context.} In December 1968, \carnet{1968}{13} relates supports and bordism. On page~64, Quillen takes on a space $X$ a family of supports $\Phi$: a family of closed sets, stable under finite unions and under passage to closed subsets. On page~65, he forms the space $X \cup_\Phi \{\infty\}$, in which $X$ is open and the neighbourhoods of $\infty$ are the sets $(X \setminus Z) \cup \{\infty\}$, $Z \in \Phi$. For $\Phi$ the family of closed sets, $\infty$ is an isolated point; for $\Phi$ the family of compact sets, one recovers the Alexandroff compactification $X^+$. A crossed-out passage on page~64 notes that the space is Hausdorff if $\Phi$ contains a neighbourhood of each point. On pages~67--69, after his own correction (\q{Correction}: he had defined the smash product wrongly, it being the quotient of $X \times Y$ by the wedge), he states and proves that $X^+ \wedge (Y \cup_\Phi \{\infty\}) = (X \times Y) \cup_\Psi \{\infty\}$, where $\Psi$ is the family of closed subsets of $X \times Y$ contained in some $X \times F$, $F \in \Phi$, and proper over $F$.

\emph{Intuition.} This is the argument of the page. A neighbourhood of $\infty$ in the smash product is the image of an open set containing the wedge $X^+ \times \{\infty\} \cup \{\infty\} \times (Y \cup_\Phi \{\infty\})$; its complement is a closed set $C$ of $X \times Y$, closed in the whole product. Since $X^+$ is compact, the tube lemma keeps $C$ away from $X^+ \times \{\infty\}$, that is, places it inside some $X \times F$; and $C$ is closed in $X^+ \times F$ (figure~\ref{fig:quillen-1968-13}).


\begin{figure}[htbp]
\centering
\begin{tikzpicture}[every node/.style={font=\small}]
\fill[black!8] (0,0.7) rectangle (6,2.3);
\draw[black!30] (0,0.7) -- (6,0.7) (0,2.3) -- (6,2.3);
\draw[black!55] (0,0) -- (6,0) (0,0) -- (0,4.2);
\draw[line width=1.8pt] (0,4.2) -- (6,4.2) -- (6,0);
\fill[frise!20] (1.0,1.3) .. controls (1.4,0.8) and (3.2,0.8) .. (4.2,1.2) .. controls (4.9,1.5) and (4.6,2.1) .. (3.6,2.0) .. controls (2.6,1.9) and (1.6,2.15) .. (1.2,1.9) .. controls (0.9,1.7) and (0.85,1.5) .. cycle;
\draw[frise] (1.0,1.3) .. controls (1.4,0.8) and (3.2,0.8) .. (4.2,1.2) .. controls (4.9,1.5) and (4.6,2.1) .. (3.6,2.0) .. controls (2.6,1.9) and (1.6,2.15) .. (1.2,1.9) .. controls (0.9,1.7) and (0.85,1.5) .. cycle;
\node[frise] at (2.8,1.45) {$C$};
\draw[decorate, decoration={brace, amplitude=4pt}] (-0.12,0.7) -- (-0.12,2.3) node[midway, left=5pt] {$F$};
\fill (6,0) circle (1.6pt) node[below] {$\infty$};
\fill (0,4.2) circle (1.6pt) node[left] {$\infty$};
\node[below] at (3,-0.05) {$X$};
\node[left] at (-0.05,3.3) {$Y$};
\node[anchor=west] at (6.25,2.6) {$\{\infty\} \times (Y \cup_\Phi \{\infty\})$};
\node[above] at (3,4.25) {$X^+ \times \{\infty\}$};
\node[black!60] at (5.25,1.0) {$X \times F$};
\end{tikzpicture}
\caption{The product $X^+ \times (Y \cup_\Phi \{\infty\})$, with $X^+$ on the horizontal axis and $Y \cup_\Phi \{\infty\}$ on the vertical axis. The wedge, in thick line, is the union of the upper edge $X^+ \times \{\infty\}$ and the right edge $\{\infty\} \times (Y \cup_\Phi \{\infty\})$. The complement of a neighbourhood of $\infty$ in the smash product is a closed set $C$ of $X \times Y$ (in blue). The tube lemma, applied to the compact set $X^+ \times \{\infty\}$, places it in a strip $X \times F$ with $F \in \Phi$ (in grey), and $C$ is closed in $X^+ \times F$: it does not accumulate on the right edge.}
\label{fig:quillen-1968-13}
\end{figure}

\begin{encadre}
\begin{theoreme}[Quillen, notebook 1968-13, pp.~64--69]\label{thm:quillen-1968-13}
Let $X$, $Y$ be topological spaces and $\Phi$ a family of supports on $Y$.
\begin{enumerate}
\item[a)] $X$ is an open subset of $X \cup_\Phi \{\infty\}$ and $\{\infty\}$ is closed. If $X$ is Hausdorff, $X \cup_\Phi \{\infty\}$ is Hausdorff if and only if every point of $X$ has a neighbourhood in $\Phi$.
\item[b)] For $\Phi$ the family of compact closed sets, $X \cup_\Phi \{\infty\}$ is homeomorphic to $X^+$; for $\Phi$ the family of closed sets, $\infty$ is isolated.
\item[c)] Let $\Psi$ be the family of closed subsets $A$ of $X \times Y$ contained in some $X \times F$, $F \in \Phi$, and closed in $X^+ \times Y$. The map $X^+ \times (Y \cup_\Phi \{\infty\}) \to (X \times Y) \cup_\Psi \{\infty\}$, which is the identity on $X \times Y$ and sends the wedge to $\infty$, is a quotient map: $(X \times Y) \cup_\Psi \{\infty\}$ is the smash product $X^+ \wedge (Y \cup_\Phi \{\infty\})$ with the quotient topology.
\end{enumerate}
\end{theoreme}
\end{encadre}

\noindent Quillen assumes $X$ locally compact, which makes $X^+$ Hausdorff; the proof uses only the compactness of $X^+$, and c) holds without any hypothesis. The family $\Psi$ is taken in the form that the proof on the page uses (closed in $X^+ \times F$); the equivalence with \q{proper over $F$}, which requires $X^+$ Hausdorff, is not formalized.

\begin{proof}
a) Separating $x \in X$ from $\infty$ requires an open set $U \ni x$ disjoint from some $X \setminus Z$, $Z \in \Phi$, that is, $U \subset Z$: $Z$ is a neighbourhood of $x$. Conversely, the interior of $Z$ and $(X \setminus Z) \cup \{\infty\}$ separate $x$ and $\infty$.

b) An open set containing $\infty$ has, in both topologies, a complement in $X$ that is closed and compact, or contained in a compact closed set, which amounts to the same thing. For the closed sets, $X$ itself is in $\Phi$, and $\{\infty\}$ is open.

c) The map is surjective. For a set $U$ not containing $\infty$, its preimage is the image of $U \cap (X \times Y)$ under the open embedding $X \times Y \to X^+ \times (Y \cup_\Phi \{\infty\})$. For $U \ni \infty$, the preimage is the complement of the image of the closed set $C = (X \times Y) \setminus U$, and everything comes down to seeing that this image is closed if and only if $C \in \Psi$. If it is closed, the tube lemma, applied to the compact set $X^+ \times \{\infty\}$, gives $F \in \Phi$ with $C \subset X \times F$, and the preimage under $\mathrm{id} \times \iota : X^+ \times Y \to X^+ \times (Y \cup_\Phi \{\infty\})$ shows that $C$ is closed in $X^+ \times Y$. Conversely, the image of $C$ is contained in the closed set $X^+ \times F$, which does not meet $\{\infty\}$, and it is there the image of a closed set under an embedding.
\end{proof}

\begin{sloppypar}
\noindent\emph{Verification in Lean.} \texttt{SupportFamily} and \texttt{Extension} define the family and the space; a) is \texttt{isOpenEmbedding\_ι}, \texttt{isClosed\_infty} and \texttt{t2Space\_iff}; b) is \texttt{homeomorphOnePoint}, to mathlib's compactification \texttt{OnePoint X}, \texttt{isOpen\_infty\_closedSupports} and \texttt{homeomorphSum}, to the sum of $X$ and a point; c) is \texttt{isQuotientMap\_smashMap}, with the family \texttt{smashSupports} and the key lemma \texttt{isClosed\_image\_e\_iff}, in \href{\leanblobmain/Notebook1968_13.lean}{\texttt{Notebook1968\_13.\allowbreak{}lean}}.
\end{sloppypar}

Not formalized: $\Gamma(X^+, j_!F) = \Gamma_\Phi(X, F)$ (p.~65), the exact sequence that follows from it (p.~66), and the application to bordism (pp.~66--69).

\litterature{the space is found in the literature, the smash lemma is not. G.~Beer and M.~C.~Vipera associate to every family of ideals of subsets with a closed base of a space an extension of that space, of which the Alexandroff compactification is the prototype \cite{beervipera2012}; a family of supports generates such an ideal, and the one-point extension corresponding to it is $X \cup_\Phi \{\infty\}$. G.~Beer was already studying the one-point extension of a space equipped with a bornology \cite{beer2008}, a setting that goes back to S.-T.~Hu \cite{hu1949}. These references rest on the zbMATH Open reviews (Zbl~1259.54008 and 1171.54012); the full text has not been read, and it has not been checked that the separation criterion a) appears there in this form. Supports and $\Gamma_\Phi$ are due to Godement (1958) and to Bredon (\emph{Sheaf Theory}), whose book has not been combed page by page. The smash lemma has not been found in print.}

\section[Carnet 1968-14 : préimages compactes des fermés loin de $A$]{Carnet 1968-14 : applications dont les préimages des fermés loin de $A$ sont compactes\protect\piste\protect\verifie{Notebook1968_14}}\label{sec:quillen-1968-14}

\emph{Le contexte.} Le \carnet{1968}{14}, de décembre 1968, cherche une description géométrique de l'homologie, par classes de bordisme d'applications de variétés orientées. Pour l'homologie d'un espace $X$ relative à un fermé $A$, les cycles proposés aux pages~26--27 sont des applications $f : V \to X \setminus A$ telles que $f^{-1}(B)$ soit compact pour tout fermé $B$ de $X$ disjoint de $A$. Quillen note que, si $X$ est compact, cette condition revient à dire que $f$ est propre comme application vers $X \setminus A$ : un fermé de $X$ disjoint de $A$ est alors un compact de $X \setminus A$, et réciproquement.

\emph{La page.} À la page~27, il énonce une proposition sur les limites inductives prises sur les compacts de $A$. Le cœur de la preuve est le point suivant : si $f$ vérifie la condition, l'intersection de l'adhérence de $f(V)$ avec $A$ est compacte. Il raisonne par l'absurde : un fermé non compact contient une suite discrète $(a_n)$ ; on choisit une suite exhaustive $(Q_n)$ de compacts de $V$, puis des points $x_n \notin Q_n$ dont les images $f(x_n)$ sont assez proches des $a_n$ pour former un ensemble discret. Cet ensemble est fermé, disjoint de $A$, et sa préimage n'est contenue dans aucun $Q_n$ : elle n'est pas compacte. En haut de la page~28, une remarque tire de la preuve l'équivalence de la condition avec deux propriétés : $f$ est propre, et l'adhérence de $f(V)$ dans $X$ est compacte. C'est cette remarque que l'on démontre ici. Ses espaces sont des variétés ; la preuve de la page se sert d'une exhaustion de $V$ par des compacts, donc de la $\sigma$-compacité de $V$.

\emph{L'intuition.} Une application propre vers $X \setminus A$ peut encore \q{partir à l'infini} de deux façons, vues depuis $X$ : en s'approchant de $A$, ce que la condition autorise, ou en sortant de tout compact de $X$, ce qu'elle interdit. La condition dit exactement que $f$ ne fuit que vers $A$ (figure~\ref{fig:quillen-1968-14}). La compacité de l'adhérence de $f(V)$ exclut la seconde fuite, la propreté vers $X \setminus A$ dit que la première est la seule.


\begin{figure}[htbp]
\centering
\begin{tikzpicture}[every node/.style={font=\small}]
\draw[black!55, rounded corners=12pt] (0,0) rectangle (9,4.4);
\node[black!65] at (8.6,4.0) {$X$};
\fill[frise!14] (2.901,1.340) .. controls (4,0.9) and (6,0.8) .. (7,1.6) .. controls (7.6,2.4) and (6.5,3.3) .. (5,3.1) .. controls (4,3.0) and (3.4,3.2) .. (2.901,3.060) arc[start angle=35, end angle=-35, x radius=1.1, y radius=1.5];
\draw[frise] (2.901,1.340) .. controls (4,0.9) and (6,0.8) .. (7,1.6) .. controls (7.6,2.4) and (6.5,3.3) .. (5,3.1) .. controls (4,3.0) and (3.4,3.2) .. (2.901,3.060);
\fill[piste!20] (2,2.2) ellipse [x radius=1.1, y radius=1.5];
\draw[piste] (2,2.2) ellipse [x radius=1.1, y radius=1.5];
\draw[piste, line width=1.6pt] (2.901,3.060) arc[start angle=35, end angle=-35, x radius=1.1, y radius=1.5];
\node[piste] at (1.7,2.2) {$A$};
\node[frise] at (4.6,2.35) {$f(V)$};
\fill[black!35, opacity=0.35] (6.2,2.0) ellipse [x radius=0.8, y radius=0.6];
\draw[black!70, dashed] (6.2,2.0) ellipse [x radius=0.8, y radius=0.6];
\node at (6.35,2.0) {$B$};
\end{tikzpicture}
\caption{Les conditions du théorème~\ref{thm:quillen-1968-14}. L'image $f(V)$ (en bleu) évite $A$ (en rouge) mais s'en approche : son adhérence rencontre $A$, ici le long de l'arc en trait épais. La condition~(1) demande que la préimage de tout fermé $B$ disjoint de $A$ (en gris) soit compacte : $f$ peut fuir vers $A$, et seulement vers $A$. La condition~(2) dit la même chose autrement : $f$ est propre vers $X \setminus A$, et l'adhérence de $f(V)$ est compacte.}
\label{fig:quillen-1968-14}
\end{figure}

\begin{encadre}
\begin{theoreme}[d'après Quillen, carnet 1968-14, p.~28]\label{thm:quillen-1968-14}
Soient $X$ un espace métrisable, $A$ une partie de $X$, $V$ un espace topologique et $f : V \to X$ une application continue à valeurs dans $X \setminus A$. Les conditions suivantes sont équivalentes :
\begin{enumerate}
\item pour tout fermé $B$ de $X$ disjoint de $A$, la partie $f^{-1}(B)$ de $V$ est compacte ;
\item $f$ est propre comme application de $V$ dans $X \setminus A$, et l'adhérence de $f(V)$ dans $X$ est compacte.
\end{enumerate}
\end{theoreme}
\end{encadre}

\noindent Les hypothèses sont plus faibles que celles de la page : $X$ n'est pas supposé localement compact, $A$ n'est pas supposé fermé, et $V$ est quelconque, sans $\sigma$-compacité. La preuve ci-dessous remplace l'exhaustion de $V$ par des suites de valeurs de $f$. Seule la métrisabilité de $X$ sert, dans le sens (1) $\Rightarrow$ (2), et seulement pour pouvoir raisonner sur des suites.

\begin{proof}
\emph{(2) $\Rightarrow$ (1).} Soit $K$ l'adhérence de $f(V)$, compacte. Pour un fermé $B$ disjoint de $A$, la partie $B \cap K$ est compacte, comme fermé d'un compact, et contenue dans $X \setminus A$ ; c'est donc un compact de $X \setminus A$. Comme $f(V) \subset K$, on a $f^{-1}(B) = f^{-1}(B \cap K)$, préimage d'un compact par une application propre : elle est compacte. Ce sens n'utilise aucune hypothèse sur $X$.

\emph{Un lemme sur les suites.} Supposons (1). Soit $(v_n)$ une suite de $V$ ; montrons que la suite $(f(v_n))$ a une valeur d'adhérence dans $X$. Sinon, soit $B$ l'ensemble de ses valeurs. Il est fermé : un point $y$ de l'adhérence de $B$ qui n'est pas dans $B$ serait valeur d'adhérence de la suite, car tout voisinage de $y$, privé des valeurs en nombre fini $f(v_0), \dots, f(v_{N-1})$ (qui forment un fermé ne contenant pas $y$), rencontre encore $B$, donc contient un $f(v_n)$ avec $n \ge N$. Il est disjoint de $A$, puisque $f$ évite $A$. Par (1), $f^{-1}(B)$ est compact ; il contient tous les $v_n$, qui ont donc une valeur d'adhérence $w$. Par continuité, $f(w)$ est valeur d'adhérence de $(f(v_n))$ : contradiction.

\emph{(1) $\Rightarrow$ adhérence compacte.} On munit $X$ d'une distance. Dans un espace métrique, une partie est compacte si et seulement si toute suite de cette partie a une sous-suite convergente dans la partie. Soit $(z_n)$ une suite de l'adhérence $K$ de $f(V)$. On choisit $v_n \in V$ avec $d(z_n, f(v_n)) < 1/(n+1)$. Par le lemme, $(f(v_n))$ a une valeur d'adhérence $y$, donc une sous-suite $(f(v_{\varphi(k)}))$ qui tend vers $y$ (un espace métrique est à base dénombrable de voisinages). Comme $\varphi(k) \ge k$, la distance de $z_{\varphi(k)}$ à $f(v_{\varphi(k)})$ est majorée par $1/(k+1)$, et $z_{\varphi(k)}$ tend aussi vers $y$. Enfin $y$ est limite de points de $f(V)$, donc dans $K$.

\emph{(1) $\Rightarrow$ propreté.} L'espace $X \setminus A$ est métrisable, donc séparé et à génération compacte ; dans un tel espace d'arrivée, une application continue est propre si et seulement si la préimage de tout compact est compacte. Soit $C$ un compact de $X \setminus A$. Vu dans $X$, c'est un compact, donc un fermé ($X$ est séparé), et il est disjoint de $A$. Sa préimage est compacte par (1).
\end{proof}

\noindent La preuve de la page passe par l'intersection de l'adhérence de $f(V)$ avec $A$ ; celle-ci montre directement que l'adhérence entière est compacte, ce que la remarque de la page affirme. Le choix d'une suite discrète, que Quillen fait dans $A$ et approche par des valeurs de $f$, est fait ici directement parmi les valeurs de $f$, ce qui dispense de l'exhaustion de $V$.

\begin{exemple}
Soient $X = \mathbf R$, $A = \{0\}$ et $f$ l'inclusion de $V = \,]0, 1[$. Le fermé $B = [1/2, 1]$ ne rencontre pas $A$, et $f^{-1}(B) = [1/2, 1[$ n'est pas compact : la condition (1) échoue, et $f$ n'est pas propre vers $\mathbf R \setminus \{0\}$, bien que l'adhérence $[0, 1]$ de $f(V)$ soit compacte. Avec $V = \,]0, 1]$, au contraire, $f^{-1}(B) = B \cap [0, 1]$ est compact pour tout fermé $B$ évitant $0$ : les deux conditions sont vraies, l'application fuit vers $A = \{0\}$ et seulement vers lui. Avec $A = \emptyset$, la condition (1) appliquée à $B = X$ dit que $V$ est compact, et (2) le dit aussi.
\end{exemple}

\emph{Vérification en Lean.} Le fichier \href{\leanblobmain/Notebook1968_14.lean}{\texttt{Notebook1968\_14.\allowbreak{}lean}} démontre le théorème~\ref{thm:quillen-1968-14} sous le nom \texttt{preimage\_compact\_iff}, pour $X$ muni de la classe \texttt{MetrizableSpace} et $f$ restreinte au but $A^c$ par \texttt{Set.codRestrict}. Le lemme sur les suites est \texttt{exists\_mapClusterPt}, appuyé sur \texttt{mapClusterPt\_of\_mem\_closure\_range} ; la compacité de l'adhérence est \texttt{isCompact\_closure\_range}, la propreté \texttt{isProperMap\_codRestrict}, par le critère de mathlib \texttt{isProperMap\_iff\_isCompact\_preimage} ; la réciproque est \texttt{isCompact\_preimage\_of\_isProperMap}, énoncée pour un espace $X$ quelconque. Les axiomes utilisés sont les trois axiomes usuels de mathlib (\texttt{propext}, \texttt{Classical.choice}, \texttt{Quot.sound}).

Ne sont pas démontrés ici : la proposition de la page~27 sur les limites inductives d'homologie, dont l'énoncé est le point de départ ; la description de l'homologie par bordisme ; la question de savoir si la métrisabilité de $X$ peut être remplacée par une hypothèse plus faible (espace séquentiel, par exemple).

\litterature{la recherche n'a pas retrouvé cet énoncé imprimé sous cette forme. Il est élémentaire et voisin de résultats classiques sur les applications propres (Bourbaki, \emph{Topologie générale}, chapitre~I, \S\,10) ; le cadre, homologie relative par bordisme, se rattache aux travaux de Quillen sur le cobordisme publiés en 1969--1971. L'édition ne le donne pas pour nouveau. Ce n'est pas une piste.}


\part{Algebraic $K$-theory and general linear groups}
\section[Notebooks 1968-15 and 1968-16: perfect subgroups, the elementary group and Whitehead's lemma]{Notebooks 1968-15 and 1968-16: perfect subgroups, the elementary group and Whitehead's lemma\protect\verifie{Notebook1968_15}}\label{sec:quillen-1968-15}

\emph{Context.} \carnet{1968}{15} and \carnet{1968}{16} contain the fair copy of the \q{plus} construction. In \S~1 of the first, \q{Killing a perfect subgroup of the fundamental group}, Proposition~1.2 kills a perfect subgroup $E$ of $\pi_1 X$; its proof (PDF, pp.~28--29) first replaces $E$ by the normal subgroup $N$ it generates, still perfect because $[N, N]$ is normal and contains $[E, E] = E$. In \S~2 of the second, \q{The K-theory associated to a ring} (p.~1), the construction is applied to the subgroup $E(A)$ of $GL(A)$ generated by the elementary matrices, which is perfect (Quillen refers to Bass); page~2 uses that $E(A)$ is normal in $GL(A)$, which rests on Whitehead's lemma. The ring $A$ is arbitrary there, \q{not necessarily commutative}.

\emph{Intuition.} The elementary matrix $e_{ij}(a) = 1 + a\,e_{ij}$ is a commutator as soon as a third index $k$ is available: $[e_{ik}(a), e_{kj}(b)] = e_{ij}(ab)$. For Whitehead's lemma, the block unipotent matrices $\left(\begin{smallmatrix} 1 & B \\ 0 & 1 \end{smallmatrix}\right)$ are products of elementary matrices, one for each coefficient of $B$, and three of them build the antidiagonal matrix $\left(\begin{smallmatrix} 0 & \theta \\ -\theta^{-1} & 0 \end{smallmatrix}\right)$.

\begin{encadre}
\begin{theoreme}[Quillen, notebooks 1968-15, p.~28, and 1968-16, pp.~1--2]\label{thm:quillen-1968-15}
Let $G$ be a group and $A$ a unital ring.
\begin{enumerate}
\item[a)] The normal subgroup generated by a perfect subgroup of $G$ is perfect.
\item[b)] If $i$, $k$, $j$ are pairwise distinct, $[e_{ik}(a), e_{kj}(b)] = e_{ij}(ab)$; in particular, the subgroup $E(n, A)$ of invertible matrices generated by the $e_{ij}(a)$ is perfect as soon as $n \ge 3$.
\item[c)] (Whitehead's lemma.) For $\theta$ invertible in $M_n(A)$, the matrix $\operatorname{diag}(\theta, \theta^{-1})$ is in $E(2n, A)$; for $\theta$, $\eta$ invertible, $\operatorname{diag}(\theta\eta\theta^{-1}\eta^{-1}, 1)$ is too.
\end{enumerate}
\end{theoreme}
\end{encadre}

\begin{proof}
a) Let $H$ be perfect and $N$ the normal subgroup it generates. The commutator subgroup $[N, N]$ is normal and contains $[H, H] = H$, hence contains $N$, so $[N, N] = N$.

b) Put $X = a\,e_{ik}$ and $Y = b\,e_{kj}$. Then $X^2 = Y^2 = YX = 0$, $XY = ab\,e_{ij}$, and $XY$ annihilates $X$ and $Y$ on the right; one checks that $(1 + X)(1 + Y)(1 - X)(1 - Y) = 1 + XY$. For $n \ge 3$, every generator $e_{ij}(a)$ is the commutator of $e_{ik}(a)$ and $e_{kj}(1)$, for a third index $k$.

c) $\left(\begin{smallmatrix} 1 & B \\ 0 & 1 \end{smallmatrix}\right)\left(\begin{smallmatrix} 1 & C \\ 0 & 1 \end{smallmatrix}\right) = \left(\begin{smallmatrix} 1 & B + C \\ 0 & 1 \end{smallmatrix}\right)$, and $\left(\begin{smallmatrix} 1 & a\,e_{ij} \\ 0 & 1 \end{smallmatrix}\right)$ is elementary; likewise for lower triangular matrices. A block computation gives
\[
\begin{pmatrix} 1 & \theta \\ 0 & 1 \end{pmatrix}\begin{pmatrix} 1 & 0 \\ -\theta^{-1} & 1 \end{pmatrix}\begin{pmatrix} 1 & \theta \\ 0 & 1 \end{pmatrix} = \begin{pmatrix} 0 & \theta \\ -\theta^{-1} & 0 \end{pmatrix},
\]
and the product of this matrix with the same one for $\theta = -1$ is $\operatorname{diag}(\theta, \theta^{-1})$. Finally
\[
\operatorname{diag}(\theta\eta\theta^{-1}\eta^{-1}, 1) = \operatorname{diag}(\theta, \theta^{-1})\operatorname{diag}(\eta, \eta^{-1})\operatorname{diag}((\eta\theta)^{-1}, \eta\theta).
\]
\end{proof}

\begin{sloppypar}
\noindent\emph{Verification in Lean.} a) is \texttt{normalClosure\_isPerfect}, b) is \texttt{commutator\_transvectionUnit} and \texttt{elementary\_isPerfect}, with the ring identity \texttt{commutator\_aux}, in \href{\leanblobmain/Notebook1968_15.lean}{\texttt{Notebook1968\_15.\allowbreak{}lean}}; c) is \texttt{exists\_mem\_elementary\_diag} and \texttt{exists\_mem\_elementary\_commutator}, with \texttt{upper\_mem\_elementary}, \texttt{lower\_mem\_elementary} and \texttt{upper\_lower\_upper}, in \href{\leanblobmain/Notebook1968_16.lean}{\texttt{Notebook1968\_16.\allowbreak{}lean}}. The matrices are the units of the ring \texttt{Matrix n n A}, for an arbitrary ring $A$.
\end{sloppypar}

\begin{sloppypar}\noindent Not formalized: the stable statements ($E(A)$ is normal in $GL(A)$, and $E(A)$ is the commutator subgroup $[GL(A), GL(A)]$), which require the inductive limit over $n$, and the \q{plus} construction itself.
\end{sloppypar}

\litterature{found in the literature: Milnor, \emph{Introduction to Algebraic K-Theory} (1971), \S~3 and Lemma~3.1; Rosenberg, \emph{Algebraic K-Theory and its Applications}, \S~2.1. mathlib has perfect groups (\texttt{Group.IsPerfect}) but none of these three statements. They are proposed for mathlib, and the edition does not present them as new.}

\section[Notebook 1971-4: a monoid and its Grothendieck group]{Notebook 1971-4: the group completion of a commutative monoid is a flat functor\protect\verifie{Notebook1971_4}}\label{sec:quillen-1971-4}

\emph{Context.} \carnet{1971}{4} (20 March -- 11 April 1971) studies torsors under a small category $\mathcal{C}$ over a space or a topos, with a view to Mather's theorem. On PDF page~24 (its \q{8}, in the second pagination), Quillen says that a functor $u : \mathcal{C}_1 \to \mathcal{C}_2$ is \emph{strongly cofinal} if $(u_!, u^*)$ is a morphism of presheaf topoi, which amounts to requiring that, for every object $Y$ of $\mathcal{C}_2$, the functor $X \mapsto \operatorname{Hom}(Y, uX)$ be a pro-object of $\mathcal{C}_1$, in other words that its category of elements be cofiltered. He notes that such a $u$ induces isomorphisms on the $R^q\varprojlim$. On page~25 (its \q{9}), he tests the definition on the universal map from a commutative monoid $M$ to its group completion $G$, viewed as a functor between one-object categories: the condition becomes that the category whose objects are the elements of $G$, and whose arrows are given by translations by $M$, is cofiltered, and he checks the two axioms in a few lines.

\emph{Intuition.} An arrow $g \to g'$ is an $m \in M$ such that $mg = g'$. Two elements of $G$ can be written with a common denominator, $g_i = m_i\,b^{-1}$, and are thus translates of $b^{-1}$: this is the axiom on products. Two parallel arrows $m_1, m_2 : g \to g'$ have the same image in $G$; by the construction of the Grothendieck group, $c\,m_1 = c\,m_2$ for some $c \in M$, and precomposing with the arrow $c : c^{-1}g \to g$ equalizes them: this is the axiom on equalizers.

\begin{encadre}
\begin{theoreme}[Quillen, notebook 1971-4, p.~25]\label{thm:quillen-1971-4}
Let $M$ be a commutative monoid, $G(M)$ its Grothendieck group and $u : M \to G(M)$ the canonical map. The induced functor between one-object categories is representably flat: for the unique object $X$, the category $X/u$ is cofiltered. In particular, it is final.
\end{theoreme}
\end{encadre}

\begin{proof}
The objects of $X/u$ are the elements of $G(M)$, and an arrow $g \to g'$ is an $m \in M$ such that $u(m)\,g = g'$. The category is nonempty (it contains $1$).

Every $x \in G(M)$ satisfies $u(b)\,x = u(a)$ for some $a, b \in M$. For $g_1, g_2$, choose $u(b_i)\,g_i = u(a_i)$ and put $g = u(b_1 b_2)^{-1}$; then $u(b_2 a_1)\,g = g_1$ and $u(b_1 a_2)\,g = g_2$.

If $m_1, m_2 : g \to g'$, then $u(m_1)\,g = u(m_2)\,g$, so $u(m_1) = u(m_2)$, so $c\,m_1 = c\,m_2$ for some $c \in M$, since $G(M)$ is the localization of $M$ at all of $M$. The arrow $c : u(c)^{-1} g \to g$ equalizes $m_1$ and $m_2$.

A representably flat functor is final (mathlib, \texttt{final\_of\_representablyFlat}).
\end{proof}

\begin{sloppypar}
\noindent\emph{Verification in Lean.} In \href{\leanblobmain/Notebook1971_4.lean}{\texttt{Notebook1971\_4.\allowbreak{}lean}}, the functor is \texttt{toGroupFunctor}, induced by \texttt{Algebra.GrothendieckGroup.of}; the theorem is \texttt{representablyFlat\_toGroupFunctor}, in the sense of mathlib's \texttt{CategoryTheory.RepresentablyFlat} (all the categories \texttt{StructuredArrow X F} cofiltered), and \texttt{final\_toGroupFunctor}. Quillen's two axioms are isolated as statements about the Grothendieck group: \texttt{exists\_common} (products) and \texttt{exists\_mul\_eq\_mul} (equalizers), with \texttt{exists\_mul\_eq\_of}. The monoid is written multiplicatively, as on the page.
\end{sloppypar}

Not formalized: the equivalence of the definitions on page~24 (morphism of topoi, pro-objects, cofiltered categories of elements), the remark on the $R^q\varprojlim$, and the statements on torsors, which mathlib does not have over a small category.

\litterature{folklore of group completion, found in the literature: flat functors and filtered categories of elements, SGA~4, exposés~I and~IV; Mac~Lane--Moerdijk, \emph{Sheaves in Geometry and Logic} (1992), VII.6. mathlib has \texttt{RepresentablyFlat} and the Grothendieck group, but not this example. It is proposed for mathlib, and the edition does not present it as new.}

\section[Notebook 1971-7: the monoid of injections of an infinite set]{Notebook 1971-7: the monoid of injections of an infinite set has trivial group completion\protect\verifie{Notebook1971_7}}\label{sec:quillen-1971-7}

\emph{Context.} \carnet{1971}{7} (1--23 May 1971) looks for small categories with the homotopy type of $BG^+$, for $G$ a group of diffeomorphisms with compact support, the infinite symmetric group or $GL(R)$. The method, in an entry dated 5 May 1971, is to build a category cofibred over a monoid $M$ of \q{embeddings}, whose fibres are the group, and to show that the one-object category of $M$ is homotopically trivial. For $S$ an infinite set and $M$ the monoid of injective maps $S \to S$, Quillen first shows, on page~14, that the fundamental group of this category is trivial, that is, that the group generated by $M$ is trivial; he then shows that $M$ has the integral homology of a point, by Mather's trick of infinite sums, and deduces on pages~17--19 that the category cofibred over finite sets and injections, with fibres the symmetric groups, has the homotopy type of $B\Sigma_\infty^+$.

\emph{The page.} Quillen's argument for the fundamental group runs as follows. Two injections that agree on a subset of $S$ of the same cardinality as $S$ have the same image in the group, since one can precompose both with an injection onto that subset. An injection $u$ whose image has a complement as large as $S$ defines a \q{transport} $u_*(\theta)$ of each $\theta \in M$: it is $u\theta u^{-1}$ on the image of $u$ and the identity off it. Then $u_*(\theta)\, u = u\, \theta$, while $u_*(\theta)$ agrees with the identity on a large set, hence becomes trivial in the group; so $\theta$ does too.

\emph{Intuition.} An injection with a large co-image behaves like an invertible element: conjugation by it is defined, as the transport, even though $u$ itself has no inverse. The transport of any $\theta$ is \q{small}, since it fixes everything outside the image of $u$, and in a group the relation $u_*(\theta)\,u = u\,\theta$ then forces $\theta$ to vanish.

\begin{encadre}
\begin{theoreme}[Quillen, notebook 1971-7, p.~14]\label{thm:quillen-1971-7}
Let $S$ be an infinite set and $M$ the monoid, under composition, of injective maps $S \to S$. Every monoid homomorphism from $M$ to a group is trivial. Equivalently, the group generated by $M$, the fundamental group of the one-object category of $M$, is trivial.
\end{theoreme}
\end{encadre}

\begin{proof}
Since $S$ is infinite, $S \sqcup S$ is in bijection with $S$; composing a bijection with the two inclusions gives injections $t, s : S \to S$ with disjoint images. For $f \in M$, let $c(f)$ be the map equal to $t f t^{-1}$ on the image of $t$ and to the identity elsewhere. It is injective: on the image of $t$ it is a composite of injections and takes values in the image of $t$, and off that image it is the identity. One has $c(f)\, t = t\, f$ and, since $s$ takes values off the image of $t$, $c(f)\, s = s$. Let $\varphi : M \to G$ be a homomorphism to a group. The second relation gives $\varphi(c(f))\,\varphi(s) = \varphi(s)$, so $\varphi(c(f)) = 1$; the first then gives $\varphi(t) = \varphi(t)\,\varphi(f)$, so $\varphi(f) = 1$.
\end{proof}

\noindent The proof keeps Quillen's transport but replaces his first step (injections agreeing on a large set have the same image) by the second injection $s$, which kills the transport directly. It is the shortcut of the published proof cited below.

\emph{Verification in Lean.} The file \href{\leanblobmain/Notebook1971_7.lean}{\texttt{Notebook1971\_7.\allowbreak{}lean}} defines the monoid $M$ as the submonoid \texttt{injMonoid S} of injective maps in \texttt{Function.End S}, and the transport as \texttt{transport t f}, the extension of $t \circ f$ along $t$ by the identity. Its properties are \texttt{transport\_comp\_self}, \texttt{transport\_comp\_of\_disjoint} and \texttt{transport\_injective}; the two injections with disjoint images are \texttt{exists\_disjoint\_injections}. Theorem~\ref{thm:quillen-1971-7} is \texttt{monoidHom\_eq\_one}, for an arbitrary infinite type, and \texttt{monoidHom\_eq\_one'} as an equality of homomorphisms. The axioms used are the three usual axioms of mathlib.

Not formalized here are the vanishing of the higher homology of $M$ (Mather's trick, pages~15--16) and the homotopy statements of pages~17--19.

\litterature{the statement is in print: S.~Schwede, \q{On the homotopy groups of symmetric spectra}, \emph{Geom.\ Topol.}~12 (2008), Lemma~5.2, proves that the classifying space of the monoid of injections of $\mathbf N$ is contractible, with a proof credited to N.~Strickland that uses the same transport ($c_t(f)\,t = t f$) together with a second injection; the paper notes before the lemma that $\operatorname{Hom}(M, A) = 0$ for every abelian group $A$. The search did not locate the case of an arbitrary infinite set stated separately, which follows by the same proof. The edition does not present the statement as new.}

\section[Notebook 1973-6: commutators of elementary matrices; a Sylow reduction]{Notebook 1973-6: commutators of elementary matrices, and a Sylow reduction\protect\verifie{Notebook1973_6}}\label{sec:quillen-1973-6}

\emph{Context.} \carnet{1973}{6} holds Quillen's notes of Serre's course on trees, in the autumn of 1973, followed by his own pages on shift equivalence. Two short arguments of the course are checked here.

\emph{The page, p.~30.} To show that $SL_3(\mathbf Z)$ has Serre's property (FA), the page writes two commutator relations between the elementary matrices $\bar e_{ij} = 1 + e_{ij}$, for $i, j, k$ distinct: $(\bar e_{ij}, \bar e_{jk}) = \bar e_{ik}$ and $(\bar e_{ij}, \bar e_{kj}) = \bar e_{ik}^{-1}$, with $(a, b) = aba^{-1}b^{-1}$. It then places the six matrices on a hexagon, in the order $\bar e_{12}, \bar e_{13}, \bar e_{23}, \bar e_{21}, \bar e_{31}, \bar e_{32}$, and notes that each commutes with its neighbours and is the commutator of its two neighbours. The first relation is the usual Steinberg relation. The second cannot hold as written: $\bar e_{ij}$ and $\bar e_{kj}$ lie in the same column, $e_{ij}e_{kj} = e_{kj}e_{ij} = 0$, and they commute. The relation that holds is $(\bar e_{ij}, \bar e_{ki}) = \bar e_{kj}^{-1}$; the hexagon needs only the first one.

\emph{The page, p.~45.} In the proof of Brown's theorem that an integer dividing the length of every orbit of a finite group divides the Euler characteristic, the page reduces to a Sylow $p$-subgroup $H$: if $p^k$ divides $[G : G_x]$, it divides $[H : H_x]$, because $[G:G_x]$ divides $[G : H_x] = [G : H][H : H_x]$ and $[G : H]$ is prime to $p$.

\begin{encadre}
\begin{theoreme}[After Quillen, notebook 1973-6, p.~30]\label{thm:quillen-1973-6-a}
Let $R$ be a commutative ring, $i, j, k$ distinct indices, $a, b \in R$, and $t_{ij}(a) = 1 + a e_{ij}$. Then $(t_{ij}(a), t_{jk}(b)) = t_{ik}(ab)$; $t_{ij}(a)$ commutes with $t_{kj}(b)$ and with $t_{ik}(b)$, so that $(t_{ij}(a), t_{kj}(b)) = 1$, which is not $t_{ik}(-1)$ if $R \ne 0$; and $(t_{ij}(a), t_{ki}(b)) = t_{kj}(-ba)$. For $3 \times 3$ matrices, each matrix of the hexagon commutes with the next one and is the commutator, in one order or the other, of its two neighbours.
\end{theoreme}
\end{encadre}

\begin{proof}
Put $x = ae_{ij}$ and $y = be_{jk}$. Then $x^2 = y^2 = yx = 0$ and $xy = ab\,e_{ik}$, $xyx = 0$; expanding, $(1+x)(1+y)(1-x)(1-y) = 1 + xy$. For $y = be_{ki}$ instead, $xy = 0$, $yx = ba\,e_{kj}$ and $yxy = 0$, and the same expansion gives $1 - yx$. Two elementary matrices in the same row or the same column have products $0$ in both orders, so they commute. The hexagon is the six instances of the first relation: $\bar e_{12} = (\bar e_{13}, \bar e_{32})$, $\bar e_{13} = (\bar e_{12}, \bar e_{23})$, $\bar e_{23} = (\bar e_{21}, \bar e_{13})$, and so on, and consecutive matrices share a row or a column.
\end{proof}

\begin{encadre}
\begin{theoreme}[Quillen, notebook 1973-6, p.~45]\label{thm:quillen-1973-6-b}
Let $H$ be a Sylow $p$-subgroup of a finite group $G$ and $K$ a subgroup. If $p^k$ divides $[G : K]$, then $p^k$ divides $[H : H \cap K]$.
\end{theoreme}
\end{encadre}

\begin{proof}
$[G : K]$ divides $[G : H \cap K] = [H : H \cap K]\,[G : H]$, and $[G:H]$ is prime to $p$.
\end{proof}

\emph{Verification in Lean.} The file \href{\leanblobmain/Notebook1973_6.lean}{\texttt{Notebook1973\_6.\allowbreak{}lean}} proves Theorem~\ref{thm:quillen-1973-6-a} with mathlib's \texttt{Matrix.transvection}, as \texttt{transvection\_commutator}, \texttt{transvection\_commute\_col}, \texttt{transvection\_commute\_row}, \texttt{transvection\_commutator\_col}, \texttt{not\_second\_relation}, \texttt{transvection\_commutator'} and, for the hexagon, \texttt{hexagon\_commute} and \texttt{hexagon\_commutator}; the commutators are written out as products, the inverse of $t_{ij}(a)$ being $t_{ij}(-a)$. The first relation was already in the library for units over any ring (Section~\ref{sec:quillen-1968-15}). Theorem~\ref{thm:quillen-1973-6-b} is \texttt{pow\_dvd\_relIndex\_of\_sylow}. The property (FA), the trees and Brown's theorem are not formalised. The axioms used are the three usual axioms of mathlib.

\litterature{Both statements are classical: the commutator relations are Steinberg's relations for $SL_n$, and the hexagon argument is in Serre's \emph{Arbres, amalgames, $SL_2$} (I, \S 6); the Sylow reduction is the first step of Brown's proof (Invent.\ Math.\ 27, 1974). The point of the section is the second relation of p.~30, which fails as transcribed and is replaced by the one that holds. The edition does not present anything here as new.}

\section[Notebook 1973-8: cancelling a free summand along a unimodular vector]{Notebook 1973-8: $A \oplus M \cong A \oplus M'$ with $f(1, m)$ of unit first coordinate gives $M \cong M'$\protect\verifie{Notebook1973_8Splitting}}\label{sec:quillen-1973-8-splitting}

\emph{Context.} On pages~26--28 of \carnet{1973}{8} Quillen reviews the Jacobson radical and Bass's cancellation for semi-local rings, which he reduces to the splitting argument: an isomorphism $A \oplus M \cong A \oplus M'$ sending a vector $(1, m)$ to a vector whose first coordinate is a unit induces $M \cong M'$.

\begin{encadre}
\begin{theoreme}[Quillen, notebook 1973-8, p.~27]\label{thm:quillen-1973-8-splitting}
Let $A$ be a commutative ring, $M, M'$ modules and $f : A \times M \to A \times M'$ an isomorphism. If the first coordinate of $f(1, m)$ is a unit for some $m \in M$, then $M \cong M'$.
\end{theoreme}
\end{encadre}

\begin{proof}
Write $f(1, m) = u(1, m')$. The map $\pi_m(a, n) = n - am$ is onto $M$ with kernel $A(1, m)$, and similarly $\pi_{m'}$. Then $\varphi = \pi_{m'} \circ f \circ \iota_2$ and $\psi = \pi_m \circ f^{-1} \circ \iota_2$ are inverse: $f(0, n) - (0, \varphi(n)) = b(1, m')$ with $b$ the first coordinate of $f(0, n)$, so $f^{-1}(0, \varphi(n)) = (0, n) - bu^{-1}(1, m)$, whose image under $\pi_m$ is $n$; symmetrically for $\varphi\psi$.
\end{proof}

\emph{Verification in Lean.} The file \href{\leanblobmain/Notebook1973_8Splitting.lean}{\texttt{Notebook1973\_8Splitting.\allowbreak{}lean}} proves \texttt{proj\_symm\_proj} and \texttt{nonempty\_equiv}, with the explicit inverse pair. Bass's lemma (a unit in $a + I$ when $Aa + I = A$, semi-local $A$) and the noncommutative case are not formalised. The axioms used are the three usual axioms of mathlib.

\litterature{Classical: Bass (Publ.\ Math.\ IHÉS 22, 1964), stable range one for semi-local rings. The edition does not present it as new.}

\section[Notebook 1973-12: the kernel of $GL_m(A) \to GL_m(A/I)$ for $I^2 = 0$]{Notebook 1973-12: the kernel of $GL_m(A) \to GL_m(A/I)$ for $I^2 = 0$\protect\verifie{Notebook1973_12}}\label{sec:quillen-1973-12}

\emph{Context.} On page~48 of \carnet{1973}{12} Quillen shows that $K_*(A) \otimes \mathbf Z[1/p] \cong K_*(A/I) \otimes \mathbf Z[1/p]$ for a nilpotent two-sided ideal $I$ with $p^rI = 0$, and deduces that $K_*(A) \otimes \mathbf Z[1/n] \cong K_*(A[t]) \otimes \mathbf Z[1/n]$ when $nA = 0$.

\emph{The page.} After reducing to $I^2 = 0$ he writes the exact sequence $0 \to N \to GL_m(A) \to GL_m(A/I) \to 0$, with $N$ isomorphic to the additive group of matrices with entries in $I$; hence $p^rN = 0$, the homology of $N$ vanishes after inverting $p$, and $GL_m(A) \to GL_m(A/I)$ is a homology isomorphism after inverting $p$.

\begin{encadre}
\begin{theoreme}[Quillen, notebook 1973-12, p.~48]\label{thm:quillen-1973-12}
Let $J$ be a two-sided ideal of a ring $R$ with $xy = 0$ for $x, y \in J$ (for $R = M_m(A)$: $J = M_m(I)$ with $I^2 = 0$). Then the units $u$ with $u - 1 \in J$ form a subgroup $N$, and $u \mapsto u - 1$ is an isomorphism from $N$ onto the additive group of $J$; if $nJ = 0$ then $u^n = 1$ on $N$. If $uv - 1$ and $vu - 1$ lie in $J$, then $u$ is a unit, so units lift modulo $J$.
\end{theoreme}
\end{encadre}

\begin{proof}
For $x, y \in J$, $(1 + x)(1 + y) = 1 + x + y$ and $(1 + x)(1 - x) = 1$. If $uv = 1 + x$ and $vu = 1 + y$ with $x, y \in J$, then $v(1 - x)$ is a right inverse and $(1 - y)v$ a left inverse of $u$, and they coincide. If $I^2 = 0$, a product of two matrices with entries in $I$ is $0$.
\end{proof}

\emph{Verification in Lean.} The file \href{\leanblobmain/Notebook1973_12.lean}{\texttt{Notebook1973\_12.\allowbreak{}lean}} proves the theorem as \texttt{kernel}, \texttt{kernelEquiv}, \texttt{pow\_eq\_one\_of\_kernel} and \texttt{isUnit\_of\_lift}, for a \texttt{TwoSidedIdeal} of square zero, and the matrix case as \texttt{matrix\_sq\_zero}, with mathlib's \texttt{TwoSidedIdeal.matrix}. The identification of $N$ with the kernel of reduction modulo $I$ is the definition $u - 1 \in M_m(I)$. The homology and the $K$-theory are not formalised. The axioms used are the three usual axioms of mathlib.

\litterature{Classical (for instance Weibel, \emph{The $K$-book}, for the nilpotent extension argument). The edition does not present it as new.}

\section[Notebook 1974-17: the commutator trick of the Whitehead lemma]{Notebook 1974-17: $(\alpha, \beta) = (\alpha\sigma\alpha^{-1}\sigma^{-1}, \beta\tau\beta^{-1}\tau^{-1})$\protect\verifie{Notebook1974_17}}\label{sec:quillen-1974-17}

\emph{Context.} On page~3 of \carnet{1974}{17}, to prove that the commutator subgroup of $GL(A)$ is perfect, Quillen moves two matrices of $GL_n(A)$ into a block of $GL_{3n}(A)$ and rewrites their commutator as a commutator of two commutators with even permutation matrices; a second proof places the commutator in the normal subgroup of $GL_{2n}(A)$ generated by permutations.

\begin{encadre}
\begin{theoreme}[Quillen, notebook 1974-17, p.~3]\label{thm:quillen-1974-17}
Let $A$ be any ring and $\alpha, \beta \in GL_n(A)$, embedded in the first of three diagonal blocks of $GL_{3n}(A)$. Let $\sigma$ and $\tau$ be the permutation matrices of the cyclic permutations of the three blocks which send the first block to the second and to the third; they are even. Then
\[
(\alpha, \beta) = (\alpha\sigma\alpha^{-1}\sigma^{-1},\ \beta\tau\beta^{-1}\tau^{-1}).
\]
Moreover, in any group, if $\sigma\beta\sigma^{-1}$ commutes with $\alpha$ then $(\alpha, \beta)$ lies in the normal subgroup generated by $\sigma$; so $(GL_n, GL_n)$ lies in the normal subgroup of $GL_{2n}(A)$ generated by the swap of the two blocks.
\end{theoreme}
\end{encadre}

\begin{proof}
In any group, $(xa, yb) = (x, y)$ when $a$ commutes with $y$ and $b$, and $b$ with $x$. Here $x = \alpha$, $y = \beta$, $a = \sigma\alpha^{-1}\sigma^{-1}$ (in the second block), $b = \tau\beta^{-1}\tau^{-1}$ (in the third block), and matrices in different blocks commute. For the second statement, $\sigma$ is trivial modulo its normal closure, so $\beta \equiv \sigma\beta\sigma^{-1}$ commutes with $\alpha$ there.
\end{proof}

\emph{Verification in Lean.} The file \href{\leanblobmain/Notebook1974_17.lean}{\texttt{Notebook1974\_17.\allowbreak{}lean}} proves \texttt{commutator\_mul\_mul}, \texttt{commute\_blockUnit}, \texttt{blockPerm\_conj}, \texttt{sign\_blockPerm\_finRotate}, \texttt{commutator\_eq}, \texttt{commutator\_mem\_normalClosure} and \texttt{commutator\_mem\_normalClosure\_swap}, with matrices indexed by $\iota \times \mathrm{Fin}\,r$. The block swap of $GL_{2n}$ is a product of $n$ transpositions, even only for $n$ even; the page's normal subgroup generated by the alternating group is reached by composing the swap with an even correction inside the second block, which is not formalised, nor is the perfectness of the alternating groups or anything about $BGL(A)^+$. The axioms used are the three usual axioms of mathlib.

\litterature{Classical (the Whitehead lemma; Milnor, \emph{Introduction to algebraic $K$-theory}, \S3); the proof through alternating groups was not compared with a printed source; the edition does not present it as new.}

\section[Notebook 1982-5: complementary minors, the $K_1$ pairing, parametrices]{Notebook 1982-5: complementary minors, the pairing $g \mapsto \tau(g^{-1}, g)$, parametrices\protect\verifie{Notebook1982_5}}\label{sec:quillen-1982-5}

\emph{Context.} In \carnet{1982}{5}, among determinant-line computations, three algebraic steps recur. On page~15 the determinant of the compression of $T^{-1}$ to a subspace is compared with that of the complementary compression of $T$ (the same identity reappears in \carnet{1985}{1}, pp.~7--9). On pages~55--56 a skew cyclic $1$-cocycle $\tau$ on an algebra is evaluated on $(g^{-1}, g)$ for invertible $g$. On pages~41 and~60 a parametrix $P$ of an operator $D$, with errors $K_0 = 1 - PD$ and $K_1 = 1 - DP$, is improved to one with errors $K_0^n$ and $K_1^n$.

\begin{encadre}
\begin{theoreme}[Quillen, notebook 1982-5, pp.~15, 41, 55--56, 60]\label{thm:quillen-1982-5}
\begin{enumerate}
\item (Jacobi.) Over a commutative ring, if $M = \begin{psmallmatrix} A & B\\ C & D\end{psmallmatrix}$ is invertible, then $\det M \cdot \det\bigl((M^{-1})_{22}\bigr) = \det A$.
\item Let $A$ be a ring, $k$ an abelian group and $\tau : A \times A \to k$ with $\tau(x, y) = -\tau(y, x)$ and $\tau(ax, y) - \tau(x, ya) + \tau(xy, a) = 0$. Then $g \mapsto \tau(g^{-1}, g)$ is a homomorphism $A^\times \to k$. The coboundaries $\tau(x, y) = \varphi(xy) - \varphi(yx)$ satisfy the hypotheses and give the zero homomorphism.
\item In a ring, if $PD = 1 - K_0$ and $DP = 1 - K_1$, then $P_n = \sum_{k=1}^n (-1)^{k-1}\binom nk (PD)^{k-1}P$ satisfies $P_nD = 1 - K_0^n$ and $DP_n = 1 - K_1^n$; $P_2 = P + K_0P$. If moreover $DQ = 1 - p_1$, then $P - Q = Pp_1 - K_0Q$.
\end{enumerate}
\end{theoreme}
\end{encadre}

\begin{proof}
(1) From $MM^{-1} = 1$ in blocks, $M\begin{psmallmatrix} 1 & (M^{-1})_{12}\\ 0 & (M^{-1})_{22}\end{psmallmatrix} = \begin{psmallmatrix} A & 0\\ C & 1\end{psmallmatrix}$; take determinants of the triangular blocks. (2) The cocycle identity at $(a, x, y) = (h^{-1}, g^{-1}, gh)$ reads $\tau((gh)^{-1}, gh) - \tau(g^{-1}, g) + \tau(h, h^{-1}) = 0$, and $\tau(h, h^{-1}) = -\tau(h^{-1}, h)$. (3) $P_nD = \sum (-1)^{k-1}\binom nk (PD)^k = 1 - (1 - PD)^n$, and $D(PD)^{k-1}P = (DP)^k$ gives the other side; the last identity is $P(1 - DQ) - (1 - PD)Q$.
\end{proof}

\emph{Verification in Lean.} The file \href{\leanblobmain/Notebook1982_5.lean}{\texttt{Notebook1982\_5.\allowbreak{}lean}} proves (1) as \texttt{complementary\_minor}, (2) as \texttt{pairing\_mul}, \texttt{pairing} (a monoid homomorphism to \texttt{Multiplicative k}), \texttt{coboundary\_isCocycle} and \texttt{pairing\_coboundary}, and (3) as \texttt{parametrix\_mul}, \texttt{mul\_parametrix}, \texttt{Pn\_two} and \texttt{sub\_eq\_of\_parametrices}. No linearity of $\tau$ is used, nor commutativity in (3). The Fredholm determinants and the trace-class setting are not formalised. The axioms used are the three usual axioms of mathlib.

\litterature{All classical: Jacobi's complementary minor theorem (Gantmacher, \emph{Theory of matrices}, I, ch.~I, \S 4; mathlib has the Schur complement but not this statement), the pairing of odd cyclic cocycles with $K_1$ (Connes, \emph{Non-commutative differential geometry}, 1985), and the iteration of a parametrix in Fredholm theory. The edition does not present them as new.}

\section[Notebook 1988-7: the idempotent of an index problem]{Notebook 1988-7: the idempotent of an index problem; two involutions\protect\verifie{Notebook1988_7}}\label{sec:quillen-1988-7}

\emph{Context.} On page~30 of \carnet{1988}{7} Quillen writes the idempotent representing the index of $p$ when $qp = 1 - x$ and $pq = 1 - y$ with $x$, $y$ small (of trace class), in a form adapted to the traces of odd powers; on page~33 he compares two involutions $F$ and $\varepsilon$.

\begin{encadre}
\begin{theoreme}[Quillen, notebook 1988-7, pp.~30, 33]\label{thm:quillen-1988-7}
In a ring, let $qp = 1 - x$ and $pq = 1 - y$. Then $px = yp$ and $xq = qy$. For $n \in \mathbf N$ put $\tilde q = (1 + x + \cdots + x^{2n-1})q$; then $\tilde qp = 1 - x^{2n}$, $p\tilde q = 1 - y^{2n}$, the row $(\tilde q, x^n)$ times the column $(p; x^n)$ is $1$, and
\[
e = \begin{pmatrix} p\\ x^n\end{pmatrix}\begin{pmatrix}\tilde q & x^n\end{pmatrix} = \begin{pmatrix} 1 - y^{2n} & px^n\\ x^n\tilde q & x^{2n}\end{pmatrix}
\]
is idempotent; $e - \operatorname{diag}(1, 0)$ has its entries in the ideal generated by $x$ and $y$. If $F^2 = \varepsilon^2 = 1$, then $(F - \varepsilon)^2 = 2 - F\varepsilon - \varepsilon F$ commutes with $F$ and with $\varepsilon$.
\end{theoreme}
\end{encadre}

\begin{proof}
$px = p - pqp = yp$. The geometric sum gives $\tilde qp = (\sum_{k<2n}x^k)(1 - x) = 1 - x^{2n}$, and $p\,x^k = y^kp$ gives $p\tilde q = (\sum_{k<2n} y^k)(1 - y) = 1 - y^{2n}$. A column times a row whose product the other way is $1$ is idempotent. The last identities are expansions.
\end{proof}

\emph{Verification in Lean.} The file \href{\leanblobmain/Notebook1988_7.lean}{\texttt{Notebook1988\_7.\allowbreak{}lean}} proves the theorem as \texttt{mul\_x}, \texttt{x\_mul}, \texttt{qt\_mul}, \texttt{mul\_qt}, \texttt{row\_mul\_col}, \texttt{isIdempotentElem\_e}, \texttt{e\_entries}, \texttt{sub\_sq\_eq}, \texttt{commute\_sub\_sq\_left} and \texttt{commute\_sub\_sq\_right}, in any ring. The trace formula for the index is not formalised. The axioms used are the three usual axioms of mathlib.

\litterature{Classical: the idempotent of the boundary map $K_1 \to K_0$ (Milnor, \emph{Introduction to algebraic $K$-theory}, \S 3) in the form used for Fredholm modules (Connes, 1985). The edition does not present it as new.}

\section[Notebook 1989-1: the idempotent-lifting polynomial]{Notebook 1989-1: the idempotent-lifting polynomial; a unit of $\mathbf Q[x]/(x^{2n-1} - 1)$\protect\verifie{Notebook1989_1}}\label{sec:quillen-1989-1}

\emph{Context.} On page~9 of \carnet{1989}{1} Quillen lifts idempotents modulo a nilpotent ideal with the normalised primitive of $(t - t^2)^n$; on page~72 he needs an odd sum of powers of a root of unity to be invertible.

\begin{encadre}
\begin{theoreme}[Quillen, notebook 1989-1, pp.~9, 72]\label{thm:quillen-1989-1}
\begin{enumerate}
\item Let $f(x) = \int_0^x(t - t^2)^n\,dt\big/\int_0^1(t - t^2)^n\,dt \in \mathbf R[x]$. Then $f(0) = 0$, $f(1) = 1$, $f(x) + f(1 - x) = 1$, and $(x^2 - x)^{n+1}$ divides $f^2 - f$.
\item In a commutative $\mathbf Q$-algebra, if $x^{2n-1} = 1$ with $n \ge 1$, then $\sum_{j<n}x^{2j+1}$ is a unit.
\end{enumerate}
\end{theoreme}
\end{encadre}

\begin{proof}
(1) The primitive $F$ of $(x - x^2)^n$ vanishing at $0$ is divisible by $x^{n+1}$; $F(x) + F(1 - x)$ has zero derivative, hence equals $F(1) = \int_0^1(t - t^2)^n\,dt > 0$. So $f = F/F(1)$ satisfies $f(1 - x) = 1 - f(x)$, and $f^2 - f = -f(x)f(1 - x)$ is divisible by $x^{n+1}(1 - x)^{n+1}$. (2) Put $m = 2n - 1$ and $N = \sum_{k<m}x^k$. Then $xN = N$, so $N^2 = mN$; $(1 + x)\sum_{j<n}x^{2j} = \sum_{k<2n}x^k = 1 + N$, and $(1 + N)(1 - N/(m + 1)) = 1$. Finally $x$ is a unit since $x^m = 1$.
\end{proof}

\emph{Verification in Lean.} The file \href{\leanblobmain/Notebook1989_1.lean}{\texttt{Notebook1989\_1.\allowbreak{}lean}} proves (1) as \texttt{liftPoly\_eval\_zero}, \texttt{liftPoly\_eval\_one}, \texttt{liftPoly\_symm} and \texttt{liftPoly\_sq\_sub\_dvd}, with an explicit primitive \texttt{F} and the fundamental theorem of calculus for $F(1) > 0$, and (2) as \texttt{isUnit\_sum\_odd\_pow}. The page's argument for (2) goes through the roots of unity; the one here is algebraic. The value $F(1) = (n!)^2/(2n + 1)!$ is \S\ref{sec:quillen-1988-9}, (1); the use of $f$ to lift idempotents is not formalised. The axioms used are the three usual axioms of mathlib.

\litterature{Classical: lifting idempotents modulo a nilpotent ideal by a polynomial is standard (see e.g.\ Lam, \emph{A first course in noncommutative rings}, \S21); (2) is elementary. The edition does not present them as new.}

\section[Notebook 1989-2: Whitehead matrices and triples]{Notebook 1989-2: Whitehead's identity and the category of triples\protect\verifie{Notebook1989_2Whitehead}}\label{sec:quillen-1989-2-whitehead}

\emph{Context.} On pages~28 and~30 of \carnet{1989}{2} Quillen reduces the action of inner automorphisms on cyclic homology to elementary $2 \times 2$ matrices; on page~84 he composes homotopy equivalence data and checks associativity.

\begin{encadre}
\begin{theoreme}[Quillen, notebook 1989-2, pp.~28, 30, 84]\label{thm:quillen-1989-2-whitehead}
Let $A$ be any ring, $z, a \in A$ and $u$ a unit. Then:
\begin{enumerate}
\item $u_z(a) = \begin{pmatrix}a & az\\ 0 & 0\end{pmatrix}$ is multiplicative and equals $\begin{pmatrix}1 & -z\\ 0 & 1\end{pmatrix}\begin{pmatrix}a & 0\\ 0 & 0\end{pmatrix}\begin{pmatrix}1 & z\\ 0 & 1\end{pmatrix}$;
\item $\operatorname{diag}(u, u^{-1})\operatorname{diag}(a, 0)\operatorname{diag}(u, u^{-1})^{-1} = \operatorname{diag}(uau^{-1}, 0)$ and $\begin{pmatrix}1 & u\\ 0 & 1\end{pmatrix}\begin{pmatrix}1 & 0\\ -u^{-1} & 1\end{pmatrix}\begin{pmatrix}1 & u\\ 0 & 1\end{pmatrix} = \begin{pmatrix}0 & u\\ -u^{-1} & 0\end{pmatrix}$, whose square is $-1$;
\item for $d \in A$, if $u, v$ and $u', v'$ commute with $d$ and $vu = 1 + dh + hd$, $v'u' = 1 + dh' + h'd$, then $(u'u, vv', h + vh'u)$ has the same property, and this composition is associative with unit $(1, 1, 0)$.
\end{enumerate}
\end{theoreme}
\end{encadre}

\begin{proof}
Matrix multiplication, and for (3) $vv'u'u = v(1 + dh' + h'd)u = vu + d(vh'u) + (vh'u)d$ since $u$, $v$ commute with $d$.
\end{proof}

\emph{Verification in Lean.} The file \href{\leanblobmain/Notebook1989_2Whitehead.lean}{\texttt{Notebook1989\_2Whitehead.\allowbreak{}lean}} proves \texttt{uz\_mul}, \texttt{uz\_conj}, \texttt{diag\_conj}, \texttt{whitehead}, \texttt{whitehead\_sq}, \texttt{comp\_spec}, \texttt{id\_spec} and \texttt{comp\_assoc}, in one ring (the page's triples live in endomorphism rings of complexes; the grading is not formalised). These are short identities, collected because the pages rely on them. The axioms used are among the three usual axioms of mathlib.

\litterature{Classical (Whitehead's lemma; Connes, 1985; Gugenheim, Lambe and Stasheff for the perturbation context); the edition does not present them as new.}

\section[Notebook 1990-5: full idempotents; the splitting lemma]{Notebook 1990-5: full idempotents; the splitting lemma; $e = (1 + \varepsilon)/2$\protect\verifie{Notebook1990_5}}\label{sec:quillen-1990-5}

\emph{Context.} On page~26 of \carnet{1990}{5} Quillen characterises the full idempotents, on pages~27--28 he forms the idempotent $(1 + \varepsilon)/2$ of a grading, and on page~67 he uses the splitting lemma for a short exact sequence.

\begin{encadre}
\begin{theoreme}[Quillen, notebook 1990-5, pp.~26--28, 67]\label{thm:quillen-1990-5}
\begin{enumerate}
\item Let $e$ be an idempotent of a ring $R$. Then $ReR = R$ (some $\sum_ia_ieb_i = 1$) iff some $(Re)^n$ maps onto $R$ as left $R$-modules; every $R$-linear $\varphi : Re \to R$ is $x \mapsto x\varphi(e)$.
\item If $i : E' \to E$ is injective, $p : E \to E''$ surjective, $pi = 0$, and $ik + hp = 1$, then $ki = 1$, $ph = 1$ and $kh = 0$.
\item In a $\mathbf Q$-algebra with $\varepsilon^2 = F^2 = 1$ and $\varepsilon F = -F\varepsilon$, $e = (1 + \varepsilon)/2$ is idempotent, $FeF = 1 - e$ and $eFe = 0$.
\end{enumerate}
\end{theoreme}
\end{encadre}

\begin{proof}
(1) If $\sum a_ieb_i = 1$, $(v_i) \mapsto \sum v_ib_i$ is onto, the image of $(ra_ie)_i$ being $r$. Conversely $\varphi(e) = \varphi(e\cdot e) = e\varphi(e)$, so $\varphi(re) = re\varphi(e)$; writing $1 = f(v)$ with $v_i = r_ie$ gives $1 = \sum_ir_ieb_i$. (2) Apply $ik + hp = 1$ to $i(x)$, to $h(y)$, and use $pi = 0$ with the injectivity of $i$ and surjectivity of $p$. (3) Expansion.
\end{proof}

\emph{Verification in Lean.} The file \href{\leanblobmain/Notebook1990_5.lean}{\texttt{Notebook1990\_5.\allowbreak{}lean}} proves \texttt{hom\_apply\_eq\_mul}, \texttt{full\_iff\_exists\_surjective}, \texttt{splitting} and \texttt{half\_one\_add}. The page works with right modules $eR$; the statement is the same in the opposite ring. The isomorphism $\mathrm{End}(eR) \cong eRe$ of the page is not formalised beyond the description of the maps $Re \to R$. Only $pi = 0$ is used of the exactness. The axioms used are the three usual axioms of mathlib.

\litterature{Classical (Morita theory; the splitting lemma); the edition does not present it as new.}

\section[Notebook 1995-2: when $za + 1 - a$ is a unit of $R[z, z^{-1}]$]{Notebook 1995-2: when is $za + (1 - a)$ a unit of $R[z, z^{-1}]$?\protect\verifie{Notebook1995_2}}\label{sec:quillen-1995-2}

\emph{Context.} In 1995 \carnet{1995}{2} works on the $K$-theory of nonunital rings: homotopy idempotents and their refinement to $A_\infty$-idempotents, perfect complexes that become contractible modulo an ideal, and a $K_0$ built from complexes of length one. In the entry of 1 July 1995, looking for a \q{good} $K_0$ of a nonunital ring inside a $K_1$ of Laurent polynomials, Quillen turns to the maps $K_0 R \to K_1(R[z, z^{-1}]) \to K_0 R$ of Bass's fundamental theorem and of the Atiyah--Bott proof of periodicity. The first map sends an idempotent $p$ to the unit $zp + 1 - p$.

\emph{The page.} Quillen studies linear families $az + b$ that are invertible over $R[z, z^{-1}]$, and splits $R$ by an idempotent given by a residue. On page~65 (his page~329), a change of coordinates on the projective line reduces such a family to the form $za + (1 - a)$, and he notes that this element is invertible over $R[z, z^{-1}]$ if and only if $a(1 - a)$ is nilpotent. His proof of the direct implication multiplies by $z^{-1}a + (1 - a)$; for the converse he passes to the subring of Laurent polynomials invariant under $z \mapsto z^{-1}$, which is a polynomial ring in $z - 2 + z^{-1}$. The page argues for a matrix $a$, that is, in a ring that need not be commutative.

\emph{Intuition.} When $a$ is idempotent, $za + 1 - a$ acts as $z$ on the image of $a$ and as $1$ on its kernel, and its inverse is $z^{-1}a + 1 - a$. The criterion says that it is enough for $a$ to be idempotent modulo nilpotents. The key identity is
\[
  \bigl(za + (1 - a)\bigr)\bigl(z^{-1}a + (1 - a)\bigr) = 1 + (z - 2 + z^{-1})\,a(1 - a),
\]
since $a^2 + (1 - a)^2 = 1 - 2a(1 - a)$. When $a(1 - a)$ is nilpotent the right-hand side is $1$ plus a nilpotent, hence a unit. Conversely, over a field the units of $k[z, z^{-1}]$ are the monomials $cz^n$, so $za + (1 - a)$ can only be a unit if $a = 0$ or $a = 1$, that is $a(1 - a) = 0$.

\begin{encadre}
\begin{theoreme}[Quillen, notebook 1995-2, p.~65]\label{thm:quillen-1995-2}
Let $R$ be a commutative ring and $a \in R$. The Laurent polynomial $za + (1 - a)$ is a unit of $R[z, z^{-1}]$ if and only if $a(1 - a)$ is nilpotent.
\end{theoreme}
\end{encadre}

\begin{proof}
If $a(1 - a)$ is nilpotent, the identity above shows that $(za + 1 - a)(z^{-1}a + 1 - a)$ is a unit, and so is each factor. Conversely, let $g$ be an inverse of $za + (1 - a)$, and choose $N \ge 0$ such that $z^N g$ is a polynomial $G$. Then $(aX + 1 - a)\,G = X^N$ in $R[X]$. A nilpotent element lies in every prime ideal, and conversely, so let $\mathfrak p$ be a prime ideal and reduce modulo $\mathfrak p$. In the domain $D = R/\mathfrak p$, write $\bar a X + \bar b$ for the reduction of $aX + (1 - a)$, and show by induction on $N$ that $(\bar aX + \bar b)\bar G = X^N$ forces $\bar a = 0$ or $\bar b = 0$. For $N = 0$, $\bar aX + \bar b$ is a unit of $D[X]$, hence of degree $0$, so $\bar a = 0$. For $N \ge 1$ and $\bar b \neq 0$, the constant term gives $\bar b\,\bar G(0) = 0$, so $\bar G = X\bar G_1$, and cancelling $X$ gives $(\bar aX + \bar b)\bar G_1 = X^{N-1}$. Hence $\bar a(1 - \bar a) = 0$, that is $a(1 - a) \in \mathfrak p$.
\end{proof}

\noindent The first implication is the one on the page. For the converse we do not use the invariant subring: the proof reduces modulo the prime ideals, where a unit $\bar a z + \bar b$ must be a monomial.

\emph{Verification in Lean.} The file \href{\leanblobmain/Notebook1995_2.lean}{\texttt{Notebook1995\_2.\allowbreak{}lean}} defines \texttt{bott a} as \texttt{C a * T 1 + C (1 - a)} in \texttt{R[T;T⁻¹]}. The identity above is \texttt{bott\_mul\_invert}, which uses the involution \texttt{LaurentPolynomial.invert}. The two implications are \texttt{isUnit\_bott\_of\_isNilpotent} and \texttt{isNilpotent\_of\_isUnit\_bott}. The second goes to polynomials through \texttt{exists\_T\_pow} and \texttt{Polynomial.toLaurent\_injective}, and the lemma on domains is \texttt{eq\_zero\_or\_eq\_zero\_of\_mul\_eq\_X\_pow}. Theorem~\ref{thm:quillen-1995-2} is \texttt{isUnit\_bott\_iff}. Only the commutative case is formalised; the matrix case of the page is not. The axioms used are the three usual axioms of mathlib.

\litterature{Bass's map $p \mapsto zp + 1 - p$ comes from Atiyah--Bott, \q{On the periodicity theorem for complex vector bundles}, \emph{Acta Math.}~112 (1964), and from Bass, \emph{Algebraic $K$-theory} (1968). Over a commutative ring the criterion follows from the known description of the units of $R[z, z^{-1}]$, as recalled: they are monomials $uz^n$ plus nilpotents, locally for a decomposition of $1$ into orthogonal idempotents. The search (zbMATH titles on units of Laurent polynomial rings) did not locate the criterion in this form, and mathlib has no description of the units of Laurent polynomials. The statement is elementary, and the edition does not present it as new.}


\part{Groups: commutators, perfect subgroups and permutations}
\section[Notebooks 1971-2 and 1975-10: $\Sigma_{n-1}$ is self-normalizing]{Notebooks 1971-2 and 1975-10: $\Sigma_{n-1}$ is its own normalizer in $\Sigma_n$\protect\verifie{Notebook1971_2Normalizer}}\label{sec:quillen-1971-2-normalizer}

\emph{Context.} On page~14 of \carnet{1971}{2} Quillen proves that the transfer $H_*(\Sigma_{m\ell+1}, \mathrm{sgn}/\ell) \to H_*(\Sigma_{m\ell}, \mathrm{sgn}/\ell)$ is an isomorphism, by a double-coset argument in which every intersection $H \cap xHx^{-1}$, $H = \Sigma_{n-1}$, is a $\Sigma_{n-2}$ unless $x$ normalizes $H$, which, he notes, does not happen outside $H$, since $\Sigma_{n-1}$ is its own normalizer in $\Sigma_n$. Page~7 of \carnet{1975}{10} opens with the same remark, for the universal $\Sigma_n$-space with orbit types $\Sigma_n$ and $\Sigma_n/\Sigma_{n-1}$.

\begin{encadre}
\begin{theoreme}[Quillen, notebook 1971-2, p.~14; notebook 1975-10, p.~7]\label{thm:quillen-1971-2-normalizer}
Let $\alpha$ be a finite set with at least three elements and $a \in \alpha$. The normalizer in the symmetric group of the stabilizer of $a$ is the stabilizer of $a$. For two elements the statement fails.
\end{theoreme}
\end{encadre}

\begin{proof}
Let $g$ normalize $H = \mathrm{Stab}(a)$, and suppose $a' = g^{-1}a \ne a$. Choose $c \notin \{a, a'\}$. The transposition $(a'\,c)$ fixes $a$, so $g(a'\,c)g^{-1}$ fixes $a$; but it sends $a$ to $g(c) \ne a$. For two elements the stabilizer is trivial and its normalizer is the whole group.
\end{proof}

\emph{Verification in Lean.} The file \href{\leanblobmain/Notebook1971_2Normalizer.lean}{\texttt{Notebook1971\_2Normalizer.\allowbreak{}lean}} proves the theorem as \texttt{normalizer\_stabilizer}, for \texttt{Equiv.Perm α} with $|\alpha| \ge 3$, and the failure for two elements as \texttt{normalizer\_stabilizer\_two}. The homology and the double cosets are not formalised. The axioms used are the three usual axioms of mathlib.

\litterature{Classical: a point stabilizer of a $2$-transitive group is self-normalizing. The edition does not present it as new.}

\section[Notebook 1971-3: the swindle that makes $g$ a commutator]{Notebook 1971-3: the swindle that makes a compactly supported homeomorphism a commutator\protect\verifie{Notebook1971_3}}\label{sec:quillen-1971-3}

\emph{Context.} On page~32 of \carnet{1971}{3} Quillen starts to read Mather's theorem that the group $G$ of compactly supported homeomorphisms of $\mathbf R^n$ has the homology of a point. He first shows that $G$ is perfect.

\emph{The page.} Take $g$ supported in a ball $B_0$, balls $B_1, B_2, \ldots$ disjoint and shrinking to a point, $z \in G$ moving each $B_i$ onto $B_{i+1}$, $\sigma(x) = zxz^{-1}$, and the infinite product $T = \sigma(g)\sigma^2(g)\cdots$, which converges in $G$. Then $\sigma^{-1}(T) = gT$, and the page concludes that $g$ is a commutator, with the formula $g = T^{-1}z^{-1}Tz$. From $z^{-1}Tz = gT$ one gets $g = z^{-1}TzT^{-1} = (z^{-1}, T)$ instead; the page's expression is $T^{-1}gT$, the conjugate of $g$ by $T$. It is also a commutator, and it equals $g$ only when $g$ and $T$ commute. The conclusion of the page stands.

\begin{encadre}
\begin{theoreme}[After Quillen, notebook 1971-3, pp.~32--33]\label{thm:quillen-1971-3}
Let $g, z, T$ be elements of a group with $T = z(gT)z^{-1}$. Then $g = (z^{-1}, T)$, so $g$ is a commutator, and $T^{-1}z^{-1}Tz = T^{-1}gT$, which equals $g$ if and only if $g$ commutes with $T$. In the symmetric group on three letters there are $g, z, T$ with $T = z(gT)z^{-1}$ and $T^{-1}z^{-1}Tz \ne g$.
\end{theoreme}
\end{encadre}

\begin{proof}
The hypothesis is $z^{-1}Tz = gT$, whence $g = z^{-1}TzT^{-1}$ and $T^{-1}z^{-1}Tz = T^{-1}(gT)$. For the example take $z = (1\,2)$, $T = (0\,1)$ and $g = z^{-1}TzT^{-1}$, a $3$-cycle, which does not commute with the transposition $T$.
\end{proof}

\emph{Verification in Lean.} The file \href{\leanblobmain/Notebook1971_3.lean}{\texttt{Notebook1971\_3.\allowbreak{}lean}} proves Theorem~\ref{thm:quillen-1971-3} as \texttt{eq\_commutator}, \texttt{mem\_commutator}, \texttt{page\_formula\_eq}, \texttt{page\_formula\_eq\_iff} and \texttt{page\_formula\_ne} (the example is checked by \texttt{decide}). It also proves the abelian form of the question of page~34, \texttt{hom\_eq\_one}: if $x \mapsto x\sigma(x)^{-1}$ is onto for an endomorphism $\sigma$, a homomorphism to an abelian group invariant under $\sigma$ is trivial. The construction of $T$ (convergence of the infinite product) is not formalised. The axioms used are the three usual axioms of mathlib.

\litterature{The swindle is classical (Mather, Topology 10, 1971, for the acyclicity of the group of compactly supported homeomorphisms of $\mathbf R^n$). The section records only that the formula of the page is a conjugate of the commutator that the hypothesis gives. The edition does not present anything here as new.}

\section[Notebook 1971-8: $G \times G$ modulo the diagonal]{Notebook 1971-8: $G \times G$ modulo the normal closure of the diagonal\protect\verifie{Notebook1971_8}}\label{sec:quillen-1971-8}

\emph{Context.} \carnet{1971}{8} opens, on 2 June 1971, with problems left by Quillen's work on the group completion of a topological monoid $M$. He studies a space built from $M \times M$ with the diagonal action of $M$, which should be the group completion, and local coefficient systems on it.

\emph{The page.} On pages~7 and~8 he shows that for a local system compatible with the action of $M$, a loop of $M$ embedded diagonally acts trivially. The action of $\pi_1(M) \times \pi_1(M)$ on a fibre therefore kills the diagonal subgroup, and he concludes that the action is abelian, since the quotient of $G \times G$ by the normal subgroup generated by the diagonal $\Delta G$ is the abelianization $G_{\mathrm{ab}}$. He deduces that the identity component of the space has an abelian fundamental group.

\emph{Intuition.} Modulo the diagonal, $(x, y)$ is equivalent to $(xy^{-1}, 1)$, so the quotient is a quotient of $G$. The diagonal is not normal in $G \times G$: conjugating $(b, b)$ by $(a, 1)$ gives $(aba^{-1}, b)$, and dividing by $(b, b)$ gives $(aba^{-1}b^{-1}, 1)$. So the normal subgroup generated by the diagonal contains the commutators in the first factor, and the quotient of $G$ obtained is exactly $G_{\mathrm{ab}}$.

\begin{encadre}
\begin{theoreme}[Quillen, notebook 1971-8, p.~8]\label{thm:quillen-1971-8}
Let $G$ be a group and $N$ the normal subgroup of $G \times G$ generated by the diagonal $\{(g, g)\}$. The map $(x, y) \mapsto [x]\,[y]^{-1}$ induces an isomorphism from $(G \times G)/N$ onto the abelianization $G_{\mathrm{ab}} = G/[G, G]$.
\end{theoreme}
\end{encadre}

\begin{proof}
The map $\varphi : (x, y) \mapsto [x][y]^{-1}$ is a homomorphism because $G_{\mathrm{ab}}$ is abelian, it is onto since $\varphi(x, 1) = [x]$, and its kernel is a normal subgroup containing the diagonal, hence containing $N$. Conversely let $(x, y)$ be in the kernel, so that $c = xy^{-1}$ lies in $[G, G]$. The set of $c \in G$ with $(c, 1) \in N$ is a subgroup of $G$, and it contains every commutator, by the computation above: $(a, 1)(b, b)(a, 1)^{-1}(b, b)^{-1} = (aba^{-1}b^{-1}, 1)$. Hence $(c, 1) \in N$, and $(x, y) = (c, 1)(y, y) \in N$.
\end{proof}

\emph{Verification in Lean.} The file \href{\leanblobmain/Notebook1971_8.lean}{\texttt{Notebook1971\_8.\allowbreak{}lean}} defines \texttt{diagClosure G} with \texttt{Subgroup.normalClosure} and the homomorphism \texttt{toAb G} to \texttt{Abelianization G}. The computation with commutators is \texttt{commutator\_le\_comap}; the kernel and the surjectivity are \texttt{ker\_toAb} and \texttt{toAb\_surjective}. Theorem~\ref{thm:quillen-1971-8} is the isomorphism \texttt{quotientEquiv}, with its formula \texttt{quotientEquiv\_mk}. The topological statements of the page, on local systems and on the fundamental group of the space, are not formalised. The axioms used are the three usual axioms of mathlib.

\litterature{The group-theoretic statement is elementary and the page states it as known. Its consequence on the page is, in published form, part of the group completion theorem: the fundamental group of a group-like $H$-space is abelian (McDuff--Segal, \emph{Invent.\ Math.}~31 (1976), as recalled). We did not find the statement in mathlib. The edition does not present it as new.}

\section[Notebook 1971-10: commutators of displaced elements; Artin--Schreier]{Notebook 1971-10: commutators of displaced elements, and the Artin--Schreier equation\protect\verifie{Notebook1971_10}}\label{sec:quillen-1971-10}

\emph{Context.} \carnet{1971}{10} opens, on 6~June 1971, with a proof that the commutator subgroup of the group $G$ of compactly supported diffeomorphisms of $\mathbf R$ is perfect, and continues with Kummer theory for $BGL(k)^+$ over an algebraically closed field $k$ of characteristic $p$.

\emph{The page, p.~1.} Quillen calls a diffeomorphism elementary when it is supported on two disjoint intervals and inverse to itself under a transport from one to the other; the elementary elements generate a normal subgroup $E \subseteq (G, G)$. Given $g_1, g_2$ supported in an interval $K$, he takes $u$, $v$ moving $K$ off itself in opposite directions; then $g_1 \cdot u g_1^{-1}u^{-1}$ and $g_2 \cdot v g_2^{-1} v^{-1}$ are elementary and their commutator is still $(g_1, g_2)$. Hence $(G,G) \subseteq (E,E)$, so $E = (G,G)$ and $E$ is perfect. The geometry only provides elements that commute: $a' = ug_1^{-1}u^{-1}$ is supported in $uK$, $b' = vg_2^{-1}v^{-1}$ in $vK$, and $K$, $uK$, $vK$ are disjoint.

\emph{The page, p.~8.} In a direct proof of Lang's theorem for $GL_n$, the page solves $x^q - x + a = 0$ in $k$, $q$ a power of $p$.

\begin{encadre}
\begin{theoreme}[After Quillen, notebook 1971-10, p.~1]\label{thm:quillen-1971-10}
Let $G$ be a group. \emph{(i)} If $a'$ commutes with $b$ and $b'$, and $b'$ commutes with $a$, then $(aa', bb') = (a, b)$. \emph{(ii)} Let $E$ be a subgroup such that every pair $g_1, g_2$ admits $a'$, $b'$ as in (i), for $a = g_1$, $b = g_2$, with $g_1a'$ and $g_2b'$ in $E$. Then $(G, G) \subseteq (E, E)$; if moreover $E \subseteq (G, G)$, then $E = (G, G)$ and $(E, E) = E$.
\end{theoreme}
\end{encadre}

\begin{proof}
(i) Since $a'$ commutes with $bb'$, $aa'bb'a'^{-1}a^{-1}b'^{-1}b^{-1} = abb'a^{-1}b'^{-1}b^{-1}$, and since $b'$ commutes with $a$ this is $aba^{-1}b^{-1}$. (ii) By (i) every commutator $(g_1, g_2)$ is the commutator $(g_1a', g_2b')$ of two elements of $E$. If $E \subseteq (G,G)$, then $(G,G) \subseteq (E,E) \subseteq E \subseteq (G,G)$.
\end{proof}

\begin{encadre}
\begin{proposition}[Notebook 1971-10, p.~8]\label{prop:quillen-1971-10}
Over an algebraically closed field $k$, every $y$ is of the form $x^q - x$ for $q \ge 2$. If $k$ has characteristic $p$ and $p \mid q$, the equation $x^q - x = y$ has exactly $q$ solutions.
\end{proposition}
\end{encadre}

\begin{proof}
$X^q - X - y$ has degree $q \ge 1$, hence a root. In characteristic $p$ with $p \mid q$ its derivative is $-1$, so it is separable and has $q$ distinct roots.
\end{proof}

\emph{Verification in Lean.} The file \href{\leanblobmain/Notebook1971_10.lean}{\texttt{Notebook1971\_10.\allowbreak{}lean}} proves Theorem~\ref{thm:quillen-1971-10} as \texttt{commutator\_mul\_mul}, \texttt{commutator\_le\_commutator} and \texttt{eq\_commutator\_of\_le}, for any group and any subgroup $E$ (normality is not needed), and Proposition~\ref{prop:quillen-1971-10} as \texttt{exists\_pow\_sub\_eq} and \texttt{card\_rootSet\_pow\_sub}. The diffeomorphisms, their supports and Lang's theorem are not formalised. The axioms used are the three usual axioms of mathlib.

\litterature{The displacement argument is the standard one for groups of compactly supported homeomorphisms and diffeomorphisms (Epstein, Compositio Math.\ 22, 1970; Mather, Comment.\ Math.\ Helv.\ 49, 1974, who proves the stronger statement that the group is perfect). The Artin--Schreier statement is classical. The edition does not present either as new.}

\section[Notebook 1972-1: a transposition acts on $A_n$ by an outer automorphism]{Notebook 1972-1: a transposition acts on $A_n$ by an outer automorphism\protect\verifie{Notebook1972_1}}\label{sec:quillen-1972-1}

\emph{Context.} On page~25 of \carnet{1972}{1} Quillen asks whether a transposition $\sigma$ acts on $A_n \subset \Sigma_n$ by an outer automorphism, in view of the extension $A_n \to \Sigma_n \to \mathbf Z/2$ and its cohomology.

\emph{The page.} He answers yes: if $\sigma x\sigma^{-1} = yxy^{-1}$ for all $x \in A_n$, then $y^{-1}\sigma$ lies in the \q{center of $A_n$}, hence is the identity. Since $y \in A_n$ and $\sigma$ is odd, $y^{-1}\sigma$ is not in $A_n$; what the argument uses is that the centralizer of $A_n$ in $\Sigma_n$ is trivial, which holds for $n \ge 4$. For $n = 3$ the answer is still yes, because $A_3$ is abelian and $\sigma$ inverts it.

\begin{encadre}
\begin{theoreme}[After Quillen, notebook 1972-1, p.~25]\label{thm:quillen-1972-1}
For $n \ge 4$, a permutation of $n$ letters that commutes with every even permutation is the identity. For $n \ge 3$ and a transposition $\sigma$, there is no $y \in A_n$ with $\sigma x \sigma^{-1} = yxy^{-1}$ for every $x \in A_n$.
\end{theoreme}
\end{encadre}

\begin{proof}
Let $g$ commute with $A_n$ and suppose $g(a) \ne a$. Choose $c, d$ distinct and outside $\{a, g(a)\}$, and let $k = (g(a)\,d)(g(a)\,c)$, an even permutation fixing $a$ and sending $g(a)$ to $c$. Then $kg(a) = c$ while $gk(a) = g(a)$, a contradiction. If $\sigma$ acted as $y$, then $y^{-1}\sigma$ would commute with $A_n$, so $\sigma = y$ would be even. For $n = 3$, every $y \in A_3$ commutes with the $3$-cycle $c$, while $\sigma c\sigma^{-1} = c^{-1} \ne c$.
\end{proof}

\emph{Verification in Lean.} The file \href{\leanblobmain/Notebook1972_1.lean}{\texttt{Notebook1972\_1.\allowbreak{}lean}} proves the theorem as \texttt{eq\_one\_of\_commute\_alternating} (the argument of mathlib's \texttt{alternatingGroup.center\_eq\_bot}, for a permutation that need not be even), \texttt{not\_inner\_swap} for $|\alpha| \ge 4$, and \texttt{not\_inner\_swap\_three} for $n = 3$ by computation. The cohomological question of the page is not formalised. The axioms used are the three usual axioms of mathlib.

\litterature{Classical: for $n \ne 6$, $\mathrm{Aut}(A_n) = \Sigma_n$ for $n \ge 4$, and conjugation by an odd permutation is outer. Mathlib has the triviality of the center of $A_n$ for $n \ge 4$, not of its centralizer in $\Sigma_n$. The edition does not present it as new.}

\section[Notebook 1974-16: perfect subgroups]{Notebook 1974-16: the largest perfect subgroup, and perfect subgroups of a central extension\protect\verifie{Notebook1974_16}}\label{sec:quillen-1974-16}

\emph{Context.} \carnet{1974}{16} reviews the plus construction. A subgroup $H$ is perfect when $(H, H) = H$.

\emph{The page, p.~5.} Quillen takes the largest perfect subgroup $N$ of $\pi_1X$, with the remark that a group generated by perfect subgroups is perfect, and builds $X \to X^+$ with it.

\emph{The page, p.~6.} For a fibration $F \to E \to B$ with an acyclic map of fibres, $\pi_1F$ is a central extension of its image in $\pi_1E$, and he uses the known fact that a perfect subgroup $N$ of the quotient is the image of a unique perfect subgroup, the commutator subgroup of the preimage of $N$.

\begin{encadre}
\begin{theoreme}[Quillen, notebook 1974-16, pp.~5--6]\label{thm:quillen-1974-16}
\emph{(i)} In a group, the subgroup generated by a family of perfect subgroups is perfect; so the subgroup generated by all perfect subgroups is the largest perfect subgroup, and it is normal. \emph{(ii)} Let $f \colon E \to Q$ be a surjective homomorphism whose kernel is central, and $N$ a perfect subgroup of $Q$. Then $M = (f^{-1}N, f^{-1}N)$ is perfect, $f(M) = N$, and $M$ is the only perfect subgroup of $E$ with image $N$.
\end{theoreme}
\end{encadre}

\begin{proof}
(i) If $K$ is generated by perfect $H_i$, then $H_i = (H_i, H_i) \subseteq (K, K) \subseteq K$, so $K \subseteq (K, K)$. A conjugate of a perfect subgroup is perfect, so the largest one is stable under conjugation. (ii) $f(M) = (N, N) = N$ since $f$ is onto. Let $M'$ be any subgroup with $f(M') = N$. Every $x \in f^{-1}N$ is $mc$ with $m \in M'$ and $c$ in the kernel, hence central, and $(m_1c_1, m_2c_2) = (m_1, m_2)$; so $M \subseteq (M', M')$. With $M' = M$ this shows that $M$ is perfect. If $M'$ is perfect, then also $M' = (M', M') \subseteq (f^{-1}N, f^{-1}N) = M$, so $M' = M$.
\end{proof}

\emph{Verification in Lean.} The file \href{\leanblobmain/Notebook1974_16.lean}{\texttt{Notebook1974\_16.\allowbreak{}lean}} proves (i) as \texttt{commutator\_iSup\_eq}, \texttt{commutator\_perfectRadical}, \texttt{le\_perfectRadical} and \texttt{perfectRadical\_normal}, and (ii) as \texttt{map\_commutator\_comap}, \texttt{commutator\_comap\_perfect} and \texttt{eq\_commutator\_comap}, with the key step \texttt{commutator\_comap\_le}; a subgroup is perfect when its commutator subgroup with itself is itself. The uniqueness does not use the surjectivity of $f$. The plus construction is not formalised. The axioms used are the three usual axioms of mathlib.

\litterature{Both statements are classical (the perfect radical is in the textbooks on the plus construction, for instance Berrick, \emph{An Approach to Algebraic $K$-Theory}, 1982; the statement on central extensions belongs to the theory of universal central extensions). Mathlib has perfect groups (\texttt{Group.IsPerfect}) but not these two statements. The edition does not present them as new.}

\section[Notebooks 1990-2, 1990-3: rigid triple; Green operator]{Notebooks 1990-2 and 1990-3: the rigid triple in $S_n$; the Green operator of $\mathbf Z/n$\protect\verifie{Notebook1990_3}}\label{sec:quillen-1990-3}

\emph{Context.} Page~2 of \carnet{1990}{3} opens with a rigid triple in the symmetric group: an $n$-cycle, a transposition and their product, an $(n - 1)$-cycle. Page~18, and page~11 of \carnet{1990}{2}, invert $1 - \kappa$ for $\kappa^n = 1$ on the complement of the invariants.

\begin{encadre}
\begin{theoreme}[Quillen, notebooks 1990-2, p.~11, and 1990-3, pp.~2, 18]\label{thm:quillen-1990-3}
\begin{enumerate}
\item Let $\gamma$ be a cycle moving all $n \ge 3$ points of a finite set, and $x \ne y$. Then $\gamma\circ(x\ y)$ is a cycle moving $n - 1$ points if and only if $\gamma x = y$ or $\gamma y = x$.
\item In a $\mathbf Q$-algebra, let $\kappa^n = 1$, $n \ge 1$, $P = \frac1n\sum_{j<n}\kappa^j$ and $G = \frac1n\sum_{j<n}(\frac{n-1}2 - j)\kappa^j$. Then $G(1 - \kappa) = (1 - \kappa)G = 1 - P$, $P(1 - \kappa) = 0$, $P^2 = P$ and $PG = GP = 0$.
\end{enumerate}
\end{theoreme}
\end{encadre}

\begin{proof}
(1) If $\gamma x = y$, then $\gamma(x\ y) = (y\ \gamma y)\gamma$, and swapping $y$ with its image splits it off the cycle, which remains a cycle on the other points as $\gamma^2y \ne y$ (here $n \ge 3$). Conversely, a cycle moving $n - 1$ points has a fixed point $z$; since $\gamma$ has none, $z \in \{x, y\}$, and $\gamma y = x$ or $\gamma x = y$. (2) The sum $\sum_{j<m}(c - j)\kappa^j(1 - \kappa)$ telescopes to $c - (c - m)\kappa^m - \sum_{j<m}\kappa^{j+1}$; take $m = n$, $c = \frac{n-1}2$. $P\kappa = P$ gives the rest, with $\sum_{j<n}(\frac{n-1}2 - j) = 0$.
\end{proof}

\emph{Verification in Lean.} The file \href{\leanblobmain/Notebook1990_3.lean}{\texttt{Notebook1990\_3.\allowbreak{}lean}} proves (1) as \texttt{isCycle\_mul\_swap\_iff}, the direct implication through mathlib's \texttt{Equiv.Perm.IsCycle.swap\_mul} and \texttt{support\_swap\_mul\_eq}, and (2) as \texttt{green\_mul\_one\_sub}, \texttt{green\_commute}, \texttt{proj\_mul\_one\_sub}, \texttt{proj\_idem}, \texttt{proj\_mul\_green} and \texttt{green\_mul\_proj}, in any $\mathbf Q$-algebra. Rigidity of the triple, and the application of $G$ to cyclic homology, are not formalised. The axioms used are the three usual axioms of mathlib.

\litterature{Classical: (1) is the cut-and-join rule for a transposition times a cycle; (2) is the standard Green operator of a finite cyclic group. The edition does not present them as new.}


\part{$p$-groups, nilpotent groups and their automorphisms}
\section[Notebook 1971-2: unique $p$-divisibility from the Adams operation]{Notebook 1971-2: unique $p$-divisibility from an automorphism killed by $(\psi - p)\cdots(\psi - p^n)$\protect\verifie{Notebook1971_2}}\label{sec:quillen-1971-2}

\emph{Context.} On page~12 of \carnet{1971}{2}, on 3~February 1971, Quillen proves that the higher $K$-groups of a perfect ring of characteristic $p$ are uniquely $p$-divisible.

\emph{The page.} The Frobenius of the ring induces an automorphism of $K_i$, which is the Adams operation $\Psi^p$, and every $x \in K_i$, $i > 0$, is killed by $(\Psi^p - p)\cdots(\Psi^p - p^n)$ for some $n$ (taken as known on the page). With $F_n$ the set of elements killed by this product, $F_{n-1} \subseteq F_n$ and $\Psi^p = p^n$ on $F_n/F_{n-1}$; since $\Psi^p$ is an automorphism, each quotient is uniquely $p$-divisible, hence so is the union.

\begin{encadre}
\begin{theoreme}[Quillen, notebook 1971-2, p.~12]\label{thm:quillen-1971-2}
Let $M$ be an abelian group, $\psi$ an automorphism of $M$ and $p$ a natural number. If every element of $M$ is killed by $(\psi - p)(\psi - p^2)\cdots(\psi - p^n)$ for some $n$, then $M$ is uniquely $p$-divisible: $x \mapsto px$ is bijective.
\end{theoreme}
\end{encadre}

\begin{proof}
Let $P_n = (\psi - p)\cdots(\psi - p^n)$ and $F_n = \ker P_n$. Since all these operators commute with $\psi$ and $\psi^{-1}$, $F_n$ is stable under $\psi$ and $\psi^{-1}$, $F_{n-1} \subseteq F_n$, and $(\psi - p^n)F_n \subseteq F_{n-1}$. We show by induction on $n$ that multiplication by $p$ is bijective on $F_n$. If $x \in F_n$ and $px = 0$, then $y = (\psi - p^n)x \in F_{n-1}$ and $py = 0$, so $y = 0$ by induction; then $\psi x = p^nx = 0$ and $x = 0$. If $x \in F_n$, write $(\psi - p^n)x = p^n w$ with $w \in F_{n-1}$, by induction; then $\psi x = p^n(x + w)$ and $x = p\,\psi^{-1}(p^{n-1}(x + w))$, with $\psi^{-1}(p^{n-1}(x+w)) \in F_n$.
\end{proof}

\emph{Verification in Lean.} The file \href{\leanblobmain/Notebook1971_2.lean}{\texttt{Notebook1971\_2.\allowbreak{}lean}} proves Theorem~\ref{thm:quillen-1971-2} as \texttt{bijective\_nsmul}, with $\psi$ and its inverse given as elements of the ring \texttt{Module.End ℤ M}, and the induction step as \texttt{bijOn\_ker}. The proof works inside $F_n$ instead of on the quotients $F_n/F_{n-1}$. The $K$-theoretic input (the Frobenius is $\Psi^p$, and the local vanishing of the product) is not formalised. The axioms used are the three usual axioms of mathlib.

\litterature{The theorem on $K$-groups of perfect rings is folklore in the later literature, and its algebraic step is elementary; we did not find this step stated as a lemma. The edition does not present it as new.}

\section[Notebook 1971-11: automorphisms of order $p$ trivial on $\Omega_1$]{Notebook 1971-11: automorphisms of order $p$ of an abelian $p$-group that are trivial on the elements of order $p$\protect\piste\protect\verifie{Notebook1971_11}}\label{sec:quillen-1971-11}

\emph{Context.} In 1971, Quillen studies the cohomology of finite groups and the role of their elementary abelian subgroups. \carnet{1971}{11} contains, from page~68 on, a series dated 24 August 1971, titled \q{$p$ odd}, annotated in the margin as a revision of Thompson's theorems. One reads there a lemma, then two propositions on the maximal normal elementary abelian subgroups of a $p$-group, which use it (pp.~69--71).

\emph{Notation.} For an abelian group $B$, written additively, and a prime $p$, we write $\Omega_1(B)$ for the subgroup of elements $b$ such that $pb = 0$. An \emph{abelian $p$-group} is an abelian group in which every element is killed by a power of $p$.

\emph{The page.} On page~68, Quillen states the following lemma: if $\theta$ is an automorphism of an abelian $p$-group $B$, with $\theta^p = 1$, which is the identity on $\Omega_1(B)$, then $\theta - 1$ sends $B$ into $\Omega_1(B)$. His proof, on page~69, passes to the dual group $\widehat B$, writes $\theta = 1 + N$ with $N$ having image in $p\widehat B$, and compares the terms of the binomial expansion of $(1 + N)^p = 1$; a first attempt, by a $p$-adic filtration, is crossed out. The hypothesis $p \ge 3$ serves to make the last term $N^p$ small enough. Page~69 then applies the lemma to an element of order $p$ that normalizes an abelian subgroup.

\emph{Intuition.} Put $\delta = \theta - 1$. The hypothesis says that $\delta$ kills $\Omega_1(B)$, and the conclusion that $p\delta = 0$. Since $\delta$ commutes with multiplication by $p$, it lowers the order: if $p^{k+1} b = 0$, then $p^k \delta(b) = \delta(p^k b) = 0$, since $p^k b \in \Omega_1(B)$. The powers of $\delta$ thus get smaller and smaller, and in the expansion of $(1 + \delta)^p = 1$ the term $p\delta$ can be compensated only by vanishing terms.

\begin{encadre}
\begin{theoreme}[Quillen, notebook 1971-11, p.~68]\label{thm:quillen-1971-11}
Let $p \ge 3$ be a prime, $B$ an abelian group in which every element is killed by a power of $p$, and $\theta$ an endomorphism of $B$ such that $\theta^p = 1$ and $\theta(b) = b$ for all $b \in \Omega_1(B)$. Then $p(\theta(b) - b) = 0$ for all $b \in B$, that is, $(\theta - 1)B \subset \Omega_1(B)$.
\end{theoreme}
\end{encadre}

\noindent The statement is that of the page, in a slightly more general form: $B$ is not assumed finite, and $\theta$ is an endomorphism, automatically bijective since $\theta^p = 1$.

\begin{proof}
Let $\delta = \theta - 1$. We show, by induction on $k$, that $p\delta(c) = 0$ for every $c$ such that $p^k c = 0$. For $k = 0$, $c = 0$. Suppose the result true for $k$, and let $c$ be such that $p^{k+1} c = 0$. By the lowering of the order, $p^k \delta(c) = 0$; the induction hypothesis applied to $\delta(c)$ gives $p\delta^2(c) = 0$. Then $\delta^2(c) \in \Omega_1(B)$, so $\delta^3(c) = 0$, then $\delta^m(c) = 0$ for all $m \ge 3$, and $p\delta^m(c) = 0$ for all $m \ge 2$.

Since $\theta = 1 + \delta$ and $1$ commutes with $\delta$, the binomial formula gives
\[
c = \theta^p(c) = \sum_{m = 0}^{p} \binom{p}{m}\, \delta^m(c) = c + p\,\delta(c) + \sum_{m = 2}^{p - 1} \binom{p}{m}\, \delta^m(c) + \delta^p(c).
\]
For $2 \le m \le p - 1$, the coefficient $\binom{p}{m}$ is divisible by $p$, and $p\delta^m(c) = 0$: these terms vanish. The last one, $\delta^p(c)$, vanishes because $p \ge 3$. What remains is $p\delta(c) = 0$.
\end{proof}

\noindent The proof is more direct than that of the page: it does not pass to the dual and proceeds by induction on the order of the elements. The hypothesis $p \ge 3$ serves at the same place: for $p = 2$, the last term is $\delta^2(c)$, which the induction cannot cancel.

\begin{contrexemple}[the edition's statement]\label{ce:quillen-1971-11}
For $p = 2$, the lemma is false. On $B = \mathbf Z/8$, let $\theta = -1$. Then $\theta^2 = 1$, and $\theta$ fixes $\Omega_1(B) = \{0, 4\}$, since $-4 = 4$. But $2(\theta(1) - 1) = -4 \ne 0$.
\end{contrexemple}

\noindent On $\mathbf Z/4$, on the other hand, $\theta = -1$ satisfies the conclusion ($2 \cdot (-2) = 0$): the first cyclic counterexample requires order~$8$. An exhaustive search, over the groups $\mathbf Z/9$, $\mathbf Z/27$, $\mathbf Z/9 \times \mathbf Z/3$, $\mathbf Z/9 \times \mathbf Z/9$, $\mathbf Z/27 \times \mathbf Z/3$ for $p = 3$ and $\mathbf Z/25$ for $p = 5$, found no automorphism contradicting the lemma; for $p = 2$, it finds some on $\mathbf Z/8$ and $\mathbf Z/8 \times \mathbf Z/2$, and none on $\mathbf Z/4$, $\mathbf Z/4 \times \mathbf Z/2$, $\mathbf Z/4 \times \mathbf Z/4$. The title \q{$p$ odd} of the page is therefore necessary (figure~\ref{fig:quillen-1971-11}). The hypothesis that $B$ is a $p$-group is necessary too: the ring $\mathbf Z[\zeta_p]$ of cyclotomic integers, with $\theta$ multiplication by $\zeta_p$, has $\Omega_1 = 0$ and $\theta^p = 1$, without $p(\theta - 1)$ being zero.


\begin{figure}[htbp]
\centering
\begin{tikzpicture}[>=stealth, every node/.style={font=\small}]
\begin{scope}
\foreach \k in {0,...,8} { \coordinate (n\k) at ({90-40*\k}:1.45); }
\foreach \a/\b in {1/4,4/7,7/1,2/8,8/5,5/2} { \draw[->, frise, shorten >=4pt, shorten <=4pt] (n\a) -- (n\b); }
\foreach \k in {1,2,4,5,7,8} { \fill (n\k) circle (1.7pt); \node at ({90-40*\k}:1.8) {$\k$}; }
\foreach \k in {0,3,6} { \fill[piste] (n\k) circle (2.2pt); \node[piste] at ({90-40*\k}:1.8) {$\k$}; }
\node at (0,-2.35) {$p = 3$, $B = \mathbf Z/9$, $\theta = 4$};
\end{scope}
\begin{scope}[xshift=5.6cm]
\foreach \k in {0,...,7} { \coordinate (m\k) at ({90-45*\k}:1.45); }
\foreach \k in {1,2,3} { \pgfmathtruncatemacro{\l}{8-\k} \draw[<->, frise, shorten >=4pt, shorten <=4pt] (m\k) -- (m\l); }
\foreach \k in {1,2,3,5,6,7} { \fill (m\k) circle (1.7pt); \node at ({90-45*\k}:1.8) {$\k$}; }
\foreach \k in {0,4} { \fill[piste] (m\k) circle (2.2pt); \node[piste] at ({90-45*\k}:1.8) {$\k$}; }
\node at (0,-2.35) {$p = 2$, $B = \mathbf Z/8$, $\theta = -1$};
\end{scope}
\end{tikzpicture}
\caption{The action of $\theta$ (arrows) and the subgroup $\Omega_1(B)$ (in red). On the left, an example of the edition: multiplication by $4$ on $\mathbf Z/9$ satisfies $\theta^3 = 1$, fixes $\Omega_1(B) = \{0, 3, 6\}$ and cyclically permutes $1, 4, 7$ and $2, 8, 5$; $\theta - 1$ is multiplication by $3$, with values in $\Omega_1(B)$, as the theorem requires. On the right, counterexample~\ref{ce:quillen-1971-11}: $\theta = -1$ fixes $\Omega_1(B) = \{0, 4\}$, but $\theta(1) - 1 = 6$ is not in $\Omega_1(B)$.}
\label{fig:quillen-1971-11}
\end{figure}

\emph{Verification in Lean.} The file \href{\leanblobmain/Notebook1971_11.lean}{\texttt{Notebook1971\_11.\allowbreak{}lean}} proves Theorem~\ref{thm:quillen-1971-11} under the name \texttt{nsmul\_sub\_one\_eq\_zero}, for $\theta$ in the ring \texttt{AddMonoid.End B} of endomorphisms of $B$, and under the name \texttt{addEquiv\_nsmul\_sub\_one\_eq\_zero} for an automorphism \texttt{B ≃+ B} whose $p$-th iterate is the identity, as on the page. The lowering of the order is \texttt{sub\_one\_nsmul\_eq\_zero\_of\_nsmul}, the binomial expansion evaluated at a point is \texttt{pow\_apply\_eq\_sum}. Counterexample~\ref{ce:quillen-1971-11} is \texttt{zmod8\_neg}, checked by computation. The axioms used are the three usual axioms of mathlib.

\litterature{the search did not locate a printed reference stating this lemma in this form. It is of the type of lemmas that serve Thompson's replacement theorems, which the page says it revises, and the study of elementary abelian subgroups in Quillen's 1971--1978 work on the cohomology of finite groups; the expositions of these theorems in textbooks on finite groups have not been combed. The statement is elementary, and the edition does not present it as new.}

\subsection*{The quotient of a $p$-central $p$-group by $\Omega_1$\protect\verifie{Notebook1971_11Central}}\label{sec:quillen-1971-11-a}

\emph{Notation.} For a group $P$, we write $\Omega_1(P)$ for the subgroup generated by the elements $x$ such that $x^p = 1$; it is always normal. We say that $P$ is \emph{$p$-central} if every element of order $p$ of $P$ is central, that is, if $\Omega_1(P) \subset Z(P)$. In this case, $\Omega_1(P)$ is abelian of exponent $p$ and consists exactly of the elements $x$ such that $x^p = 1$.

\emph{The page.} The series of 24 August 1971 ends, on pages~70--71 of \carnet{1971}{11}, with a proposition: for $p$ odd, if $P$ is a $p$-central $p$-group, the quotient $P/\Omega_1(P)$ is also $p$-central. Quillen's proof puts $A = \Omega_1(P)$ and chooses a subgroup $B \supset A$, normal in $P$, such that $B/A$ is a maximal normal elementary abelian subgroup of $P/A$. Raising to the $p$-th power sends $B$ into $A$; it induces an injective map from $B/A$ to $A$, compatible with conjugation since $A$ is central, from which he deduces that $B/A$ is central in $P/A$. He concludes with the preceding proposition (pp.~69--70): a maximal normal elementary abelian subgroup is a maximal elementary abelian subgroup ($p$ odd), so that $B/A$ contains all the elements of order $p$ of $P/A$. It is this intermediate proposition that uses the lemma of Theorem~\ref{thm:quillen-1971-11}.

\emph{Intuition.} An element $x$ whose class modulo $A$ has order dividing $p$ satisfies $x^p \in A$. For $g \in P$, the conjugate $u = gxg^{-1}$ has the same $p$-th power as $x$, since $x^p$ is central. If $u$ and $x$ commute modulo $A$, the group they generate has class at most~$2$, and for $p$ odd raising to the $p$-th power is a homomorphism there: $(ux^{-1})^p = u^p x^{-p} = 1$, so the commutator $gxg^{-1}x^{-1}$ is in $A$. It remains to see that $u$ and $x$ commute modulo $A$; this is the statement itself in a maximal subgroup containing $x$, whence an induction on the order.

\begin{encadre}
\begin{theoreme}[Quillen, notebook 1971-11, pp.~70--71]\label{thm:quillen-1971-11-a}
Let $p$ be an odd prime and $P$ a finite $p$-group in which every element $z$ with $z^p = 1$ is central. Then every element $y$ of $P/\Omega_1(P)$ with $y^p = 1$ is central; in other words, $\Omega_1(P/\Omega_1(P)) \subset Z(P/\Omega_1(P))$.
\end{theoreme}
\end{encadre}

\begin{proof}
Put $A = \Omega_1(P)$. Since $A = \{x \in P \mid x^p = 1\}$, we must show: if $(x^p)^p = 1$, then $(gxg^{-1}x^{-1})^p = 1$ for all $g \in P$. We argue by induction on the order of $P$; the hypothesis passes to subgroups, since an element central in $P$ is central in every subgroup containing it.

If $x$ generates $P$, the group is commutative and the commutator is trivial. Otherwise, $\langle x \rangle$ is contained in a maximal subgroup $M$, normal because $P$ is nilpotent. The conjugate $u = gxg^{-1}$ is in $M$ and satisfies $u^p = g x^p g^{-1} = x^p$, whence $(u^p)^p = 1$. The induction hypothesis, applied in $M$ to the elements $u$ and $x^{-1}$, gives $d^p = 1$ for $d = x^{-1} u x u^{-1}$; thus $d$ is central. We have $x^{-1} u = d\, u\, x^{-1}$, and an immediate induction gives, for all $n$,
\[
(u x^{-1})^n = u^n\, x^{-n}\, d^{\binom n2}.
\]
For $n = p$ odd, $p$ divides $\binom p2$, so $d^{\binom p2} = 1$ and $(ux^{-1})^p = u^p x^{-p} = 1$. Now $u x^{-1} = gxg^{-1}x^{-1}$.

Finally, if $y = xA$ satisfies $y^p = 1$, then $x^p \in A$, that is $(x^p)^p = 1$, and for all $g$ the commutator $gxg^{-1}x^{-1}$ is in $A$: $y$ commutes with $gA$.
\end{proof}

\noindent The proof keeps Quillen's first step (raising to the $p$-th power is an injective homomorphism compatible with conjugation, which forces the commutators into $A$), but replaces the appeal to the proposition of pages~69--70 by an induction on the order of $P$. It therefore does not use the lemma of Theorem~\ref{thm:quillen-1971-11}. The hypothesis $p$ odd serves at the same place as in Quillen: the divisibility of $\binom p2$ by $p$, which makes raising to the $p$-th power a homomorphism in class~$2$.

\emph{Verification in Lean.} The file \href{\leanblobmain/Notebook1971_11Central.lean}{\texttt{Notebook1971\_11Central.\allowbreak{}lean}} defines $\Omega_1(P)$ under the name \texttt{omega1}, as the subgroup generated by the elements $x$ such that $x^p = 1$, and shows that it is normal (\texttt{omega1\_normal}) and, under the hypothesis of $p$-centrality, that it consists exactly of these elements (\texttt{mem\_omega1\_iff}). The formula $(uv)^n = u^n v^n d^{\binom n2}$ is \texttt{mul\_pow\_eq\_mul\_pow\_mul\_pow\_choose}, and its case $n = p$ odd is \texttt{mul\_pow\_eq\_of\_commutator\_pow}. The induction on the order is \texttt{conj\_commutator\_pow\_eq\_one\_aux}, its statement \texttt{conj\_commutator\_pow\_eq\_one}. Theorem~\ref{thm:quillen-1971-11-a} is \texttt{pCentral\_quotient}, for elements, and \texttt{omega1\_quotient\_le\_center}, in the form of the page $\Omega_1(P) \le Z(P) \Rightarrow \Omega_1(P/\Omega_1(P)) \le Z(P/\Omega_1(P))$. The axioms used are the three usual axioms of mathlib.

Not proved here is the proposition of pages~69--70 (a maximal normal elementary abelian subgroup of a $p$-group, $p$ odd, is a maximal elementary abelian subgroup), through which Quillen's proof passes and which uses the lemma of Theorem~\ref{thm:quillen-1971-11}.

\litterature{the statement is known before the notebook: it is the theorem that J.~González-Sánchez and Th.~Weigel (2011, introduction) call Blackburn--Buckley--Thompson. It is published in this form by J.~Buckley, \q{Finite groups whose minimal subgroups are normal}, \emph{Math.~Z.} 116 (1970), 15--17, and follows from a result of Thompson that B.~Huppert reproduces as early as 1967, with a proof attributed to Blackburn (\emph{Endliche Gruppen}~I, Springer, III, Hilfssatz~12.2). Theorem~B of González-Sánchez and Weigel generalizes it to $p^i$-central pro-$p$-groups of height~$k$. The notebook reproves it in 1971, as a lemma in a revision of Thompson's theorems; only the proof by induction given here is the edition's.}

\section[Notebook 1971-12: a group with $FV = VF = p^d$ and $F = p^{d+r}$]{Notebook 1971-12: an abelian group with $FV = VF = p^d$ and $F = p^{d+r}$\protect\verifie{Notebook1971_12}}\label{sec:quillen-1971-12}

\emph{Context.} In September 1971, \carnet{1971}{12} considers the $K$-groups $K_nA$ of a ring $A$ of characteristic $p$, with the map $F$ induced by the Frobenius and its transfer $V$. On page~22 he finds $VF = p^d$, where $p^d$ is the rank of $A$ over its subring of $p$-th powers, and he argues on pages~22 and~23 that $FV$ is also multiplication by $p^d$. He wants to combine this with a conjecture, that $K_nA$ has high $\gamma$-filtration for $n$ large, which would make $F$ act through high powers of $p$.

\emph{The page.} He isolates the essential case as a lemma on abelian groups (page~23, his page~2). Let $M$ be an abelian group with two endomorphisms $F$, $V$ such that $FV = VF = p^d$ and $F = p^{d+r}$ with $r > 0$. Then $M = M' \oplus M''$, where $M' = p^dM$ is the largest $p$-divisible subgroup and $M''$, the subgroup of elements killed by $p^d$, is the $p$-torsion subgroup. So $M$ is the sum of a $\mathbf Z[1/p]$-module and a group of exponent $p^d$. The proof, on page~24, shows that $p^d(x - p^rVx) = 0$, splits $M$ by the operator $1 - p^rV$, which is an automorphism of $M''$, and iterates $x = p^rVx$ on its kernel.

\emph{Intuition.} Since $F$ is multiplication by $p^{d+r}$, the relation $FV = p^d$ reads $p^d\,(1 - p^rV) = 0$: up to elements killed by $p^d$, every $x$ equals $p^rVx$, hence $p^{2r}V^2x$, and so on, so that $x$ is divisible by arbitrarily high powers of $p$ up to such elements. The relation $VF = p^d$ reads $p^d = V \circ p^{d+r}$: what is killed by $p^{d+r}$ is already killed by $p^d$, so the torsion cannot grow.

\begin{encadre}
\begin{theoreme}[Quillen, notebook 1971-12, p.~23]\label{thm:quillen-1971-12}
Let $M$ be an abelian group, $p \ge 1$, $d \ge 0$ and $r > 0$ integers, and $F$, $V$ endomorphisms of $M$ such that $FV = VF = p^d$ and $F = p^{d+r}$. Then:
\begin{enumerate}
\item[\emph{(a)}] $M = p^dM \oplus M[p^d]$, where $M[p^d]$ is the subgroup of elements killed by $p^d$;
\item[\emph{(b)}] every element of $M$ killed by a power of $p$ is killed by $p^d$;
\item[\emph{(c)}] $p^dM = p^sM$ for all $s \ge d$, and $p^dM$ contains every subgroup $D$ such that $D \subset pD$.
\end{enumerate}
\end{theoreme}
\end{encadre}

\begin{proof}
For all $x$, $p^dx = FVx = p^{d+r}Vx$, so $t(x) = x - p^rVx$ is killed by $p^d$. By induction on $n$, $x = t_n + p^{nr}V^nx$ with $p^dt_n = 0$: apply the identity to $V^nx$ and multiply by $p^{nr}$. (b) If $p^{d+r}x = 0$, in particular if $p^sx = 0$ with $s \le d + r$, then $p^dx = VFx = V(p^{d+r}x) = 0$; if $p^sx = 0$ with $s > d + r$, apply this to $p^{s-d-r}x$ to get $p^{s-r}x = 0$, and conclude by induction on $s$. (c) Multiply $x = t_n + p^{nr}V^nx$ by $p^d$: $p^dx = p^{d+nr}V^nx$, which lies in $p^sM$ as soon as $nr \ge s - d$. If $D \subset pD$ then $D \subset p^dD \subset p^dM$. (a) Take $n = d$: since $dr \ge d$, $x = t_d + p^{dr}V^dx$ lies in $M[p^d] + p^dM$. If $x = p^dy$ and $p^dx = 0$, then $p^{2d}y = 0$, so $p^dy = 0$ by (b), that is $x = 0$.
\end{proof}

\noindent The proof of the page splits $M$ as the kernel of $1 - p^rV$ plus $M[p^d]$, using that $1 - p^rV$ is invertible on $M[p^d]$, where $p^rV$ is nilpotent. The proof above iterates the same identity on every element and does not use that inverse. Neither proof uses that $p$ is prime.

\emph{Verification in Lean.} The file \href{\leanblobmain/Notebook1971_12.lean}{\texttt{Notebook1971\_12.\allowbreak{}lean}} takes two additive endomorphisms \texttt{F}, \texttt{V} of an additive commutative group \texttt{M}, with the three hypotheses written pointwise. The identity is \texttt{nsmul\_sub\_eq\_zero} and its iteration \texttt{exists\_decomp}. Part~(b) is \texttt{nsmul\_eq\_zero\_of\_pow\_nsmul}, part~(c) is \texttt{exists\_nsmul\_eq} and \texttt{le\_range\_of\_divisible}, and part~(a) is \texttt{isCompl\_range\_ker}, which says that the image and the kernel of multiplication by $p^d$ are complementary subgroups. The application to $K$-theory is not formalised; on the page it depends on a conjecture that Quillen names as such. The axioms used are the three usual axioms of mathlib.

\litterature{The search did not locate the lemma as a stated result; it is elementary, of the kind used for Frobenius and Verschiebung on groups killed by, or divisible by, powers of $p$. The consequence drawn on the page, that $K_nA$ would be the sum of a $\mathbf Z[1/p]$-module and a group of exponent $p^d$, is conditional on the conjecture. The edition does not present the lemma as new.}

\section[Notebook 1971-12: Kaloujnine's theorem]{Notebook 1971-12: Kaloujnine's theorem $[A_r, A_s] \le A_{r+s}$\protect\verifie{Notebook1971_12Kaloujnine}}\label{sec:quillen-1971-12-kaloujnine}

\emph{Context.} On pages~21 and~34 of \carnet{1971}{12} Quillen states Kaloujnine's theorem on the automorphisms that act trivially on the successive quotients of a sequence of normal subgroups, and proves it by the three-subgroup lemma in the semidirect product of $G$ by its automorphisms.

\begin{encadre}
\begin{theoreme}[Kaloujnine; Quillen, notebook 1971-12, p.~34]\label{thm:quillen-1971-12-kaloujnine}
Let $(G_i)_{i \ge 0}$ be normal subgroups of a group $G$, and for $r \ge 0$ let $A_r$ be the set of automorphisms $\theta$ of $G$ with $\theta(G_i) = G_i$ and $\theta(x)x^{-1} \in G_{i+r}$ for all $i$ and all $x \in G_i$. Then $A_r$ is a subgroup and $[A_r, A_s] \le A_{r+s}$.
\end{theoreme}
\end{encadre}

\begin{proof}
Work in the group of permutations of the set $G$, which contains the left translations $L_g$ and the automorphisms, with $\theta L_g\theta^{-1} = L_{\theta(g)}$; hence $[L_x, \theta] = L_{x\theta(x)^{-1}}$, a translation by an element of $G_{i+r}$ when $x \in G_i$ and $\theta \in A_r$. For $\theta \in A_r$, $\varphi \in A_s$ and $g \in G_i$, the Hall--Witt identity writes $[L_g, [\theta, \varphi]]$ as a product of conjugates of $[\varphi, [L_{g^{-1}}, \theta^{-1}]]$ and $[\theta^{-1}, [\varphi^{-1}, L_g]]$, both translations by elements of $G_{i+r+s}$; as $G_{i+r+s}$ is normal and stable under $\theta$, $\varphi$, the result is a translation by an element of $G_{i+r+s}$, which is $[\theta,\varphi](g)g^{-1}$ up to inversion.
\end{proof}

\emph{Verification in Lean.} The file \href{\leanblobmain/Notebook1971_12Kaloujnine.lean}{\texttt{Notebook1971\_12Kaloujnine.\allowbreak{}lean}} defines \texttt{A} and proves \texttt{commutator\_le}. The semidirect product of the page is replaced by the permutation group of $G$, in which it embeds; no monotonicity of the $G_i$ is used. The corollary of p.~21 (nilpotency of the automorphisms of a nilpotent group trivial on the abelianisation) is not formalised. The axioms used are the three usual axioms of mathlib.

\litterature{Classical: Kaloujnine (1953); P.~Hall, Illinois J. Math. 2 (1958). The edition does not present it as new.}

\section[Notebook 1971-12: Kaloujnine's corollary]{Notebook 1971-12: automorphisms trivial on $G^{\mathrm{ab}}$ form a nilpotent group\protect\verifie{Notebook1971_12Nilpotent}}\label{sec:quillen-1971-12-nilpotent}

\emph{Context.} On page~21 of \carnet{1971}{12}, after Kaloujnine's theorem (section~\ref{sec:quillen-1971-12-kaloujnine}), Quillen draws its classical corollary for nilpotent groups.

\begin{encadre}
\begin{theoreme}[Quillen, notebook 1971-12, p.~21]\label{thm:quillen-1971-12-nilpotent}
Let $G$ be a group, $\Gamma_0 = G$, $\Gamma_{i+1} = [\Gamma_i, G]$ its lower central series.
\begin{enumerate}
\item For an antitone sequence of normal subgroups $G_i$ with $G_0 = G$ and some $G_c = 1$, the group $A_1$ of automorphisms preserving the $G_i$ and trivial on each $G_i/G_{i+1}$ is nilpotent.
\item Conjugation by an element of $\Gamma_k$ lies in $A_{k+1}$ (for the lower central series); hence $[\Gamma_j, \Gamma_i] \le \Gamma_{i+j+1}$.
\item An automorphism $\theta$ lies in $A_1$ for the lower central series iff $\theta(g)g^{-1} \in [G, G]$ for all $g$.
\item Hence, if $G$ is nilpotent, the automorphisms inducing the identity on $G^{\mathrm{ab}}$ form a nilpotent group.
\end{enumerate}
\end{theoreme}
\end{encadre}

\begin{proof}
(1) By Kaloujnine's theorem the $n$-th term of the lower central series of $A_1$ lies in $A_{n+1}$, and $A_r = 1$ once $G_r = 1$. (2) By induction on $k$: $c_g \in A_1$ since $gxg^{-1}x^{-1} \in [\Gamma_i, G]$; then $c_{[x,g]} = [c_x, c_g] \in [A_{k+1}, A_1] \le A_{k+2}$. (3) By induction on $i$ over the generators $[p, q]$ of $\Gamma_{i+1}$: modulo $\Gamma_{i+2}$, $\theta(p)p^{-1}$ is central and $\theta(q)q^{-1} \in \Gamma_1$ commutes with $p$ by~(2), so $\theta[p, q] \equiv [p, q]$; the elements with $\theta(x)x^{-1} \in \Gamma_{i+2}$ form a subgroup.
\end{proof}

\emph{Verification in Lean.} The file \href{\leanblobmain/Notebook1971_12Nilpotent.lean}{\texttt{Notebook1971\_12Nilpotent.\allowbreak{}lean}} proves \texttt{isNilpotent\_A\_one}, \texttt{conj\_mem\_A}, \texttt{commutator\_lowerCentralSeries\_le}, \texttt{mem\_A\_one\_iff} and \texttt{isNilpotent\_trivialOnAb}, importing Kaloujnine's theorem. Triviality on $G^{\mathrm{ab}}$ is stated as $\theta(g)g^{-1} \in [G, G]$. The axioms used are the three usual axioms of mathlib.

\litterature{Classical (Kaloujnine; P.~Hall); $[\Gamma_i, \Gamma_j] \le \Gamma_{i+j}$ is textbook but absent from mathlib, derived here from Kaloujnine's theorem. The edition does not present it as new.}


\part{Representations and group actions}
\section[Notebook 1968-1: a simple module is determined by its trace]{Notebook 1968-1: a simple module is determined by its trace\protect\verifie{Notebook1968_1}}\label{sec:quillen-1968-1}

\emph{Context.} \carnet{1968}{1}, dated January 1968, is mostly about representations of semisimple Lie algebras. On PDF page~4, Quillen isolates a basic statement: an irreducible representation of a finite-dimensional algebra over $\mathbb{C}$ is determined by its trace function $a \mapsto \operatorname{tr}(a \mid S)$. His sketch: the trace vanishes on the radical, whose elements act nilpotently; by Wedderburn, the semisimple quotient is a product of simple algebras, and the trace of $S$ is nonzero on the factor of $S$ and zero on the others. A remark that the trace equals $\dim S$ at $1$ is crossed out, rightly so in characteristic $p$, where $\dim S$ can vanish in the field.

\emph{Intuition.} If $S \not\simeq T$, it suffices to find an element of $A$ acting on $S$ as a projector of rank one and on $T$ as zero: its traces are $1$ and $0$, in every characteristic. Jacobson's density theorem provides such an element, because the endomorphisms of $S \times T$ are diagonal and scalar.

\begin{encadre}
\begin{theoreme}[Quillen, notebook 1968-1, p.~4]\label{thm:quillen-1968-1}
Let $k$ be an algebraically closed field of arbitrary characteristic, $A$ a $k$-algebra, and $S$, $T$ two simple $A$-modules of finite dimension over $k$. If $\operatorname{tr}(a \mid S) = \operatorname{tr}(a \mid T)$ for all $a \in A$, then $S \simeq T$. More precisely, if $S \not\simeq T$, there exists $a \in A$ with $\operatorname{tr}(a \mid S) = 1$ and $\operatorname{tr}(a \mid T) = 0$. When $A$ is finite-dimensional, as in the notebook, the finite-dimensionality hypothesis on $S$ and $T$ is automatic.
\end{theoreme}
\end{encadre}

\begin{lemme}[Schur]\label{lem:quillen-1968-1-schur}
An $A$-linear map between non-isomorphic simple modules is zero. If $k$ is algebraically closed and $S$ is of finite dimension over $k$, every $A$-linear endomorphism of $S$ is a scalar multiplication.
\end{lemme}

\begin{proof}
The kernel and the image are submodules, so a nonzero map between simple modules is bijective. An endomorphism $g$ of $S$ has an eigenvalue $c \in k$; $g - c$ is $A$-linear and not injective, hence zero.
\end{proof}

\begin{proof}[Proof of the theorem]
The module $S \times T$ is semisimple, being the sum of the images of the two injections. By Lemma~\ref{lem:quillen-1968-1-schur}, if $S \not\simeq T$, an $A$-linear endomorphism $g$ of $S \times T$ has the form $g(s, t) = (c\,s, d\,t)$ with $c, d \in k$. Let $v \neq 0$ in $S$, $\varphi$ a linear form with $\varphi(v) = 1$, and $\psi = \varphi(\cdot)\,v$, of trace $1$. The map $f(s, t) = (\psi(s), 0)$ commutes with all these $g$; since $S \times T$ is finitely generated over its endomorphism ring (it already is over $k$), the density theorem gives $a \in A$ acting as $f$. Then $\operatorname{tr}(a \mid S) = \operatorname{tr}\psi = 1$ and $\operatorname{tr}(a \mid T) = 0$, which contradicts the equality of traces. If $A$ is finite-dimensional, a simple module is cyclic, hence finite-dimensional over $k$.
\end{proof}

\begin{remarque}
The proof avoids the radical and Wedderburn--Artin, which Quillen uses; the two routes give the same element, a matrix unit $e_{11}$ of the block of $S$. The finiteness of $A$ is not used, only that of $S$ and $T$.
\end{remarque}

\begin{contrexemple}
The hypothesis \q{algebraically closed} cannot be dropped in characteristic $p$. Let $L/k$ be a purely inseparable extension of degree $p$ and $A = L \times L$; the two simple modules $L \times 0$ and $0 \times L$ are not isomorphic, and their traces, $(x, y) \mapsto \operatorname{Tr}_{L/k}(x)$ and $(x, y) \mapsto \operatorname{Tr}_{L/k}(y)$, are both zero. Likewise, over an algebraically closed field of characteristic $p$, the trace does not determine semisimple modules: $S^{\oplus p}$ and the zero module have the same trace.
\end{contrexemple}

\begin{sloppypar}
\noindent\emph{Verification in Lean.} The theorem is \texttt{nonempty\_linearEquiv\_of\_trace\_eq}, and \texttt{nonempty\_linearEquiv\_of\_trace\_eq'} for $A$ finite-dimensional; the element with traces $1$ and $0$ is \texttt{exists\_trace\_eq\_one\_trace\_eq\_zero}. Lemma~\ref{lem:quillen-1968-1-schur} is \texttt{eq\_zero\_of\_isEmpty\_linearEquiv} and \texttt{exists\_eq\_smul\_of\_isAlgClosed}, the semisimplicity of $S \times T$ is \texttt{isSemisimpleModule\_prod}, in \href{\leanblobmain/Notebook1968_1.lean}{\texttt{Notebook1968\_1.\allowbreak{}lean}}. Density is mathlib's \texttt{Module.Finite.toModuleEnd\_moduleEnd\_surjective}.
\end{sloppypar}

Not formalized: the counterexample above, and the linear independence of the traces of pairwise non-isomorphic simple modules, which follows from the same argument.

\litterature{found in the literature: Curtis--Reiner, \emph{Representation Theory of Finite Groups and Associative Algebras} (1962), the chapter on characters of algebras; Lam, \emph{A First Course in Noncommutative Rings}, the chapter on representations. mathlib has the density theorem and Wedderburn--Artin over an algebraically closed field, but not this statement; it could be proposed for mathlib, and the edition does not present it as new.}

\section[Carnet 1968-13 : quand la restriction des représentations est pleine]{Carnet 1968-13 : quand la restriction des représentations est pleine\protect\verifie{Notebook1968_13Rep}}\label{sec:quillen-1968-13-rep}

\emph{Le contexte.} Le \carnet{1968}{13} s'ouvre, à l'automne 1968, sur un théorème de Whitehead pour les petites catégories : un foncteur $f : \mathcal{C}_1 \to \mathcal{C}_2$ devrait être une équivalence faible dès que le retrait le long de $f$ induit une équivalence entre catégories dérivées de systèmes de coefficients localement constants. Dans la preuve de l'implication \q{(i) $\Rightarrow$ (ii)} (PDF, p.~4), Quillen se ramène au cas d'un homomorphisme de groupes et a besoin d'un petit lemme, qu'il énonce entre parenthèses : la restriction des modules le long d'un homomorphisme de groupes est pleinement fidèle exactement quand l'homomorphisme est surjectif, et une équivalence exactement quand il est bijectif. Il justifie la surjectivité en comparant les invariants du module de permutation $\mathbb{Z}[G/f(H)]$ sous les deux groupes, l'injectivité en remarquant que le noyau agit trivialement sur tout module restreint. Sur la page, les lettres $G$ et $H$ sont échangées entre l'énoncé et l'argument des rangs.

\emph{L'intuition.} Une flèche entre restrictions commute à l'action de l'image $f(H)$, pas forcément à celle de $G$ ; il suffit d'un module où l'écart se voit. Dans $k[G/f(H)]$, la classe de base est fixée par $f(H)$ et par lui seul : la projection sur la droite qu'elle engendre commute à $H$, et elle ne commute à $G$ que si $G$ fixe cette classe, c'est-à-dire si $f(H) = G$. Pour l'injectivité, la représentation régulière $k[H]$ est fidèle, alors que le noyau de $f$ agit trivialement sur tout ce qui vient de $G$.

\begin{encadre}
\begin{theoreme}[Quillen, carnet 1968-13, p.~4]\label{thm:quillen-1968-13-rep}
Soient $f : H \to G$ un homomorphisme de groupes et $k$ un anneau commutatif non nul. Le foncteur de restriction $\operatorname{Rep}_k(G) \to \operatorname{Rep}_k(H)$ est plein si et seulement si $f$ est surjectif ; c'est une équivalence si et seulement si $f$ est bijectif.
\end{theoreme}
\end{encadre}

\begin{proof}
La restriction est toujours fidèle : elle ne change pas l'application linéaire sous-jacente. Si $f$ est surjectif, une application $H$-linéaire entre restrictions commute à tout $g = f(h)$, donc est $G$-linéaire.

Supposons la restriction pleine. Soient $V = k[G/f(H)]$, $c_0$ la classe de $f(H)$ et $p : V \to V$ l'application $x \mapsto x(c_0)\,c_0$. Pour $h \in H$, $f(h)$ et $f(h)^{-1}$ fixent $c_0$, donc $p(f(h)x) = f(h)\,p(x)$ : $p$ est $H$-linéaire. Elle provient donc d'une application $G$-linéaire, qui est encore $p$. Pour $g \in G$, on a $p(g c_0) = g\,p(c_0) = g c_0 \neq 0$ puisque $k$ est non nul, donc $g c_0 = c_0$, c'est-à-dire $g \in f(H)$.

Si $f$ est bijectif, d'inverse $f^{-1}$, toute représentation $\rho$ de $H$ est la restriction de $\rho \circ f^{-1}$ ; la restriction est alors pleine, fidèle et essentiellement surjective. Réciproquement, supposons-la essentiellement surjective et soit $h \in \ker f$. La représentation régulière $k[H]$ est isomorphe à la restriction d'une représentation $W$ de $G$ ; $h$ agit sur $W$ par $f(h) = 1$, donc trivialement, et par transport il fixe $\delta_1 \in k[H]$. Or $h\,\delta_1 = \delta_h$, d'où $h = 1$.
\end{proof}

\begin{sloppypar}
\noindent\emph{Vérification en Lean.} Dans \href{\leanblobmain/Notebook1968_13Rep.lean}{\texttt{Notebook1968\_13Rep.\allowbreak{}lean}}, la restriction est le foncteur \texttt{Rep.resFunctor f} de mathlib. Le critère de plénitude est \texttt{full\_resFunctor\_iff}, dont le sens direct est \texttt{surjective\_of\_full}, avec la projection \texttt{proj} sur la classe de base du module \texttt{cosetRep} ; le critère d'équivalence est \texttt{isEquivalence\_resFunctor\_iff}, avec \texttt{injective\_of\_essSurj} (testé sur \texttt{regRep}) et \texttt{essSurj\_of\_bijective}. Les représentations vivent dans l'univers $\max(u, v_1, v_2)$, assez grand pour contenir les objets tests.
\end{sloppypar}

Le cas de Quillen est $k = \mathbb{Z}$ ; l'argument par les rangs des invariants est remplacé par la projection, qui lui est équivalente (une flèche depuis la représentation triviale est un invariant). N'est pas formalisé : le théorème de Whitehead lui-même, qui demande les catégories dérivées de coefficients localement constants.

\litterature{folklore, retrouvé : l'inflation le long d'un quotient $G \to G/N$ identifie les $G/N$-modules aux $G$-modules sur lesquels $N$ agit trivialement, sous-catégorie pleine. mathlib a le sens \q{surjectif $\Rightarrow$ plein} (\texttt{Rep.full\_res}, et \texttt{Action.full\_res} pour les actions) mais ni la réciproque ni le critère d'équivalence. Ils sont proposés pour mathlib, et l'édition ne les donne pas pour nouveaux.}

\section[Notebook 1971-12: marks along a normal $p$-step]{Notebook 1971-12: marks along a normal $p$-step\protect\verifie{Notebook1971_12Marks}}\label{sec:quillen-1971-12-marks}

\emph{Context.} On pages~27 to~33 of \carnet{1971}{12} Quillen studies the Burnside ring of a finite group $G$ and its marks $\varphi_H(X) = \lvert X^H\rvert$ modulo a prime $p$. On pages~28 and~30 he uses that the marks of two subgroups agree modulo $p$ when one is normal in the other with a $p$-group as quotient; iterating the step identifies the marks modulo $p$ with those of subgroups of index prime to $p$ in their normaliser.

\emph{Intuition.} If $H \lhd H'$, the group $H'/H$ acts on $X^H$, and its fixed points there are $X^{H'}$. A $p$-group acting on a finite set has a number of fixed points congruent to the size of the set modulo $p$.

\begin{encadre}
\begin{theoreme}[After Quillen, notebook 1971-12, pp.~28, 30]\label{thm:quillen-1971-12-marks}
Let a group $G$ act on a finite set $X$ and let $N \lhd G$ with $G/N$ a $p$-group. Then $G/N$ acts on $X^N$ with fixed points $X^G$, and $\lvert X^G\rvert \equiv \lvert X^N\rvert \pmod p$.
\end{theoreme}
\end{encadre}

\begin{proof}
For $x \in X^N$, $n \in N$ and $g \in G$, $n g x = g (g^{-1}ng) x = gx$, so $G$ preserves $X^N$, and $N$ acts trivially on it, which gives the action of $G/N$; its fixed points are the points fixed by all of $G$. The congruence is the fixed-point congruence for the $p$-group $G/N$ on the finite set $X^N$.
\end{proof}

\emph{Verification in Lean.} The file \href{\leanblobmain/Notebook1971_12Marks.lean}{\texttt{Notebook1971\_12Marks.\allowbreak{}lean}} constructs the action (\texttt{quotientAction}), the bijection $(X^N)^{G/N} \simeq X^G$ (\texttt{fixedPointsQuotientEquiv}) and proves the congruence as \texttt{card\_fixedPoints\_modEq}, from mathlib's \texttt{IsPGroup.card\_modEq\_card\_fixedPoints}. $G$ need not be finite. The companion fact of page~28, $(G/H)^H \simeq N_G(H)/H$, is mathlib's \texttt{Sylow.fixedPointsMulLeftCosetsEquivQuotient}. The Burnside ring is not formalised. The axioms used are the three usual axioms of mathlib.

\litterature{Classical (A.~Dress, \emph{A characterisation of solvable groups}, Math.~Z.~110, 1969, on the prime ideals of the Burnside ring; tom Dieck, \emph{Transformation groups and representation theory}, LNM~766, ch.~1). The edition does not present it as new.}

\section[Notebook 1973-6: the index over double cosets]{Notebook 1973-6: the index of a subgroup as a sum over double cosets\protect\verifie{Notebook1973_6Doset}}\label{sec:quillen-1973-6-doset}

\emph{Context.} On page~61 of \carnet{1973}{6} Quillen counts the index of a subgroup $\Gamma'$ of finite index of $\Gamma$ through the orbits of a finite subgroup $\Gamma_P$ on $\Gamma/\Gamma'$, a sum of ratios of orders of stabilisers.

\begin{encadre}
\begin{theoreme}[Quillen, notebook 1973-6, p.~61]\label{thm:quillen-1973-6-doset}
Let $K$ be a subgroup of finite index of a group $G$ and $H$ a finite subgroup. Then
\[ [G : K] = \sum_{HgK} \frac{\lvert H\rvert}{\lvert H \cap gKg^{-1}\rvert}, \]
the sum over the double cosets $HgK$, one representative $g$ for each; each term is the index $[H : H \cap gKg^{-1}]$, and in this form $H$ need not be finite.
\end{theoreme}
\end{encadre}

\begin{proof}
$H$ acts on $G/K$; its orbits are the double cosets, and the stabiliser of $gK$ is $H \cap gKg^{-1}$, since $hgK = gK$ iff $g^{-1}hg \in K$. Sum the orbit--stabiliser formula over the orbits.
\end{proof}

\emph{Verification in Lean.} The file \href{\leanblobmain/Notebook1973_6Doset.lean}{\texttt{Notebook1973\_6Doset.\allowbreak{}lean}} proves \texttt{stabilizer\_eq}, \texttt{index\_eq\_sum} and \texttt{index\_eq\_sum\_index}, the orbits being indexed by mathlib's quotient by the orbit relation. The count itself is mathlib's \texttt{MulAction.card\_eq\_sum\_card\_group\_div\_card\_stabilizer}; the file identifies the stabilisers. The axioms used are the three usual axioms of mathlib.

\litterature{Classical (orbit--stabiliser; Mackey's double cosets); the edition does not present it as new.}


\part{Number theory and congruences}
\section[Notebook 1969-4: $(p^j)!$ divided by its power of $p$]{Notebook 1969-4: $(p^j)!$ divided by its power of $p$, modulo $p$\protect\verifie{Notebook1969_4}}\label{sec:quillen-1969-4}

\emph{Context.} \carnet{1969}{4}, of March 1969, contains among other things a study of formal group laws of height one and of characteristic numbers in $K$-theory. On page~44 Quillen recalls a theorem of Adams on the Chern character and states a lemma: if $f$ is a root of $X^{p-1} = p$, the power series $(1 - e^{-fX})/(fX)$ is $p$-integral and congruent modulo $p$ to $\sum_{j \ge 0} (-1)^j X^{p^j - 1}$.

\emph{The page.} The proof is on page~45. The coefficient of $X^{n-1}$ is, up to sign, $f^{n-1}/n!$, of valuation $(n-1)/(p-1) - v(n!)$. Quillen notes that this is positive unless $n = p^j$, in which case $v(n!) = (n-1)/(p-1)$. It remains to find the residue modulo $p$ of $(p^j)!/p^r$, with $r = (p^j - 1)/(p - 1)$. He splits the integers up to $p^j$ into those prime to $p$ and those divisible by $p$, cites Wilson for the first, and finds by induction that $(p^j)!/p^{1 + p + \dots + p^{j-1}} \equiv (-1)^j$ modulo $p$.

\emph{Intuition.} Among the integers from $1$ to $pm$, the multiples of $p$ are $p, 2p, \dots, mp$, with product $p^m\,m!$. The others fall into $m$ blocks $pk + 1, \dots, pk + p - 1$, each congruent modulo $p$ to $1, \dots, p - 1$, with product $(p-1)! \equiv -1$ by Wilson's theorem. Hence
\[
  (pm)! = p^m \cdot m! \cdot U_m, \qquad U_m \equiv (-1)^m \pmod p.
\]
With $m = p^{j-1}$ this passes from $(p^{j-1})!$ to $(p^j)!$, and each step contributes a sign $(-1)^{p^{j-1}}$, which is $-1$ modulo $p$ (for $p = 2$, $-1 = 1$).

\begin{encadre}
\begin{theoreme}[Quillen, notebook 1969-4, p.~45]\label{thm:quillen-1969-4}
Let $p$ be a prime and $j \ge 0$. The $p$-adic valuation of $(p^j)!$ is $r = (p^j - 1)/(p - 1) = 1 + p + \dots + p^{j-1}$, and
\[
  \frac{(p^j)!}{p^{\,r}} \equiv (-1)^j \pmod p.
\]
\end{theoreme}
\end{encadre}

\begin{proof}
By induction on $m$, $(pm)! = p^m\,m!\,U_m$, where $U_m$ is the product of the integers up to $pm$ that are prime to $p$: passing from $m$ to $m + 1$ multiplies the factorial by $(pm + 1)\cdots(pm + p - 1)$ and by $p(m + 1)$. Each block $(pk + 1)\cdots(pk + p - 1)$ is congruent to $(p - 1)! \equiv -1$ modulo $p$, so $U_m \equiv (-1)^m$. We show by induction on $j$ that $(p^j)! = p^{r_j} a_j$ with $r_j = 1 + \dots + p^{j-1}$ and $a_j \equiv (-1)^j$. For $j = 0$, $a_0 = 1$. For the step, $(p^{j+1})! = p^{p^j}\,(p^j)!\,U_{p^j} = p^{p^j + r_j}\,a_j U_{p^j}$, and $a_jU_{p^j} \equiv (-1)^j(-1)^{p^j} = (-1)^{j+1}$ modulo $p$, since $p^j$ is odd when $p$ is odd and $-1 = 1$ modulo $2$. As $a_j$ is not divisible by $p$, the valuation of $(p^j)!$ is $r_j$, and $r_j = (p^j - 1)/(p - 1)$ is the sum of the geometric progression.
\end{proof}

\noindent The page writes the sign of one step as $(-1)^{p^{j-1}}$ and concludes $(-1)^j$, which supposes $p$ odd. The statement also holds for $p = 2$, where it says that $(2^j)!/2^{2^j - 1}$ is odd.

\emph{Verification in Lean.} The file \href{\leanblobmain/Notebook1969_4.lean}{\texttt{Notebook1969\_4.\allowbreak{}lean}} defines \texttt{block p m}, the product of $pm + 1, \dots, pm + p - 1$, and \texttt{primeToPart p m}, the number $U_m$. The splitting is \texttt{factorial\_mul}, and the congruences are \texttt{block\_zmod}, from \texttt{ZMod.wilsons\_lemma} of mathlib, and \texttt{primeToPart\_zmod}. The induction is \texttt{exists\_factorial\_pow\_eq}. The two parts of Theorem~\ref{thm:quillen-1969-4} are \texttt{padicValNat\_factorial\_pow} and \texttt{factorial\_pow\_div\_zmod}; in the second the quotient is the quotient of natural numbers, and the congruence is an equality in \texttt{ZMod p}. The lemma on the power series $(1 - e^{-fX})/(fX)$, and the claim that the valuation is positive when $n$ is not a power of $p$, are not formalised. The axioms used are the three usual axioms of mathlib.

\litterature{The valuation is Legendre's formula, which mathlib has in general (\texttt{sub\_one\_mul\_padicValNat\_factorial}). The residue of the part of $n!$ prime to $p$ is classical (generalisations of Wilson's theorem; see for instance Granville, \q{Binomial coefficients modulo prime powers}, 1997, as recalled). The statement is formalised because the lemma of the page rests on it. The edition does not present it as new.}

\section[Notebook 1970-5: the index of $N_n$ in $O_n(\mathbf F_q)$ is odd]{Notebook 1970-5: the index of $N_n(\mathbf F_q)$ in $O_n(\mathbf F_q)$ is odd when $4 \mid q - 1$\protect\verifie{Notebook1970_5}}\label{sec:quillen-1970-5}

\emph{Context.} On pages~65 to~68 of \carnet{1970}{5} Quillen studies the mod~$2$ cohomology of the orthogonal groups $O_n(\mathbf F_q)$, $q$ odd, through the subgroups $N_n(\mathbf F_q) = (\Sigma_m \wr O_2(\mathbf F_q)) \times (\mathbf Z/2)^\varepsilon$, $n = 2m + \varepsilon$.

\emph{The page.} On page~67 he writes the quotients of the orders,
\[
\frac{|O_{2m+1}|}{|N_{2m+1}|} = q^{m^2}\prod_{j=1}^m \frac{q^{2j} - 1}{2j(q - 1)}, \qquad \frac{|O_{2m}|}{|N_{2m}|} = q^{m(m-1)}\,\frac{q^m - 1}{m(q-1)}\prod_{j=1}^{m-1}\frac{q^{2j}-1}{2j(q-1)},
\]
says that they are prime to~$2$ when $4 \mid q - 1$, and concludes that $N_n(\mathbf F_q)$ then contains a Sylow $2$-subgroup of $O_n(\mathbf F_q)$, a result he attributes to Chevalley.

\begin{encadre}
\begin{theoreme}[Quillen, notebook 1970-5, p.~67]\label{thm:quillen-1970-5}
Let $q \equiv 1 \pmod 4$, $q > 1$. Then $v_2(q^k - 1) = v_2(q - 1) + v_2(k)$ for $k \ge 1$; the integers $\prod_{j=1}^m (q^{2j} - 1)$ and $2^m\,m!\,(q-1)^m$ have the same $2$-adic valuation; and, for $m \ge 1$, so do $(q^m - 1)\prod_{j=1}^{m-1}(q^{2j}-1)$ and $m(q-1)\cdot 2^{m-1}(m-1)!\,(q-1)^{m-1}$. For $q = 3$, $m = 1$ the first equality fails.
\end{theoreme}
\end{encadre}

\begin{proof}
The first formula is the $2$-adic lifting of the exponent, valid when $4 \mid q - 1$. It gives $v_2(q^{2j} - 1) = v_2(q-1) + 1 + v_2(j) = v_2(2j(q-1))$ for each factor, and $v_2(q^m - 1) = v_2(m(q-1))$; summing, and using $\sum_{j \le m} v_2(j) = v_2(m!)$, gives both equalities. For $q = 3$, $v_2(3^2 - 1) = 3$ while $v_2(2 \cdot 1 \cdot 2) = 2$.
\end{proof}

\emph{Verification in Lean.} The file \href{\leanblobmain/Notebook1970_5.lean}{\texttt{Notebook1970\_5.\allowbreak{}lean}} proves the theorem as \texttt{padicValNat\_pow\_sub\_one} (from mathlib's \texttt{Int.two\_pow\_sub\_pow'}), \texttt{padicValNat\_factor}, \texttt{padicValNat\_odd\_index}, \texttt{padicValNat\_even\_index} and \texttt{not\_padicValNat\_odd\_index\_three}. It compares valuations of numerator and denominator, and does not claim that the quotients are integers, which is the group theory of the page. The orders of the orthogonal groups and the Sylow statement are not formalised; the powers of $q$ are odd. The axioms used are the three usual axioms of mathlib.

\litterature{The arithmetic is the lifting of the exponent lemma; the Sylow statement is in Carter and Fong, \emph{The Sylow $2$-subgroups of the finite classical groups} (J.\ Algebra~1, 1964), and the page refers to Chevalley. The edition does not present it as new.}

\section[Notebook 1970-11: the 2-parts of $|GL_n(\mathbf F_q)|$ and $|N|$]{Notebook 1970-11: the 2-parts of $|GL_n(\mathbf F_q)|$ and of $|N|$ for $q \equiv 3 \bmod 4$\protect\verifie{Notebook1970_11}}\label{sec:quillen-1970-11}

\emph{Context.} In \carnet{1970}{11} Quillen computes the mod~$\ell$ cohomology of $GL_n(\mathbf F_q)$ through subgroups of index prime to $\ell$. Page~17 takes up the exceptional case $\ell = 2$, $q \equiv 3 \pmod 4$, with the subgroup $N = (\Sigma_m \wr (\mathrm{Gal} \ltimes \mathbf F_{q^2}^\times)) \times (\mathbf F_q^\times)^e$ of $GL_n(\mathbf F_q)$, $n = 2m + e$, $e \in \{0, 1\}$.

\emph{The page.} He writes $|GL_n(\mathbf F_q)| = q^{n(n-1)/2}\prod_{i=1}^n (q^i - 1)$ and $|N| = m!\,(2(q^2 - 1))^m (q - 1)^e$, and says that they have the same power of~$2$, since $v_2(q^i - 1) = 1$ for $i$ odd and $v_2(q^{2i} - 1) = v_2(i) + v_2(q^2 - 1)$.

\begin{encadre}
\begin{theoreme}[Quillen, notebook 1970-11, p.~17]\label{thm:quillen-1970-11}
Let $q \equiv 3 \pmod 4$. Then $v_2(q^k - 1) = 1$ for $k$ odd, and $v_2(q^{2i} - 1) = v_2(i) + v_2(q^2 - 1)$ for $i \ge 1$ (this for every odd $q > 1$). For $n = 2m + e$, $e \in \{0,1\}$, the integers $q^{n(n-1)/2}\prod_{i=1}^n (q^i - 1)$ and $m!\,(2(q^2-1))^m (q-1)^e$ have the same $2$-adic valuation.
\end{theoreme}
\end{encadre}

\begin{proof}
For $k$ odd, $q^k \equiv 3 \pmod 4$. For the even exponents, the $2$-adic lifting of the exponent gives $v_2(q^{2i} - 1) + 1 = v_2(q + 1) + v_2(q - 1) + v_2(2i)$, and $v_2(q+1) + v_2(q-1) = v_2(q^2 - 1)$. In $\prod_{i=1}^{n}$ there are $m + e$ odd exponents and the even exponents $2, 4, \ldots, 2m$, so its valuation is $m + e + v_2(m!) + m\,v_2(q^2 - 1)$; this is the valuation of $m!\,(2(q^2 - 1))^m(q-1)^e$, since $v_2(q - 1) = 1$. The power of $q$ is odd.
\end{proof}

\emph{Verification in Lean.} The file \href{\leanblobmain/Notebook1970_11.lean}{\texttt{Notebook1970\_11.\allowbreak{}lean}} proves the theorem as \texttt{padicValNat\_pow\_odd\_sub\_one}, \texttt{padicValNat\_pow\_even\_sub\_one} and \texttt{padicValNat\_card\_GL\_eq}, the products written over $i < n$ with exponent $i + 1$. The orders of $GL_n(\mathbf F_q)$ and of $N$, and the cohomology, are not formalised. The axioms used are the three usual axioms of mathlib.

\litterature{The valuations are the lifting of the exponent lemma, and the comparison is the standard count showing that $N$ has odd index (compare Quillen, \emph{On the cohomology and $K$-theory of the general linear groups over a finite field}, Ann.\ of Math.\ 96, 1972). The edition does not present it as new.}

\section[Notebooks 1970-11 and 1974-2: the digit-sum bound]{Notebooks 1970-11 and 1974-2: a nonzero multiple of $p^d - 1$ has digit sum at least $d(p-1)$\protect\verifie{Notebook1970_11}}\label{sec:quillen-1970-11-digits}

\emph{Context.} On pages~20 to~22 of \carnet{1970}{11} Quillen looks for the first invariant of $\mathbf F_q^\times$, $q = p^r$, in the cohomology of the additive group $\mathbf F_q$ with coefficients in $\bar{\mathbf F}_p$. A monomial has degree $\sum_a (2n_a - \varepsilon_a)$ and weight $\sum_a n_ap^a$, with $\varepsilon_a \in \{0, 1\}$ and $\varepsilon_a = 0$ when $n_a = 0$, and it is invariant when $p^r - 1$ divides the weight. Pages~21 and~22 are crossed out by a diagonal stroke and marked to be ignored; the statement they prove holds. In \carnet{1974}{2}, page~26, the same count appears for the characters of $\mathbf F_{p^d}^\times$ in $\Lambda^i \otimes \Gamma_j$: an invariant needs $a_0, \ldots, a_{d-1} \ge 0$ with $\sum a_i = i + j$ and $\sum a_ip^i \equiv 0 \pmod{p^d - 1}$, and Quillen notes that this forces $i + j \ge d(p-1)$.

\emph{The page.} The argument of 1970-11 is a cyclic carrying: if $n_a \ge p$, replace $n_a$ by $n_a - p$ and $n_{a+1}$ by $n_{a+1} + 1$, the index read modulo $r$; the weight is unchanged modulo $p^r - 1$ and the degree drops. At the end all $n_a \le p - 1$, the weight is then $p^r - 1$, so every $n_a = p - 1$, and the first invariant is in degree $r(2p - 3)$.

\begin{encadre}
\begin{theoreme}[Quillen, notebook 1970-11, pp.~21--22; notebook 1974-2, p.~26]\label{thm:quillen-1970-11-digits}
Let $p \ge 2$ and $d \ge 0$. \emph{(i)} If $a_0, \ldots, a_{d-1} \ge 0$ are not all zero and $p^d - 1$ divides $\sum a_ip^i$, then $\sum a_i \ge d(p - 1)$. \emph{(ii)} If moreover $\varepsilon_i \in \{0, 1\}$, with $\varepsilon_i = 0$ when $a_i = 0$, then $\sum (2a_i - \varepsilon_i) \ge d(2p - 3)$, with equality for $a_i = p - 1$, $\varepsilon_i = 1$.
\end{theoreme}
\end{encadre}

\begin{proof}
(i) By induction on $\sum a_i$. If some $a_i \ge p$, the carry described above keeps the family nonzero and $\sum a_ip^i$ modulo $p^d - 1$ (it changes by $p^{i+1} - p^{i+1}$, or by $1 - p^d$ when $i = d - 1$), and lowers $\sum a_i$ by $p - 1$; the claim for the new family gives it for the old. If all $a_i \le p - 1$, then $0 < \sum a_ip^i \le \sum (p-1)p^i = p^d - 1$, so $\sum a_ip^i = p^d - 1$; then $\sum (p - 1 - a_i)p^i = 0$, every $a_i = p - 1$ and $\sum a_i = d(p-1)$. (ii) $\sum(2a_i - \varepsilon_i) \ge 2\sum a_i - d \ge 2d(p-1) - d$.
\end{proof}

\emph{Verification in Lean.} The second part of the file \href{\leanblobmain/Notebook1970_11.lean}{\texttt{Notebook1970\_11.\allowbreak{}lean}} proves (i) as \texttt{le\_sum\_of\_dvd}, by the carrying of the page (\texttt{carry}, \texttt{val\_carry}, \texttt{sum\_carry}) and the base case \texttt{base}, with the digits indexed by \texttt{Fin d}, and (ii) as \texttt{le\_degree} and \texttt{degree\_extremal}. The cohomology and the characters are not formalised. The axioms used are the three usual axioms of mathlib.

\litterature{The bound is an elementary fact on base-$p$ digits: the least digit sum of a nonzero multiple of $p^d - 1$ is $d(p-1)$; the character count of 1974-2 is that of Quillen, \emph{On the cohomology and $K$-theory of the general linear groups over a finite field} (Ann.\ of Math.\ 96, 1972). The edition does not present it as new.}

\section[Notebook 1971-1: two binomial congruences]{Notebook 1971-1: two binomial congruences and a product in characteristic $p$\protect\verifie{Notebook1971_1}}\label{sec:quillen-1971-1}

\emph{Context.} Around the new year 1971, \carnet{1971}{1} computes the cohomology of a metacyclic $p$-group $\langle x, y\rangle$ with $xyx^{-1} = y^{1+p}$ by its Hochschild--Serre spectral sequence.

\emph{The page, p.~3.} To rule out a central extension, Quillen computes modulo $p^3$, for $p$ odd, $(1 + p + \beta p^2)^p \equiv 1 + p(p + \beta p^2) + p\frac{p-1}{2}(p + \beta p^2)^2 \equiv 1 + p^2$.

\emph{The page, pp.~6--7.} Restricting a representation to $\langle y\rangle$, he uses $(1 + p)^a \equiv 1 + ap \pmod{p^2}$, and then computes a total Chern class as $\prod_{a=0}^{p-1}(1 + (1 + ap)c) = (1 + c)^p = 1 + c^p$ in characteristic~$p$.

\begin{encadre}
\begin{theoreme}[Quillen, notebook 1971-1, pp.~3, 6--7]\label{thm:quillen-1971-1}
\emph{(i)} For every odd natural number $p$ and every integer $\beta$, $(1 + p + \beta p^2)^p \equiv 1 + p^2 \pmod{p^3}$; for $p = 2$ the left side is $\equiv 1 \pmod 8$ and the congruence fails. \emph{(ii)} $(1 + p)^a \equiv 1 + ap \pmod{p^2}$ for all natural numbers $p$, $a$. \emph{(iii)} In a commutative ring of prime characteristic $p$, $\prod_{a=0}^{p-1}(1 + (1 + ap)c) = 1 + c^p$.
\end{theoreme}
\end{encadre}

\begin{proof}
In $\mathbf Z$, $x^3$ divides $(1+x)^n - 1 - nx - \binom n2x^2$, by induction on $n$. With $x = p + \beta p^2$: $x^3$ and $\binom p2 x^2 = p\frac{p-1}{2}x^2$ are divisible by $p^3$ when $p$ is odd, and $px \equiv p^2$. For $p = 2$, $(3 + 4\beta)^2 = 8(1 + 3\beta + 2\beta^2) + 1$. (ii) is the same expansion modulo $x^2$. (iii) In characteristic $p$, $1 + ap = 1$ and $(1 + c)^p = 1 + c^p$.
\end{proof}

\emph{Verification in Lean.} The file \href{\leanblobmain/Notebook1971_1.lean}{\texttt{Notebook1971\_1.\allowbreak{}lean}} proves (i) as \texttt{pow\_congr\_odd} and \texttt{not\_pow\_congr\_two}, (ii) as \texttt{one\_add\_pow\_congr}, (iii) as \texttt{prod\_one\_add}, with the expansions \texttt{three\_terms} and \texttt{two\_terms}. The number $p$ of (i) need not be prime. The cohomology and the Chern classes are not formalised. The axioms used are the three usual axioms of mathlib.

\litterature{Elementary binomial congruences. The edition does not present them as new.}

\section[Notebook 1971-14: the functional equation of the zeta function of a curve]{Notebook 1971-14: $Z(z) = q^{g-1}z^{2g-2}Z(1/qz)$ from Riemann--Roch\protect\verifie{Notebook1971_14}}\label{sec:quillen-1971-14}

\emph{Context.} On pages~23--25 of \carnet{1971}{14} Quillen rewrites the zeta function of a curve $X$ of genus $g$ over $\mathbf F_q$ as a sum over $\mathrm{Pic}(X)$, splits it into a polynomial part (degrees $0$ to $2g - 2$) and a geometric tail, continues the tail analytically, and derives the functional equation from Riemann--Roch by $L \mapsto K - L$.

\begin{encadre}
\begin{theoreme}[After Quillen, notebook 1971-14, pp.~24--25]\label{thm:quillen-1971-14}
Let $K$ be a field, $q, z \in K$ with $q \ne 0$, $z \ne 0$, $z \ne 1$, $qz \ne 1$, and $g \ge 1$, $h \ge 0$ integers. Let $I$ be a finite set with maps $\deg, h^0 : I \to \mathbf N$ and an involution $\sigma$ such that $\deg i + \deg\sigma(i) = 2g - 2$, $h^0(i) = \deg i + 1 - g + h^0(\sigma(i))$, and exactly $h$ elements of each degree $0, \dots, 2g - 2$. Put
\[Z(w) = \sum_{i \in I}w^{\deg i}(q^{h^0(i)} - 1) + h\Bigl(\frac{q^gw^{2g-1}}{1 - qw} - \frac{w^{2g-1}}{1 - w}\Bigr).\]
Then $Z(z) = q^{g-1}z^{2g-2}Z(1/(qz))$.
\end{theoreme}
\end{encadre}

For a curve, $I$ is the set of classes of degree $0$ to $2g - 2$, $\sigma(L) = K - L$, $h = |\mathrm{Pic}^0(X)|$, and $Z$ is $(q - 1)$ times the zeta function in $z = q^{-s}$: the page's polynomial part and tail.

\begin{proof}
Since $\sum_{d < 2g - 1}w^d + w^{2g-1}/(1 - w) = 1/(1 - w)$, $Z(w) = B(w) + hS(w)$ with $B(w) = \sum_iw^{\deg i}q^{h^0(i)}$ and $S(w) = q^gw^{2g-1}/(1 - qw) - 1/(1 - w)$. Riemann--Roch and $\deg\sigma(i) = 2g - 2 - \deg i$ give, term by term after $i \mapsto \sigma(i)$, $q^{g-1}z^{2g-2}B(1/qz) = B(z)$; a direct computation gives $q^{g-1}z^{2g-2}S(1/qz) = S(z)$.
\end{proof}

\emph{Verification in Lean.} The file \href{\leanblobmain/Notebook1971_14.lean}{\texttt{Notebook1971\_14.\allowbreak{}lean}} defines \texttt{Z} and proves \texttt{functional\_equation}. The page continues an infinite sum over all classes analytically; the proof above uses only finite sums, the split into $B$ and $S$ being the edition's. The identification of $Z$ with the zeta function (counting effective divisors in a class) and Riemann--Roch are hypotheses, not formalised. The axioms used are the three usual axioms of mathlib.

\litterature{Classical: F.\,K.~Schmidt (1931); the proof through $\mathrm{Pic}$ and Riemann--Roch is the standard one (e.g.\ Rosen, \emph{Number theory in function fields}, ch.~5). The edition does not present it as new.}

\section[Notebook 1972-6: denominators of $B_{2s}/4s$]{Notebook 1972-6: the denominators of $B_{2s}/4s$ and the valuations of $p^t - 1$\protect\verifie{Notebook1972_6}}\label{sec:quillen-1972-6}

\emph{Context.} On pages~34 to~40 and~49 to~55 of \carnet{1972}{6} Quillen studies the image of the $J$-homomorphism in the $K$-theory of $\mathbf Z$. Its order in degree $4s - 1$ is the denominator $m(2s)$ of $B_{2s}/4s$, Adams' number, which he compares with the orders $p^{2s} - 1$ of $K_{4s-1}(\mathbf F_p)$.

\emph{The page.} On page~40, under a heading on number theory, he writes $m(2s) = \prod_\ell \ell^{m_\ell(2s)}$, with $m_\ell(t) = v_\ell(p^t - 1)$ for $\ell$ odd and $p$ a generator of $\mathbf Z_\ell^\times$, equal to $0$ if $\ell - 1 \nmid t$ and to $v_\ell(t) + 1$ if $\ell - 1 \mid t$, and $m_2(t) = v_2(3^t - 1)$, equal to $1$ for $t$ odd and to $v_2(t) + 2$ for $t$ even. He tabulates $m(2s)$ for $s = 1, \ldots, 5$: $24$, $240$, $504$, $480$, $264$. His Bernoulli numbers are the topologists' ones, $B_s$ for our $B_{2s}$, with $B_2 = 1/6$.

\begin{encadre}
\begin{theoreme}[Quillen, notebook 1972-6, p.~40]\label{thm:quillen-1972-6}
\emph{(i)} $v_2(3^t - 1) = 1$ for $t$ odd and $v_2(3^t - 1) = v_2(t) + 2$ for $t \ge 2$ even. \emph{(ii)} Let $\ell$ be an odd prime and $p \ge 2$ an integer whose class modulo $\ell$ has order $\ell - 1$ and with $p^{\ell - 1} \not\equiv 1 \pmod{\ell^2}$, that is, a generator of $(\mathbf Z/\ell^2)^\times$. For $t \ge 1$, $v_\ell(p^t - 1) = 0$ if $\ell - 1 \nmid t$ and $v_\ell(p^t - 1) = v_\ell(t) + 1$ if $\ell - 1 \mid t$. \emph{(iii)} The denominators of $B_{2s}/4s$ for $s = 1, \ldots, 5$ are $24$, $240$, $504$, $480$ and $264$.
\end{theoreme}
\end{encadre}

\begin{proof}
(i) For $t$ odd, $3^t \equiv 3 \pmod 4$, so $3^t - 1 \equiv 2 \pmod 4$. For $t$ even, the $2$-adic lifting of the exponent gives $v_2(3^t - 1) + 1 = v_2(4) + v_2(2) + v_2(t)$. (ii) If $\ell \mid p^t - 1$, the order $\ell - 1$ of $p$ modulo $\ell$ divides $t$. If $t = (\ell - 1)s$, put $x = p^{\ell - 1}$: then $v_\ell(x - 1) = 1$ by hypothesis, and the lifting of the exponent for odd $\ell$ gives $v_\ell(x^s - 1) = 1 + v_\ell(s) = 1 + v_\ell(t)$, since $\ell \nmid \ell - 1$. (iii) $B_2 = 1/6$, $B_4 = -1/30$, $B_6 = 1/42$, $B_8 = -1/30$, $B_{10} = 5/66$, computed from the recursion, and $B_{10}/20 = 1/264$.
\end{proof}

\emph{Verification in Lean.} The file \href{\leanblobmain/Notebook1972_6.lean}{\texttt{Notebook1972\_6.\allowbreak{}lean}} proves (i) as \texttt{padicValNat\_three\_pow\_sub\_one\_odd} and \texttt{padicValNat\_three\_pow\_sub\_one\_even}, (ii) as \texttt{padicValNat\_pow\_sub\_one\_of\_not\_dvd} and \texttt{padicValNat\_pow\_sub\_one\_of\_dvd}, from mathlib's lifting the exponent lemmas (\texttt{padicValNat.pow\_two\_sub\_one}, \texttt{padicValNat.pow\_sub\_pow}), and (iii) as \texttt{den\_bernoulli\_div}, with mathlib's \texttt{bernoulli}. The generator hypothesis is stated through its two consequences used in the proof. The identity $m(2s) = \prod_\ell \ell^{m_\ell(2s)}$ for every $s$ (von Staudt--Clausen and the numerators) is not formalised; for $s \le 5$ it can be checked on the table. The axioms used are the three usual axioms of mathlib.

\litterature{All of this is classical: Adams, \emph{On the groups $J(X)$}~IV (Topology~5, 1966), for $m(2s)$ and its $\ell$-adic valuations, and the lifting of the exponent. The edition does not present it as new.}

\section[Notebook 1972-8: $1 - \sigma$ on $k'^\times$; weights on $\Lambda^q$]{Notebook 1972-8: $z \mapsto \sigma(z)/z$ on $k'^\times$ has kernel $k^\times$ and cokernel of order $q - 1$; weights on $\Lambda^q$\protect\verifie{Notebook1972_8}}\label{sec:quillen-1972-8}

\emph{Context.} In \carnet{1972}{8} Quillen computes $K_*$ of a twisted polynomial ring over finite fields, where $1 - \sigma$ on $k'^\times$ appears (p.~20); and uses a weight argument on $\Lambda^q$ for the dual numbers (p.~30). The orders of the image of $J$ listed on p.~12 are those of section~\ref{sec:quillen-1972-6}.

\begin{encadre}
\begin{theoreme}[Quillen, notebook 1972-8, pp.~20, 30]\label{thm:quillen-1972-8}
\begin{enumerate}
\item Let $k$ be a subfield of a finite field $K$, $|k| = q$, and $\sigma(z) = z^q$. The homomorphism $z \mapsto \sigma(z)/z = z^{q-1}$ of $K^\times$ has kernel $k^\times$, of order $q - 1$, and its image has index $q - 1$.
\item For a module $M$ over a commutative ring and $\lambda$ a scalar, multiplication by $\lambda$ induces multiplication by $\lambda^q$ on $\Lambda^qM$.
\end{enumerate}
\end{theoreme}
\end{encadre}

\begin{proof}
(1) $K^\times$ is cyclic of order $q^d - 1$, divisible by $q - 1$; in a cyclic group of order $N$ the $e$-th power map with $e \mid N$ has kernel of order $e$ and image of index $e$. Elements of $k^\times$ satisfy $z^{q-1} = 1$, and both groups have $q - 1$ elements. (2) On generators, $\lambda v_1 \wedge \dots \wedge \lambda v_q = \lambda^q v_1 \wedge \dots \wedge v_q$.
\end{proof}

\emph{Verification in Lean.} The file \href{\leanblobmain/Notebook1972_8.lean}{\texttt{Notebook1972\_8.\allowbreak{}lean}} proves \texttt{oneSubFrob\_apply}, \texttt{card\_ker}, \texttt{index\_range}, \texttt{ker\_eq} and \texttt{exteriorPower\_map\_smul}; the cyclic-group counts are mathlib's. The $K$-theoretic statements these serve are not formalised. The axioms used are the three usual axioms of mathlib.

\litterature{Classical: the cokernel of $1 - \sigma$ in Quillen's computation of $K_*(\mathbf F_q)$ (Ann.\ of Math.\ 96, 1972). The edition does not present them as new.}

\section[{Notebook 1987-5: $\mathbf Z[1/p] + \mathbf Z_p = \mathbf Q_p$}]{Notebook 1987-5: $\mathbf Z[1/p] + \mathbf Z_p = \mathbf Q_p$ and $\mathbf Z[1/p] \cap \mathbf Z_p = \mathbf Z$\protect\verifie{Notebook1987_5}}\label{sec:quillen-1987-5}

\emph{Context.} On page~64 of \carnet{1987}{5} Quillen uses the decomposition of the $p$-adic numbers into $p$-adic integers and rationals with a power of $p$ as denominator, the two meeting in $\mathbf Z$ (the local form of $\mathbf Q + \prod_p\mathbf Z_p = \mathbf A_f$, $\mathbf Q \cap \prod_p\mathbf Z_p = \mathbf Z$).

\begin{encadre}
\begin{theoreme}[Quillen, notebook 1987-5, p.~64]\label{thm:quillen-1987-5}
Let $p$ be a prime. Every $x \in \mathbf Q_p$ is $m/p^k + z$ with $m \in \mathbf Z$, $k \in \mathbf N$, $z \in \mathbf Z_p$; and if $m/p^k \in \mathbf Z_p$ then $p^k \mid m$, so that $\mathbf Z[1/p] \cap \mathbf Z_p = \mathbf Z$.
\end{theoreme}
\end{encadre}

\begin{proof}
Choose $k$ with $p^kx \in \mathbf Z_p$ and an integer $m \equiv p^kx \pmod{p^k\mathbf Z_p}$; then $x - m/p^k \in \mathbf Z_p$. Conversely $\lvert m\rvert_p = \lvert m/p^k\rvert_p\,p^{-k} \le p^{-k}$.
\end{proof}

\emph{Verification in Lean.} The file \href{\leanblobmain/Notebook1987_5.lean}{\texttt{Notebook1987\_5.\allowbreak{}lean}} proves the theorem as \texttt{exists\_eq\_add}, \texttt{dvd\_of\_norm\_div\_le\_one} and \texttt{inter\_eq\_int}, with mathlib's \texttt{PadicInt.appr}. The adelic statement is not formalised. The axioms used are the three usual axioms of mathlib.

\litterature{Classical and elementary (e.g.\ the strong approximation for $\mathbf Q$, Cassels and Fröhlich, \emph{Algebraic number theory}, ch.~II). The edition does not present it as new.}


\part{Combinatorics, $q$-analogues and posets}
\section[Notebook 1970-14: the Euler characteristic of the Tits building]{Notebook 1970-14: the Euler characteristic of the Tits building of $GL_n(\mathbf F_q)$\protect\verifie{Notebook1970_14}}\label{sec:quillen-1970-14}

\emph{Context.} On page~4 of \carnet{1970}{14} Quillen computes the Euler characteristic of the poset of proper nonzero subspaces of $\mathbf F_q^n$ (the Tits building) by counting its simplices, the flags $0 < W_1 < \cdots < W_d < V$. A flag has a type, the composition $(s_0, \ldots, s_d)$ of $n$ given by the successive dimensions, and the flags of a given type are counted by the $q$-multinomial coefficient. The page writes the alternating sum, lists the values $q + 1$, $1 - q^3$, $1 + q^6$ for $n = 2, 3, 4$, and concludes that the building should be a bouquet of $q^{n(n-1)/2}$ spheres of dimension $n - 2$, which it then proves by induction.

\emph{Intuition.} Write $[m]! = \prod_{b=1}^m (q^b - 1)$ and let $F(n)$ be the sum over all compositions $c$ of $n$ of $(-1)^{\ell(c)}[n]!/\prod_a [c_a]!$. Splitting off the first part $s$ gives $F(n) = -\sum_{s=1}^n \binom{n}{s}_q F(n - s)$, and the $q$-binomial theorem at $t = -1$ shows that $(-1)^n q^{n(n-1)/2}$ satisfies the same recursion.

\begin{encadre}
\begin{theoreme}[Quillen, notebook 1970-14, p.~4]\label{thm:quillen-1970-14}
Let $q$ be an element of a field with $q^b \ne 1$ for $b \ge 1$, and $[m]! = \prod_{b=1}^m (q^b - 1)$. For $n \ge 1$,
\[
\sum_{d=1}^{n-1} (-1)^{d+1} \sum_{\substack{s_0, \ldots, s_d > 0\\ s_0 + \cdots + s_d = n}} \frac{[n]!}{\prod_{a=0}^d [s_a]!} = 1 + (-1)^n q^{n(n-1)/2}.
\]
\end{theoreme}
\end{encadre}

\begin{proof}
First, $\binom{n}{k}_q [k]!\,[n-k]! = [n]!$ for the Gaussian binomial defined by the $q$-Pascal rule, by induction on $n$: the two terms of the rule contribute $[n]!(q^{k+1} - 1)$ and $q^{k+1}[n]!(q^{n-k} - 1)$, whose sum is $[n+1]!$. In particular $\binom{n}{k}_q = \binom{n}{n-k}_q$. The compositions of $n + 1$ are the lists $(s + 1, c)$ with $c$ a composition of $n - s$, $0 \le s \le n$, so $F(n+1) = -\sum_{s=0}^{n}\binom{n+1}{s+1}_q F(n - s)$. By strong induction $F(m) = (-1)^m q^{\binom m2}$: the identity $\sum_{k=0}^{N}(-1)^k q^{\binom k2}\binom Nk_q = 0$ for $N \ge 1$, read backwards with the symmetry, isolates the term $k = 0$. Finally the left side of the statement is $F(n)$ minus the term of the one-part composition, which is $-1$.
\end{proof}

\emph{Verification in Lean.} The file \href{\leanblobmain/Notebook1970_14.lean}{\texttt{Notebook1970\_14.\allowbreak{}lean}} proves the theorem as \texttt{euler\_characteristic}, the sum running over mathlib's \texttt{Composition n} with at least two blocks, and the full sum as \texttt{sum\_compositions}; the product formula is \texttt{qBinomial\_mul\_qf}, in any commutative ring, for the Gaussian binomial of the library's \texttt{Upstream/QBinomial} (written for notebook 1970-15). The hypothesis on $q$ covers $\mathbf Q(q)$, hence the identity of rational functions of the page, and every prime power in $\mathbf Q$. The topology (the building, its order complex and the count of flags of a type) is not formalised. The axioms used are the three usual axioms of mathlib.

\litterature{Classical: the Solomon--Tits theorem (L.~Solomon, 1969) gives the homotopy type of the building, a bouquet of $q^{n(n-1)/2}$ spheres, and the Euler characteristic follows; the alternating sum of $q$-multinomials is the $q$-analogue of the count of the faces of a simplex. The edition does not present it as new.}

\section[Notebook 1970-15: the Möbius function of the lattice of subspaces]{Notebook 1970-15: the Möbius function of the lattice of subspaces\protect\verifie{Notebook1970_15}}\label{sec:quillen-1970-15}

\emph{Context.} In September 1970, \carnet{1970}{15} compares the cohomology of a group in a topos with that of its points. On pages~21--22, Quillen inverts Euler characteristics over the isotropy types of an action, with Möbius coefficients, and checks the formula for an elementary abelian group acting on a point. In rank~$3$, the check reduces to the identity $1 - (q^2 + q + 1) + q(q^2 + q + 1) - q^3 = 0$: the alternating sum, over the subspaces $W$ of a space of dimension $n$ over $\mathbf F_q$, of $(-1)^{\dim W} q^{\binom{\dim W}{2}}$ is zero for $n \ge 1$, since $(-1)^k q^{\binom{k}{2}}$ is the value of the Möbius function of the lattice of subspaces on an interval of rank~$k$. Figure~\ref{fig:quillen-1970-15} shows the case $n = 3$, $q = 2$.

\emph{Intuition.} The number of subspaces of dimension $k$ is the Gaussian binomial coefficient $\left[\begin{smallmatrix} n \\ k \end{smallmatrix}\right]_q$, and these coefficients satisfy Gauss's binomial theorem $\prod_{i < n}(1 + q^i t) = \sum_k q^{\binom{k}{2}} \left[\begin{smallmatrix} n \\ k \end{smallmatrix}\right]_q t^k$. For $t = -1$, the factor $i = 0$ of the product vanishes.


\begin{figure}[htbp]
\centering
\begin{tikzpicture}[x=1.05cm, y=1.25cm, every node/.style={font=\small}]
\foreach \i in {0,...,6} {
  \coordinate (d\i) at ({\i-3},1);
  \coordinate (p\i) at ({\i-3},2);
  \draw[black!45] (0,0) -- (d\i);
  \draw[black!45] (p\i) -- (0,3);
}
\foreach \i/\j in {0/1,0/3,0/5,1/0,1/3,1/4,2/2,2/3,2/6,3/0,3/1,3/2,4/1,4/4,4/6,5/0,5/5,5/6,6/2,6/4,6/5} { \draw[black!35] (d\i) -- (p\j); }
\fill (0,0) circle (2pt); \fill (0,3) circle (2pt);
\foreach \i in {0,...,6} { \fill (d\i) circle (2pt); \fill (p\i) circle (2pt); }
\node[below=2pt] at (0,0) {$0$};
\node[above=2pt] at (0,3) {$\mathbf F_2^3$};
\foreach \y/\t in {0/{1 \times 1}, 1/{7 \times (-1)}, 2/{7 \times 2}, 3/{1 \times (-8)}} {
  \node[anchor=west] at (4.1,\y) {$k = \y$: \quad $\t$};
}
\node[anchor=west, black!65] at (4.1,3.65) {dimension, number $\times$ value};
\end{tikzpicture}
\caption{The lattice of subspaces of $\mathbf F_2^3$: the zero subspace, seven lines, seven planes and the whole space; each plane contains three lines. On the right, for each dimension $k$, the number $\left[\begin{smallmatrix} 3 \\ k \end{smallmatrix}\right]_2$ of subspaces and the value $(-1)^k q^{\binom{k}{2}}$ of the Möbius function between $0$ and a subspace of dimension $k$, for $q = 2$. The sum $1 - 7 + 14 - 8$ is zero: this is the identity of page~22 for $q = 2$.}
\label{fig:quillen-1970-15}
\end{figure}

\begin{encadre}
\begin{theoreme}[Quillen, notebook 1970-15, pp.~21--22]\label{thm:quillen-1970-15}
Let $R$ be a commutative ring and $q \in R$. Define $\left[\begin{smallmatrix} n \\ k \end{smallmatrix}\right]_q$ by $\left[\begin{smallmatrix} n \\ 0 \end{smallmatrix}\right] = 1$, $\left[\begin{smallmatrix} 0 \\ k + 1 \end{smallmatrix}\right] = 0$ and $\left[\begin{smallmatrix} n + 1 \\ k + 1 \end{smallmatrix}\right] = \left[\begin{smallmatrix} n \\ k \end{smallmatrix}\right] + q^{k+1} \left[\begin{smallmatrix} n \\ k + 1 \end{smallmatrix}\right]$. Then, for all $t \in R$,
\[
\prod_{i = 0}^{n - 1} (1 + q^i t) = \sum_{k = 0}^{n} q^{\binom{k}{2}} \begin{bmatrix} n \\ k \end{bmatrix}_q t^k,
\qquad\text{and}\qquad
\sum_{k = 0}^{n} (-1)^k q^{\binom{k}{2}} \begin{bmatrix} n \\ k \end{bmatrix}_q = 0 \quad (n \ge 1).
\]
For $n = 3$, this is the identity of page~22.
\end{theoreme}
\end{encadre}

\begin{proof}
By induction on $n$, for all $t$ at once. Write $\prod_{i \le n}(1 + q^i t) = (1 + t) \prod_{i < n}(1 + q^i (qt))$ and apply the induction hypothesis to $qt$. The coefficient of $t^{k+1}$ in the product with $1 + t$ is $q^{\binom{k}{2} + k}\left[\begin{smallmatrix} n \\ k \end{smallmatrix}\right] + q^{\binom{k+1}{2} + k + 1}\left[\begin{smallmatrix} n \\ k + 1 \end{smallmatrix}\right]$, that is $q^{\binom{k+1}{2}}\left[\begin{smallmatrix} n + 1 \\ k + 1 \end{smallmatrix}\right]$ by Pascal's recurrence, since $\binom{k+1}{2} = \binom{k}{2} + k$ and $\left[\begin{smallmatrix} n \\ n + 1 \end{smallmatrix}\right] = 0$. For $t = -1$ and $n \ge 1$, the factor $1 + q^0 \cdot (-1)$ vanishes.
\end{proof}

\begin{sloppypar}
\noindent\emph{Verification in Lean.} The coefficients are \texttt{qBinom}, with \texttt{qBinom\_eq\_zero\_of\_lt}; the binomial theorem is \texttt{prod\_one\_add\_eq\_sum}, the alternating sum is \texttt{sum\_alternating\_eq\_zero}, and the check of the page is \texttt{quillen\_check\_rank\_three}, in \href{\leanblobmain/Notebook1970_15.lean}{\texttt{Notebook1970\_15.\allowbreak{}lean}}. mathlib has no Gaussian binomial coefficients; the file defines them.
\end{sloppypar}

Not formalized: the counting of subspaces of $\mathbf F_q^n$ by the $\left[\begin{smallmatrix} n \\ k \end{smallmatrix}\right]_q$, and Möbius inversion over isotropy types. The proposition of page~14 (a pro-$\ell$-group has cohomology of finite type if and only if it has an open subgroup of finite cohomology) would require the cohomology of profinite groups, absent from mathlib; it is not attempted.

\litterature{found in the literature: P.~Hall (1936); Rota, \q{On the foundations of combinatorial theory I} (1964); Stanley, \emph{Enumerative Combinatorics}~I, 3.10. The statement is classical; the edition does not present it as new.}

\section[Notebook 1971-15: the poset of infinite coinfinite subsets is coned off]{Notebook 1971-15: every finite family of infinite, coinfinite subsets of $\mathbf N$ is coned off\protect\verifie{Notebook1971_15}}\label{sec:quillen-1971-15}

\emph{Context.} On pages~3--4 of \carnet{1971}{15}, a fair copy proving that the group of all permutations of a countable set is acyclic, Quillen shows that the poset $X$ of subsets $x \subset \mathbf N$ with $x$ and $\mathbf N - x$ infinite, ordered by $x \le y$ iff $x = y$ or $x \subset y$ with $y - x$ infinite, is contractible: every finite subset of $X$ is coned off by a map $x \mapsto x \cup z'$.

\begin{encadre}
\begin{theoreme}[Quillen, notebook 1971-15, pp.~3--4]\label{thm:quillen-1971-15}
Let $F$ be a finite subset of $X$. There is $c \in X$ such that, for all $x, y \in F$: $x \cup c \in X$, $x \le x \cup c$, $c \le x \cup c$, and $x \le y$ implies $x \cup c \le y \cup c$.
\end{theoreme}
\end{encadre}

\begin{proof}
By induction on $F$ there is an infinite $z$ with, for each $x \in F$, $z \subset x$ or $z \cap x = \emptyset$ (keep whichever of $z \cap x$, $z - x$ is infinite). Split $z = c \sqcup c'$ into two infinite parts (even and odd terms of an enumeration). Then $c \in X$ ($c' \subset \mathbf N - c$). If $z \subset x$, $x \cup c = x$ and $x - c \supset c'$; if $z \cap x = \emptyset$, $(x \cup c) - x = c$, $(x \cup c) - c = x$, and $c' $ lies outside $x \cup c$. Monotonicity is checked in the same way, case by case.
\end{proof}

\emph{Verification in Lean.} The file \href{\leanblobmain/Notebook1971_15.lean}{\texttt{Notebook1971\_15.\allowbreak{}lean}} proves \texttt{exists\_infinite\_subset\_or\_disjoint}, \texttt{exists\_split} and \texttt{cone}. Contractibility of the nerve of $X$ (finite subcomplexes coned off), the spectral sequence of pp.~4--7 and the acyclicity of the permutation group are not formalised. The axioms used are the three usual axioms of mathlib.

\litterature{Posets of large subsets are used for acyclicity by de la Harpe and McDuff (Comment.\ Math.\ Helv.\ 58, 1983); this exact poset and coning argument were not located. The statement is elementary; the edition does not present it as new.}

\section[Notebook 1972-8: the $2 \times m$ chessboard complex]{Notebook 1972-8: the $2 \times m$ chessboard complex is disconnected for $m = 2$ and connected for $m \ge 3$\protect\verifie{Notebook1972_8Chess}}\label{sec:quillen-1972-8-chess}

\emph{Context.} On pages~46--47 of \carnet{1972}{8} Quillen asks for the homotopy type of the complex of partial injections $\sigma \hookrightarrow C$, $\sigma \subset \{1, \dots, n\}$, and treats $n = 2$: disconnected when $|C| = 2$, a bouquet of circles when $|C| > 2$.

\begin{encadre}
\begin{theoreme}[Quillen, notebook 1972-8, p.~47]\label{thm:quillen-1972-8-chess}
Let $G_m$ be the graph on $\{1, 2\} \times C$, $|C| = m$, in which $(i, c)$ and $(j, c')$ are joined iff $i \ne j$ and $c \ne c'$. It is $(m - 1)$-regular with $m(m - 1)$ edges; it is disconnected for $m = 2$ and connected for $m \ge 3$. Hence for $m \ge 3$ it has the homotopy type of a bouquet of $m^2 - 3m + 1$ circles; for $m = 3$, one circle.
\end{theoreme}
\end{encadre}

\begin{proof}
Each vertex has $m - 1$ neighbours, so $2|E| = 2m(m - 1)$. For $m = 2$, $i + c \bmod 2$ is constant along edges, so $(1, 1)$ and $(1, 2)$ are not joined by a path. For $m \ge 3$, two vertices $(1, c)$, $(1, c')$ are joined through $(2, d)$ with $d \ne c, c'$, and $(2, c)$ is joined to some $(1, c'')$. A connected graph is homotopy equivalent to a bouquet of $|E| - |V| + 1$ circles.
\end{proof}

\emph{Verification in Lean.} The file \href{\leanblobmain/Notebook1972_8Chess.lean}{\texttt{Notebook1972\_8Chess.\allowbreak{}lean}} defines \texttt{chess} and proves \texttt{degree\_chess}, \texttt{card\_edgeFinset\_chess}, \texttt{not\_connected\_chess\_two} and \texttt{connected\_chess}. The homotopy type of a graph is not formalised (only the counts and connectivity). The \q{$\vee^4 S^1$} written beside a $3 \times 3$ array on the page does not fit $n = 2$, $|C| = 3$, which gives one hexagon; the reading of the exponent is doubtful. The axioms used are the three usual axioms of mathlib.

\litterature{The chessboard complex: Björner, Lovász, Vrećica and Živaljević (J.\ London Math.\ Soc.\ 49, 1994), who show that the $m \times n$ complex is $(\nu - 2)$-connected, $\nu = \min(m, n, \lfloor(m + n + 1)/3\rfloor$). The edition does not present the case $n = 2$ as new.}

\section[Notebook 1973-5: a partial order is the intersection of its linear extensions]{Notebook 1973-5: a partial order is the intersection of its linear extensions, and matrices supported on it\protect\verifie{Notebook1973_5}}\label{sec:quillen-1973-5}

\emph{Context.} \carnet{1973}{5} (October 1973) works towards the $K$-spectrum of an exact category. On 8 and 9 October Quillen reads Wagoner's paper and compares the constructions of Wagoner and Volodin, both built from cosets of families of unipotent subgroups of $GL_n(A)$. Volodin's family consists of the finite intersections of the conjugates $\pi U \pi^{-1}$ of the upper unitriangular group by permutation matrices. On page~24 Quillen asserts that these are exactly the groups attached to the semi-standard flags of Wagoner, then marks the assertion \q{WRONG} with a counterexample on three letters. Pages~28--30 find the correct description.

\emph{The page.} On page~28, for a set $X$ with a partial order $\le$ and two elements with $y \not\le x$, Quillen adds to $\le$ all the pairs $(u, v)$ such that $u \le x$ and $y \le v$, and checks in four cases each that the result is again transitive and antisymmetric; every mixed case would force $y \le x$. He concludes that a maximal partial order is linear, that every partial order is refined by a linear order (Zorn), and that every partial order is the intersection of the linear orders refining it. On page~29 he applies this to matrices: for a partial order $R$ on $\{1, \ldots, n\}$, the matrices whose $(i, j)$ entry vanishes unless $i \mathrel{R} j$ form a subring $M_R$, the intersection of a family of relations gives the intersection of the subrings, and therefore every $M_R$ is a finite intersection of conjugates of the ring of upper triangular matrices by permutations. Page~30 counts the partial orders on two and three points and observes that the orders coming from flags are those in which incomparability is an equivalence relation, which explains the failure of page~24.

\emph{Intuition.} If $a \not\le b$, the enlargement of page~28 with $x = b$, $y = a$ forces $b$ below $a$; any linear order refining it puts $b$ strictly below $a$. So a pair that is not in the order is excluded by some linear extension.

\begin{encadre}
\begin{theoreme}[Quillen, notebook 1973-5, p.~28]\label{thm:quillen-1973-5}
Let $\le$ be a partial order on a set $X$.
\begin{enumerate}
\item If $y \not\le x$, the relation $u \preceq v \iff (u \le v \text{ or } (u \le x \text{ and } y \le v))$ is a partial order containing $\le$, with $x \preceq y$.
\item For all $a, b \in X$: $a \le b$ if and only if $a \le' b$ for every linear order $\le'$ containing $\le$.
\end{enumerate}
\end{theoreme}
\end{encadre}

\begin{proof}
(1) Reflexivity is clear. If $u \preceq v \preceq w$, either both relations come from $\le$, or exactly one is of the second kind, and then $u \le x$, $y \le w$ by transitivity; if both are of the second kind, $y \le v \le x$, which is excluded. If $u \preceq v \preceq u$, either $u \le v \le u$, or one of the relations is of the second kind and one finds again $y \le x$ by composing the other three inequalities.

(2) One direction is trivial. If $a \not\le b$, apply (1) with $x = b$, $y = a$ and extend $\preceq$ to a linear order $\le'$ (Szpilrajn). Then $b \le' a$; if also $a \le' b$, antisymmetry gives $a = b$, hence $a \le b$, a contradiction.
\end{proof}

\begin{encadre}
\begin{theoreme}[Quillen, notebook 1973-5, p.~29]\label{thm:quillen-1973-5-b}
Let $A$ be a ring, $n$ a finite set and $R$ a relation on $n$; write $M_R$ for the set of matrices $(a_{ij})$ over $A$ with $a_{ij} = 0$ unless $i \mathrel{R} j$.
\begin{enumerate}
\item If $R$ is reflexive and transitive, $M_R$ is a subring of $M_n(A)$.
\item For a family $(R_\nu)$, $M_{\bigcap_\nu R_\nu} = \bigcap_\nu M_{R_\nu}$.
\item If $R$ is a partial order, $M_R = \bigcap M_{\le'}$, the intersection over the linear orders $\le'$ containing $R$.
\end{enumerate}
\end{theoreme}
\end{encadre}

\begin{proof}
(1) The identity is supported on the diagonal, and sums and negatives preserve the support. If $i \not\mathrel{R} j$, each term $a_{ik} b_{kj}$ of $(ab)_{ij}$ vanishes: if $i \mathrel{R} k$ then $k \not\mathrel{R} j$ by transitivity, so $b_{kj} = 0$; otherwise $a_{ik} = 0$. (2) is immediate from the definition, and (3) follows from (2) and Theorem~\ref{thm:quillen-1973-5}.
\end{proof}

\noindent For a linear order on $\{1, \ldots, n\}$, $M_{\le'}$ is the conjugate of the upper triangular matrices by the permutation that sorts $\le'$; on the unipotent groups, (3) says that the group of matrices $I + N$ with $N$ supported off the diagonal on $R$ is a finite intersection of conjugates of the unitriangular group, which is the description of Volodin's family the page arrives at. This conjugation is not formalized.

\emph{Verification in Lean.} The file \href{\leanblobmain/Notebook1973_5.lean}{\texttt{Notebook1973\_5.\allowbreak{}lean}} defines the enlarged relation under the name \texttt{adjoin} and proves part~(1) of Theorem~\ref{thm:quillen-1973-5} as \texttt{isPartialOrder\_adjoin}, \texttt{le\_adjoin} and \texttt{adjoin\_self}; part~(2) is \texttt{rel\_iff\_forall\_linear}, through \texttt{exists\_linear\_not\_rel}, using mathlib's extension theorem \texttt{extend\_partialOrder}. In Theorem~\ref{thm:quillen-1973-5-b}, $M_R$ is \texttt{matSupp}, part~(1) is the subring \texttt{matSubring}, part~(2) \texttt{matSupp\_iInf} and part~(3) \texttt{matSupp\_eq\_iInter}. The axioms used are the three usual axioms of mathlib.

\litterature{the extension of a partial order to a linear order is E.~Szpilrajn, \q{Sur l'extension de l'ordre partiel}, \emph{Fund.\ Math.} 16 (1930), 386--389, and the description of a partial order as the intersection of its linear extensions is the starting point of the order dimension of B.~Dushnik and E.~W.~Miller, \q{Partially ordered sets}, \emph{Amer.\ J.\ Math.} 63 (1941), 600--610. Rings of matrices supported on a preorder are the structural matrix rings, or finite incidence algebras. The statements are classical; mathlib has the extension theorem but not the intersection. The notebook uses them to correct its own assertion on Volodin's subgroups, which A.~Suslin's comparison of the $K$-theories (\emph{Comm.\ Algebra} 9, 1981) handles through the same family.}

\section[Notebook 1975-21: Birkhoff duality on morphisms]{Notebook 1975-21: Birkhoff duality on morphisms\protect\verifie{Notebook1975_21}}\label{sec:quillen-1975-21}

\emph{Context.} On pages~37 to~50 of \carnet{1975}{21} Quillen describes filtrations of a vector space indexed by a finite poset $J$ as lattice homomorphisms from the lattice $\mathrm{Cl}(J)$ of closed (lower) subsets of $J$ to the lattice of subspaces. On pages~41 to~48 he works out « some sorites » on finite distributive lattices: every such lattice is $\mathrm{Cl}$ of its poset of irreducibles, and a morphism $u$ between two of them comes from a map between the posets of irreducibles in the other direction, constructed through the left adjoint $v(y) = \inf\{x \mid y \le u(x)\}$.

\emph{Intuition.} For $b \in \beta$, the lower sets $Z$ of $\alpha$ with $b \in u(Z)$ are closed under finite intersections, since $u$ preserves them; their intersection is a principal set ${\downarrow}f(b)$, and $u(Z) = f^{-1}(Z)$.

\begin{encadre}
\begin{theoreme}[After Quillen, notebook 1975-21, pp.~41--48]\label{thm:quillen-1975-21}
Let $\alpha$ be a finite partial order and $\beta$ a preorder. Every map $u$ from the lower sets of $\alpha$ to the lower sets of $\beta$ preserving $\emptyset$, the whole set, unions and intersections is $Z \mapsto f^{-1}(Z)$ for a unique monotone $f \colon \beta \to \alpha$.
\end{theoreme}
\end{encadre}

\begin{proof}
Every lower set $Z$ is the finite union of the ${\downarrow}a$, $a \in Z$, so $b \in u(Z)$ if and only if $b \in u({\downarrow}a)$ for some $a \in Z$; applied to $Z = \alpha$, the set $S_b = \{a \mid b \in u({\downarrow}a)\}$ is not empty. If $a, a' \in S_b$, then $b \in u({\downarrow}a \cap {\downarrow}a')$, so some $c \le a, a'$ lies in $S_b$; a minimal element $f(b)$ of the finite set $S_b$ is therefore its least element. Then $b \in u(Z)$ if and only if $f(b) \in Z$. Monotonicity: if $b' \le b$, then $b' \in u({\downarrow}f(b))$, a lower set, so $f(b') \le f(b)$. Uniqueness: if $f^{-1} = g^{-1}$ on lower sets, $b \in g^{-1}({\downarrow}g(b)) = f^{-1}({\downarrow}g(b))$ gives $f(b) \le g(b)$, and symmetrically.
\end{proof}

\emph{Verification in Lean.} The file \href{\leanblobmain/Notebook1975_21.lean}{\texttt{Notebook1975\_21.\allowbreak{}lean}} defines the preimage as a bounded lattice homomorphism (\texttt{comap}) and proves the theorem as \texttt{exists\_unique\_comap}, with \texttt{comap\_injective}; the map $f$ is \texttt{dual u}, the least element of $S_b$. Only $\alpha$ is finite. Mathlib has Birkhoff's representation on objects (\texttt{Mathlib/Order/Birkhoff.lean}, with « TODO: extend to morphisms »); the categorical equivalence between finite posets and finite distributive lattices is not formalised. The axioms used are the three usual axioms of mathlib.

\litterature{Classical (G.~Birkhoff, 1937; the duality between finite ordered sets and finite distributive lattices, on objects and morphisms, is in Davey and Priestley, \emph{Introduction to Lattices and Order}, ch.~5). A gap in mathlib, and a possible contribution to it. The edition does not present it as new.}

\section[Notebook 1975-21: splitting a filtration]{Notebook 1975-21: splitting a filtration indexed by a finite poset\protect\verifie{Notebook1975_21Filtration}}\label{sec:quillen-1975-21-filtration}

\emph{Context.} On pages~37--38 of \carnet{1975}{21} Quillen splits a lattice homomorphism from the order ideals of a finite poset to the subspaces of a vector space by choosing complements; Birkhoff duality on morphisms, from the same notebook, is \S\ref{sec:quillen-1975-21}.

\begin{encadre}
\begin{theoreme}[Quillen, notebook 1975-21, pp.~37--38]\label{thm:quillen-1975-21-filtration}
Let $J$ be a finite poset, $L$ a modular complete lattice (the submodules of a module) and $W$ a lattice homomorphism from the order ideals of $J$ to $L$ with $W(\varnothing) = 0$. For each $j$, let $V_j$ satisfy $V_j + W(\{<j\}) = W(\{\le j\})$ and $V_j \cap W(\{<j\}) = 0$. Then $W(Z) = \sum_{j \in Z}V_j$ for every order ideal $Z$, and the $V_j$ are independent.
\end{theoreme}
\end{encadre}

\begin{proof}
By induction on $Z$: for $j$ maximal in $Z$, $Z' = Z \setminus \{j\}$ is an order ideal, $Z = Z' \cup \{\le j\}$ and $\{<j\} \subset Z'$, so $W(Z) = W(Z') + V_j$. Moreover $\{\le j\} \cap Z' \subset \{<j\}$, so $V_j \cap W(Z') = V_j \cap W(\{\le j\}) \cap W(Z') \subset V_j \cap W(\{<j\}) = 0$; in a modular lattice, this disjointness and the independence over $Z'$ give the independence over $Z$.
\end{proof}

\emph{Verification in Lean.} The file \href{\leanblobmain/Notebook1975_21Filtration.lean}{\texttt{Notebook1975\_21Filtration.\allowbreak{}lean}} proves \texttt{iSup\_eq} and \texttt{iSupIndep\_V} for any modular complete lattice, and \texttt{submodule}, the page's case. The equality $W(Z) = \sum V_j$ uses only the joins; the independence uses the meets. The existence of the complements over a field is mathlib's. The axioms used are the three usual axioms of mathlib.

\litterature{Elementary (distributive filtrations split); the edition does not present it as new.}

\section[Notebook 1986-7: the strict-sum lemma]{Notebook 1986-7: the strict-sum lemma\protect\verifie{Notebook1986_7Strict}}\label{sec:quillen-1986-7-strict}

\emph{Context.} On page~56 of \carnet{1986}{7}, on the way to the Durfee and Jacobi identities, Quillen sums $q^{i_1 + \dots + i_s}$ over the strictly increasing sequences of positive integers.

\begin{encadre}
\begin{theoreme}[Quillen, notebook 1986-7, p.~56]\label{thm:quillen-1986-7-strict}
Let $D_s \in R[[q]]$, over any commutative ring, have as $N$-th coefficient the number of $s$-element sets of positive integers with sum $N$, so that $D_s = \sum_{0<i_1<\dots<i_s}q^{i_1+\dots+i_s}$. Then $(1 - q^{s+1})D_{s+1} = q^{s+1}D_s$, $D_0 = 1$, and
\[
(1 - q)(1 - q^2)\cdots(1 - q^s)\,D_s = q^{s(s+1)/2},
\]
the product being a unit: $D_s = q^{s(s+1)/2}/\prod_{i=1}^s(1 - q^i)$.
\end{theoreme}
\end{encadre}

\begin{proof}
Let $E_s$ count the $s$-element sets of naturals ($0$ allowed). Removing $0$ gives $E_{s+1} = D_{s+1} + D_s$ coefficientwise; adding $1$ to every element gives $D_s(N + s) = E_s(N)$, and $D_s(N) = 0$ for $N < s$. So $D_{s+1} = q^{s+1}(D_{s+1} + D_s)$; induct on $s$.
\end{proof}

\emph{Verification in Lean.} The file \href{\leanblobmain/Notebook1986_7Strict.lean}{\texttt{Notebook1986\_7Strict.\allowbreak{}lean}} proves \texttt{count\_succ}, \texttt{count\_shift}, \texttt{one\_sub\_mul\_D}, \texttt{prod\_mul\_D} and \texttt{isUnit\_prod}; the sets are \texttt{Finset}s of $\{0, \dots, N\}$. The Durfee identity of pp.~55--56 is the object of section~\ref{sec:quillen-1986-7-durfee}. The axioms used are the three usual axioms of mathlib.

\litterature{Classical (partitions into $s$ distinct parts; Euler); the edition does not present it as new.}

\section[Notebook 1986-7: the Durfee identity]{Notebook 1986-7: $\prod_{i>0}(1 - q^i)^{-1} = \sum_s q^{s^2}/(q)_s^2$\protect\verifie{Notebook1986_7Durfee}}\label{sec:quillen-1986-7-durfee}

\emph{Context.} On pages~55--56 of \carnet{1986}{7}, after Quillen's notes of a conversation with Graeme Segal, the Poincaré series of a conjectured decomposition of $H^*(BU)$ leads to the identity of the Durfee square, of which he records two proofs.

\begin{encadre}
\begin{theoreme}[after Quillen, notebook 1986-7, pp.~55--56]\label{thm:quillen-1986-7-durfee}
Let $R$ be a commutative ring, $[n, s]_q$ the Gaussian binomial coefficients and $(q)_s = (1 - q)\cdots(1 - q^s)$. For all $m, n$,
\[
\sum_{s} q^{s^2}[m, s]_q\,[n, s]_q = [m + n, m]_q
\]
in $R[q]$. Consequently, in $R[[q]]$, for $N \le M$ the $N$-th coefficients of $1/(q)_M$ and of $\sum_{s\le M} q^{s^2}/(q)_s^2$ agree, and the number of partitions of $N$ is the $N$-th coefficient of $\sum_{s \le N} q^{s^2}/(q)_s^2$: that is, $\prod_{i>0}(1 - q^i)^{-1} = \sum_{s \ge 0} q^{s^2}/(q)_s^2$.
\end{theoreme}
\end{encadre}

\begin{proof}
The finite identity (the Durfee square of a partition in an $m \times n$ box) follows by induction on $m + n$ from the two $q$-Pascal recursions: both sides satisfy $F(m + 1, n + 1) = F(m, n + 1) + q^{m+1}F(m + 1, n)$, the left one after a telescoping of the sum. From $[a + b, b]_q\,(q)_a(q)_b = (q)_{a+b}$, the power series $[a + s, s]_q$ agrees with $1/(q)_s$ below degree $a + 1$; take $m = n = M$. Finally $1/(q)_M$ is the generating function of the partitions with parts at most $M$, which below degree $M + 1$ are all partitions.
\end{proof}

\emph{Verification in Lean.} The file \href{\leanblobmain/Notebook1986_7Durfee.lean}{\texttt{Notebook1986\_7Durfee.\allowbreak{}lean}} proves \texttt{qBinomial\_pascal'}, \texttt{durfee\_succ\_succ}, \texttt{sum\_durfee}, \texttt{qBinomial\_mul\_qfac}, \texttt{coeff\_inverse\_qfac} and \texttt{card\_partition}, with the $q$-binomials of the project's \texttt{Upstream/QBinomial} and mathlib's generating function of restricted partitions. The infinite product is handled degree by degree; neither the Jacobi triple product proof nor the cohomological conjecture is formalised. The axioms used are the three usual axioms of mathlib.

\litterature{Classical (the Durfee square; Hardy and Wright, ch.~XIX; the finite form is a case of the $q$-Vandermonde identity); the edition does not present it as new.}

\section[Notebook 1987-14: cyclic embeddings]{Notebook 1987-14: $k(k - 1)$ cyclic embeddings of $\{1, \dots, k - 1\}$ in $\{1, \dots, k\}$\protect\verifie{Notebook1987_14Cyclic}}\label{sec:quillen-1987-14-cyclic}

\emph{Context.} On page~45 of \carnet{1987}{14}, to see that a natural map between cyclic tensor powers is unique up to a scalar, Quillen notes that the cyclic-order-preserving embeddings of a cyclic set with $k - 1$ elements into one with $k$ elements form a torsor under $\mathbf Z/k \times \mathbf Z/(k - 1)$.

\begin{encadre}
\begin{theoreme}[Quillen, notebook 1987-14, p.~45]\label{thm:quillen-1987-14-cyclic}
Let $k \ge 2$, and give $\mathbf Z/(k-1)$ and $\mathbf Z/k$ their cyclic orders. Every injective map $f : \mathbf Z/(k-1) \to \mathbf Z/k$ preserving the cyclic order is $f(x) = \iota(x + s) + r$ for a unique pair $(r, s) \in \mathbf Z/k \times \mathbf Z/(k-1)$, where $\iota$ is the inclusion of $\{0, \dots, k-2\}$; every such map preserves the cyclic order. Hence there are $k(k - 1)$ of them.
\end{theoreme}
\end{encadre}

\begin{proof}
Rotations and $\iota$ preserve the cyclic order. Given $f$, it misses one value $p$; after the rotation by $r = p - (k - 1)$ it lands in $\{0, \dots, k - 2\}$ and becomes a cyclic bijection $h$ of $\mathbf Z/(k-1)$. Such an $h$ is a rotation: if $h(x + 1) \ne h(x) + 1$, the preimage $y$ of $h(x) + 1$ lies cyclically after $x + 1$, and its image would lie strictly between $h(x)$ and $h(x) + 1$. The pair $(r, s)$ is recovered from the missing value and from $f(0)$.
\end{proof}

\emph{Verification in Lean.} The file \href{\leanblobmain/Notebook1987_14Cyclic.lean}{\texttt{Notebook1987\_14Cyclic.\allowbreak{}lean}} proves \texttt{isCyclic\_emb}, \texttt{rotation\_of\_cyclic}, \texttt{emb\_injective}, \texttt{exists\_emb} and \texttt{card\_isCyclic}, with $k = n + 2$, maps $\mathrm{Fin}(n+1) \to \mathrm{Fin}(n+2)$ and mathlib's circular order on $\mathrm{Fin}$. For $k = 1$ the count is $1$ (the empty map), not $k(k-1)$. The use of the torsor for maps of derived tensor powers is not formalised. The axioms used are the three usual axioms of mathlib.

\litterature{Elementary (the findings call it so); the edition does not present it as new.}


\part{Polynomials, power series and symmetric functions}
\section[Notebook 1969-1: a Lagrange identity for push-forwards]{Notebook 1969-1: a Lagrange interpolation identity for push-forwards\protect\verifie{Notebook1969_1}}\label{sec:quillen-1969-1}

\emph{Context.} \carnet{1969}{1}, of January 1969, works on complex cobordism: projective bundles, Riemann--Roch, and the formal group law $F$ of a complex-oriented cohomology theory. On pages~49 to~54 Quillen computes the push-forward $f_*\xi^q$ along the projective bundle $f : \mathbf P E \to X$ of a sum $E = L_1 + \dots + L_n$ of line bundles, where $\xi$ is the first Chern class of $\mathcal O(1)$. With $X_i$ the Chern class of the dual of $L_i$ and $I$ the inverse of the formal group law, he finds (p.~52) that $f_*\xi^q$ is the sum over $i$ of $X_i^q / \prod_{j \ne i} F(X_i, IX_j)$.

\emph{The page.} On page~54 (his page~9) he takes the law $F(X, Y) = X + Y - aXY$, for which $F(X, IY) = (X - Y)/(1 - aY)$, and evaluates the sum. He writes it as a Lagrange interpolation of the polynomial $a^{n-1}Z^q$ at the point $Z = 1/a$, and concludes that $f_*\xi^q = a^{n-1-q}$ for $0 \le q \le n - 1$. For $a = 0$ this is ordinary cohomology and for $a = 1$ it is $K$-theory, and he notes that the old formulas are recovered.

\emph{Intuition.} A polynomial $P$ of degree less than $n$ is determined by its values at $n$ distinct points $X_1, \dots, X_n$:
\[
  P(Z) = \sum_{i=1}^n P(X_i) \prod_{j \ne i} \frac{Z - X_j}{X_i - X_j}.
\]
Take $P = Z^q$ and $Z = 1/a$, and multiply by $a^{n-1}$: each factor $1/a - X_j$ becomes $1 - aX_j$, and the left-hand side becomes $a^{n-1-q}$. The argument divides by $a$. For $a = 0$ the identity still holds, with the convention $0^0 = 1$: the sum of the $X_i^q / \prod_{j \ne i}(X_i - X_j)$ is the coefficient of $Z^{n-1}$ in the interpolation polynomial of $Z^q$, that is $1$ for $q = n - 1$ and $0$ for $q < n - 1$.

\begin{encadre}
\begin{theoreme}[Quillen, notebook 1969-1, p.~54]\label{thm:quillen-1969-1}
Let $K$ be a field, $X_1, \dots, X_n$ distinct elements of $K$, $a \in K$ and $0 \le q \le n - 1$. Then
\[
  \sum_{i=1}^n X_i^q \prod_{j \ne i} \frac{1 - aX_j}{X_i - X_j} = a^{\,n-1-q},
\]
where $a^0 = 1$ also for $a = 0$.
\end{theoreme}
\end{encadre}

\begin{proof}
Suppose $a \ne 0$. The polynomial $Z^q$ has degree $q < n$, so the interpolation formula applies; at $Z = a^{-1}$ it reads $a^{-q} = \sum_i X_i^q \prod_{j \ne i} (a^{-1} - X_j)/(X_i - X_j)$. For each $i$ the product has $n - 1$ factors and $(1 - aX_j) = a\,(a^{-1} - X_j)$, so the sum of the statement is $a^{n-1}$ times this one, that is $a^{n-1}a^{-q} = a^{n-1-q}$, since $q \le n - 1$. Suppose $a = 0$. The interpolation formula, read on the coefficient of $Z^{n-1}$, says that for a polynomial $P$ of degree less than $n$ this coefficient is $\sum_i P(X_i)/\prod_{j \ne i}(X_i - X_j)$, because each basis polynomial $\prod_{j \ne i}(Z - X_j)/(X_i - X_j)$ has degree $n - 1$ and leading coefficient $1/\prod_{j \ne i}(X_i - X_j)$. For $P = Z^q$ the coefficient is $1$ if $q = n - 1$ and $0$ otherwise, which is $0^{\,n-1-q}$.
\end{proof}

\noindent The case $a = 0$ is the one of ordinary cohomology. The page obtains it by putting $a = 0$ in the final formula; the interpolation at $1/a$ does not cover it, and the proof above treats it separately.

\emph{Verification in Lean.} The file \href{\leanblobmain/Notebook1969_1.lean}{\texttt{Notebook1969\_1.\allowbreak{}lean}} states the identity for a finite set \texttt{s} of indices and a map \texttt{v} injective on \texttt{s}, with $n$ the cardinality of \texttt{s}. Theorem~\ref{thm:quillen-1969-1} is \texttt{sum\_pow\_mul\_prod\_eq\_pow}. The case $a \ne 0$ uses \texttt{Lagrange.eq\_interpolate} of mathlib, evaluated at $a^{-1}$. The case $a = 0$ is \texttt{sum\_pow\_mul\_prod\_eq\_zero\_pow}, from \texttt{Lagrange.coeff\_eq\_sum}. The push-forward, the projective bundle and the formal group law are not formalised: only the algebraic identity is. The axioms used are the three usual axioms of mathlib.

\litterature{The push-forward formula is in Quillen's published papers on cobordism and formal groups (\emph{Bull.\ Amer.\ Math.\ Soc.}~75 (1969), in residue form, and \emph{Adv.\ Math.}~7 (1971), \S 1, as recalled), and the identity is a classical consequence of Lagrange interpolation. It is formalised here because the page rests on it. The edition does not present it as new.}

\section[Notebook 1970-1: the polynomial of a finite additive subgroup is additive]{Notebook 1970-1: the polynomial of a finite additive subgroup of a field is additive\protect\verifie{Notebook1970_1}}\label{sec:quillen-1970-1}

\emph{Context.} In \carnet{1970}{1} Quillen studies the mod~$2$ homology of the symmetric groups through the Dickson invariants. The Stiefel--Whitney classes $w_{2^a - 2^i}$ of the regular representation of $(\mathbf Z/2)^a$ are the coefficients of $\prod_{v \in V}(X - v)$, where $V$ is the dual of $(\mathbf Z/2)^a$ inside the cohomology.

\emph{The page.} On pages~11 and~15--16 he uses that this product is $\sum_i w_{2^a - 2^i}X^{2^i}$, an additive function of $X$, and defines with it the operation $\Gamma(e(L)) = \sum_i w_{2^a-2^i}\,e(L)^{2^i}$ on Euler classes of line bundles.

\begin{encadre}
\begin{theoreme}[After Quillen, notebook 1970-1, pp.~11, 15--16]\label{thm:quillen-1970-1}
Let $V$ be a finite additive subgroup of a field $K$ and $P_V = \prod_{v \in V}(X - v)$. Then $P_V$ is monic of degree $|V|$, $P_V(X + v) = P_V(X)$ for $v \in V$, and $P_V(x + y) = P_V(x) + P_V(y)$ for all $x, y \in K$.
\end{theoreme}
\end{encadre}

\begin{proof}
Translation by $v \in V$ permutes the factors. For fixed $x$, the polynomial $R(Y) = P_V(Y + x) - P_V(Y) - P_V(x)$ has degree less than $|V|$, since $P_V(Y + x)$ and $P_V(Y)$ are monic of the same degree, and it vanishes at every $v \in V$: $P_V(v + x) = P_V(x)$ and $P_V(v) = 0$. So $R = 0$.
\end{proof}

\emph{Verification in Lean.} The file \href{\leanblobmain/Notebook1970_1.lean}{\texttt{Notebook1970\_1.\allowbreak{}lean}} proves the theorem as \texttt{subgroupPoly\_monic}, \texttt{natDegree\_subgroupPoly}, \texttt{subgroupPoly\_comp\_X\_add\_C} and \texttt{subgroupPoly\_add}, for an \texttt{AddSubgroup} with a \texttt{Fintype} instance. No hypothesis on the characteristic is used. That the coefficients sit only in degrees $p^i$ (Ore's $p$-polynomials), and the cohomology, are not formalised. The axioms used are the three usual axioms of mathlib.

\litterature{Classical: Ore, \emph{On a special class of polynomials} (Trans.\ AMS 35, 1933), and Dickson's invariants; mathlib has the case $V = \mathbf F_q$ through $X^q - X$, not the general one. The edition does not present it as new.}

\section[Notebook 1972-9: additive polynomials]{Notebook 1972-9: an additive polynomial in characteristic $p$ has only monomials $X^{p^i}$\protect\verifie{Notebook1972_9}}\label{sec:quillen-1972-9}

\emph{Context.} On pages~66 and~67 of \carnet{1972}{9} Quillen considers additive polynomial maps between vector spaces in characteristic $p$, and says that they are sums of $p^i$-homogeneous pieces. The one-variable statement behind this is that a polynomial with $f(X + Y) = f(X) + f(Y)$ involves only the monomials $X^{p^i}$.

\begin{encadre}
\begin{theoreme}[After Quillen, notebook 1972-9, p.~66]\label{thm:quillen-1972-9}
Let $K$ be a field of prime characteristic $p$ and $f \in K[X]$ with $f(X + Y) = f(X) + f(Y)$ in $K[X, Y]$. If the coefficient $f_n$ of $X^n$ is nonzero, then $n$ is a power of $p$.
\end{theoreme}
\end{encadre}

\begin{proof}
The coefficient of $Y^a$ in $f(X + Y)$ is the Hasse derivative $D^{(a)}f$; the identity says that $D^{(a)}f$ is the constant $f_a$ for $a \ge 1$, and that $f_0 = 0$. Taking the coefficient of $X^{n-a}$ gives $\binom na f_n = 0$ for $0 < a < n$. If $n$ is not a power of $p$, some $\binom na$ with $0 < a < n$ is prime to $p$ (Lucas), so $f_n = 0$.
\end{proof}

\emph{Verification in Lean.} The file \href{\leanblobmain/Notebook1972_9.lean}{\texttt{Notebook1972\_9.\allowbreak{}lean}} proves the theorem as \texttt{eq\_pow\_of\_coeff\_ne\_zero}, with the identity written in $K[X][Y]$ as \texttt{IsAdditivePoly} and the steps \texttt{hasseDeriv\_eq\_C}, \texttt{coeff\_zero\_eq\_zero} and \texttt{choose\_mul\_coeff}; Lucas's theorem is mathlib's \texttt{Choose.eq\_pow\_multiplicity\_of\_choose\_modEq\_zero\_nat}. The several-variable form and the inverse up to Frobenius of the page are not formalised. The axioms used are the three usual axioms of mathlib.

\litterature{Classical: Ore's $p$-polynomials (Trans.\ AMS 35, 1933). The edition does not present it as new.}

\section[1972 mixed papers: square-zero symmetric functions; the cone product]{1972 mixed papers: $(\sum\varepsilon_i)^k = k!\,e_k$ when $\varepsilon_i^2 = 0$; the cone product; the sum map\protect\verifie{Notebook1972_Mixed}}\label{sec:quillen-1972-mixed}

\emph{Context.} In the \href{https://www.claymath.org/library/Quillen/Working_papers/quillen\%201972/1972-mixed\%20papers.pdf}{file 1972-mixed papers} Quillen tests the cohomology of symmetric products of $\mathbf{CP}^1$ with classes $\varepsilon_i$ of square zero (p.~40), composes homotopies into a cone product (p.~17; the same algebra in \carnet{1971}{14}, pp.~10--11), and counts the $\ell$-torsion of a Brauer group as the kernel of a sum map (p.~33).

\begin{encadre}
\begin{theoreme}[Quillen, 1972 mixed papers, pp.~17, 33, 40]\label{thm:quillen-1972-mixed}
\begin{enumerate}
\item In a commutative ring, if $\varepsilon_i^2 = 0$ for $i \in s$, then $(\sum_{i\in s}\varepsilon_i)^k = k!\,e_k(\varepsilon)$, the elementary symmetric function; in particular $(\sum\varepsilon_i)^2 = 2\sum_{i<j}\varepsilon_i\varepsilon_j$ and $(\sum_{i\le n}\varepsilon_i)^n = n!\prod\varepsilon_i$.
\item In a ring, if $dh + hd = u - v$, $dh' + h'd = u' - v'$ and $d$ commutes with $v$ and $u'$, then $d(hu' + vh') + (hu' + vh')d = uu' - vv'$.
\item The kernel of the sum map $K^s \to K$, $s \ge 1$, has dimension $s - 1$.
\end{enumerate}
\end{theoreme}
\end{encadre}

\begin{proof}
(1) By induction on $s$: $(\varepsilon_a + S)^{k+1} = S^{k+1} + (k + 1)\varepsilon_aS^k$ since $\varepsilon_a^2 = 0$, and $e_{k+1}$ of $s \cup \{a\}$ is $e_{k+1}(s) + \varepsilon_ae_k(s)$. (2) Expand: $[d, h]u' = (u - v)u'$, $v[d, h'] = v(u' - v')$. (3) The map is onto; rank--nullity.
\end{proof}

\emph{Verification in Lean.} The file \href{\leanblobmain/Notebook1972_Mixed.lean}{\texttt{Notebook1972\_Mixed.\allowbreak{}lean}} proves \texttt{sum\_pow\_of\_sq\_eq\_zero}, \texttt{sum\_pow\_card}, \texttt{cone\_product} (ungraded commutators: $[d, h] = dh + hd$ for $h$ odd is written out) and \texttt{finrank\_ker\_sum}. The cohomology of symmetric products and the Brauer group are not formalised. Of the page's other short statements, the group of power series with constant term $1$ is the kernel of a homomorphism and the lifting-the-exponent identity of p.~30 is mathlib's \texttt{padicValNat.pow\_sub\_pow}; they are not repeated. The axioms used are the three usual axioms of mathlib.

\litterature{Elementary or classical (Macdonald, Proc.\ Cambridge Philos.\ Soc.\ 58, 1962, for the symmetric products); the edition does not present them as new.}

\section[Notebook 1973-8: $t + t^2 + t^4 + \texorpdfstring{\cdots}{...}$ inverts $t + t^2$]{Notebook 1973-8: $\chi(t) = t + t^2 + t^4 + \cdots$ inverts $t \mapsto t + t^2$\protect\verifie{Notebook1973_8}}\label{sec:quillen-1973-8}

\emph{Context.} On pages~61 and~62 of \carnet{1973}{8} Quillen recalls Wu's determination of the Stiefel--Whitney classes of a manifold by the Steenrod squares, and rewrites it with the antipode $\chi(Sq)$ of the total square. On a class $t$ of degree one, $Sq(t) = t + t^2$, and the page gives $\chi(Sq)(t) = t + t^2 + t^4 + \cdots$, so that $Sq(\chi(Sq)\,t) = t$.

\emph{Intuition.} In characteristic $2$ squaring is additive, so the square of $\sum_k t^{2^k}$ is $\sum_k t^{2^{k+1}}$: adding the series to its square leaves only $t$.

\begin{encadre}
\begin{theoreme}[Quillen, notebook 1973-8, p.~62]\label{thm:quillen-1973-8}
In the ring of formal power series $R[[t]]$ over a commutative ring of characteristic~$2$, the series $\chi = \sum_{k \ge 0} t^{2^k}$ satisfies $\chi + \chi^2 = t$, and it is the only series without constant term with this property.
\end{theoreme}
\end{encadre}

\begin{proof}
In characteristic $2$, $f^2$ is obtained from $f$ by replacing $t$ by $t^2$ and squaring the coefficients; the coefficients of $\chi$ are $0$ and $1$. So the coefficient of $t^n$ in $\chi + \chi^2$ is $[n = 2^k] + [n = 2^{k+1}]$: it is $1$ for $n = 1$, $1 + 1 = 0$ for $n = 2^{k}$ with $k \ge 1$, and $0$ otherwise. For uniqueness, $(f - g)(1 + f + g) = (f + f^2) - (g + g^2)$, and $1 + f + g$ is invertible when $f, g$ have no constant term.
\end{proof}

\emph{Verification in Lean.} The file \href{\leanblobmain/Notebook1973_8.lean}{\texttt{Notebook1973\_8.\allowbreak{}lean}} proves the identity as \texttt{chi\_add\_chi\_sq}, from mathlib's Frobenius formula for power series (\texttt{PowerSeries.map\_frobenius\_expand}), and uniqueness as \texttt{eq\_chi\_of\_add\_sq} (\texttt{eq\_of\_add\_sq\_eq} in any commutative ring). The Steenrod algebra and Wu's formula are not formalised. The axioms used are the three usual axioms of mathlib.

\litterature{Classical (Milnor, \emph{The Steenrod algebra and its dual}, 1958, for the conjugation; the series is the compositional inverse of $t + t^2$ over $\mathbf F_2$). The edition does not present it as new.}

\section[Notebook 1974-3: the Bézout partner of small degree]{Notebook 1974-3: the unique Bézout partner of degree less than $\deg f$\protect\verifie{Notebook1974_3}}\label{sec:quillen-1974-3}

\emph{Context.} In February 1974, on pages~4 to~6 of \carnet{1974}{3}, Quillen proves that the ordered set of unimodular subspaces of $k[t]^2$ is contractible (a statement he notes is equivalent to Nagao's theorem), by retracting the subspaces whose lines have denominators of degree $n + 1$ onto those of degree at most $n$. A line other than the base line is spanned by $(f, g)$, with $f, g \in k[t]$ coprime and $f \ne 0$.

\emph{The page.} On page~5 he proves that there is a unique pair $(f', g')$ with $fg' - f'g = 1$ and $\deg f' < \deg f$: Bézout gives some pair, adding a multiple of $(f, g)$ reduces $\deg f'$, and a second pair would give $f \mid \bar f'$ with $\deg \bar f' < \deg f$. A clause \q{and in fact $\deg f' = n - 1$} is struck out; it is false in general.

\begin{encadre}
\begin{theoreme}[Quillen, notebook 1974-3, p.~5]\label{thm:quillen-1974-3}
Let $k$ be a field and $f, g \in k[t]$ coprime with $f \ne 0$. There is a unique pair $(f', g')$ with $fg' - f'g = 1$ and $\deg f' < \deg f$. For $f = t^2$, $g = 1$ this pair is $(-1, 0)$, so $\deg f' = 0 \ne \deg f - 1$.
\end{theoreme}
\end{encadre}

\begin{proof}
Write $af + bg = 1$ and $-b = qf + f'$ with $\deg f' < \deg f$; then $(f', a - qg)$ works. If $(f'_1, g'_1)$ also works, $f(g'_1 - g') = (f'_1 - f')g$, so $f$ divides $f'_1 - f'$, whose degree is less than $\deg f$; hence $f'_1 = f'$, and then $g'_1 = g'$ since $f \ne 0$.
\end{proof}

\emph{Verification in Lean.} The file \href{\leanblobmain/Notebook1974_3.lean}{\texttt{Notebook1974\_3.\allowbreak{}lean}} proves the theorem as \texttt{exists\_unique\_bezout}, with \texttt{IsCoprime} and \texttt{Polynomial.degree}, and the example as \texttt{not\_degree\_eq\_pred}. The poset and its contractibility are not formalised. The axioms used are the three usual axioms of mathlib.

\litterature{Elementary (the normalised Bézout identity; Nagao's theorem itself is in Serre, \emph{Arbres, amalgames, $SL_2$}, chapter~II). The edition does not present it as new.}

\section[Notebook 1982-8: the symmetrised cyclic products vanish]{Notebook 1982-8: the symmetrisation of $1/((z_1 - z_2)\cdots(z_n - z_1))$ vanishes for $n \ge 3$\protect\verifie{Notebook1982_8}}\label{sec:quillen-1982-8}

\emph{Context.} On page~63 of \carnet{1982}{8} Quillen computes $\log\langle 0|S|0\rangle = \log\det(1 - K) = -\sum_n \operatorname{tr}(K^n)/n$ for a fermion in an abelian background field, and needs that only $n = 2$ contributes. The trace of $K^n$ is an integral against the cyclic function $\varphi = 1/((z_1 - z_2)(z_2 - z_3)\cdots(z_n - z_1))$ with a symmetric numerator, so what matters is the symmetrisation of $\varphi$. The page observes that the reversal $z_i \mapsto z_{n+1-i}$ only gives the odd case and does not see at first why $n = 4$ vanishes; it then argues with the Vandermonde product: $\omega\varphi$ would be an antisymmetric polynomial of degree less than $\binom n2$, for $n \ge 3$. In the non-abelian case only the cyclic symmetry is left and no vanishing follows.

\emph{Intuition.} The proof here is a different one, by partial fractions. Inserting a point $y$ between $a$ and $b$ in a cyclic product multiplies it by $(a - b)/((a - y)(y - b)) = 1/(a - y) - 1/(b - y)$; summed over the gaps of a cycle, these differences telescope to $0$.

\begin{encadre}
\begin{theoreme}[Quillen, notebook 1982-8, p.~63]\label{thm:quillen-1982-8}
Let $z_1, \ldots, z_n$ be distinct elements of a field, $n \ge 3$. Then
\[
\sum_{\sigma \in \mathfrak S_n} \frac{1}{(z_{\sigma 1} - z_{\sigma 2})(z_{\sigma 2} - z_{\sigma 3})\cdots(z_{\sigma n} - z_{\sigma 1})} = 0.
\]
For $n = 2$ the sum is $-2/(z_1 - z_2)^2$.
\end{theoreme}
\end{encadre}

\begin{proof}
For a path $a, c_1, \ldots, c_k, b$ write $P(a; c; b) = 1/((a - c_1)\cdots(c_k - b))$. By induction on $k$, the sum of $P$ over the $k + 1$ ways of inserting $y$ into $c$ is $P(a; c; b)\,(1/(a - y) - 1/(b - y))$. Order the permutations by inserting $z_1$ into the orderings of the others: inserting it in front or at the end gives the same cycle, and the other insertions sum to $0$ by the identity with $a = b$; so the sum is the sum over the orderings with $z_1$ in front. Insert $z_2$ into the orderings of $z_3, \ldots, z_n$ (at least one point, as $n \ge 3$): the insertions run over all the gaps of the cycle starting at $z_1$, and the identity with $a = b = z_1$ gives $0$.
\end{proof}

\emph{Verification in Lean.} The file \href{\leanblobmain/Notebook1982_8.lean}{\texttt{Notebook1982\_8.\allowbreak{}lean}} proves the theorem as \texttt{sum\_permutations\_cyc}, the sum running over \texttt{List.permutations} of a list without repetition (its $n!$ orderings), and the case $n = 2$ as \texttt{sum\_permutations\_cyc\_two}; the telescoping identity is \texttt{sum\_insert\_path}. The page's Vandermonde argument is not formalised (mathlib has no « antisymmetric implies divisible by the Vandermonde product »), nor the traces of $K^n$. The axioms used are the three usual axioms of mathlib.

\litterature{Classical in substance: the vanishing of the connected current correlators of order $n \ge 3$ of a free fermion in two dimensions (the bosonization behind the page), and, in the same algebraic form, the « $U(1)$ decoupling » identity for Parke--Taylor factors (Kleiss and Kuijf, 1989), proved by the same partial fractions. The edition does not present it as new.}

\section[Notebook 1985-8: the series of $(1 - x)^{-1/2}$]{Notebook 1985-8: $\sum_k \binom{2k}{k}x^k/4^k = (1 - x)^{-1/2}$\protect\verifie{Notebook1985_8}}\label{sec:quillen-1985-8}

\emph{Context.} On page~37 of \carnet{1985}{8} Quillen uses the generating series $\sum_{k \ge 0}\binom{2k}{k}x^k/4^k = (1 - x)^{-1/2}$, with $\binom{2k}{k}/4^k = \prod_{j=1}^k(1 - 1/2j)$.

\begin{encadre}
\begin{theoreme}[Quillen, notebook 1985-8, p.~37]\label{thm:quillen-1985-8}
As formal power series over $\mathbf Q$, $\bigl(\sum_{k \ge 0}\binom{2k}{k}x^k/4^k\bigr)^2(1 - x) = 1$; equivalently $\sum_{i+j=n}\binom{2i}{i}\binom{2j}{j} = 4^n$ for every $n$. Moreover $\binom{2k}{k}/4^k = \prod_{j=1}^k(1 - 1/2j)$.
\end{theoreme}
\end{encadre}

\begin{proof}
Let $c_k = \binom{2k}{k}$ and $s_n = \sum_{i+j=n}c_ic_j$. By symmetry $2\sum_{i+j=n}ic_ic_j = ns_n$. Applied to $n + 1$, and with $(i+1)c_{i+1} = 2(2i+1)c_i$, it gives $(n+1)s_{n+1} = 2\sum_{i+j=n}2(2i+1)c_ic_j = 4(ns_n + s_n)$, so $s_{n+1} = 4s_n$ and $s_n = 4^n$. The coefficient of $x^n$ in the square of the series is $s_n/4^n = 1$. The product formula follows from $c_{k+1}/4c_k = (2k+1)/(2k+2)$.
\end{proof}

\emph{Verification in Lean.} The file \href{\leanblobmain/Notebook1985_8.lean}{\texttt{Notebook1985\_8.\allowbreak{}lean}} proves the theorem as \texttt{centralSeries\_sq\_mul}, with mathlib's \texttt{PowerSeries.invUnitsSub}, \texttt{sum\_centralBinom\_mul} and \texttt{centralBinom\_div\_four\_pow}. The analytic statement for $|x| < 1$ is not formalised. The axioms used are the three usual axioms of mathlib.

\litterature{Classical (Graham, Knuth and Patashnik, \emph{Concrete Mathematics}, (5.39) and the binomial series). Mathlib has the central binomial coefficients and bounds on them, not the convolution identity. The edition does not present it as new.}

\section[Notebook 1986-7: triangular idempotents; antisymmetrisation]{Notebook 1986-7: triangular idempotents are orthogonal; antisymmetrisation over the Vandermonde\protect\verifie{Notebook1986_7}}\label{sec:quillen-1986-7}

\emph{Context.} On page~50 of \carnet{1986}{7} Quillen has idempotents $e_0, \ldots, e_n$ summing to $1$ with $e_te_s = 0$ for $t < s$, wants $e_ie_j = 0$ for all $i \ne j$, and looks for a better argument than an induction. On page~25 he computes the Gysin map of the full flag bundle $(BU_1)^n \to BU_n$ by antisymmetrising monomials against the Vandermonde product.

\begin{encadre}
\begin{theoreme}[Quillen, notebook 1986-7, pp.~25, 50]\label{thm:quillen-1986-7}
\begin{enumerate}
\item In a ring, let $(e_i)$ be idempotents indexed by a finite linearly ordered set, with $\sum_i e_i = 1$ and $e_te_s = 0$ for $t < s$. Then $e_ie_j = 0$ for all $i \ne j$: the $e_i$ form a complete family of orthogonal idempotents.
\item Over a commutative ring, $\sum_\sigma \operatorname{sgn}(\sigma)\prod_i u_{\sigma i}^{a_i} = \det(u_i^{a_j})$; it vanishes when two exponents coincide, and for $a_i = i$ it equals $\prod_{i<j}(u_j - u_i)$.
\end{enumerate}
\end{theoreme}
\end{encadre}

\begin{proof}
(1) Multiplying $\sum_k e_k = 1$ on the right by $e_j$ and using $e_ke_j = 0$ for $k < j$ gives $\sum_{k > j} e_ke_j = 0$. For $j < i$ multiply on the left by $e_i$: the terms $k > i$ vanish, the term $k = i$ is $e_ie_j$, and the terms $j < k < i$ vanish by induction on $i$. (2) The first identity is the Leibniz formula; then a determinant with two equal columns vanishes, and the last case is Vandermonde's determinant.
\end{proof}

\emph{Verification in Lean.} The file \href{\leanblobmain/Notebook1986_7.lean}{\texttt{Notebook1986\_7.\allowbreak{}lean}} proves (1) as \texttt{sum\_gt\_mul\_eq\_zero}, \texttt{mul\_eq\_zero\_of\_gt} and \texttt{completeOrthogonalIdempotents} (mathlib's structure), and (2) as \texttt{antisymm\_eq\_det}, \texttt{antisymm\_eq\_zero} and \texttt{antisymm\_vandermonde}. The argument for (1) is still an induction. The Gysin map itself, and the divisibility of an antisymmetric polynomial by the Vandermonde product, are not formalised. The axioms used are the three usual axioms of mathlib.

\litterature{Elementary: (1) is folklore of idempotent decompositions; (2) is the classical alternant (Macdonald, \emph{Symmetric functions and Hall polynomials}, I, \S 3). The edition does not present them as new.}

\section[Notebook 1986-7: inverse Toeplitz determinants; Newton]{Notebook 1986-7: the inverse of a unipotent Toeplitz matrix; Newton's identity for the Chern character\protect\verifie{Notebook1986_7Series}}\label{sec:quillen-1986-7-series}

\emph{Context.} Computing degeneracy classes, Quillen inverts the total Chern class, $c(Q) = 1/c(S)$, through determinants of Toeplitz matrices (pages~28--29 of \carnet{1986}{7}); on page~45 he recovers the Chern character from the logarithmic derivative of $c_t$.

\begin{encadre}
\begin{theoreme}[Quillen, notebook 1986-7, pp.~28--29, 45]\label{thm:quillen-1986-7-series}
Let $R$ be a commutative ring.
\begin{enumerate}
\item If $(1 + c_1T + c_2T^2 + \cdots)(1 + b_1T + b_2T^2 + \cdots) = 1$ in $R[[T]]$, then $(-1)^nb_n$ is the determinant of the $n \times n$ matrix with entries $c_{j-i+1}$ for $i \le j + 1$ (with $c_0 = 1$) and $0$ below the subdiagonal: $c_1$ on the diagonal, $1$ under it.
\item For $c_t = \prod_j(1 + x_jt)$ in $R[[t]]$, $t\,\frac{d}{dt}c_t = c_t\sum_{d\ge1}(-1)^{d-1}p_dt^d$, with $p_d = \sum_j x_j^d$.
\end{enumerate}
\end{theoreme}
\end{encadre}

\begin{proof}
(1) The lower triangular Toeplitz matrices $T_c = (c_{i-j})$ and $T_b = (b_{i-j})$ of size $n + 1$ are inverse to each other, since their product is the Toeplitz matrix of the product series. As $\det T_c = 1$, $T_b$ is the adjugate of $T_c$, so $b_n$ is the cofactor of the entry $(0, n)$, that is $(-1)^n$ times the minor obtained by deleting the first row and the last column; this minor is the transpose of the matrix of the statement. (2) For one factor, $(1 + xt)\sum_{d\ge1}(-1)^{d-1}x^dt^d = xt$, which is $t$ times the derivative of $1 + xt$; the general case follows by the Leibniz rule, by induction on the number of factors.
\end{proof}

\emph{Verification in Lean.} The file \href{\leanblobmain/Notebook1986_7Series.lean}{\texttt{Notebook1986\_7Series.\allowbreak{}lean}} proves (1) as \texttt{neg\_one\_pow\_mul\_coeff\_eq\_det}, through \texttt{toeplitz\_mul\_toeplitz} and mathlib's \texttt{adjugate\_fin\_succ\_eq\_det\_submatrix}, and (2) as \texttt{X\_mul\_derivative\_prod}; no division by integers is needed in this form. The Durfee identity and the strict-sum lemma of pages~55--56 are not formalised. The axioms used are the three usual axioms of mathlib.

\litterature{Classical: (1) is the determinantal expression of the inverse of a power series (Wronski; in the language of symmetric functions, $h_n = \det(e_{1-i+j})$, Macdonald, \emph{Symmetric functions and Hall polynomials}, I, \S2); (2) is Newton's identity. The edition does not present them as new.}

\section[Notebook 1987-11: the binomial series at a nilpotent]{Notebook 1987-11: the binomial series $(1 - \alpha)^{-1/2}$ at a nilpotent element\protect\verifie{Notebook1987_11}}\label{sec:quillen-1987-11}

\emph{Context.} On page~77 of \carnet{1987}{11} Quillen normalises an involution by the series $\sum_k\binom{2k}{k}\alpha^k/4^k$ for $(1 - \alpha)^{-1/2}$, with $\alpha$ small.

\begin{encadre}
\begin{theoreme}[after Quillen, notebook 1987-11, p.~77]\label{thm:quillen-1987-11}
Let $A$ be a $\mathbf Q$-algebra and $\alpha \in A$ with $\alpha^N = 0$. Then $s = \sum_{k<N}\binom{2k}{k}4^{-k}\alpha^k$ satisfies $s^2(1 - \alpha) = 1$. If moreover $u^2 = 1 - \alpha$ and $u$ commutes with $\alpha$, then $F = us$ satisfies $F^2 = 1$.
\end{theoreme}
\end{encadre}

\begin{proof}
Let $P \in \mathbf Q[X]$ be the truncation in degrees $< N$ of $\sum_k\binom{2k}k X^k/4^k$. Its square times $1 - X$ agrees with $1$ in degrees $< N$, since the formal identity $(\sum_k\binom{2k}kX^k/4^k)^2(1 - X) = 1$ holds (\S\ref{sec:quillen-1985-8}) and coefficients of degree $< N$ of a product depend only on the truncations. Hence $X^N$ divides $P^2(1 - X) - 1$; evaluating at $\alpha$ gives the first claim. For the second, $u$, $s$ and $\alpha$ commute, and $F^2 = u^2s^2 = (1 - \alpha)s^2 = 1$.
\end{proof}

\emph{Verification in Lean.} The file \href{\leanblobmain/Notebook1987_11.lean}{\texttt{Notebook1987\_11.\allowbreak{}lean}} proves \texttt{sq\_mul\_one\_sub} and \texttt{involution}, importing \texttt{Notebook1985\_8} for the formal identity; the algebra need not be commutative. The analytic case ($\alpha$ of norm $< 1$) is not formalised. The axioms used are the three usual axioms of mathlib.

\litterature{Classical (the binomial series); the edition does not present it as new.}


\part{Rings, ideals and modules}
\section[Notebook 1971-12: a ring is not a finite union of cosets]{Notebook 1971-12: a ring is not a finite union of cosets of ideals with infinite quotients\protect\verifie{Notebook1971_12Cover}}\label{sec:quillen-1971-12-cover}

\emph{Context.} On pages~17 to~20 of \carnet{1971}{12} Quillen shows that a closed subset of small dimension of an affine space over a ring can be avoided by a section, choosing the coordinates $a_1, a_2, \ldots$ one at a time so that each hyperplane $X_1 = a_1$ cuts down the dimension.

\emph{The page.} Page~19 treats the case where all residue fields of the commutative ring $A$ are infinite. If no good $a_1$ existed, $A$ would be a finite union of cosets $a_i + \mathfrak q_i$ of prime ideals. The page shows that this is impossible: one may assume the $\mathfrak q_i$ maximal; each $A/\mathfrak q_i$ is then infinite, so one can choose a class modulo each $\mathfrak q_i$ avoiding the finitely many $a_i$ attached to it, and the Chinese remainder theorem gives an element outside all the cosets.

\emph{Intuition.} B.~H.~Neumann's lemma says that if a group is a finite union of cosets of subgroups, one of the subgroups has finite index. Here every quotient $A/\mathfrak q_i$ is infinite, so no such covering exists; the hypothesis on residue fields serves only to make the quotients infinite, since a finite nonzero quotient would have a finite residue field.

\begin{encadre}
\begin{theoreme}[After Quillen, notebook 1971-12, p.~19]\label{thm:quillen-1971-12-cover}
\emph{(i)} In any ring $A$, a finite family of cosets $a_i + I_i$ of left ideals with $A/I_i$ infinite does not cover $A$. \emph{(ii)} If $A$ is commutative and every residue field $A/\mathfrak m$, $\mathfrak m$ maximal, is infinite, a finite family of cosets $a_i + \mathfrak q_i$ of proper ideals does not cover $A$.
\end{theoreme}
\end{encadre}

\begin{proof}
(i) If the cosets covered $A$, Neumann's lemma applied to the additive group would give an $I_i$ of finite index, against the hypothesis. (ii) A proper ideal $\mathfrak q$ lies in a maximal ideal $\mathfrak m$, and $A/\mathfrak q$ maps onto the infinite field $A/\mathfrak m$; apply (i).
\end{proof}

\emph{Verification in Lean.} The file \href{\leanblobmain/Notebook1971_12Cover.lean}{\texttt{Notebook1971\_12Cover.\allowbreak{}lean}} proves (i) as \texttt{not\_cover\_of\_infinite}, from mathlib's \texttt{AddSubgroup.exists\_finiteIndex\_of\_leftCoset\_cover}, and (ii) as \texttt{not\_cover}. The conclusion is stated as the existence of $x$ with $x - a_i \notin I_i$ for all $i$. The page's route through maximal ideals and the Chinese remainder theorem is replaced by Neumann's lemma. The geometric statement the page proves with it is not formalised. The axioms used are the three usual axioms of mathlib.

\litterature{Neumann's lemma (1954) is in mathlib, and the corollary for vector spaces over infinite fields too (\texttt{Subspace.biUnion\_ne\_univ\_of\_ne\_top}); the form for rings is a routine consequence, of the same kind as prime avoidance. The edition does not present it as new.}

\section[Notebook 1971-13: good submodules pass to quotients]{Notebook 1971-13: good submodules of a module over $A[z]$ pass to quotients\protect\verifie{Notebook1971_13}}\label{sec:quillen-1971-13}

\emph{Context.} In November 1971, \carnet{1971}{13} works towards the homotopy property $K_*(A) \cong K_*(A[z])$ of the $K$-theory of coherent modules, first for a field (3 November), then for a noetherian ring (12 November). For a finitely generated $A[z]$-module $E$, Quillen calls a finitely generated $A$-submodule $L \subset E$ \emph{good} if $L$ generates $E$ over $A[z]$ and $z^{-1}L \subset L$, where $z^{-1}L$ is the set of $e \in E$ such that $ze \in L$. Such an $L$ gives a resolution $0 \to A[z] \otimes_A z^{-1}L \to A[z] \otimes_A L \to E \to 0$ (page~15), and the point is to replace $E$ by the pair $(L, z^{-1}L)$; for this he needs to know how good submodules behave in exact sequences.

\emph{The page.} On page~13 (entry of 10 November 1971), Quillen states a lemma: if $E' \subset E$ is an $A[z]$-submodule, $E'' = E/E'$, $L$ is good for $E$ and $L' = L \cap E'$ generates $E'$, then $L'$ is good for $E'$ and the image $L''$ of $L$ is good for $E''$. The first half and the generation of $E''$ are immediate. For the inclusion $z^{-1}L'' \subset L''$, he writes $ze = l + e'$ with $l \in L$ and $e' \in E'$, uses the generation of $E'$ by $L'$ to write $e' = l' + z e'_1$, and concludes from $z(e - e'_1) = l + l' \in L$. Page~14 draws corollaries for filtered modules.

\emph{Intuition.} The condition $z^{-1}L \subset L$ says that $L$ is closed under \q{division by $z$} inside $E$. To show that $L + E'$ is closed under it as well, one has to absorb the part of $e'$ that is not divisible by $z$, and this is where $L \cap E'$ must generate $E'$: modulo $z E'$, every element of $E'$ is in $L \cap E'$.

\begin{encadre}
\begin{theoreme}[Quillen, notebook 1971-13, p.~13]\label{thm:quillen-1971-13}
Let $A$ be a commutative ring, $E$ an $A$-module, $z$ an $A$-linear endomorphism of $E$, $L \subset E$ a submodule such that $z^{-1}L \subset L$, and $E' \subset E$ a submodule such that $zE' \subset E'$ and $E' \subset \sum_{n \ge 0} z^n (L \cap E')$. If $e \in E$ and $ze \in L + E'$, then $e \in L + E'$. Equivalently, the image $L''$ of $L$ in $E/E'$ satisfies $z^{-1}L'' \subset L''$ for the endomorphism of $E/E'$ induced by $z$.
\end{theoreme}
\end{encadre}

\noindent An $A[z]$-module is the same thing as an $A$-module with an endomorphism $z$; the statement drops the finiteness hypotheses of the page, which play no role in the lemma.

\begin{proof}
First, $E' \subset (L \cap E') + zE'$: an element $z^n l'$ with $l' \in L \cap E'$ is $l'$ if $n = 0$, and $z(z^{n-1}l')$ with $z^{n-1} l' \in E'$ if $n \ge 1$, since $E'$ is stable under $z$. Write $ze = l + e'$ with $l \in L$, $e' \in E'$, and then $e' = l' + z e'_1$ with $l' \in L \cap E'$, $e'_1 \in E'$. Then $z(e - e'_1) = l + l' \in L$, so $e - e'_1 \in z^{-1}L \subset L$, and $e = (e - e'_1) + e'_1 \in L + E'$.
\end{proof}

\noindent The proof is the one on the page; only the first step, which the page states in one line (\q{as $L'$ generates $E'$}), is spelled out.

\emph{Verification in Lean.} The file \href{\leanblobmain/Notebook1971_13.lean}{\texttt{Notebook1971\_13.\allowbreak{}lean}} takes $z$ in \texttt{Module.End A E} and writes the generation hypothesis as \texttt{GeneratesUnder z (L ⊓ E') E'}. The first step is \texttt{le\_sup\_map\_of\_generates}; Theorem~\ref{thm:quillen-1971-13} is \texttt{mem\_sup\_of\_smul\_mem\_sup}, and its form in the quotient, with \texttt{Submodule.mkQ} and \texttt{Submodule.mapQ}, is \texttt{comap\_map\_mkQ\_le}. The easy half of the page's lemma, $z^{-1}(L \cap E') \cap E' \subset L \cap E'$, is \texttt{comap\_inf\_le}. The axioms used are the three usual axioms of mathlib.

\litterature{the homotopy property $G(A) \cong G(A[z])$ for a noetherian ring is published in Quillen, \q{Higher algebraic $K$-theory~I}, \emph{Lecture Notes in Math.}~341 (1973); the search did not locate the lemma of page~13 in this form, and the 1973 proof was not reread to see whether it passes through good submodules. The statement is elementary, and the edition does not present it as new.}

\section[Notebook 1972-2: splitting finite-rank endomorphisms; an integral over a triangle]{Notebook 1972-2: splitting an endomorphism of a free module of infinite rank; an integral over a triangle\protect\verifie{Notebook1972_2}}\label{sec:quillen-1972-2}

\emph{Context.} In February 1972, \carnet{1972}{2} looks for a category with the homotopy type of $BGL(k)^+$ and studies what happens on passing to an infinite-dimensional space $W$. For a free module $R^{(S)}$ on an infinite set $S$, Quillen considers the group $\operatorname{Aut}^f(R^{(S)})$ of automorphisms $1 + K$ with $K$ of finite rank, and conjectures (page~18) that the inclusion $GL_\infty(R) \subset \operatorname{Aut}^f(R^{(\infty)})$ induces an isomorphism on cohomology. Over a field, every such $1 + K$ comes from an automorphism of a finite-dimensional summand (page~17). The entry of 19 February asks how far this holds over a ring. The second statement belongs to the last block of the notebook, on the Godbillon--Vey class and a $2$-cocycle on the group of compactly supported diffeomorphisms of $\mathbf R$; on page~46, Quillen computes an elementary integral on the way to the cocycle of pages~49--50.

\emph{The page (pp.~20--21).} Let $R$ be a ring and $K$ an endomorphism of $R^{(S)}$ of finite rank, that is, with image in $R^T$ for a finite $T \subset S$. Quillen asks when there is a decomposition $R^{(S)} = V \oplus Q$ with $V$ finitely generated, $K(V) \subset V$ and $K(Q) = 0$; then $1 + K$ is the automorphism $(1 + K|_V) \oplus 1$, which comes from $\operatorname{Aut}(V)$. If $\operatorname{Im} K$ is finitely generated, he picks basis vectors $e_i$, $i \in T_1$, whose images generate it, writes $K e_j = \sum_{i \in T_1} r_{ji} K e_i$ for the other $j$, and changes coordinates by the vectors $e_j - \sum_i r_{ji} e_i$, which lie in the kernel; the span of $e_j$, $j \in T_1 \cup T$, and the span of the new vectors with $j \notin T_1 \cup T$, then do the job. Conversely, $\operatorname{Im} K = K(V)$ is finitely generated. He concludes that over a non-noetherian ring some finite-rank maps $R^{(S)} \to R$ have kernels containing no summand of finite corank, and on page~22 passes to a noetherian ring.

\emph{Intuition.} A summand $Q$ in the kernel of $K$ with a finitely generated complement $V$ forces $\operatorname{Im} K = K(V)$ to be finitely generated. Conversely, once a finite part $R^{(U)}$ of the basis already has the full image under $K$, every vector outside $R^{(U)}$ can be corrected by a vector of $R^{(U)}$ with the same image, and the corrected vectors span a complement in the kernel. The finite-rank hypothesis is not needed: a finitely generated image lies in some $R^T$ automatically.

\begin{encadre}
\begin{theoreme}[after Quillen, notebook 1972-2, pp.~20--21]\label{thm:quillen-1972-2}
Let $R$ be a ring, $S$ a set and $K$ an endomorphism of the free left module $R^{(S)}$. There is a decomposition $R^{(S)} = V \oplus Q$ with $V$ finitely generated, $K(Q) = 0$ and $K(V) \subset V$ if and only if $\operatorname{Im} K$ is finitely generated; $V$ can then be taken of the form $R^{(U)}$ with $U \subset S$ finite. In particular, if $R$ is left noetherian, every endomorphism of finite rank splits in this way; if $R$ is not left noetherian and $I \subset R$ is a left ideal that is not finitely generated, the map $R^{(I)} \to R$, $e_i \mapsto i$, has no such $Q$ in its kernel.
\end{theoreme}
\end{encadre}

\noindent The statement drops the page's hypothesis that $K$ is of finite rank, and makes the page's remark on non-noetherian rings explicit with the ideal $I$.

\begin{proof}
If $R^{(S)} = V \oplus Q$ with $V$ finitely generated and $K(Q) = 0$, then $\operatorname{Im} K = K(V) + K(Q) = K(V)$ is finitely generated. Conversely, let $g_1, \dots, g_n$ generate $\operatorname{Im} K$ and choose $x_k$ with $K x_k = g_k$. Let $U \subset S$ be the finite union of the supports of the $g_k$ and $x_k$, $V = R^{(U)}$ and $W = R^{(S \setminus U)}$, so that $R^{(S)} = V \oplus W$, $\operatorname{Im} K \subset V$ (hence $K(V) \subset V$) and $\operatorname{Im} K = K(V)$. The map $K|_V \colon V \to \operatorname{Im} K$ is onto and $W$ is free, so $K|_W \colon W \to \operatorname{Im} K$ lifts to $h \colon W \to V$ with $K h = K|_W$. Put $Q = \{w - h(w) : w \in W\}$. Then $K(Q) = 0$; if $w - h(w) \in V$ then $w \in V \cap W = 0$; and $v + w = (v + h(w)) + (w - h(w))$, so $R^{(S)} = V \oplus Q$. If $R$ is noetherian and $\operatorname{Im} K \subset R^T$ with $T$ finite, $\operatorname{Im} K$ is finitely generated. For the last assertion, the map $e_i \mapsto i$ has image $I$, which is not finitely generated.
\end{proof}

\noindent The lift $h$ replaces the page's explicit change of coordinates: with $h(e_j) = \sum_i r_{ji} e_i$ one recovers the vectors $e_j - \sum_i r_{ji} e_i$.

\emph{The page (p.~46).} In the course of the cocycle computation, Quillen integrates $(1 + ax + by)^{-2}$ over the triangle $x, y \ge 0$, $x + y \le 1$, first in $y$, then in $x$, and finds the symmetric expression below; the function $\log(1+x)/x$ appears in the cocycle of page~50.

\begin{encadre}
\begin{theoreme}[Quillen, notebook 1972-2, p.~46]\label{thm:quillen-1972-2-integral}
For real numbers $a, b > -1$ with $a \ne 0$, $b \ne 0$ and $a \ne b$,
\[
\int_0^1 \!\! \int_0^{1-x} \frac{dy\,dx}{(1 + ax + by)^2} = \frac{\log(1+a)}{a(b-a)} + \frac{\log(1+b)}{b(a-b)}.
\]
\end{theoreme}
\end{encadre}

\noindent The page states no hypothesis; $a, b > -1$ is the condition for $1 + ax + by$ to be positive on the closed triangle, whose vertices give the values $1$, $1 + a$, $1 + b$, and the other three exclude the zero denominators of the formula.

\begin{proof}
On the triangle, $1 + ax + by = (1 - x - y) + x(1 + a) + y(1 + b) > 0$. For $0 \le x \le 1$, an antiderivative in $y$ is $-1/(b(1 + ax + by))$, so the inner integral is $\frac1b\bigl(\frac{1}{1+ax} - \frac{1}{1 + b + (a-b)x}\bigr)$. An antiderivative of this in $x$ is $\frac1b\bigl(\frac{\log(1+ax)}{a} - \frac{\log(1 + b + (a-b)x)}{a-b}\bigr)$, and its variation from $0$ to $1$ is $\frac{\log(1+a)}{ab} - \frac{\log(1+a) - \log(1+b)}{b(a-b)}$, which is the right-hand side since $\frac{1}{ab} - \frac{1}{b(a-b)} = \frac{1}{a(b-a)}$.
\end{proof}

\noindent The proof is the computation of the page, with the positivity made explicit.

\emph{Verification in Lean.} The file \href{\leanblobmain/Notebook1972_2.lean}{\texttt{Notebook1972\_2.\allowbreak{}lean}} writes $R^{(S)}$ as \texttt{S →₀ R} and $R^{(U)}$ as \texttt{Finsupp.supported R R U}, over any \texttt{Ring R}. The lifting step, for any decomposition $M = V \oplus W$ with $W$ projective and $K(V) = K(M)$, is \texttt{exists\_isCompl\_le\_ker}; the converse is \texttt{fg\_range\_of\_isCompl}. Theorem~\ref{thm:quillen-1972-2} is \texttt{exists\_splitting\_iff}, with the noetherian case \texttt{exists\_splitting\_of\_isNoetherianRing} and the counterexample \texttt{not\_exists\_splitting\_of\_not\_isNoetherianRing}. Theorem~\ref{thm:quillen-1972-2-integral} is \texttt{integral\_simplex}, stated with iterated interval integrals, the inner one being \texttt{integral\_inner}; the integral over the triangle as a subset of $\mathbf R^2$ for Lebesgue measure, which equals the iterated one by Fubini's theorem, is not formalized. The axioms used are the three usual axioms of mathlib.

\litterature{the splitting lemma was not found in this form; it is elementary, and was not searched for specifically. The integral is elementary calculus; the formula was checked numerically for $a = 0.7$, $b = 2.3$. The edition presents neither as new.}

\section[Notebook 1972-9: a finite extension]{Notebook 1972-9: $k[t, x]$ is finite over $k[t, tx + x^q]$\protect\verifie{Notebook1972_9Finite}}\label{sec:quillen-1972-9-finite}

\emph{Context.} On page~71 of \carnet{1972}{9} Quillen deduces that the family $f + t\cdot\mathrm{id}$ is proper from the finiteness of $k[t, x]$ over the subalgebra generated by $t$ and $tx + x^q$.

\begin{encadre}
\begin{theoreme}[Quillen, notebook 1972-9, p.~71]\label{thm:quillen-1972-9-finite}
Let $k$ be a commutative ring and $q \ge 2$. Then $k[t, x]$ is a finite module over its subalgebra $k[t, y]$, $y = tx + x^q$.
\end{theoreme}
\end{encadre}

\begin{proof}
$x$ is a root of the monic polynomial $Z^q + tZ - y$ with coefficients in $k[t, y]$, so $k[t, y][x]$ is finite over $k[t, y]$; and $k[t, y][x] = k[t, x]$, since $t$ lies in $k[t, y]$.
\end{proof}

\emph{Verification in Lean.} The file \href{\leanblobmain/Notebook1972_9Finite.lean}{\texttt{Notebook1972\_9Finite.\allowbreak{}lean}} proves \texttt{isIntegral\_X\_one} and \texttt{finite}, in \texttt{MvPolynomial (Fin 2) k}. The hypothesis $q \ge 2$ makes the polynomial monic; for $q = 1$ the statement fails ($y = (1 + t)x$), which is not formalised. The axioms used are the three usual axioms of mathlib.

\litterature{Elementary (integral extensions); the edition does not present it as new.}

\section[Notebook 1972-9.5: $\bigcap\pi^nE$ is the largest divisible part]{Notebook 1972-9.5: for $E$ torsion-free over a discrete valuation ring, $\bigcap_n\pi^nE$ is $\pi$-divisible\protect\verifie{Notebook1972_9_5}}\label{sec:quillen-1972-9-5}

\emph{Context.} On page~3 of \carnet{1972}{9.5}, sorites for his $K'$-theory of modules, Quillen considers a module $E$ over a discrete valuation ring with uniformiser $\pi$ and the submodule $W = \bigcap_n\pi^nE$, which he shows is $\pi$-divisible, hence a vector space over the fraction field, before asking whether $E/W$ is a lattice.

\begin{encadre}
\begin{theoreme}[Quillen, notebook 1972-9.5, p.~3]\label{thm:quillen-1972-9-5}
Let $E$ be a module over a commutative ring and $\pi$ a scalar acting injectively on $E$ (for instance $E$ torsion-free over a discrete valuation ring with uniformiser $\pi$). Every $x \in W = \bigcap_n\pi^nE$ is $\pi y$ with $y \in W$; and every subset $D$ of $E$ with $D \subset \pi D$ is contained in $W$.
\end{theoreme}
\end{encadre}

\begin{proof}
Write $x = \pi^{n+1}e_{n+1}$ for each $n$; then $\pi(\pi^ne_{n+1}) = x = \pi e_1$, so $\pi^ne_{n+1} = e_1$ for all $n$, and $y = e_1 \in W$. For the second claim, $D \subset \pi^nD \subset \pi^nE$ by induction on $n$.
\end{proof}

\emph{Verification in Lean.} The file \href{\leanblobmain/Notebook1972_9_5.lean}{\texttt{Notebook1972\_9\_5.\allowbreak{}lean}} proves \texttt{divisible\_iInf} and \texttt{mem\_iInf\_of\_divisible}, with $W$ written elementwise. Injectivity of $\pi$ is the only hypothesis used. The transversality sorites of page~5 is the statement of section~\ref{sec:quillen-1972-11-transverse}; the coherence sorites of pages~31--33 and the lattice question are not formalised. The axioms used are the three usual axioms of mathlib.

\litterature{Standard commutative algebra (Bourbaki, \emph{Alg.\ comm.}~VII); the edition does not present it as new.}

\section[Notebook 1972-11: extensions of finitely presented modules]{Notebook 1972-11: an extension of finitely presented modules is finitely presented\protect\verifie{Notebook1972_11}}\label{sec:quillen-1972-11}

\emph{Context.} In December 1972, \carnet{1972}{11} works on the foundations of higher algebraic $K$-theory: a first proof of the resolution theorem (1~December), the embedding of an exact category into an abelian category of left exact functors (17~December), the $K'$-groups of graded and filtered rings (21~December). The run of 23~December, headed with Gersten's theorem and coherence, aims at the $K$-theory of the free algebra $k\langle X_1, \dots, X_n\rangle = T(V)$: Quillen shows that every left ideal of $T(V)$ is free, and wants the theorem on filtered rings with \q{noetherian} replaced by \q{coherent}. On the way he writes out a sorites on finitely presented modules and coherent rings (pp.~35--37).

\emph{The page.} On page~36 (his page~3 of the run), for a short exact sequence $0 \to M' \to M \to M'' \to 0$ of left modules over a ring $A$, Quillen states and proves: (a)~if $M''$ is finitely presented and $M$ finitely generated, then $M'$ is finitely generated; (b)~if $M$ is finitely presented and $M'$ finitely generated, then $M''$ is finitely presented; (c)~if $M'$ and $M''$ are finitely presented, so is $M$. All three come from one diagram: a free module $A^p$ mapping onto $M''$, its kernel $K$, and the pull-back $X = M \times_{M''} A^p$, an extension of $A^p$ by $M'$, hence isomorphic to $M' \oplus A^p$, which maps onto $M$ with kernel $K$. He then proves the equivalence of three definitions of a coherent ring (pp.~36--37), using (a), (b) and (c).

\begin{encadre}
\begin{theoreme}[Quillen, notebook 1972-11, p.~36, (c)]\label{thm:quillen-1972-11}
Let $A$ be a ring and $0 \to M' \xrightarrow{f} M \xrightarrow{g} M'' \to 0$ an exact sequence of left $A$-modules. If $M'$ and $M''$ are finitely presented, then $M$ is finitely presented.
\end{theoreme}
\end{encadre}

\begin{proof}
Since $M''$ is finitely generated, there is a surjection $\pi \colon A^n \to M''$; since $A^n$ is free, it lifts to $\sigma \colon A^n \to M$ with $g\sigma = \pi$. The map $\varphi \colon M' \oplus A^n \to M$, $(m', v) \mapsto f(m') + \sigma(v)$, is onto: given $m$, choose $v$ with $\pi(v) = g(m)$; then $g(m - \sigma v) = 0$, so $m - \sigma v = f(m')$ for some $m'$. The second projection sends $\ker\varphi$ onto $\ker\pi$ (if $\pi(v) = 0$, then $\sigma(v) = f(m')$ for some $m'$, and $(-m', v) \in \ker\varphi$), and injectively, since $f$ is injective. Now $\ker\pi$ is finitely generated, by~(a), because $M''$ is finitely presented; hence so is $\ker\varphi$. As $M' \oplus A^n$ is finitely presented, (b) applied to $\varphi$ shows that $M$ is finitely presented.
\end{proof}

\noindent This is the notebook's proof, with $\varphi$ in place of the pull-back $X$: both are the same extension of $A^n$ by $M'$, split because $A^n$ is free.

\emph{Verification in Lean.} The file \href{\leanblobmain/Notebook1972_11.lean}{\texttt{Notebook1972\_11.\allowbreak{}lean}} proves Theorem~\ref{thm:quillen-1972-11} under the name \texttt{finitePresentation\_of\_exact}, for modules over an arbitrary (not necessarily commutative) ring, with \texttt{Module.FinitePresentation} as the notion of finite presentation and \texttt{Function.Exact} for exactness at $M$. The kernel step is \texttt{fg\_ker\_coprod}. Parts (a) and (b) of the page are already in mathlib, as \texttt{Module.FinitePresentation.fg\_ker} and \texttt{Module.finitePresentation\_of\_surjective}; part (c) is there, as far as the search went, only for split sequences (\texttt{Module.finitePresentation\_of\_split\_exact}) and for a projective quotient. The general statement may be worth proposing to mathlib.

\litterature{the three statements are standard and older than the notebook: they are the usual sorites on finitely presented modules that precedes the theory of coherent rings (S.~U. Chase, \q{Direct products of modules}, \emph{Trans.\ Amer.\ Math.\ Soc.} 97 (1961), 457--473, for coherent rings; and the textbook treatments of finitely presented modules, not re-read for this section). The edition does not present them as new; what the Lean file adds is the general case of~(c) to the library.}

\section[Notebook 1972-11: transverse subobjects]{Notebook 1972-11: transverse subobjects\protect\verifie{Notebook1972_11}}\label{sec:quillen-1972-11-transverse}

\emph{Context.} In the first pages of \carnet{1972}{11} Quillen filters objects in two directions, and asks that the two filtrations be transverse. On page~9 (his page~5) he recalls what this means for two subobjects $A$, $B$ of an object $X$ of an abelian category, with the diagram of the nine quotients of $X$ by $A$, $B$, $A \cap B$ and $A + B$.

\emph{The page.} He lists four conditions and says that they are equivalent: $A + B = X$; a condition on the map $X \to X/A \times X/B$; $B \to X/A$ is onto; $A \to X/B$ is onto. When they hold, $A$ and $B$ are called transverse in $X$. In the transcription the second condition is written as an isomorphism $X \xrightarrow{\sim} X/A \times X/B$.

\emph{Intuition.} The kernel of $X \to X/A \times X/B$ is $A \cap B$. The map is onto exactly when every pair of classes $(x \bmod A,\ y \bmod B)$ has a common representative, that is, when $x - y$ always lies in $A + B$. So surjectivity is the condition $A + B = X$, and injectivity is the independent condition $A \cap B = 0$. With an isomorphism the second condition would say that $X = A \oplus B$, which is stronger than the three others: in $X = \mathbf Z$, $A = B = X$ are transverse and the map to $X/A \times X/B = 0$ is not injective.

\begin{encadre}
\begin{theoreme}[After Quillen, notebook 1972-11, p.~9]\label{thm:quillen-1972-11-transverse}
Let $A$ and $B$ be submodules of a module $X$. The following are equivalent: \emph{(i)} $A + B = X$; \emph{(ii)} the map $X \to X/A \times X/B$ is surjective; \emph{(iii)} $B \to X/A$ is surjective; \emph{(iv)} $A \to X/B$ is surjective. The map of \emph{(ii)} is injective if and only if $A \cap B = 0$, and bijective if and only if $X = A \oplus B$.
\end{theoreme}
\end{encadre}

\begin{proof}
(i) $\Rightarrow$ (iii): write $x = a + b$; then $x$ and $b$ have the same class modulo $A$. (iii) $\Rightarrow$ (i): if $b \in B$ has the class of $x$ modulo $A$, then $x = (x - b) + b$. (iv) is (iii) with $A$ and $B$ exchanged. (i) $\Rightarrow$ (ii): given $x$ and $y$, write $x - y = a + b$; then $z = x - a = y + b$ has the class of $x$ modulo $A$ and of $y$ modulo $B$. (ii) $\Rightarrow$ (i): given $x$, take $z$ with the class of $x$ modulo $A$ and of $0$ modulo $B$; then $x = (x - z) + z$ with $x - z \in A$ and $z \in B$. The kernel of the map of (ii) is $A \cap B$.
\end{proof}

\emph{Verification in Lean.} The section \texttt{Transverse} of the file \href{\leanblobmain/Notebook1972_11.lean}{\texttt{Notebook1972\_11.\allowbreak{}lean}} proves the equivalences as \texttt{sup\_eq\_top\_iff\_surjective\_right}, \texttt{sup\_eq\_top\_iff\_surjective\_left} and \texttt{sup\_eq\_top\_iff\_surjective\_prod}, with \texttt{A.mkQ.prod B.mkQ} for the map to the product, and the last sentence as \texttt{injective\_prod\_iff} and \texttt{bijective\_prod\_iff}, with \texttt{IsCompl A B}. The statement is proved for modules over a ring, not in an abelian category. The axioms used are the three usual axioms of mathlib.

\litterature{This is the elementary form of the Chinese remainder theorem for two submodules. Mathlib has the equivalence of (i) and (iii) (\texttt{Submodule.map\_mkQ\_eq\_top}); we did not find there the statement on the map to the product. We say \q{after} because the theorem replaces the isomorphism of the transcription by a surjection; the handwriting has not been re-read against the facsimile on this point, and the sign may be an arrow rather than an isomorphism. The edition does not present the statement as new.}

\section[Notebook 1987-12: Tor of cyclic modules]{Notebook 1987-12: $\operatorname{Tor}_1(A/I, A/J) \cong (I \cap J)/IJ$\protect\verifie{Notebook1987_12}}\label{sec:quillen-1987-12}

\emph{Context.} On page~81 of \carnet{1987}{12} Quillen computes $\operatorname{Tor}_1$ of two cyclic modules; on page~63 of \carnet{1990}{5} he uses the case $J = I$, $\operatorname{Tor}_1(R/I, R/I) \cong I/I^2$.

\begin{encadre}
\begin{theoreme}[Quillen, notebook 1987-12, p.~81; notebook 1990-5, p.~63]\label{thm:quillen-1987-12}
Let $I$, $J$ be ideals of a commutative ring $A$. The kernel of $(A/I) \otimes_A J \to (A/I) \otimes_A A$ is isomorphic to $(I \cap J)/IJ$. For $J = I$ the map is zero, and $(A/I) \otimes_A I \cong I/I^2$.
\end{theoreme}
\end{encadre}

\begin{proof}
Every element of $(A/I) \otimes J \cong J/IJ$ is $1 \otimes y$, zero iff $y \in IJ$; its image $\bar y \in A/I$ vanishes iff $y \in I$. So $x \mapsto 1 \otimes x$ maps $I \cap J$ onto the kernel, with kernel $IJ$.
\end{proof}

\emph{Verification in Lean.} The file \href{\leanblobmain/Notebook1987_12.lean}{\texttt{Notebook1987\_12.\allowbreak{}lean}} proves \texttt{range\_g}, \texttt{ker\_g}, \texttt{kernel\_equiv}, \texttt{lTensor\_self\_eq\_zero} and \texttt{tensor\_self\_equiv} (with mathlib's \texttt{Ideal.Cotangent} for $I/I^2$). The identification of this kernel with $\operatorname{Tor}_1$ (long exact sequence) and the two-sided noncommutative version of 1990-5 are not formalised. The axioms used are the three usual axioms of mathlib.

\litterature{Classical; mathlib has the vanishing criterion \texttt{injective\_lTensor\_quotient\_iff\_inf\_eq\_mul}, not the isomorphism. The edition does not present it as new.}

\section[Notebook 1991-4: $J \cap K = JK + KJ$ for comaximal ideals]{Notebook 1991-4: $J \cap K = JK + KJ$ for comaximal two-sided ideals\protect\verifie{Notebook1991_4}}\label{sec:quillen-1991-4}

\emph{Context.} \carnet{1991}{4}, of the summer of 1991, belongs to the period of the work of Cuntz and Quillen on quasi-free algebras. On pages~20 to~34 Quillen constructs an explicit lifting for the surjection of a tensor algebra $R = T(V)$ onto a product of two algebras. The kernel $I$ of this surjection is the intersection of two two-sided ideals of $R$ whose sum is $R$.

\emph{The page.} On page~20 (his page~424) he recalls, to compute $I$, that for two-sided ideals $J$, $K$ of a ring $R$ with $R = J + K$ one has $J \cap K = JK + KJ$. The inclusion $\supset$ is clear, and for the other he multiplies $J \cap K$ by $R$ written in terms of $J$ and $K$. The ring $T(V)$ is not commutative, so $JK$ and $KJ$ are different ideals and both are needed.

\emph{Intuition.} In a commutative ring, comaximal ideals satisfy $J \cap K = JK$: write $1 = j + k$ and $x = xj + xk$; if $x$ lies in both ideals, each term is a product of an element of $J$ and an element of $K$. Without commutativity the same line gives $xj \in KJ$ and $xk \in JK$, in this order, whence the sum of the two products.

\begin{encadre}
\begin{theoreme}[Quillen, notebook 1991-4, p.~20]\label{thm:quillen-1991-4}
Let $J$ and $K$ be two-sided ideals of a ring $R$ such that $J + K = R$. Then $J \cap K = JK + KJ$.
\end{theoreme}
\end{encadre}

\begin{proof}
The products $JK$ and $KJ$ lie in $J$ and in $K$, because both ideals are two-sided. Conversely write $1 = j + k$ with $j \in J$, $k \in K$, and let $x \in J \cap K$. Then $x = xk + xj$, with $xk \in JK$ since $x \in J$, and $xj \in KJ$ since $x \in K$.
\end{proof}

\emph{Verification in Lean.} The file \href{\leanblobmain/Notebook1991_4.lean}{\texttt{Notebook1991\_4.\allowbreak{}lean}} proves Theorem~\ref{thm:quillen-1991-4} as \texttt{inf\_eq\_mul\_sup\_mul}, for left ideals \texttt{J K : Ideal R} of a semiring with the instances \texttt{J.IsTwoSided} and \texttt{K.IsTwoSided}, and the commutative case as \texttt{inf\_eq\_mul\_of\_comm}. The computation of the kernel $I$ in the tensor algebra, for which the page uses the statement, is not formalised. The axioms used are the three usual axioms of mathlib.

\litterature{The statement is standard in noncommutative ring theory (the Chinese remainder theorem for two-sided ideals), and the page recalls it as known. Mathlib has the commutative case (\texttt{Ideal.mul\_eq\_inf\_of\_coprime}); we did not find the two-sided one. The edition does not present it as new.}


\part{Homological algebra and complexes}
\section[Notebook 1973-6: shift equivalences]{Notebook 1973-6: inverses of finite filtration, and two stabilisation lemmas\protect\verifie{Notebook1973_6}}\label{sec:quillen-1973-6-shift}

\emph{Context.} After Serre's course, the last pages of \carnet{1973}{6} take up Williams' problem on shift equivalence. Quillen considers objects $A$ of a category with an endomorphism $f_A$, and maps $r \colon A \to B$ commuting with the endomorphisms that admit $r' \colon B \to A$, also commuting, with $r'r = f_A^m$ and $rr' = f_B^m$, which he calls isomorphisms of filtration $\le m$.

\emph{The pages.} On page~64 he checks that such an $r$ has filtration $\le m + 1$ as well, with inverse $f_Ar'$; that filtrations add under composition; and that two inverses $r'$, $r''$ of filtrations $m$, $n$ satisfy $f_A^mr'' = f_A^nr'$, the usual uniqueness of an inverse. On page~67, if $f_A = r'r$ and $f_B = rr'$, then $r$ is an isomorphism of filtration $\le 1$. For the examples he uses that a self-map $f$ of a finite set is a bijection of $f^nA$ for $n$ large, and that on a finitely generated free abelian group $f$ is injective on $f^nA$ for $n$ large, since the kernels of $f^n$ stabilise.

\begin{encadre}
\begin{theoreme}[Quillen, notebook 1973-6, pp.~64, 67]\label{thm:quillen-1973-6-shift}
In any category, with composition written from right to left: \emph{(i)} if $r'$ is an inverse of filtration $m$ of $r$, then $f_Ar'$ is one of filtration $m+1$; \emph{(ii)} if $r'$, $s'$ are inverses of filtration $m$, $n$ of $r \colon A \to B$, $s \colon B \to C$, then $r's'$ is one of filtration $m + n$ of $sr$; \emph{(iii)} two inverses $r'$, $r''$ of $r$ of filtrations $m$, $n$ satisfy $f_A^mr'' = f_A^nr'$; \emph{(iv)} if $f_A = r'r$ and $f_B = rr'$, then $r'$ is an inverse of filtration $1$ of $r$. Moreover \emph{(v)} a self-map of a finite set is a bijection of the image of some iterate $f^n$ onto itself, and \emph{(vi)} an endomorphism of a Noetherian module is injective on the image of some $f^n$.
\end{theoreme}
\end{encadre}

\begin{proof}
(i)--(iv) are the computations of the page, for instance $r's'sr = r'f_B^nr = r'rf_A^n = f_A^{m+n}$ and $f_A^mr'' = r'rr'' = r'f_B^n = f_A^nr'$. (v) The images of the iterates decrease, so their cardinalities stabilise; once $f^{n+1}A = f^nA$, $f$ maps the finite set $f^nA$ onto itself, hence bijectively. (vi) For $n$ large, $\ker f^n \cap f^nA = 0$ (Fitting); if $f(x) = f(y)$ with $x, y \in f^{n}A$, then $x - y \in \ker f \cap f^{n}A = 0$.
\end{proof}

\emph{Verification in Lean.} The section \texttt{ShiftEquivalence} of \href{\leanblobmain/Notebook1973_6.lean}{\texttt{Notebook1973\_6.\allowbreak{}lean}} proves (i)--(iv) as \texttt{IsARInv.succ}, \texttt{IsARInv.comp}, \texttt{IsARInv.unique} and \texttt{isARInv\_one\_of\_factor}, for a structure \texttt{IsARInv} holding the four equations, with diagrammatic composition $\gg$ and the iterate \texttt{iPow}; the section \texttt{Examples} proves (v) as \texttt{exists\_bijOn\_range\_iterate} and (vi) as \texttt{exists\_injOn\_range\_pow}, from mathlib's \texttt{Module.End.eventually\_disjoint\_ker\_pow\_range\_pow}. The page's goal, that shift equivalence implies strong shift equivalence over $\mathbf Z$, is not formalised. The axioms used are the three usual axioms of mathlib.

\litterature{Shift and strong shift equivalence are Williams' (Ann.\ of Math.\ 98, 1973); the checks are elementary, and their equivalence over $\mathbf Z$ is in Williams' paper. The edition does not present them as new.}

\section[Notebook 1984-16: a homotopy to an invertible operator]{Notebook 1984-16: a contracting homotopy to an invertible operator kills cohomology\protect\verifie{Notebook1984_16}}\label{sec:quillen-1984-16}

\emph{Context.} On pages~19 and~20 of \carnet{1984}{16} Quillen shows that the Witten cohomology of the $S^1$-invariant forms on a manifold with a free circle action vanishes, the algebraic heart of the localisation theorem of Atiyah and Bott.

\emph{The page.} For $d_a = d + a\iota_X$, $a \ne 0$, and a connection form $\theta$, the anticommutator $d_a\theta + \theta d_a$ is multiplication by $d\theta + a$, which is invertible because $d\theta$ is nilpotent; so $\theta$, rescaled, is a contracting homotopy.

\begin{encadre}
\begin{theoreme}[Quillen, notebook 1984-16, pp.~19--20]\label{thm:quillen-1984-16}
Let $D$, $T$, $N$ be endomorphisms of a module over a commutative ring $R$ and $a$ a unit of $R$, with $D^2 = 0$, $DT + TD = N + a$, $N$ nilpotent and $ND = DN$. Then $\ker D = \operatorname{im} D$. More generally, $\delta^2 = 0$ and $\delta K + K\delta = u$ with $u$ invertible and commuting with $\delta$ give $\ker\delta = \operatorname{im}\delta$.
\end{theoreme}
\end{encadre}

\begin{proof}
If $\delta x = 0$ then $ux = \delta Kx$, so $x = u^{-1}\delta Kx = \delta(u^{-1}Kx)$. In the first statement $u = N + a$ is invertible, a unit plus a commuting nilpotent element, and commutes with $D$.
\end{proof}

\emph{Verification in Lean.} The file \href{\leanblobmain/Notebook1984_16.lean}{\texttt{Notebook1984\_16.\allowbreak{}lean}} proves the theorem as \texttt{ker\_eq\_range\_witten} and \texttt{ker\_eq\_range\_of\_homotopy}, in \texttt{Module.End R M}, with mathlib's \texttt{IsNilpotent.isUnit\_add\_right\_of\_commute}. The forms, the circle action and the identity $[d_a, \theta] = d\theta + a$ are not formalised. The axioms used are the three usual axioms of mathlib.

\litterature{Atiyah and Bott, \emph{The moment map and equivariant cohomology} (Topology~23, 1984); the algebraic step is standard. The edition does not present it as new.}

\section[Notebook 1988-10: the contracting homotopy for T(V)]{Notebook 1988-10: $b'h + sp = 1$ for the tensor algebra\protect\verifie{Notebook1988_10Bar}}\label{sec:quillen-1988-10-bar}

\emph{Context.} On pages~54--55 of \carnet{1988}{10} Quillen proves that the bimodules of noncommutative differentials of a tensor algebra are free, first through derivations, then by an explicit contracting homotopy.

\begin{encadre}
\begin{theoreme}[Quillen, notebook 1988-10, pp.~54--55]\label{thm:quillen-1988-10-bar}
Let $R = T(V)$ over a commutative ring $k$, with $V$ free. Let $b' : R^{\otimes 4} \to R^{\otimes 3}$ be the bar differential, $p : R^{\otimes 3} \to R \otimes V \otimes R$ the bimodule map $p(x, v_1\cdots v_n, y) = \sum_i(xv_1\cdots v_{i-1}, v_i, v_{i+1}\cdots v_ny)$, and $s$, $h$ the bimodule maps with $s(1, v, 1) = (1, v, 1)$ and $h(1, v_1\cdots v_n, 1) = (1, v_1\cdots v_n, 1, 1) - \sum_i(v_1\cdots v_{i-1}, v_i, v_{i+1}\cdots v_n, 1)$. Then $pb' = 0$, $ps = 1$ and $b'h + sp = 1$; hence $R^{\otimes 4} \to R^{\otimes 3} \to R \otimes V \otimes R \to 0$ is exact.
\end{theoreme}
\end{encadre}

\begin{proof}
On a basis element $(x, w, y)$, the terms of $b'h$ coming from the sum telescope: $b'\bigl(\sum_i(xv_{<i}, v_i, v_{>i}, y)\bigr) = (xw, 1, y) - (x, w, y) + sp(x, w, y)$, by induction on the length of $w$; and $b'(x, w, 1, y) = (xw, 1, y)$. The identity $pb' = 0$ reduces to $p(x, yz, w) = p(x, y, zw) + p(xy, z, w)$.
\end{proof}

\emph{Verification in Lean.} The file \href{\leanblobmain/Notebook1988_10Bar.lean}{\texttt{Notebook1988\_10Bar.\allowbreak{}lean}} proves \texttt{pw\_append}, \texttt{p\_b'}, \texttt{p\_s}, \texttt{b'\_gw}, \texttt{b'\_h\_add\_s\_p}, \texttt{exact\_at\_middle} and \texttt{p\_surjective}. $V$ is free on a type of letters, so $R^{\otimes n}$ is the free module on $n$-tuples of words; the maps are defined on that basis, outer factors included. The first proof (through derivations) and the corollary for $R^{\otimes 3} \to R^{\otimes 2} \to R \otimes V$ are not formalised. The axioms used are the three usual axioms of mathlib.

\litterature{Classical (freeness of $\Omega^1$ of a tensor algebra; Cuntz and Quillen, 1995); the explicit homotopy as written was not located; the edition does not present it as new.}

\section[Notebook 1989-2: the perturbation lemma]{Notebook 1989-2: the perturbation lemma, algebraic form\protect\verifie{Notebook1989_2}}\label{sec:quillen-1989-2}

\emph{Context.} On pages~76--80 of \carnet{1989}{2} Quillen transfers a perturbation $\theta$ of a differential along a homotopy equivalence $(u, v, h)$, with $vu - 1 = [d, h]$, to $\theta' = u(1 - \theta h)^{-1}\theta v$: the basic perturbation lemma.

\begin{encadre}
\begin{theoreme}[after Quillen, notebook 1989-2, pp.~76--80]\label{thm:quillen-1989-2}
In a ring, let $d$, $d_M$, $\delta$, $h$, $i$, $p$ satisfy $d^2 = d_M^2 = 0$, $(d + \delta)^2 = 0$, $di = id_M$, $pd = d_Mp$, $ip = 1 + dh + hd$, and let $1 - \delta h$ be invertible; put $X = (1 - \delta h)^{-1}\delta$. Then $dX + Xd + X^2 + X(dh + hd)X = 0$, and with $d_M' = d_M + pXi$, $i' = i + hXi$, $p' = p + pXh$, $h' = h + hXh$:
\[ d_M'^2 = 0,\quad (d + \delta)i' = i'd_M',\quad p'(d + \delta) = d_M'p',\quad i'p' = 1 + (d + \delta)h' + h'(d + \delta). \]
If moreover $pi = 1$, $hi = 0$, $ph = 0$ and $h^2 = 0$, then $p'i' = 1$.
\end{theoreme}
\end{encadre}

\begin{proof}
$(1 - \delta h)X = \delta$, and $1 - h\delta$ is invertible with inverse $1 + h(1 - \delta h)^{-1}\delta$, with $X(1 - h\delta) = \delta$. Multiplying the first identity by $1 - \delta h$ on the left and $1 - h\delta$ on the right gives $d\delta + \delta d + \delta^2 = (d + \delta)^2 - d^2 = 0$. The four identities reduce to it: each difference is $p(\cdot)i$, $h(\cdot)i$, $p(\cdot)h$ or $h(\cdot)h$ of the first, plus multiples of $(1 - \delta h)X - \delta$ and $X(1 - h\delta) - \delta$.
\end{proof}

\emph{Verification in Lean.} The file \href{\leanblobmain/Notebook1989_2.lean}{\texttt{Notebook1989\_2.\allowbreak{}lean}} proves \texttt{key}, \texttt{sq\_eq\_zero}, \texttt{chain\_left}, \texttt{chain\_right}, \texttt{homotopy} and \texttt{retract}. Everything is stated in one ring (the endomorphisms of the sum of the two complexes), with operators composed without signs; the side conditions are used only for $p'i' = 1$. The page's $\theta'$ is $pXi$. The axioms used are the three usual axioms of mathlib.

\litterature{Classical (Brown, 1965; Gugenheim, 1972; Huebschmann--Kadeishvili, 1991); the edition does not present it as new.}

\section[Notebook 1989-2: the cyclic bar construction]{Notebook 1989-2: the cyclic bar construction of a group\protect\verifie{Notebook1989_2Simplicial}}\label{sec:quillen-1989-2-simplicial}

\emph{Context.} On pages~94--95 of \carnet{1989}{2} Quillen compares the cyclic bar construction of a group (attributed there, after R.~Cohen, to Waldhausen) with a simplicial set on $G \times G^n$.

\begin{encadre}
\begin{theoreme}[Quillen, notebook 1989-2, pp.~94--95]\label{thm:quillen-1989-2-simplicial}
Let $G$ be a group. On $G^{n+1}$ let $d_i(g_0, \dots, g_n) = (\dots, g_ig_{i+1}, \dots)$ for $i < n$ and $d_n(g_0, \dots, g_n) = (g_ng_0, g_1, \dots, g_{n-1})$; on $G \times G^n$ let $d_0(x, g) = (g_1^{-1}xg_1, g_2, \dots, g_n)$, $d_i$ multiply $g_ig_{i+1}$, and $d_n$ drop $g_n$. Then $\Phi(x, g_1, \dots, g_n) = ((g_1\cdots g_n)^{-1}x, g_1, \dots, g_n)$ is a bijection commuting with every face.
\end{theoreme}
\end{encadre}

\begin{proof}
The inverse is $(g_0, g) \mapsto (g_1\cdots g_n\,g_0, g)$. For $d_0$: $(g_2\cdots g_n)^{-1}g_1^{-1}xg_1 = (g_1\cdots g_n)^{-1}xg_1$; inner faces keep the total product; for $d_n$: $g_n(g_1\cdots g_n)^{-1}x = (g_1\cdots g_{n-1})^{-1}x$.
\end{proof}

\emph{Verification in Lean.} The file \href{\leanblobmain/Notebook1989_2Simplicial.lean}{\texttt{Notebook1989\_2Simplicial.\allowbreak{}lean}} proves \texttt{$\Phi$\_bijective}, \texttt{$\Phi$\_face\_zero}, \texttt{$\Phi$\_face\_inner} and \texttt{$\Phi$\_face\_last}, with tuples \texttt{Fin n} $\to G$, products \texttt{Fin.partialProd} and inner faces \texttt{Fin.contractNth}. The simplicial identities and degeneracies are not formalised. The axioms used are the three usual axioms of mathlib.

\litterature{Classical (Waldhausen's cyclic bar construction; Goodwillie, Burghelea); the edition does not present it as new.}

\section[Notebook 1990-2: endomorphisms of a split extension]{Notebook 1990-2: endomorphisms of a split extension of complexes\protect\verifie{Notebook1990_2Split}}\label{sec:quillen-1990-2-split}

\emph{Context.} On pages~64--67 of \carnet{1990}{2} Quillen classifies the lifts of an endomorphism $L$ of an exact sequence $0 \to X \to E \to Y \to 0$ of complexes, split as graded modules, by homotopies $LS = [d, g]$; on page~71 he uses annihilating idempotents.

\begin{encadre}
\begin{theoreme}[Quillen, notebook 1990-2, pp.~64--67, 71]\label{thm:quillen-1990-2-split}
In a ring, let $ph = 1 = ki$, $ik + hp = 1$, $kh = 0$, $di = id$, $dp = pd$, and $S = kdh$. Then $pi = 0$, $iS = [d, h]$, $Sp = -[d, k]$. If $Li = iL$, $pL = Lp$, $dL = Ld$: $[L, S] = [d, kLh]$; if $LS = [d, g]$, then $J = Lh - ig$, $I = Lk + gp$ satisfy $pJ = L$, $iI + Jp = L$, $Ii = L$, $[d, J] = [d, I] = 0$; conversely if $pJ = L$ and $[d, J] = 0$, $g = kLh - kJ$ satisfies $ig = Lh - J$ and $LS = [d, g]$. Moreover, for idempotents $e$, $f$ with $ef = fe = 0$, $1 - e - f$ is an idempotent with $(1 - e - f)e = (1 - e - f)f = 0$.
\end{theoreme}
\end{encadre}

\begin{proof}
Expansions using the relations: e.g. $iS = (1 - hp)dh = dh - hdph = dh - hd$, and $[d, J] = L[d, h] - i[d, g] = LiS - iLS = 0$.
\end{proof}

\emph{Verification in Lean.} The file \href{\leanblobmain/Notebook1990_2Split.lean}{\texttt{Notebook1990\_2Split.\allowbreak{}lean}} proves \texttt{p\_mul\_i}, \texttt{iS\_eq}, \texttt{Sp\_eq}, \texttt{commutator\_L\_S}, \texttt{of\_homotopy}, \texttt{homotopy\_of} and \texttt{annihilating}. Everything lives in one ring (endomorphisms of $X \oplus E \oplus Y$) with ungraded commutators; the graded signs are not formalised. The page's hypothesis that $L$ is a map of complexes is the hypothesis $dL = Ld$. The axioms used are the three usual axioms of mathlib.

\litterature{Classical (extensions of complexes classified by maps $Y \to \Sigma X$); the edition does not present it as new.}


\part{Hochschild and cyclic cohomology}
\section[{Notebook 1986-2: the cocycle $\mathrm{tr}(a_0[F, a_1])$}]{Notebook 1986-2: the cocycle $\operatorname{tr}(a_0[F, a_1])$ and $\operatorname{tr}(F[F,a][F,b]) = 2\operatorname{tr}(a[F,b])$\protect\verifie{Notebook1986_2}}\label{sec:quillen-1986-2}

\emph{Context.} On pages~65 and~70 of \carnet{1986}{2} Quillen considers, for an operator $F$, the bilinear form $\tau(a_0, a_1) = \operatorname{tr}(a_0[F, a_1])$, the character of a $1$-summable Fredholm module in Connes' theory.

\emph{The page.} He notes that $\tau$ is a Hochschild cocycle, that it is cyclic when $\operatorname{tr}[F, a] = 0$ for all $a$, in particular when $F^2 = 1$, and that then $\operatorname{tr}(F[F, a_0][F, a_1]) = 2\operatorname{tr}(a_0[F, a_1])$; he relates it to the $2$-form $\operatorname{tr} F(dF)^2$ on involutions.

\begin{encadre}
\begin{theoreme}[After Quillen, notebook 1986-2, pp.~65, 70]\label{thm:quillen-1986-2}
Let $F$, $a_0$, $a_1$, $a_2$ be square matrices over a commutative ring and $\tau(a_0, a_1) = \operatorname{tr}(a_0[F, a_1])$. Then $\tau(a_0a_1, a_2) - \tau(a_0, a_1a_2) + \tau(a_2a_0, a_1) = 0$ and $\tau(a_0, a_1) = -\tau(a_1, a_0)$. If $F^2 = 1$, then $\operatorname{tr}(F[F, a][F, b]) = 2\operatorname{tr}(a[F, b])$.
\end{theoreme}
\end{encadre}

\begin{proof}
With $[F, a_1a_2] = [F, a_1]a_2 + a_1[F, a_2]$ the first sum is $\operatorname{tr}(a_2a_0[F,a_1]) - \operatorname{tr}(a_0[F, a_1]a_2) = 0$ by cyclicity of the trace; cyclicity also gives the antisymmetry, since traces of commutators vanish in finite dimension. If $F^2 = 1$, $F[F, a] = a - FaF$, and $\operatorname{tr}((a - FaF)(Fb - bF)) = \operatorname{tr}(aFb) - \operatorname{tr}(abF) - \operatorname{tr}(Fab) + \operatorname{tr}(FaFbF) = 2\operatorname{tr}(aFb) - 2\operatorname{tr}(abF)$.
\end{proof}

\emph{Verification in Lean.} The file \href{\leanblobmain/Notebook1986_2.lean}{\texttt{Notebook1986\_2.\allowbreak{}lean}} proves the theorem as \texttt{tau\_hochschild}, \texttt{tau\_cyclic} and \texttt{trace\_F\_comm\_comm}, for \texttt{Matrix n n R}. In finite dimension the cyclicity needs no hypothesis on $F$; the infinite-dimensional setting of the page (trace-class commutators), where the hypothesis matters, and the $2$-form on involutions are not formalised. The axioms used are the three usual axioms of mathlib.

\litterature{Connes, \emph{Non-commutative differential geometry} (Publ.\ IHÉS 62, 1985), for the character of a Fredholm module and the normalisation $\operatorname{tr}(F[F,a][F,b])$. The edition does not present it as new.}

\section[Notebook 1987-13: kernel and image of the norm]{Notebook 1987-13: $\operatorname{Ker}N = \operatorname{Im}(1 - \sigma)$, $\operatorname{Im}N = \operatorname{Ker}(1 - \sigma)$\protect\verifie{Notebook1987_13}}\label{sec:quillen-1987-13}

\emph{Context.} On pages~41--42 of \carnet{1987}{13} Quillen identifies the cyclic coinvariants and invariants of the tensor algebra through the norm of the cyclic group.

\begin{encadre}
\begin{theoreme}[Quillen, notebook 1987-13, pp.~41--42]\label{thm:quillen-1987-13}
Let $\sigma$ be a linear operator with $\sigma^n = 1$ on a module over a commutative ring in which $n$ is invertible, and $N = 1 + \sigma + \dots + \sigma^{n-1}$. Then $\operatorname{Ker}N = \operatorname{Im}(1 - \sigma)$ and $\operatorname{Im}N = \operatorname{Ker}(1 - \sigma)$.
\end{theoreme}
\end{encadre}

\begin{proof}
$(1 - \sigma)N = N(1 - \sigma) = 1 - \sigma^n = 0$. If $\sigma v = v$ then $v = N(v/n)$. In any ring $(1 - \sigma)\sum_{i<n}\sum_{j<i}\sigma^j = n - N$, so $Nv = 0$ gives $v = (1 - \sigma)\bigl(\tfrac1n\sum_{i<n}\sum_{j<i}\sigma^jv\bigr)$.
\end{proof}

\emph{Verification in Lean.} The file \href{\leanblobmain/Notebook1987_13.lean}{\texttt{Notebook1987\_13.\allowbreak{}lean}} proves \texttt{one\_sub\_mul\_Q}, \texttt{ker\_norm} and \texttt{range\_norm}. The application to $P/[P, P]$ for $P = T(V)$ is not formalised. The axioms used are the three usual axioms of mathlib.

\litterature{Classical (averaging); mathlib has the analogues for the regular representation of a finite cyclic group (\texttt{FiniteCyclic}), not for an arbitrary module with $n$ invertible. The edition does not present it as new.}

\section[Notebook 1987-13: the cup product of Hochschild cochains]{Notebook 1987-13: $b(\varphi \cup \psi) = b\varphi \cup \psi + (-1)^p\varphi \cup b\psi$\protect\verifie{Notebook1987_13Cup}}\label{sec:quillen-1987-13-cup}

\emph{Context.} On pages~71--74 of \carnet{1987}{13} Quillen writes the cup product of Hochschild cochains, checks its Leibniz rule, and computes the cup square of a derivation and the first-order deformations given by two derivations.

\begin{encadre}
\begin{theoreme}[Quillen, notebook 1987-13, pp.~72--74]\label{thm:quillen-1987-13-cup}
Let $A$ be a ring, and for functions $\varphi$ of $p$ arguments in $A$ with values in $A$ let $(\varphi \cup \psi)(a_1, \dots, a_{p+q}) = \varphi(a_1, \dots, a_p)\,\psi(a_{p+1}, \dots, a_{p+q})$ and $b$ the Hochschild coboundary. Then:
\begin{enumerate}
\item $b(\varphi \cup \psi) = b\varphi \cup \psi + (-1)^p\varphi \cup b\psi$, with no linearity needed;
\item for a derivation $X$ and $f = X \circ X$, $f(a_1a_2) = f(a_1)a_2 + 2X(a_1)X(a_2) + a_1f(a_2)$, so $X \cup X = -\tfrac12 bf$ when $2$ is invertible; and $X \cup Y$ is a cocycle for derivations $X, Y$;
\item for $X = [c, \cdot\,]$ inner and $Y$ a derivation, $s(a) = a + hcY(a)$ is a homomorphism from $A$ to $A[h]/h^2$ with the product $a_1 * a_2 = a_1a_2 + hX(a_1)Y(a_2)$; and $s(a) = a + hXY(a)$ gives $s(a_1)s(a_2) - s(a_1a_2) = -h(Xa_1\,Ya_2 + Ya_1\,Xa_2)$.
\end{enumerate}
\end{theoreme}
\end{encadre}

\begin{proof}
(1) In $b(\varphi \cup \psi)$, the faces contracting inside the first $p$ arguments, with the left action, give $b\varphi \cup \psi$ except for its last term; the face contracting $a_p a_{p+1}$ cancels it against the first term of $\varphi \cup b\psi$; the rest is $(-1)^p\varphi \cup b\psi$. (2) and (3) are the Leibniz rule applied twice, in a noncommutative ring.
\end{proof}

\emph{Verification in Lean.} The file \href{\leanblobmain/Notebook1987_13Cup.lean}{\texttt{Notebook1987\_13Cup.\allowbreak{}lean}} proves \texttt{coboundary\_cup}, \texttt{sq\_derivation\_mul}, \texttt{cup\_self\_eq}, \texttt{coboundary\_cup\_derivations}, \texttt{inner\_section} and \texttt{section\_XY}; cochains are functions of sequences depending only on their first $p$ terms, and the coefficients are in $A$ itself (the page's cochains with values in two bimodules and their tensor product are not formalised, nor the identification of the cup product with the Yoneda product). The axioms used are the three usual axioms of mathlib.

\litterature{Classical (Gerstenhaber, 1963); the edition does not present it as new.}

\section[{Notebook 1988-1: the cocycle $\mathrm{tr}(x[y, \rho])$}]{Notebook 1988-1: the cocycle $\operatorname{tr}(x[y, \rho])$ attached to an extension\protect\verifie{Notebook1988_1Cocycle}}\label{sec:quillen-1988-1-cocycle}

\emph{Context.} On pages~73 to~77 of \carnet{1988}{1}, after the lemma of Section~\ref{sec:quillen-1988-1}, Quillen builds maps from $\mathrm{Ext}$ groups to duals of Hochschild homology by cocycles $(a_0, \ldots, a_n) \mapsto \operatorname{tr}_M(a_0[a_1, h]\cdots[a_n, h])$, with $h$ a $k$-linear splitting.

\emph{The page.} For $n = 1$ and $A$ commutative, an extension of $M$ by $M$ is described by a $k$-linear endomorphism $\rho$ of $M$, and the cocycle $\varphi(x, y) = \operatorname{tr}(x[y, \rho])$ gives a map $\mathrm{Ext}^1_A(M, M) \to (\Omega^1_A)^*$: it is a derivation in its second variable, so that $x\,dy \mapsto \varphi(x, y)$ is well defined, and it does not change when $\rho$ is changed by an $A$-linear endomorphism.

\begin{encadre}
\begin{theoreme}[After Quillen, notebook 1988-1, pp.~73--77]\label{thm:quillen-1988-1-cocycle}
Let $\rho, \sigma, x, y, z$ be square matrices over a commutative ring and $\varphi_\rho(x, y) = \operatorname{tr}(x[y, \rho])$. If $xz = zx$ then $\varphi_\rho(x, yz) = \varphi_\rho(xy, z) + \varphi_\rho(xz, y)$. Moreover $\varphi_\rho(x, 1) = 0$, $\varphi_\rho(x, y) = -\varphi_\rho(y, x)$, and $\varphi_{\rho + \sigma}(x, y) = \varphi_\rho(x, y)$ if $\sigma$ commutes with $y$.
\end{theoreme}
\end{encadre}

\begin{proof}
Expanding, the first identity reduces to $\operatorname{tr}(xy\rho z) = \operatorname{tr}(xzy\rho)$ and $\operatorname{tr}(x\rho yz) = \operatorname{tr}(xz\rho y)$, which follow from the cyclicity of the trace and $zx = xz$. The other statements are immediate.
\end{proof}

\emph{Verification in Lean.} The file \href{\leanblobmain/Notebook1988_1Cocycle.lean}{\texttt{Notebook1988\_1Cocycle.\allowbreak{}lean}} proves the theorem as \texttt{phi\_mul\_right}, \texttt{phi\_one}, \texttt{phi\_antisymm} and \texttt{phi\_add\_of\_commute}, for \texttt{Matrix n n k}, the elements of $A$ acting as commuting matrices on $M = k^n$. The $\mathrm{Ext}$ groups, the Kähler differentials and the residue are not formalised. The axioms used are the three usual axioms of mathlib.

\litterature{The shape is that of Connes' character of a Fredholm module, $\operatorname{tr}(a_0[F, a_1]\cdots[F, a_n])$ (Publ.\ IHÉS 62, 1985), and the Hochschild route to the residue is Lipman's (\emph{Residues and traces of differential forms via Hochschild homology}, Contemp.\ Math.\ 61, 1987); for $n = 1$ the identities are elementary. The edition does not present them as new.}

\section[Notebook 1988-5: the cyclic counting lemma]{Notebook 1988-5: the cyclic counting lemma in trace form\protect\verifie{Notebook1988_5}}\label{sec:quillen-1988-5}

\emph{Context.} On pages~87--88 of \carnet{1988}{5} Quillen counts words in two letters up to cyclic permutation; in trace form the count is the derivative of a trace of powers.

\begin{encadre}
\begin{theoreme}[after Quillen, notebook 1988-5, pp.~87--88]\label{thm:quillen-1988-5}
For square matrices $A$, $B$ over a commutative ring, $\frac{d}{ds}\operatorname{tr}((sA + B)^{m+1}) = (m + 1)\operatorname{tr}(A(sA + B)^m)$ as polynomials in $s$.
\end{theoreme}
\end{encadre}

\begin{proof}
For a matrix $M$ of polynomials, $(M^n)' = \sum_{i<n}M^iM'M^{n-1-i}$ by the Leibniz rule; by cyclicity each term has trace $\operatorname{tr}(M'M^{n-1})$. Apply this to $M = sA + B$, with $M' = A$.
\end{proof}

\emph{Verification in Lean.} The file \href{\leanblobmain/Notebook1988_5.lean}{\texttt{Notebook1988\_5.\allowbreak{}lean}} proves \texttt{map\_derivative\_pow}, \texttt{trace\_map\_derivative\_pow} and \texttt{derivative\_trace\_pow}. The coefficientwise form, $k$ times the trace of the words of bidegree $(k, p)$ equals $k + p$ times the trace of $A$ times the words of bidegree $(k - 1, p)$, is not formalised. The axioms used are the three usual axioms of mathlib.

\litterature{Classical (cyclicity of the trace); the edition does not present it as new.}

\section[Notebook 1988-8: the cyclic word count]{Notebook 1988-8: the cyclic word count; traces on a twisted product; a gauge transformation\protect\verifie{Notebook1988_8}}\label{sec:quillen-1988-8}

\emph{Context.} On page~16 of \carnet{1988}{8} Quillen counts the words in two matrices cyclically under the trace, the identity behind the trace form of \S\ref{sec:quillen-1988-5}; on page~10 he notes that traces vanish on the twisted part of a crossed product, and on page~21 he conjugates $\mathrm{diag}(F, -F)$ by a rotation.

\begin{encadre}
\begin{theoreme}[Quillen, notebook 1988-8, pp.~10, 16, 21]\label{thm:quillen-1988-8}
\begin{enumerate}
\item For square matrices $A$, $B$ over a commutative ring, let $P_{k,l}$ be the sum of the words with $k$ letters $A$ and $l$ letters $B$. Then $(k+1)\operatorname{tr}P_{k+1,l} = (k+l+1)\operatorname{tr}(AP_{k,l})$.
\item If $\tau$ is additive with $\tau(ab) = \tau(ba)$ and $\varepsilon t = -t\varepsilon$, then $2\tau(t^n\varepsilon) = 0$ for $n \ge 1$; hence $\tau(t^n\varepsilon) = 0$ with values in a $\mathbf Q$-vector space.
\item If $F^2 = 1$ and $c$, $s$ are scalars with $c^2 + s^2 = 1$, then $g = \begin{pmatrix}c & -sF\\ sF & c\end{pmatrix}$ has inverse $\begin{pmatrix}c & sF\\ -sF & c\end{pmatrix}$ and $g\,\mathrm{diag}(F, -F)\,g^{-1} = \begin{pmatrix}(c^2-s^2)F & 2cs\\ 2cs & (s^2-c^2)F\end{pmatrix}$.
\end{enumerate}
\end{theoreme}
\end{encadre}

\begin{proof}
(1) $P_{k,l}$ is the coefficient of $s^k$ in $(sA + B)^{k+l}$; take the coefficient of $s^k$ in $\frac d{ds}\operatorname{tr}(sA + B)^{k+l+1} = (k+l+1)\operatorname{tr}(A(sA + B)^{k+l})$ (\S\ref{sec:quillen-1988-5}). (2) $\varepsilon t^n = (-1)^nt^n\varepsilon$; for $n$ odd, $\tau(t^n\varepsilon) = \tau(\varepsilon t^n) = -\tau(t^n\varepsilon)$; for $n$ even, move one $t$ around: $\tau(t^n\varepsilon) = \tau(t^{n-1}\varepsilon t) = -\tau(t^n\varepsilon)$. (3) Direct computation.
\end{proof}

\emph{Verification in Lean.} The file \href{\leanblobmain/Notebook1988_8.lean}{\texttt{Notebook1988\_8.\allowbreak{}lean}} defines $P_{k,l}$ as the coefficient, proves its recursion by the first letter (\texttt{P\_succ\_succ} and the edge cases), and proves \texttt{trace\_word\_count}, \texttt{two\_nsmul\_trace\_pow\_mul}, \texttt{trace\_pow\_mul}, \texttt{gauge\_mul\_inv} and \texttt{gauge\_conj}. The identity (1) holds over any commutative ring, without division; in (2), $\varepsilon^2 = 1$ is not needed. In (3) the page writes the factors $F$ on the off-diagonal of the result; with $g$ as above they cancel, since $F^2 = 1$. The axioms used are the three usual axioms of mathlib.

\litterature{Classical (cyclicity of the trace); the edition does not present it as new.}

\section[Notebook 1988-10: the identities of the cyclic bicomplex]{Notebook 1988-10: $Nb = b'N$, $b(1 - t) = (1 - t)b'$\protect\verifie{Notebook1988_10}}\label{sec:quillen-1988-10}

\emph{Context.} On pages~69--70 of \carnet{1988}{10}, as corollaries of an exact sequence of comodules over the tensor coalgebra, Quillen obtains the two identities that make Connes's cyclic bicomplex a double complex.

\begin{encadre}
\begin{theoreme}[Quillen, notebook 1988-10, pp.~69--70]\label{thm:quillen-1988-10}
Let $A$ be a ring and, on $A^{\otimes(n+1)}$, let $d_i(a_0, \dots, a_n) = (a_0, \dots, a_ia_{i+1}, \dots, a_n)$ for $i < n$, $d_n(a_0, \dots, a_n) = (a_na_0, a_1, \dots, a_{n-1})$, $b' = \sum_{i<n}(-1)^id_i$, $b = b' + (-1)^nd_n$, $t(a_0, \dots, a_n) = (-1)^n(a_n, a_0, \dots, a_{n-1})$ and $N = \sum_{j \le n}t^j$. Then $Nb = b'N$ and $b(1 - t) = (1 - t)b'$.
\end{theoreme}
\end{encadre}

\begin{proof}
From the relations $d_0t = \pm d_n$ and $d_{i+1}t = -td_i$ one gets $t^jd_i = \pm d_{i-j}t^j$ for $j \le i$ and $t^{i+1+k}d_i = \pm d_{n-k}t^{i+k}$; in $Nb$ and $b'N$ the double sums over $(j, i)$ then correspond term by term under the bijection $(j, i) \mapsto (j, i - j)$ for $j \le i$, $(j, i) \mapsto (j - 1, n + 1 - j + i)$ for $i < j$. The second identity follows from the same relations: $t$ carries the faces of $b$ other than $d_0$ to those of $b't$, and $d_0$ to the last face.
\end{proof}

\emph{Verification in Lean.} The file \href{\leanblobmain/Notebook1988_10.lean}{\texttt{Notebook1988\_10.\allowbreak{}lean}} proves \texttt{face\_rot\_zero}, \texttt{face\_rot\_succ}, \texttt{T\_pow\_dual\_le}, \texttt{T\_pow\_dual\_gt}, \texttt{one\_sub\_T\_b} and \texttt{b\_N}, in transposed form: on cochains, functions of sequences with values in an abelian group, where the operators act by precomposition; an identity of operators holds on all elementary tensors if and only if its transpose holds on all such functions, so the two forms are equivalent. Quillen's derivation from the cotensor exact sequence is not formalised. The axioms used are the three usual axioms of mathlib.

\litterature{Classical (Connes, 1985; Loday, \emph{Cyclic homology}); their derivation from a cotensor exact sequence was not located; the edition does not present the identities as new.}

\section[Notebook 1989-4: the cyclic operator is homotopic to 1]{Notebook 1989-4: $(bs + sb)\psi = \psi - \psi^\lambda$\protect\verifie{Notebook1989_4Cochain}}\label{sec:quillen-1989-4-cochain}

\emph{Context.} On pages~70--71 of \carnet{1989}{4} Quillen shows that the cyclic operator acts on Hochschild cochains by a map homotopic to the identity, with homotopy $s$.

\begin{encadre}
\begin{theoreme}[Quillen, notebook 1989-4, p.~71]\label{thm:quillen-1989-4-cochain}
Let $A$ be a monoid, $M$ an abelian group, and for cochains $\psi : A^{m} \to M$ let $(s\psi)(a_1, \dots, a_{m-1}) = \psi(1, a_1, \dots, a_{m-1})$, $b'$ the alternating sum of the inner faces and $b = b' + (-1)^m\,(\text{crossover } a_ma_0)$. For $\psi : A^{n+1} \to M$, $n \ge 1$, and
\[
\psi^\lambda(a_0, \dots, a_n) = (-1)^{n-1}\psi(1, a_na_0, a_1, \dots, a_{n-1}) + (-1)^n\psi(a_n, a_0, \dots, a_{n-1}),
\]
one has $b's + sb' = 1$ and $(bs + sb)\psi = \psi - \psi^\lambda$.
\end{theoreme}
\end{encadre}

\begin{proof}
In $(sb'\psi)(a)$ the face at $0$ gives $\psi(a)$ and the others give $-(b's\psi)(a)$. The crossover of $bs\psi$ is $(-1)^n\psi(1, a_na_0, a_1, \dots)$, that of $sb\psi$ is $(-1)^{n+1}\psi(a_n, a_0, \dots)$.
\end{proof}

\emph{Verification in Lean.} The file \href{\leanblobmain/Notebook1989_4Cochain.lean}{\texttt{Notebook1989\_4Cochain.\allowbreak{}lean}} proves \texttt{s\_b'}, \texttt{b'\_s\_add\_s\_b'} and \texttt{bs\_add\_sb}, with cochains on \texttt{Fin m} $\to A$; neither identity needs $\psi$ normalised. \emph{Correction}: the intermediate box of p.~70 gives the two crossover terms with the signs $(-1)^{n+1}$ and $(-1)^n$; the lines above it and the box of p.~71 give $(-1)^n$ and $(-1)^{n+1}$, which is what holds (the slip may be the transcription's). The axioms used are the three usual axioms of mathlib.

\litterature{Classical in cyclic cohomology (Connes; Loday); the edition does not present it as new.}

\section[Notebook 1989-5: Cuntz's polynomial and two homomorphisms]{Notebook 1989-5: $(1 - \lambda)W(\lambda) = 1 - N(\lambda)/(n + 1)$; the difference of two homomorphisms\protect\verifie{Notebook1989_5}}\label{sec:quillen-1989-5}

\emph{Context.} On pages~95, 107--108 and~124 of \carnet{1989}{5} Quillen uses a polynomial of Cuntz to invert $1 - \lambda$ on the non-invariant part when $\lambda^{n+1} = 1$, inverts $1 + \lambda$ for odd $n$, and studies the half-difference of two ring homomorphisms.

\begin{encadre}
\begin{theoreme}[Quillen, notebook 1989-5, pp.~95, 107, 124]\label{thm:quillen-1989-5}
\begin{enumerate}
\item With $W(x) = \frac{1}{n+1}\sum_{j \le n}(\frac n2 - j)x^j$ and $N(x) = \sum_{j \le n}x^j$, $(1 - x)W(x) = 1 - \frac{N(x)}{n + 1} + \frac{n}{2(n+1)}(x^{n+1} - 1)$ in $\mathbf Q[x]$; so $(1 - \lambda)W(\lambda) = 1 - N(\lambda)/(n + 1)$ for $\lambda^{n+1} = 1$ in any $\mathbf Q$-algebra.
\item In a ring where $2$ is invertible, if $\lambda^n = 1$ with $n$ odd, then $1 + \lambda$ is a unit with inverse $\frac12\sum_{i<n}(-\lambda)^i$.
\item For ring homomorphisms $u, v : A \to R$, $q = u - v$ and $\rho = u + v$: $q(xy) = q(x)u(y) + v(x)q(y) = u(x)q(y) + q(x)v(y)$, $2q(xy) = \rho(x)q(y) + q(x)\rho(y)$ and $q(xy) = u(x)q(y) + q(x)u(y) - q(x)q(y)$ (the page's formulas for $(u - v)/2$, $(u + v)/2$ multiplied out).
\item If $1 + X$ and $1 - X$ are units and $g = (1 + X)(1 - X)^{-1}$, then $a - g^{-1}ag = (1 + X)^{-1}\cdot 2(Xa - aX)\cdot(1 - X)^{-1}$.
\end{enumerate}
\end{theoreme}
\end{encadre}

\begin{proof}
(1) $W = W' - \frac{n}{2(n+1)}N$ with $(n + 1)W'(x) = \sum(n - j)x^j$, and $(1 - x)\sum(n - j)x^j = n + 1 - N(x)$ by induction; $(1 - x)N(x) = 1 - x^{n+1}$. (2) $(1 + \lambda)\sum_{i<n}(-\lambda)^i = 1 - (-\lambda)^n = 2$. (3) Expand. (4) $1 + X$ and $1 - X$ commute; multiply by $1 + X$ on the left and $1 - X$ on the right.
\end{proof}

\emph{Verification in Lean.} The file \href{\leanblobmain/Notebook1989_5.lean}{\texttt{Notebook1989\_5.\allowbreak{}lean}} proves \texttt{one\_sub\_mul\_S}, \texttt{one\_sub\_mul\_W'}, \texttt{W\_eq}, \texttt{one\_sub\_mul\_W}, \texttt{one\_sub\_mul\_aeval\_W}, \texttt{one\_add\_mul\_inv}, \texttt{one\_add\_inv\_mul}, \texttt{q\_mul\_left}, \texttt{q\_mul\_right}, \texttt{q\_mul\_}$\rho$, \texttt{q\_mul} and \texttt{sub\_conj}, the last ones in noncommutative rings and without dividing by $2$. These are short identities, collected because the pages rely on them. The axioms used are the three usual axioms of mathlib.

\litterature{Classical or elementary (Cuntz's homotopy polynomial; Cayley transform); the edition does not present them as new.}

\section[Notebook 1990-2: the exact order of the corrected Karoubi operator]{Notebook 1990-2: the exact order of $\tilde\kappa$; the first Green operator\protect\verifie{Notebook1990_2}}\label{sec:quillen-1990-2}

\emph{Context.} On pages~2--3 of \carnet{1990}{2} Quillen corrects the Karoubi operator $\kappa$, which satisfies $\kappa^{n(n+1)} = 1 + N$ with $N^2 = 0$ and $\kappa N = N$, to $\tilde\kappa = \kappa - cN$ of finite order; on page~99 he writes a first Green operator for $1 - \sigma$.

\begin{encadre}
\begin{theoreme}[Quillen, notebook 1990-2, pp.~2--3, 99]\label{thm:quillen-1990-2}
\begin{enumerate}
\item In a $\mathbf Q$-algebra, let $\kappa$, $N$ commute, $N^2 = 0$, $\kappa N = N$. Then $(\kappa - cN)^j = \kappa^j - jcN$ for all $j$ and $c \in \mathbf Q$. If $\kappa^m = 1 + N$ and $N \ne 0$, then $(\kappa - cN)^m = 1$ iff $mc = 1$; in particular $\tilde\kappa = \kappa - N/(n(n+1))$ has $\tilde\kappa^{n(n+1)} = 1$.
\item In any ring, $(1 - \sigma)\sum_{j<n}j\sigma^{n-1-j} = n - \sum_{i<n}\sigma^i$.
\end{enumerate}
\end{theoreme}
\end{encadre}

\begin{proof}
(1) By induction, since $\kappa^jN = N\kappa = N$ and $N^2 = 0$; then $(\kappa - cN)^m = 1 + (1 - mc)N$. (2) By induction on $n$.
\end{proof}

\emph{Verification in Lean.} The file \href{\leanblobmain/Notebook1990_2.lean}{\texttt{Notebook1990\_2.\allowbreak{}lean}} proves \texttt{pow\_sub\_smul}, \texttt{pow\_eq\_one\_iff}, \texttt{pow\_eq\_one} and \texttt{one\_sub\_mul\_green}. The canonical Green operator of p.~11 is in \S\ref{sec:quillen-1990-3}, its independence of $n$ in \S\ref{sec:quillen-1990-2-green}. The axioms used are the three usual axioms of mathlib.

\litterature{Elementary; the edition does not present it as new.}

\section[Notebook 1990-2: the Green operator does not depend on the period]{Notebook 1990-2: the canonical Green operator does not depend on the period\protect\verifie{Notebook1990_2Green}}\label{sec:quillen-1990-2-green}

\emph{Context.} On page~14 of \carnet{1990}{2} Quillen checks that the projection $P$ and the Green operator $Q$ of $1 - \kappa$ (\S\ref{sec:quillen-1990-3}) do not depend on the period $n$ chosen with $\kappa^n = 1$.

\begin{encadre}
\begin{theoreme}[Quillen, notebook 1990-2, p.~14]\label{thm:quillen-1990-2-green}
In a $\mathbf Q$-algebra, put $P_n = \frac1n\sum_{j<n}\kappa^j$ and $Q_n = \frac1n\sum_{j<n}(\frac{n-1}2 - j)\kappa^j$. If $\kappa^m = 1$ and $k \ge 1$, then $P_{km} = P_m$ and $Q_{km} = Q_m$.
\end{theoreme}
\end{encadre}

\begin{proof}
Write $j < km$ as $j = qm + r$ with $q < k$, $r < m$; then $\kappa^j = \kappa^r$, and $\sum_{q<k}(\frac{km-1}2 - qm - r) = k(\frac{m-1}2 - r)$.
\end{proof}

\emph{Verification in Lean.} The file \href{\leanblobmain/Notebook1990_2Green.lean}{\texttt{Notebook1990\_2Green.\allowbreak{}lean}} proves \texttt{proj\_mul} and \texttt{green\_mul}, for the operators of \texttt{Notebook1990\_3}, through \texttt{sum\_range\_mul\_smul}. The axioms used are the three usual axioms of mathlib.

\litterature{Elementary; the edition does not present it as new.}


\part{Noncommutative algebras: traces, extensions and separability}
\section[Notebook 1988-1: $\mathrm{Hom}_k(M, A)$ modulo commutators is the dual of $M$]{Notebook 1988-1: $\mathrm{Hom}_k(M, A)$ modulo commutators is the dual of $M$\protect\verifie{Notebook1988_1}}\label{sec:quillen-1988-1}

\emph{Context.} In February 1988, \carnet{1988}{1} turns from the cyclic homology of algebra extensions to residues and duality. Quillen wants to reconstruct Cartier's construction of the Grothendieck residue symbol and Tate's one-dimensional residue through Hochschild cochains. For an $A$-module $M$ that is finite-dimensional over the ground field $k$, he identifies $\mathrm{Ext}^1_A(M, A)$ with derivations of $A$ into the bimodule $\mathrm{Hom}_k(M, A)$ modulo inner ones. To land in linear functionals on $M$, he needs to know the commutator quotient of that bimodule.

\emph{The page.} On pages~71--72 (entry of 29 February 1988), Quillen computes $A \otimes_{A \otimes A} \mathrm{Hom}_k(M, A) = \mathrm{Hom}_k(M, A)/[A, \mathrm{Hom}_k(M, A)]$. He chooses a basis $m_j$ of $M$ with dual basis $\mu_j$, writes $\lambda = \sum_j \mu_j \otimes \lambda(m_j)$, and moves the coefficient $\lambda(m_j)$ from the right of the tensor sign to the left modulo commutators. The result is the functional $m \mapsto \sum_j \mu_j(\lambda(m_j) m)$, which he recognises as the trace of $M \xrightarrow{\lambda} A \xrightarrow{\cdot m} M$. He states it as a lemma: the map $\lambda \mapsto (m \mapsto \mathrm{tr}_M(x \mapsto \lambda(x) m))$ induces an isomorphism $\mathrm{Hom}_k(M, A)/[A, \mathrm{Hom}_k(M, A)] \cong M^*$. The following pages build from it the maps $\mathrm{Ext}^n_A(M, A) \to \mathrm{Hom}_A(\Omega^n_A, M^*)$.

\emph{Intuition.} When $M$ is finite-dimensional, $\mathrm{Hom}_k(M, A) = M^* \otimes_k A$. Dividing by commutators glues the right action of $A$ (through $M$) to the left action (on the factor $A$), which is $M^* \otimes_A A = M^*$. The trace is the intrinsic way to write this identification without choosing a basis.

\begin{encadre}
\begin{theoreme}[Quillen, notebook 1988-1, pp.~71--72]\label{thm:quillen-1988-1}
Let $k$ be a field, $A$ a $k$-algebra, and $M$ a left $A$-module that is finite-dimensional over $k$. Make $\mathrm{Hom}_k(M, A)$ an $A$-bimodule by $(a\lambda)(x) = a\lambda(x)$ and $(\lambda a)(x) = \lambda(ax)$, and let $[A, \mathrm{Hom}_k(M, A)]$ be the $k$-span of the $a\lambda - \lambda a$. The $k$-linear map
\[
  \Phi : \mathrm{Hom}_k(M, A) \to M^*, \qquad \Phi(\lambda)(m) = \mathrm{tr}_M\bigl(x \mapsto \lambda(x)\, m\bigr),
\]
is surjective with kernel $[A, \mathrm{Hom}_k(M, A)]$. Hence $\mathrm{Hom}_k(M, A)/[A, \mathrm{Hom}_k(M, A)] \cong M^*$.
\end{theoreme}
\end{encadre}

\noindent The page assumes $A$ commutative, as everywhere in the residue entries. The argument does not use it, and the statement here is for any $k$-algebra.

\begin{proof}
Write $T_{\lambda, m}$ for $x \mapsto \lambda(x) m$ and $L_a$ for $x \mapsto ax$. Then $T_{a\lambda, m} = L_a \circ T_{\lambda, m}$ and $T_{\lambda a, m} = T_{\lambda, m} \circ L_a$, so $\Phi(a\lambda) = \Phi(\lambda a)$ because $\mathrm{tr}(fg) = \mathrm{tr}(gf)$. For $\mu \in M^*$ let $\iota(\mu)(x) = \mu(x) \cdot 1$. Then $T_{\iota(\mu), m}$ is the rank-one map $x \mapsto \mu(x) m$, whose trace is $\mu(m)$, so $\Phi \circ \iota = \mathrm{id}$ and $\Phi$ is surjective. Finally let $\lambda \in \mathrm{Hom}_k(M, A)$, choose a basis $m_j$ with dual basis $\mu_j$, and put $c_j = \lambda(m_j)$. Then $\lambda = \sum_j c_j\, \iota(\mu_j)$, and $\iota(\mu_j)\, c_j = \iota(\mu_j \circ L_{c_j})$. Applying $\Phi$ and the first step gives $\Phi(\lambda) = \sum_j \mu_j \circ L_{c_j}$. So
\[
  \lambda - \iota(\Phi(\lambda)) = \sum_j \bigl(c_j\, \iota(\mu_j) - \iota(\mu_j)\, c_j\bigr) \in [A, \mathrm{Hom}_k(M, A)],
\]
and if $\Phi(\lambda) = 0$ then $\lambda$ lies in the commutator subspace.
\end{proof}

\noindent The proof is the one on the page. The page computes the quotient through $M^* \otimes_A A$ and identifies the result with the trace. Here the trace is the definition and the basis appears only in the last step.

\emph{Verification in Lean.} The file \href{\leanblobmain/Notebook1988_1.lean}{\texttt{Notebook1988\_1.\allowbreak{}lean}} defines $\Phi$ as \texttt{traceForm}, the two actions as \texttt{leftMul} and \texttt{rightMul}, the commutator subspace as \texttt{commutators}, and $\iota$ as \texttt{embed}. The first step is \texttt{traceForm\_leftMul}, the second \texttt{traceForm\_embed} and \texttt{traceForm\_surjective}, and the decomposition is \texttt{sub\_embed\_traceForm\_mem}. Theorem~\ref{thm:quillen-1988-1} is \texttt{ker\_traceForm} together with the isomorphism \texttt{quotientEquivDual}. The axioms used are the three usual axioms of mathlib.

\litterature{Lipman, \q{Residues and traces of differential forms via Hochschild homology}, \emph{Contemp. Math.}~61 (1987), Zbl~0606.14015, builds the residue from trace maps of this kind; the search did not locate the lemma there in this form. The identification $\mathrm{Hom}_k(M, A) \cong M^* \otimes_k A$ and the computation of $N \otimes_{A \otimes A^{o}} -$ on it are elementary, and the edition does not present the statement as new.}

\section[Notebooks 1988-2 and 1988-4: the GNS product; a surjective $\partial$]{Notebooks 1988-2 and 1988-4: associativity of the GNS product; a surjective $\partial$ is injective\protect\verifie{Notebook1988_2}}\label{sec:quillen-1988-2}

\emph{Context.} To a linear map $\rho : A \to B$ of algebras, page~62 of \carnet{1988}{2} and page~6 of \carnet{1988}{4} attach the product on $A \otimes B \otimes A$ of a Stinespring (GNS) construction, and use that it is associative; it is the same statement in both notebooks, proved once here. On page~45 of 1988-2 a candidate differential graded algebra is ruled out by a two-line argument.

\begin{encadre}
\begin{theoreme}[Quillen, notebook 1988-2, pp.~45, 62; notebook 1988-4, p.~6]\label{thm:quillen-1988-2}
\begin{enumerate}
\item Let $R$ be a commutative ring, $A$, $B$ two $R$-algebras and $\rho : A \to B$ an $R$-linear map. The bilinear product on $A \otimes_R B \otimes_R A$ with $(a_1 \otimes b_1 \otimes a_1')(a_2 \otimes b_2 \otimes a_2') = a_1 \otimes b_1\rho(a_1'a_2)b_2 \otimes a_2'$ is associative.
\item Let $P$ carry a left and a right action of a ring $D$, and $\partial : P \to D$ satisfy $\partial(x)\,y = x\,\partial(y)$. If $\partial$ is onto, it is one-to-one.
\end{enumerate}
\end{theoreme}
\end{encadre}

\begin{proof}
(1) Both sides are trilinear, so it suffices to compare them on pure tensors, where both are $a_1 \otimes b_1\rho(a_1'a_2)b_2\rho(a_2'a_3)b_3 \otimes a_3'$. (2) With $\partial y = 1$, $x = x\,\partial(y) = \partial(x)\,y$, so $x$ is determined by $\partial x$.
\end{proof}

\emph{Verification in Lean.} The file \href{\leanblobmain/Notebook1988_2.lean}{\texttt{Notebook1988\_2.\allowbreak{}lean}} defines the product \texttt{gnsMul} from tensor-product maps (\texttt{gnsMul\_tmul} on pure tensors) and proves (1) as \texttt{gnsMul\_assoc}, by induction on the three tensors, and (2) as \texttt{injective\_of\_surjective}. Neither multiplicativity nor unitality of $\rho$ is used, nor additivity of $\partial$. Positivity and the Hilbert-space completion are not formalised. The axioms used are the three usual axioms of mathlib.

\litterature{Elementary; the product is the algebraic core of Stinespring's dilation (and of the bimodule $A \otimes B \otimes A$ in Cuntz--Quillen's work on extensions). Not located as a stated lemma. The edition does not present it as new.}

\section[Notebook 1988-4: Cuntz--Krieger matrix units]{Notebook 1988-4: Cuntz--Krieger matrix units; the $\mathcal O_2$ relations as a unitary block\protect\verifie{Notebook1988_4}}\label{sec:quillen-1988-4}

\emph{Context.} Reading Cuntz and Krieger, Quillen writes on page~59 of \carnet{1988}{4} the matrix units of a Cuntz--Krieger algebra, and on page~60 the relations of $\mathcal O_2$ in matrix form.

\begin{encadre}
\begin{theoreme}[Quillen, notebook 1988-4, pp.~59--60]\label{thm:quillen-1988-4}
In a ring with involution:
\begin{enumerate}
\item let $(s_i)$ satisfy $s_j^*s_l = 0$ for $j \ne l$, let $P_k = s_ks_k^*$ be pairwise orthogonal idempotents, $Q_j = s_j^*s_j$, and $P_kQ_j = P_k$ when $(j, k)$ is an edge of a graph $\Gamma$. Then for $(j, k) \in \Gamma$, $(s_iP_ks_j^*)(s_lP_ns_m^*) = \delta_{jl}\delta_{kn}\,s_iP_ks_m^*$;
\item if $s_0^*s_0 = s_1^*s_1 = 1$, $s_0^*s_1 = 0$ and $s_0s_0^* + s_1s_1^* = 1$, the row $(s_0\ s_1)$ and the column $(s_0^*\ s_1^*)^t$ are inverse to each other.
\end{enumerate}
\end{theoreme}
\end{encadre}

\begin{proof}
(1) The middle factor $s_j^*s_l$ is $0$ for $j \ne l$, and $Q_j$ for $j = l$; then $P_kQ_jP_n = P_kP_n = \delta_{kn}P_k$. (2) The product row times column is $s_0s_0^* + s_1s_1^*$; the product column times row has entries $s_i^*s_j$, where $s_1^*s_0 = (s_0^*s_1)^* = 0$.
\end{proof}

\emph{Verification in Lean.} The file \href{\leanblobmain/Notebook1988_4.lean}{\texttt{Notebook1988\_4.\allowbreak{}lean}} proves \texttt{matrixUnit\_mul} (only the edge condition of the left factor is used, and no property of the involution), \texttt{row\_mul\_col} and \texttt{col\_mul\_row}. The $C^*$-algebras themselves are not formalised. The axioms used are the three usual axioms of mathlib.

\litterature{Classical (Cuntz--Krieger, \emph{Invent. Math.} 56 (1980); Cuntz, \emph{Comm. Math. Phys.} 57 (1977)); the edition does not present it as new.}

\section[Notebook 1989-4: powers to first order]{Notebook 1989-4: powers to first order; graded commutators; the modular operator\protect\verifie{Notebook1989_4}}\label{sec:quillen-1989-4}

\emph{Context.} On page~38 of \carnet{1989}{4} Quillen expands powers of $a - da\,d$ to first order; on pages~14 and~18 he uses the graded commutator $[z, \theta] = z\theta + \theta z$, and on page~34 the modular operator of a faithful trace functional on matrices.

\begin{encadre}
\begin{theoreme}[Quillen, notebook 1989-4, pp.~14, 18, 34, 38]\label{thm:quillen-1989-4}
In any ring:
\begin{enumerate}
\item if $exe = 0$ for every $x$, then $(a + e)^{n+1} = a^{n+1} + \sum_{k=0}^na^{n-k}ea^k$;
\item with $X = z\theta + \theta z$, $zX - Xz = z^2\theta - \theta z^2$; if $z^2 = 1$, $zX = Xz = \theta + z\theta z$;
\item for an invertible matrix $T$ over a commutative ring, $\operatorname{tr}(Tyx) = \operatorname{tr}(Tx\,TyT^{-1})$.
\end{enumerate}
\end{theoreme}
\end{encadre}

\begin{proof}
(1) By induction: $(a^n + S_n)(a + e) = a^{n+1} + a^ne + S_na$, since $S_ne = 0$. (2) Expansion. (3) Cyclicity of the trace.
\end{proof}

\emph{Verification in Lean.} The file \href{\leanblobmain/Notebook1989_4.lean}{\texttt{Notebook1989\_4.\allowbreak{}lean}} proves \texttt{add\_pow\_first\_order}, \texttt{commutator\_graded}, \texttt{mul\_graded\_of\_sq} and \texttt{trace\_modular}. The page applies (1) to $e = -da\,d$. The axioms used are the three usual axioms of mathlib.

\litterature{Elementary; the edition does not present it as new.}

\section[Notebook 1989-4: powers in a square-zero extension]{Notebook 1989-4: powers in the square-zero extension of a cup product\protect\verifie{Notebook1989_4Square}}\label{sec:quillen-1989-4-square}

\emph{Context.} On pages~53--54 of \carnet{1989}{4} Quillen forms, from two derivations $X$, $Y$, the square-zero extension with product $a_1 * a_2 = a_1a_2 + Xa_1\,Ya_2$, and computes the powers $a^{*n}$ (boxed).

\begin{encadre}
\begin{theoreme}[Quillen, notebook 1989-4, p.~54]\label{thm:quillen-1989-4-square}
Let $E$ be a ring, $X$, $Y$ additive maps $E \to E$ satisfying the Leibniz rule, and on $E \times E$ put $(a, m)*(b, n) = (ab, an + mb + Xa\,Yb)$. This product is associative with unit $(1, 0)$, and
\[
(a, 0)^{*(n+2)} = \Bigl(a^{n+2},\ \sum_{i+j+k=n}a^i(Xa)a^j(Ya)a^k\Bigr).
\]
\end{theoreme}
\end{encadre}

\begin{proof}
Associativity is the Leibniz rule for $X$ and $Y$. Then $Y(a^{m+1}) = \sum_{j+k=m}a^j(Ya)a^k$, and $(a, 0)*(a^{n+2}, S_n) = (a^{n+3}, aS_n + Xa\,Y(a^{n+2}))$; split $\sum_{i+j+k=n+1}$ by $i = 0$ or $i \ge 1$.
\end{proof}

\emph{Verification in Lean.} The file \href{\leanblobmain/Notebook1989_4Square.lean}{\texttt{Notebook1989\_4Square.\allowbreak{}lean}} proves \texttt{star\_assoc}, \texttt{star\_one}, \texttt{one\_star}, \texttt{Y\_pow}, \texttt{starPow\_fst} and \texttt{starPow\_snd}. On the page the second components lie in $M \otimes_A N$; here in $E$, which is enough for the identity, the page itself noting that one may suppose the cocycle a cup product of two derivations. The exponential version of p.~55 is not formalised. The axioms used are the three usual axioms of mathlib.

\litterature{Elementary; the edition does not present it as new.}

\section[Notebook 1990-4: separability makes derivations inner]{Notebook 1990-4: a separability element makes derivations inner\protect\verifie{Notebook1990_4}}\label{sec:quillen-1990-4}

\emph{Context.} On pages~13--14 of \carnet{1990}{4} Quillen shows, from a separability element $e = \sum_ix_i \otimes y_i$, that every derivation is inner.

\begin{encadre}
\begin{theoreme}[Quillen, notebook 1990-4, pp.~13--14]\label{thm:quillen-1990-4}
Let $A$ be an algebra over a commutative ring $k$, $e = \sum_ix_i\otimes y_i \in A \otimes_k A$ with $\sum_ix_iy_i = 1$ and $ae = ea$ for all $a$, and $D : A \to A$ a $k$-linear map with $D(ab) = aD(b) + D(a)b$. Then $D(a) = ac - ca$ with $c = \sum_ix_iD(y_i)$.
\end{theoreme}
\end{encadre}

\begin{proof}
Apply $x \otimes y \mapsto xD(y)$ to $ae = ea$: $\sum_iax_iD(y_i) = \sum_ix_iD(y_ia) = \sum_ix_iy_iD(a) + \sum_ix_iD(y_i)a = D(a) + ca$.
\end{proof}

\emph{Verification in Lean.} The file \href{\leanblobmain/Notebook1990_4.lean}{\texttt{Notebook1990\_4.\allowbreak{}lean}} proves \texttt{eq\_inner}, for an additive Leibniz map of any ring with the image identity as hypothesis, and \texttt{eq\_inner\_of\_comm}, from $ae = ea$ in $A \otimes_k A$. The matrix separability element is \S\ref{sec:quillen-1991-4-separable}. The axioms used are the three usual axioms of mathlib.

\litterature{Classical (separable algebras have vanishing Hochschild cohomology in degree one); the edition does not present it as new.}

\section[Notebook 1990-5: $\mathrm{tr}(g^{-1}dg)$ modulo commutators]{Notebook 1990-5: $\operatorname{tr}(g^{-1}dg)$ is additive modulo commutators\protect\verifie{Notebook1990_5Trace}}\label{sec:quillen-1990-5-trace}

\emph{Context.} On page~68 of \carnet{1990}{5} Quillen uses that $g \mapsto \operatorname{tr}(g^{-1}dg)$ is a homomorphism modulo commutators; on page~48, the derivation attached to a bimodule extension.

\begin{encadre}
\begin{theoreme}[Quillen, notebook 1990-5, pp.~48, 68]\label{thm:quillen-1990-5-trace}
\begin{enumerate}
\item In a ring with a derivation $\delta$, for units $u$, $v$: $(uv)^{-1}\delta(uv) = v^{-1}(u^{-1}\delta u)v + v^{-1}\delta v$.
\item For invertible $g_1$, $g_2 \in M_n(E)$, $E$ a ring with a derivation $d$ applied entrywise, $\operatorname{tr}((g_1g_2)^{-1}d(g_1g_2)) - \operatorname{tr}(g_1^{-1}dg_1) - \operatorname{tr}(g_2^{-1}dg_2)$ lies in the additive subgroup $[E, E]$ spanned by the $ab - ba$.
\item For a map of $R$-bimodules $\pi : M \to R$ and $\xi$ with $\pi(\xi) = 1$, $D(x) = x\xi - \xi x$ takes values in $\operatorname{Ker}\pi$ and $D(xy) = xD(y) + D(x)y$.
\end{enumerate}
\end{theoreme}
\end{encadre}

\begin{proof}
(1) Leibniz. (2) By (1) in $M_n(E)$, the difference is $\operatorname{tr}(PQ) - \operatorname{tr}(QP)$ with $P = g_2^{-1}$, $Q = g_1^{-1}dg_1\,g_2$, which is $\sum_{i,j}(P_{ij}Q_{ji} - Q_{ji}P_{ij})$. (3) Expansion.
\end{proof}

\emph{Verification in Lean.} The file \href{\leanblobmain/Notebook1990_5Trace.lean}{\texttt{Notebook1990\_5Trace.\allowbreak{}lean}} proves \texttt{logDeriv\_mul}, \texttt{trace\_mul\_sub\_mem}, \texttt{matDeriv\_mul}, \texttt{trace\_logDeriv\_mul} and \texttt{derivation\_of\_lift}. The derivation is ungraded: on the page $d$ is the differential of $\Omega E$ on elements of degree~$0$, which is all that enters. The axioms used are the three usual axioms of mathlib.

\litterature{Classical (the logarithmic derivative; $K_1 \to HC_0$ of $\Omega^1$); the edition does not present it as new.}

\section[Notebook 1991-4: the separability element of $M_n(K)$]{Notebook 1991-4: the separability element of $M_n(K)$\protect\verifie{Notebook1991_4Separable}}\label{sec:quillen-1991-4-separable}

\emph{Context.} On pages~51 to~58 of \carnet{1991}{4} Quillen works with separable algebras through an element $z = \sum x_i \otimes y_i$ of $A \otimes A$ that commutes with $A$ for the bimodule structure, $(a \otimes 1)z = z(1 \otimes a)$, and with $\sum x_iy_i = 1$. For the matrix algebra he takes $z = \frac1n\sum_{i,j} e_{ij} \otimes e_{ji}$, which is moreover symmetric.

\begin{encadre}
\begin{theoreme}[Quillen, notebook 1991-4, pp.~51--58]\label{thm:quillen-1991-4-separable}
Let $K$ be a commutative ring and $s = \sum_{i,j} e_{ij} \otimes e_{ji}$ in $M_n(K) \otimes_K M_n(K)$. Then $(a \otimes 1)s = s(1 \otimes a)$ for every matrix $a$, $\sum_{i,j} e_{ij}e_{ji} = n \cdot 1$, $s$ is fixed by the flip and $s^2 = 1$. If $n$ is invertible in $K$, $z = n^{-1}s$ satisfies $(a \otimes 1)z = z(1 \otimes a)$, $m(z) = 1$ and is symmetric.
\end{theoreme}
\end{encadre}

\begin{proof}
Both sides of the first identity are linear in $a$, and for $a = e_{kl}$ both equal $\sum_j e_{kj} \otimes e_{jl}$, by $e_{ij}e_{kl} = \delta_{jk}e_{il}$. Then $\sum_{i,j}e_{ij}e_{ji} = \sum_{i,j}e_{ii} = n \cdot 1$; the flip exchanges the two indices of summation; and $(e_{ij} \otimes e_{ji})s = e_{ii} \otimes e_{jj}$, whose sum is $1 \otimes 1$.
\end{proof}

\emph{Verification in Lean.} The file \href{\leanblobmain/Notebook1991_4Separable.lean}{\texttt{Notebook1991\_4Separable.\allowbreak{}lean}} proves the statement as \texttt{tmul\_one\_mul\_s}, \texttt{mul\_s}, \texttt{comm\_s}, \texttt{s\_mul\_s} and \texttt{separability\_idempotent}, for any finite index type, with mathlib's \texttt{Matrix.single} and \texttt{Algebra.TensorProduct}. The trace form of the page and the group-algebra analogue are not formalised. The axioms used are the three usual axioms of mathlib.

\litterature{Classical (DeMeyer and Ingraham, \emph{Separable algebras over commutative rings}, LNM~181, 1971, ch.~II). Mathlib has separable field extensions, not separability idempotents. The edition does not present it as new.}


\part{Linear algebra, matrices and determinants}
\section[1972 mixed papers: alternating perfect pairings force even dimension]{1972 mixed papers: a nondegenerate alternating form forces even dimension, in every characteristic\protect\verifie{Notebook1972_MixedAlt}}\label{sec:quillen-1972-mixed-alt}

\emph{Context.} On page~41 of the \href{https://www.claymath.org/library/Quillen/Working_papers/quillen\%201972/1972-mixed\%20papers.pdf}{file 1972-mixed papers} Quillen notes that the cup-product pairing on $H^1(\pi, \mathbf Z/p)$ of the Galois group of a $p$-adic field is a perfect alternating pairing, so that the dimension is even, which he checks against $[\mathbf Q_p(\mu_p) : \mathbf Q_p] = p - 1$. The statement below holds in every characteristic.

\begin{encadre}
\begin{theoreme}[Quillen, 1972 mixed papers, p.~41]\label{thm:quillen-1972-mixed-alt}
A nondegenerate alternating bilinear form on a finite-dimensional vector space over a field has even dimension. More generally, an alternating matrix (zero diagonal, $M^T = -M$) of odd size has determinant zero over any commutative ring.
\end{theoreme}
\end{encadre}

\begin{proof}
For the generic alternating matrix $G$ of odd size $n$ over $\mathbf Z[X_{ij}]_{i<j}$, $\det G = \det G^T = \det(-G) = -\det G$, so $2\det G = 0$ and $\det G = 0$ since $\mathbf Z[X_{ij}]$ has no $2$-torsion. Any alternating matrix is the image of $G$ under a ring homomorphism, which commutes with $\det$. The Gram matrix of a nondegenerate alternating form has nonzero determinant.
\end{proof}

\emph{Verification in Lean.} The file \href{\leanblobmain/Notebook1972_MixedAlt.lean}{\texttt{Notebook1972\_MixedAlt.\allowbreak{}lean}} proves \texttt{det\_generic\_eq\_zero}, \texttt{det\_eq\_zero\_of\_alternating} and \texttt{even\_finrank}. The argument by universality avoids the symplectic-basis theorem and the Pfaffian, neither of which the pinned mathlib has; the Galois-cohomological pairing is not formalised. The axioms used are the three usual axioms of mathlib.

\litterature{Classical (Bourbaki, \emph{Algèbre}~IX, symplectic bases); the edition does not present it as new.}

\section[Notebook 1976-1: $\det(1 + t\nu)$ a unit iff $\nu^2 = 0$]{Notebook 1976-1: for $2 \times 2$ matrices, $\det(1 + t\nu) \in K^\times$ iff $\nu^2 = 0$ iff $\nu$ is nilpotent\protect\verifie{Notebook1976_1}}\label{sec:quillen-1976-1}

\emph{Context.} On page~49 of \carnet{1976}{1} Quillen identifies the unimodular subspaces of degree at most one in $\mathbf C[t]^2$ with the matrices $1 + t\nu$, $\nu$ nilpotent, and notes that the cone of nilpotent $2 \times 2$ matrices has dimension~$2$.

\begin{encadre}
\begin{theoreme}[Quillen, notebook 1976-1, p.~49]\label{thm:quillen-1976-1}
Let $\nu$ be a $2 \times 2$ matrix over a field $K$. Then $\det(1 + t\nu) = 1 + t\operatorname{tr}\nu + t^2\det\nu$ in $K[t]$, and the following are equivalent: $\det(1 + t\nu)$ is a unit of $K[t]$; $\operatorname{tr}\nu = \det\nu = 0$; $\nu^2 = 0$; $\nu$ is nilpotent.
\end{theoreme}
\end{encadre}

\begin{proof}
The units of $K[t]$ are the nonzero constants. Cayley--Hamilton reads $\nu^2 = (\operatorname{tr}\nu)\nu - (\det\nu)1$. If $\nu^k = 0$, then $(\det\nu)^k = 0$, so $\nu^2 = (\operatorname{tr}\nu)\nu$, $\nu^k = (\operatorname{tr}\nu)^{k-1}\nu$, hence $\operatorname{tr}\nu = 0$ or $\nu = 0$, and $\nu^2 = 0$. If $\nu^2 = 0$ and $\operatorname{tr}\nu \ne 0$, Cayley--Hamilton gives $\nu = 0$, a contradiction.
\end{proof}

\emph{Verification in Lean.} The file \href{\leanblobmain/Notebook1976_1.lean}{\texttt{Notebook1976\_1.\allowbreak{}lean}} proves \texttt{det\_one\_add}, \texttt{sq\_eq}, \texttt{sq\_eq\_zero\_iff}, \texttt{isUnit\_det\_iff} and \texttt{isNilpotent\_iff}, in every characteristic. The page's Quot-scheme context and higher-degree families are not formalised. The axioms used are the three usual axioms of mathlib.

\litterature{Elementary; the edition does not present it as new.}

\section[Notebooks 1982-2 and 1982-6: the Faddeev--LeVerrier recursion]{Notebooks 1982-2 and 1982-6: the Faddeev--LeVerrier recursion, with Jacobi's formula\protect\verifie{Notebook1982_6}}\label{sec:quillen-1982-6}

\emph{Context.} On page~22 of \carnet{1982}{2} and page~17 of \carnet{1982}{6} Quillen expands the Fredholm determinant $\det(1 - tK)$ and the adjugate of $1 - tK$ in powers of $t$ and uses the recursions linking the coefficients (Plemelj--Smithies formulas). Here the finite-dimensional form is proved, over any commutative ring.

\begin{encadre}
\begin{theoreme}[Quillen, notebooks 1982-2, p.~22, and 1982-6, p.~17]\label{thm:quillen-1982-6}
Let $K$ be a square matrix over a commutative ring, $\det(1 - tK) = \sum_k \Delta_k t^k$ and $\operatorname{adj}(1 - tK) = \sum_k C_kt^k$. Then $\Delta_0 = 1$, $C_0 = 1$ and, for $k \ge 0$,
\[
C_{k+1} = KC_k + \Delta_{k+1}\cdot 1, \qquad (k + 1)\,\Delta_{k+1} = -\operatorname{tr}(KC_k).
\]
More generally (Jacobi's formula), for a square matrix $M$ over $R[t]$, $(\det M)' = \operatorname{tr}(\operatorname{adj}(M)\,M')$.
\end{theoreme}
\end{encadre}

\begin{proof}
Jacobi's formula: differentiate the Leibniz expansion of $\det M$ by the product rule; the term in which column $i$ is differentiated is the determinant of $M$ with column $i$ replaced by its derivative, which by Cramer's rule is $\sum_j \operatorname{adj}(M)_{ij} M'_{ji}$. The first recursion is the coefficient of $t^{k+1}$ in $(1 - tK)\operatorname{adj}(1 - tK) = \det(1 - tK)\cdot 1$; the second is the coefficient of $t^k$ in Jacobi's formula for $M = 1 - tK$, $M' = -K$.
\end{proof}

\emph{Verification in Lean.} The file \href{\leanblobmain/Notebook1982_6.lean}{\texttt{Notebook1982\_6.\allowbreak{}lean}} proves Jacobi's formula as \texttt{derivative\_det} and the recursions as \texttt{Delta\_zero}, \texttt{Cmat\_zero}, \texttt{Cmat\_succ} and \texttt{succ\_mul\_Delta\_succ}. Mathlib has only the derivative of $\det(1 + tM)$ at $t = 0$ (\texttt{Matrix.derivative\_det\_one\_add\_X\_smul}). The Fredholm (trace-class) version is not formalised. The axioms used are the three usual axioms of mathlib.

\litterature{Classical: Le~Verrier (1840), Faddeev and Sominsky; for Fredholm determinants, the Plemelj--Smithies formulas (Smithies, \emph{Integral equations}, 1958; Simon, \emph{Trace ideals}, ch.~5). Not in mathlib. The edition does not present it as new.}

\section[Notebook 1982-7: symmetric matrices with non-positive off-diagonal entries]{Notebook 1982-7: a symmetric matrix with non-positive off-diagonal entries and a positive vector $v$ with $Av \ge 0$ is positive semidefinite\protect\verifie{Notebook1982_7}}\label{sec:quillen-1982-7}

\emph{Context.} \carnet{1982}{7} (11 June -- 12 July 1982, his pages 700--787) works on determinants of $\bar\partial$-operators when the complex structure of the Riemann surface varies, on Riemann--Roch and the Hodge index theorem for surfaces, and on Arakelov's arithmetic surfaces (with notes of a talk by Faltings on 30 June). In the entry of 25 June (pages~32--34) Quillen wants to understand the Hodge index theorem through Mumford's result that the intersection matrix $(C_i \cdot C_j)$ of the exceptional curves of a resolved surface singularity is negative definite. He reduces it to an algebraic situation: a symmetric matrix with entries $\ge 0$ off the diagonal and a combination $D = \sum n_i C_i$ with all $n_i > 0$ and $C_i \cdot D = 0$ for every $i$.

\emph{The page.} On page~33 he passes to $A = -(C_i \cdot C_j)$ and tries an induction on the size, splitting a vector into its positive and negative parts; he notes that subtracting a multiple of $D$ destroys the induction. On page~34 he states the correct assertion: a real symmetric matrix with off-diagonal entries $\le 0$, for which some vector $v$ with strictly positive entries satisfies the hypothesis, is positive semidefinite; after normalising $v = (1, \ldots, 1)$ the proof uses that the row sums are $\ge 0$, then lowers the diagonal entries until the row sums vanish. The hypothesis is written $v^{t}Av \ge 0$; the proof uses $Av \ge 0$, and with $v^{t}Av \ge 0$ alone the assertion is false.

\emph{Intuition.} Write $x_i = v_i y_i$. The quadratic form $x^{t}Ax$ is then $\sum_i (Av)_i v_i y_i^2$ minus half of $\sum_{i,j} A_{ij} v_i v_j (y_i - y_j)^2$: the first sum is $\ge 0$ because $Av \ge 0$, the second is $\le 0$ because only the off-diagonal entries, all $\le 0$, contribute to it.

\begin{encadre}
\begin{theoreme}[Quillen, notebook 1982-7, pp.~33--34]\label{thm:quillen-1982-7}
Let $A$ be a real symmetric $n \times n$ matrix with $A_{ij} \le 0$ for $i \ne j$.
\begin{enumerate}
\item For every vector $v$ with non-zero entries and every $x$,
\[
  x^{t}Ax = \sum_i (Av)_i\,\frac{x_i^2}{v_i} - \frac12 \sum_{i,j} A_{ij}\,v_iv_j\Bigl(\frac{x_i}{v_i} - \frac{x_j}{v_j}\Bigr)^2
\]
(this holds for any symmetric matrix over a field of characteristic $\ne 2$).
\item If some $v$ with $v_i > 0$ for all $i$ satisfies $(Av)_i \ge 0$ for all $i$, then $A$ is positive semidefinite.
\item If moreover $Av = 0$ and $x^{t}Ax = 0$, then $x_i/v_i = x_j/v_j$ whenever $A_{ij} < 0$; in particular, when the graph of the non-zero off-diagonal entries is connected, $x$ is a multiple of $v$.
\item Equivalently, for an intersection matrix $C$ with $C_{ij} \ge 0$ for $i \ne j$ and a vector $D > 0$ with $CD \le 0$, one has $x^{t}Cx \le 0$ for all $x$.
\item The hypothesis $v^{t}Av \ge 0$ does not suffice: $A = \begin{pmatrix} 4 & -1 \\ -1 & -1 \end{pmatrix}$, $v = (1, 1)$ gives $v^{t}Av = 1$, but $(0, 1)A(0, 1)^{t} = -1$.
\end{enumerate}
\end{theoreme}
\end{encadre}

\begin{proof}
For (1), put $y_i = x_i/v_i$. The left side is $\sum_{i,j} A_{ij} v_iv_j y_iy_j$, and $\sum_i (Av)_i x_i^2/v_i = \sum_{i,j} A_{ij} v_iv_j y_i^2$. By symmetry of $A$, $\sum_{i,j} A_{ij} v_iv_j y_j^2$ equals the same sum, so half of $\sum_{i,j} A_{ij} v_iv_j (y_i - y_j)^2$ is $\sum_{i,j} A_{ij}v_iv_j y_i^2 - \sum_{i,j} A_{ij} v_iv_j y_iy_j$, which is the identity. In (2) the first sum is $\ge 0$; in the double sum the diagonal terms vanish and the others are products of $A_{ij} \le 0$ with non-negative factors. In (3) the first sum is $0$, so every term of the double sum, all $\le 0$, vanishes; for $A_{ij} < 0$ this forces $y_i = y_j$. (4) is (2) applied to $-C$. (5) is a computation.
\end{proof}

\noindent The page's argument (induction on the principal submatrices, then lowering the diagonal) is not the one formalised; the identity (1) replaces it. In Quillen's application, $D = \sum n_i C_i$ is the fundamental cycle and $C_i \cdot D = 0$; part (4) gives the negative semidefiniteness he wanted, and part (3) Zariski's form of the equality case.

\emph{Verification in Lean.} The file \href{\leanblobmain/Notebook1982_7.lean}{\texttt{Notebook1982\_7.\allowbreak{}lean}} proves part~(1) as \texttt{quadratic\_identity}, over any field with $2 \ne 0$; part~(2) as \texttt{posSemidef\_of\_mulVec\_nonneg}, with mathlib's \texttt{Matrix.PosSemidef}; part~(3) as \texttt{ratio\_eq\_of\_quadratic\_eq\_zero}; part~(4) as \texttt{quadratic\_nonpos\_of\_mulVec\_eq\_zero}; and part~(5) as \texttt{not\_posSemidef\_example}. The axioms used are at most the three usual axioms of mathlib.

\litterature{the statement is classical. In the theory of $M$-matrices, a real matrix with non-positive off-diagonal entries that admits a positive vector $x$ with $Ax > 0$ is a non-singular $M$-matrix (M.~Fiedler and V.~Pták, \emph{Czechoslovak Math.\ J.} 12 (1962), 382--400); the semidefinite symmetric case is in the standard accounts of $M$-matrices. In the theory of surfaces it is Zariski's lemma on the intersection form of the components of a fibre (W.~Barth, K.~Hulek, C.~Peters and A.~Van de Ven, \emph{Compact complex surfaces}, Springer, 2nd ed.\ 2004, Lemma~III.8.2, as recalled), and the negativity of exceptional configurations is due to D.~Mumford, \emph{Publ.\ Math.\ IHÉS} 9 (1961), 5--22. The page's hypothesis $v^{t}Av \ge 0$ is a slip of the pen, as part~(5) shows; nothing here is presented as new.}

\section[Notebook 1983-10: traces of odd powers of $P - Q$]{Notebook 1983-10: the trace of the odd powers of a difference of two idempotents\protect\piste\protect\verifie{Notebook1983_10}}\label{sec:quillen-1983-10}

\emph{Context.} \carnet{1983}{10} works on index formulas in cyclic cohomology, the theory that Connes was developing at the time and to which Quillen would devote several papers. The index of a Toeplitz operator, or of a graded operator, is written there as the trace of a sufficiently large odd power of an operator built from an idempotent.

\emph{The page.} On page~31 (its page~154), Quillen observes that his two index formulas, the Toeplitz form and the graded form, reduce to the same type of expression: the trace of an odd power of a difference $e_1 - e_2$ of two idempotents. He proposes to study $\operatorname{tr}((e_1 - e_2)^{2n+1})$ and to try to show that it is an integer independent of $n$. The page poses the question; it does not answer it.

\emph{Intuition.} In infinite dimension, $e_1 - e_2$ is not trace class in general, and only a sufficiently large power is; the question is then deep. In finite dimension, all the traces exist, and the expected answer is $\operatorname{tr} e_1 - \operatorname{tr} e_2$, the difference of the ranks when working over a field of characteristic zero. The classical computation puts $A = P - Q$ and $B = P + Q - 1$: these two matrices anticommute and $A^2 + B^2 = 1$ (figure~\ref{fig:quillen-1983-10}). For $m$ odd, $\operatorname{tr}(A^m B^2)$ equals its opposite, by cyclicity of the trace and anticommutation, whence $\operatorname{tr}(A^m) = \operatorname{tr}(A^{m+2})$ \emph{provided one can divide by~$2$}. One might think that characteristic~$2$ is an exception. It is not.


\begin{figure}[htbp]
\centering
\begin{tikzpicture}[scale=1.05, >=stealth, every node/.style={font=\small}]
\draw[black!40, thick] (-2.6,0) -- (2.8,0) node[right] {$\ell_P$};
\draw[black!40, thick] ({-1.9*cos(50)},{-1.9*sin(50)}) -- ({2.9*cos(50)},{2.9*sin(50)}) node[above right=-2pt] {$\ell_Q$};
\draw[frise, dashed] ({-1.1*cos(25)},{-1.1*sin(25)}) -- ({2.8*cos(25)},{2.8*sin(25)}) node[right] {$b$};
\draw[piste, dashed] ({-2.4*cos(-20)},{-2.4*sin(-20)}) -- ({2.8*cos(-20)},{2.8*sin(-20)}) node[right] {$a$};
\draw (0.6,0) arc (0:50:0.6); \node at (13:0.85) {$\varphi$};
\draw[black!45, densely dotted] (-1.2,2.0785) -- (-1.2,0);
\draw[black!45, densely dotted] (-1.2,2.0785) -- (0.5276,0.6287);
\draw[->, thick] (0,0) -- (-1.2,2.0785) node[above] {$v$};
\draw[->, thick] (0,0) -- (-1.2,0) node[above left=-1pt] {$Pv$};
\draw[->, thick] (0,0) -- (0.5276,0.6287) node[above left=-1pt] {$Qv$};
\draw[->, thick, piste] (0,0) -- (-1.7276,-0.6287) node[below left=-2pt] {$Av$};
\draw[->, thick, frise] (0,0) -- (0.5276,-1.4498) node[below] {$Bv$};
\fill (0,0) circle (1.2pt);
\end{tikzpicture}
\caption{Two orthogonal projectors $P$ and $Q$ of the Euclidean plane, onto lines $\ell_P$ and $\ell_Q$ making an angle $\varphi$ (here $50^\circ$). The matrix $A = P - Q$ is $\sin\varphi$ times the orthogonal reflection in the axis $a$, and $B = P + Q - 1$ is $\cos\varphi$ times the reflection in the axis $b$, the bisector of $\ell_P$ and $\ell_Q$. The two axes make an angle of $45^\circ$: the product of the two reflections is a quarter turn, whence $AB = -BA$, and $A^2 + B^2 = (\sin^2\varphi + \cos^2\varphi)\cdot 1 = 1$. Here $\operatorname{tr} A^{2k+1} = 0 = \operatorname{tr} P - \operatorname{tr} Q$.}
\label{fig:quillen-1983-10}
\end{figure}

\begin{encadre}
\begin{theoreme}[the edition's statement; answer, in finite dimension, to the question of notebook 1983-10, p.~31]\label{thm:quillen-1983-10}
Let $R$ be a commutative ring and $P$, $Q$ two idempotent square matrices with entries in $R$. For every integer $k \ge 0$,
\[
\operatorname{tr}\bigl((P - Q)^{2k+1}\bigr) = \operatorname{tr} P - \operatorname{tr} Q.
\]
\end{theoreme}
\end{encadre}

\begin{proof}
Put $C = (P - Q)^2 = P + Q - PQ - QP$. Expanding, with $P^2 = P$ and $Q^2 = Q$,
\[
PC = P - PQP = CP, \qquad QC = Q - QPQ = CQ :
\]
$C$ commutes with $P$ and $Q$, hence so do its powers. It follows that
\[
(P - Q)\,C = (P - Q) - (PQP - QPQ).
\]
Let $D$ be a matrix commuting with $P$ and $Q$. By cyclicity of the trace, $\operatorname{tr}(DPQP) = \operatorname{tr}(PDPQ) = \operatorname{tr}(DP^2Q) = \operatorname{tr}(DPQ)$; likewise $\operatorname{tr}(DQPQ) = \operatorname{tr}(DQP)$, and $\operatorname{tr}(DQP) = \operatorname{tr}(PDQ) = \operatorname{tr}(DPQ)$. Hence
\[
\operatorname{tr}(DPQP) = \operatorname{tr}(DQPQ).
\]
Finally $(P - Q)^{2k+1} = C^k (P - Q)$, and
\[
(P - Q)^{2k+3} = C^k (P - Q)\, C = (P - Q)^{2k+1} - \bigl(C^k PQP - C^k QPQ\bigr).
\]
The trace of the term in parentheses is zero, by the above with $D = C^k$. The traces of the odd powers of $P - Q$ are therefore all equal to $\operatorname{tr}(P - Q) = \operatorname{tr} P - \operatorname{tr} Q$.
\end{proof}

\noindent The proof uses only idempotence, the commutation of $C$ with $P$ and $Q$, and the cyclicity of the trace; it never divides by~$2$. The computation by anticommutation shows $2\operatorname{tr}(A^m B^2) = 0$; the proof above shows directly $\operatorname{tr}(A^m B^2) = 0$, since $A^m B^2 = C^{k}(PQP - QPQ)$ for $m = 2k + 1$. An exhaustive computation over all pairs of idempotents of size~$2$ and~$3$ over $\mathbf F_2$, of size~$2$ over $\mathbf Z/4$, and over $\mathbf F_3$, $\mathbf F_5$, $\mathbf Z/9$, for $k = 1, 2, 3$, confirms the identity, in characteristic~$2$ as elsewhere.

\emph{Verification in Lean.} The file \href{\leanblobmain/Notebook1983_10.lean}{\texttt{Notebook1983\_10.\allowbreak{}lean}} proves Theorem~\ref{thm:quillen-1983-10} under the name \texttt{trace\_sub\_pow\_odd}, for matrices indexed by an arbitrary finite type over an arbitrary commutative ring, idempotence being mathlib's predicate \texttt{IsIdempotentElem}. The steps are \texttt{sq\_sub\_commute\_left} and \texttt{sq\_sub\_commute\_right} (commutation of $C$), \texttt{sub\_mul\_sq\_sub} (the identity $(P - Q)C = (P - Q) - (PQP - QPQ)$) and \texttt{trace\_mul\_PQP}. Mathlib does not contain this identity: a pattern search on \texttt{Matrix.trace} and \texttt{IsIdempotentElem} returns nothing. The axioms used are the three usual axioms of mathlib.

Not proved here: the infinite-dimensional version, where $P - Q$ is not trace class and only a sufficiently large odd power is, which is the real question of the page; the interpretation of the common value as the index of a Fredholm operator.

\litterature{the infinite-dimensional version is published: J.~Avron, R.~Seiler and B.~Simon \cite{ass1994} prove, for two orthogonal projections of a Hilbert space whose difference has a trace-class odd power $(P - Q)^{2n+1}$, that the index of the pair equals $\operatorname{tr}((P - Q)^{2n+1})$, and hence does not depend on $n$; their proof rests in particular on the anticommutation of $P - Q$ and $P + Q - 1$. The finite-dimensional case is its simplest example. The form with no hypothesis on the characteristic has not been found in print; it is elementary, and the edition does not present it as new. It is not a lead.}

\section[1983, Lecture Notes 8 (Witten): two trace inequalities]{1983, Lecture Notes 8 (Witten): two trace inequalities\protect\verifie{Notebook1983_LectureNotes8}}\label{sec:quillen-1983-ln8}

\emph{Context.} The \dossier{1983-Lecture\%20Notes\%208\%20Witten}{file \q{Lecture Notes 8 Witten}} contains Quillen's notes of a lecture of Edward Witten, of December 1983, on mass inequalities in quantum chromodynamics. The principle is that the measure on gauge fields is positive, so that an inequality which holds for every fixed background field $A$ holds after integration. For a fixed $A$ the quantities compared are traces of finite matrices built from the quark propagator. The mathematics is that of the lecture.

\emph{The pages.} On page~4, the pion and rho propagators in a background field are written $\operatorname{Tr}(SS^*)$ and $\operatorname{Tr}(S\,\Gamma\,S^*\,\Gamma)$, where $S$ is a $12 \times 12$ matrix and $\Gamma = \gamma_\mu\gamma_5$. The notes say that $\Gamma^2 = 1$, so that the first is at least the second \q{by linear algebra}; this is the inequality $m_\pi \le m_\rho$ of Weingarten. On page~6, a difference of squared masses is reduced to $\operatorname{Tr}(E\,M_\mu\,E\,M_\mu^*)$ with $E \ge 0$, which the notes expand in an eigenbasis of $E$ as $\sum \lambda_i\lambda_j\,|\langle i|M_\mu|j\rangle|^2 \ge 0$.

\emph{Intuition.} On matrices, $\langle S, T\rangle = \operatorname{Tr}(ST^*)$ is a hermitian inner product, and $S \mapsto USV$ preserves its norm when $U$ and $V$ are unitary. The first inequality is then the Cauchy--Schwarz inequality between $S$ and $\Gamma S\Gamma$. It needs $\Gamma$ unitary, and $\Gamma^2 = 1$ alone does not give it: a matrix with $\Gamma^2 = 1$ is unitary exactly when it is hermitian. The gamma matrices of the lecture are hermitian up to the usual conventions, so the hypothesis is implicit on the page. For the second inequality, write $E = B^2$ with $B$ hermitian; then $\operatorname{Tr}(EMEM^*) = \operatorname{Tr}(CC^*)$ with $C = BMB$.

\begin{encadre}
\begin{theoreme}[After Quillen's notes of a lecture of Witten, 1983, p.~4]\label{thm:quillen-1983-ln8-a}
Let $S$ be a complex $n \times n$ matrix and $U$, $V$ unitary $n \times n$ matrices. Then $\operatorname{Re}\operatorname{Tr}(S\,V\,S^*\,U) \le \operatorname{Tr}(SS^*)$. In particular, if $\Gamma$ is hermitian and $\Gamma^2 = 1$, then $\operatorname{Tr}(S\,\Gamma\,S^*\,\Gamma)$ is real and
\[
  \operatorname{Tr}(S\,\Gamma\,S^*\,\Gamma) \le \operatorname{Tr}(SS^*).
\]
\end{theoreme}
\end{encadre}

\begin{proof}
For all $S$, $T$, $0 \le \operatorname{Tr}\bigl((S - T)(S - T)^*\bigr) = \operatorname{Tr}(SS^*) - 2\operatorname{Re}\operatorname{Tr}(ST^*) + \operatorname{Tr}(TT^*)$, because $\operatorname{Tr}(TS^*)$ is the conjugate of $\operatorname{Tr}(ST^*)$. Take $T = U^*SV^*$. Then $ST^* = SVS^*U$, and $TT^* = U^*SV^*VS^*U = U^*SS^*U$ has the trace of $SS^*$. Hence $2\operatorname{Re}\operatorname{Tr}(SVS^*U) \le 2\operatorname{Tr}(SS^*)$. If $\Gamma$ is hermitian with $\Gamma^2 = 1$, it is unitary, and the conjugate of $\operatorname{Tr}(S\Gamma S^*\Gamma)$ is $\operatorname{Tr}(\Gamma S\Gamma S^*) = \operatorname{Tr}(S\Gamma S^*\Gamma)$.
\end{proof}

\begin{contrexemple}[the edition's]\label{cex:quillen-1983-ln8}
The hypothesis $\Gamma^2 = 1$ alone does not suffice. For
\[
  \Gamma = \begin{pmatrix} 1 & 3 \\ 0 & -1 \end{pmatrix}, \qquad S = \begin{pmatrix} 0 & 1 \\ 1 & 0 \end{pmatrix},
\]
one has $\Gamma^2 = 1$, $\operatorname{Tr}(SS^*) = 2$ and $\operatorname{Tr}(S\,\Gamma\,S^*\,\Gamma) = 7$.
\end{contrexemple}

\begin{encadre}
\begin{theoreme}[After Quillen's notes of a lecture of Witten, 1983, p.~6]\label{thm:quillen-1983-ln8-b}
Let $E$ be a positive semidefinite complex $n \times n$ matrix and $M$ a complex $n \times n$ matrix. Then $\operatorname{Tr}(E\,M\,E\,M^*)$ is real and nonnegative.
\end{theoreme}
\end{encadre}

\begin{proof}
Let $B$ be the positive semidefinite square root of $E$; it is hermitian and $B^2 = E$. By cyclicity of the trace, $\operatorname{Tr}(B^2MB^2M^*) = \operatorname{Tr}\bigl((BMB)(BM^*B)\bigr) = \operatorname{Tr}(CC^*)$ with $C = BMB$, and $\operatorname{Tr}(CC^*) = \sum_{i,j}|C_{ij}|^2$.
\end{proof}

\noindent The notes prove the second inequality in an eigenbasis of $E$. The square root avoids the choice of a basis and gives the same sum of squares.

\emph{Verification in Lean.} The file \href{\leanblobmain/Notebook1983_LectureNotes8.lean}{\texttt{Notebook1983\_\allowbreak{}LectureNotes8.\allowbreak{}lean}} works with square complex matrices (\texttt{Matrix n n} over the complex numbers) and states the inequalities on real parts. The polarisation inequality is \texttt{re\_trace\_mul\_conjTranspose\_le}. Theorem~\ref{thm:quillen-1983-ln8-a} is \texttt{re\_trace\_unitary\_le}, \texttt{re\_trace\_involution\_le} and \texttt{im\_trace\_involution}; Counterexample~\ref{cex:quillen-1983-ln8} is \texttt{not\_re\_trace\_le\_of\_sq\_eq\_one}. Theorem~\ref{thm:quillen-1983-ln8-b} is \texttt{re\_trace\_posSemidef\_nonneg} and \texttt{im\_trace\_posSemidef}, with the square root \texttt{CFC.sqrt} of mathlib. Nothing of the physics is formalised: neither the propagators, nor the identity $S(x, 0) = \gamma_5 S(0, x)^*\gamma_5$ which puts the traces in this form, nor the integration over gauge fields. The axioms used are the three usual axioms of mathlib.

\litterature{The inequalities of the lecture are those of Weingarten (\emph{Phys.\ Rev.\ Lett.}~51 (1983)), Witten (\emph{Phys.\ Rev.\ Lett.}~51 (1983)) and Vafa--Witten, as recalled; the notes cite the paper of Vafa and Witten. The two matrix statements are the Cauchy--Schwarz inequality for the Hilbert--Schmidt inner product and the positivity of $\operatorname{Tr}(CC^*)$. The edition does not present them as new; the counterexample only records that the page's \q{$\Gamma^2 = 1$} is used together with hermiticity.}

\section[Notebook 1987-14: the kernel of the contraction by (1, ..., 1)]{Notebook 1987-14: $\dim\operatorname{Ker}(\iota_v : \Lambda^q K^k \to \Lambda^{q-1}K^k) = \binom{k-1}{q}$\protect\verifie{Notebook1987_14}}\label{sec:quillen-1987-14}

\emph{Context.} On page~49 of \carnet{1987}{14}, computing a Koszul-type complex attached to $k$ factors $A \otimes^L_P$, Quillen needs the kernel of the contraction of $\Lambda^q(K^k)$ by the vector $(1, \dots, 1)$, and identifies it with $\Lambda^q(K^{k-1})$.

\begin{encadre}
\begin{theoreme}[after Quillen, notebook 1987-14, p.~49]\label{thm:quillen-1987-14}
Let $M$ be a vector space of finite dimension $k$ over a field $K$, $d$ a nonzero linear form on $M$, and $d\lrcorner$ the contraction on the exterior algebra. Then $d\lrcorner$ maps $\Lambda^{q+1}M$ to $\Lambda^q M$, the complex $\cdots \to \Lambda^{q+2}M \to \Lambda^{q+1}M \to \Lambda^qM \to \cdots \to \Lambda^0M \to 0$ is exact, the image of $d\lrcorner$ on $\Lambda^{q+1}M$ has dimension $\binom{k-1}{q}$ and its kernel $\binom{k-1}{q+1}$. In particular this holds for the contraction with $(1, \dots, 1)$, that is with $d = \sum x_i$ on $K^k$.
\end{theoreme}
\end{encadre}

\begin{proof}
Choose $m$ with $d(m) = 1$. The identity $d\lrcorner(m \wedge x) + m \wedge (d\lrcorner x) = d(m)x = x$ shows that every $x$ with $d\lrcorner x = 0$ is $d\lrcorner(m \wedge x)$; and $d\lrcorner d\lrcorner = 0$. Rank and nullity at $\Lambda^{q+1}M$ give $\binom{k}{q+1} = \dim\operatorname{Im}_{q+1} + \dim\operatorname{Im}_q$, with $\dim\operatorname{Im}_0 = 1$; Pascal's rule gives $\dim\operatorname{Im}_q = \binom{k-1}{q}$ by induction.
\end{proof}

\emph{Verification in Lean.} The file \href{\leanblobmain/Notebook1987_14.lean}{\texttt{Notebook1987\_14.\allowbreak{}lean}} proves \texttt{contractLeft\_mem}, \texttt{ker\_contr}, \texttt{range\_contr\_zero}, \texttt{finrank\_range\_contr}, \texttt{finrank\_ker\_contr} and \texttt{finrank\_ker\_sum}, with mathlib's contraction on the Clifford algebra of the zero quadratic form and its exterior powers. Only $d \ne 0$ matters, not the value $d(v) = k$. The isomorphism with $\Lambda^q(K^{k-1})$ is given only through dimensions, and the homological computation that uses it on the page is not formalised. The axioms used are the three usual axioms of mathlib.

\litterature{Classical (exactness of the Koszul complex of a nonzero linear form); the edition does not present it as new.}

\section[Notebook 1989-3: powers of a block matrix]{Notebook 1989-3: powers of a block matrix with a zero corner; an involution\protect\verifie{Notebook1989_3}}\label{sec:quillen-1989-3}

\emph{Context.} On pages~99--100 of \carnet{1989}{3} Quillen computes the powers of a block matrix $\begin{pmatrix}A & B\\ C & 0\end{pmatrix}$ by a three-term recursion, and writes an involution in block form.

\begin{encadre}
\begin{theoreme}[Quillen, notebook 1989-3, pp.~99--100]\label{thm:quillen-1989-3}
Over any ring:
\begin{enumerate}
\item let $P_0 = 0$, $P_1 = 1$, $P_{n+2} = AP_{n+1} + BCP_n$. Then $\begin{pmatrix}A & B\\ C & 0\end{pmatrix}^{n+1} = \begin{pmatrix}P_{n+2} & P_{n+1}B\\ CP_{n+1} & CP_nB\end{pmatrix}$, with rectangular blocks;
\item if $b$, $c$ commute with $z$ and $bc = cb = 1 - z^2$, then $\begin{pmatrix}z & b\\ c & -z\end{pmatrix}^2 = 1$.
\end{enumerate}
\end{theoreme}
\end{encadre}

\begin{proof}
(1) By induction on $n$, multiplying on the left by the block matrix. (2) Direct computation.
\end{proof}

\emph{Verification in Lean.} The file \href{\leanblobmain/Notebook1989_3.lean}{\texttt{Notebook1989\_3.\allowbreak{}lean}} proves \texttt{fromBlocks\_pow\_succ} and \texttt{involution}. The page's $P_n$ is our $P_{n+1}$ (so that $P_{-1} = 0$ needs no special case). The axioms used are the three usual axioms of mathlib.

\litterature{Elementary (a noncommutative Chebyshev recursion); the edition does not present it as new.}

\section[Notebook 1989-4: faithful functionals on matrices]{Notebook 1989-4: faithful functionals on a matrix algebra\protect\verifie{Notebook1989_4Faithful}}\label{sec:quillen-1989-4-faithful}

\emph{Context.} On page~28 of \carnet{1989}{4} Quillen describes the faithful functionals on $M_n(k)$.

\begin{encadre}
\begin{theoreme}[Quillen, notebook 1989-4, p.~28]\label{thm:quillen-1989-4-faithful}
For $T \in M_n(K)$, $K$ a field, the functional $x \mapsto \operatorname{tr}(Tx)$ is faithful (for every $a \ne 0$ some $\operatorname{tr}(T\alpha a) \ne 0$) iff $T$ is invertible.
\end{theoreme}
\end{encadre}

\begin{proof}
If $T$ is invertible and $a_{ij} \ne 0$, take $\alpha = T^{-1}E_{ji}$: $\operatorname{tr}(E_{ji}a) = a_{ij}$. If $T$ is singular, a row vector $v \ne 0$ with $vT = 0$ gives the matrix $a$ with all rows $v$, $aT = 0$, and $\operatorname{tr}(T\alpha a) = \operatorname{tr}(\alpha aT) = 0$.
\end{proof}

\emph{Verification in Lean.} The file \href{\leanblobmain/Notebook1989_4Faithful.lean}{\texttt{Notebook1989\_4Faithful.\allowbreak{}lean}} proves \texttt{faithful\_iff}. The description of every functional as $\operatorname{tr}(T\,\cdot)$ is not restated. The axioms used are the three usual axioms of mathlib.

\litterature{Classical; the edition does not present it as new.}


\part{Clifford algebras, fermions and geometry}
\section[Notebook 1982-12: $(\gamma^1\cdots\gamma^m)^2$]{Notebook 1982-12: $(\gamma^1\cdots\gamma^m)^2 = (-1)^{m(m-1)/2}$\protect\verifie{Notebook1982_12}}\label{sec:quillen-1982-12}

\emph{Context.} On page~88 of \carnet{1982}{12} Quillen sets up the spinor module of an even-dimensional Euclidean space, with Clifford generators $\gamma^1, \ldots, \gamma^{2n}$ of an orthonormal basis and the chirality operator $\gamma^{2n+1} = \gamma^1\cdots\gamma^{2n}$, and uses $(\gamma^{2n+1})^2 = (-1)^{2n(2n-1)/2} = (-1)^n$ and the anticommutation of $\gamma^{2n+1}$ with the $\gamma^i$, before computing supertraces.

\begin{encadre}
\begin{theoreme}[Quillen, notebook 1982-12, p.~88]\label{thm:quillen-1982-12}
In a ring, let $e_0, e_1, \ldots$ satisfy $e_i^2 = 1$ and $e_ie_j = -e_je_i$ for $i \ne j$, and $P_m = e_0\cdots e_{m-1}$. Then $P_m^2 = (-1)^{m(m-1)/2}$; for $m = 2n$, $P_{2n}^2 = (-1)^n$ and $P_{2n}$ anticommutes with $e_i$, $i < 2n$. More precisely $e_iP_m = (-1)^{m-1}P_me_i$ for $i < m$ and $e_iP_m = (-1)^mP_me_i$ for $i \ge m$.
\end{theoreme}
\end{encadre}

\begin{proof}
The commutation rules follow by moving $e_i$ across the factors one at a time. Then $P_{m+1}^2 = P_me_mP_me_m = (-1)^mP_m^2e_m^2 = (-1)^mP_m^2$, and $\binom{m+1}{2} = \binom m2 + m$.
\end{proof}

\emph{Verification in Lean.} The file \href{\leanblobmain/Notebook1982_12.lean}{\texttt{Notebook1982\_12.\allowbreak{}lean}} proves the theorem as \texttt{prod\_mul\_prod}, \texttt{prod\_mul\_prod\_even}, \texttt{mul\_prod\_even}, \texttt{mul\_prod\_of\_lt} and \texttt{mul\_prod\_of\_le}, from the relations alone, in any ring. The traces on the spinor module (pp.~88--89) are not formalised. The axioms used are the three usual axioms of mathlib.

\litterature{Standard (Lawson and Michelsohn, \emph{Spin geometry}, ch.~I, \S 5, the volume element). The edition does not present it as new.}

\section[Notebook 1982-13: the projections $1 - e(\eta_j)\,i(X_j)$]{Notebook 1982-13: the projections $1 - e(\eta_j)\,i(X_j)$ and the joint kernel of the contractions\protect\verifie{Notebook1982_13}}\label{sec:quillen-1982-13}

\emph{Context.} \carnet{1982}{13} (22 November -- 30 December 1982) opens with the proof of Loday's conjecture on the Lie algebra homology of matrices, and continues with heat kernels, path integrals and the local index theorem, moment maps and the cohomology of gauge groups. On 23 December (pages~67--71) Quillen looks for a way to compute the equivariant cohomology of a manifold $M$ acted on by a compact connected group $G$ with differential forms. He conjectures that the complex $[\Omega(M) \otimes S(\mathfrak g^*)]^G$ does it, and reduces it to the Weil model $[\Omega(M) \otimes W(\mathfrak g)]_{\mathrm{basic}}$, provided $\Omega(M)$ is identified with the subspace of $\Lambda\mathfrak g^* \otimes \Omega(M)$ killed by all the contractions $i(X)$, $X \in \mathfrak g$, which act on both factors.

\emph{The page.} On page~69 he treats the case $\dim \mathfrak g = 1$: the map $\alpha \mapsto \alpha - \eta \otimes i(X)\alpha$ lands in the kernel of $i(X)$. On page~70, for a basis $X_j$ of $\mathfrak g$ and the dual basis $\eta_j$, he considers the operators $1 - e(\eta_j)\,i(X_j)$ on $\Lambda\mathfrak g^* \otimes \Omega(M)$, where $e$ is exterior multiplication, expands the product of two of them, and concludes that they commute, are projections, and that $i(X_j)$ kills $1 - e(\eta_j)\,i(X_j)$; their product is therefore killed by all the $i(X_j)$, and gives the identification of $\Omega(M)$ with the joint kernel.

\emph{Intuition.} Only the canonical anticommutation relations between the $e(\eta_j)$ and the $i(X_k)$ are used: $i(X_j)^2 = 0$, $e(\eta_j)\,i(X_j) + i(X_j)\,e(\eta_j) = 1$, and anticommutation of operators of different indices. The product $e_j i_j$ of two operators anticommuting with everything of index $k \ne j$ commutes with them.

\emph{Notation.} Let $A$ be a ring and $(e_j)$, $(i_j)$ families of elements of $A$ indexed by a set $J$. We say they satisfy the \emph{anticommutation relations} if $i_j^2 = 0$ and $e_j i_j + i_j e_j = 1$ for every $j$, and if, for $j \ne k$, $e_j e_k = -e_k e_j$, $i_j i_k = -i_k i_j$ and $e_j i_k = -i_k e_j$. Put $P_j = 1 - e_j i_j$.

\begin{encadre}
\begin{theoreme}[Quillen, notebook 1982-13, p.~70]\label{thm:quillen-1982-13}
Under the anticommutation relations:
\begin{enumerate}
\item each $P_j$ is idempotent, $i_j P_j = 0$, and the $P_j$ commute pairwise;
\item for distinct indices $j_1, \ldots, j_r$, the product $Q = P_{j_1} \cdots P_{j_r}$ is idempotent and $i_{j_s} Q = 0$ for every $s$;
\item on every $A$-module $V$, the image of $Q$ is $\{v \in V \mid i_{j_s} v = 0 \text{ for } s = 1, \ldots, r\}$, and $Q$ is the identity on it.
\end{enumerate}
The relations hold for $e_j$ the left multiplication by vectors $v_j$ and $i_j$ the contraction by linear forms $f_j$ on an exterior algebra $\Lambda V$, when $f_k(v_j) = \delta_{jk}$.
\end{theoreme}
\end{encadre}

\begin{proof}
From $i_j e_j = 1 - e_j i_j$ and $i_j^2 = 0$: $(e_j i_j)^2 = e_j (1 - e_j i_j) i_j = e_j i_j$, so $P_j^2 = P_j$; and $i_j P_j = i_j - (1 - e_j i_j) i_j = e_j i_j^2 = 0$. If $a$ and $b$ anticommute with $x$, then $ab\,x = -a\,x\,b = x\,ab$; hence $e_j i_j$, and $P_j$, commute with $e_k$ and $i_k$ for $k \ne j$, so $P_j$ commutes with $P_k$. In (2), $Q$ is a product of commuting idempotents, and $i_{j_s}$ commutes with the factors of index $\ne j_s$ and is killed by $P_{j_s}$. In (3), $i_{j_s} Q v = 0$; conversely, if all $i_{j_s} v = 0$, each $P_{j_s} v = v - e_{j_s} i_{j_s} v = v$, so $Qv = v$. On $\Lambda V$, $i_j^2 = 0$, left multiplications by vectors anticommute, so do contractions, and $i_f(v \wedge \omega) = f(v)\,\omega - v \wedge i_f\omega$ gives the mixed relations.
\end{proof}

\noindent The square $e_j^2 = 0$, true on the page, is not used. On $\Lambda\mathfrak g^* \otimes \Omega(M)$, with $e_j = e(\eta_j) \otimes 1$ and $i_j = i(X_j) \otimes 1 + \epsilon \otimes i(X_j)$, $\epsilon$ the parity operator on $\Lambda\mathfrak g^*$, the relations hold by the exterior-algebra case and the graded sign rule; this instance is not formalized. Nor is the final step of the page, that $\alpha \mapsto Q(1 \otimes \alpha)$ is a bijection from $\Omega(M)$ onto the joint kernel.

\emph{Verification in Lean.} The file \href{\leanblobmain/Notebook1982_13.lean}{\texttt{Notebook1982\_13.\allowbreak{}lean}} states the relations as the structure \texttt{IsCAR} and defines $P_j$ as \texttt{proj} and $Q$, for a list of indices, as \texttt{prodProj}. Part~(1) of Theorem~\ref{thm:quillen-1982-13} is \texttt{idem\_proj}, \texttt{inc\_mul\_proj} and \texttt{commute\_proj}; part~(2) is \texttt{isIdempotentElem\_prodProj} and \texttt{inc\_mul\_prodProj}, for a list without repetition; part~(3) is \texttt{prodProj\_smul\_of\_inc} and \texttt{mem\_range\_prodProj\_iff}. The exterior-algebra instance is \texttt{exterior\_isCAR}, on \texttt{ExteriorAlgebra R V} with mathlib's contraction \texttt{CliffordAlgebra.contractLeft}. The axioms used are at most the three usual axioms of mathlib.

\litterature{the Weil and Cartan models of equivariant cohomology go back to H.~Cartan's notes of 1950, reprinted in V.~Guillemin and S.~Sternberg, \emph{Supersymmetry and equivariant de Rham theory}, Springer, 1999. Their isomorphism by the operator $\exp(-\sum_j \theta^j \iota_j)$ is in V.~Mathai and D.~Quillen, \q{Superconnections, Thom classes, and equivariant differential forms}, \emph{Topology} 25 (1986), 85--110, and in J.~Kalkman, \emph{Comm.\ Math.\ Phys.} 153 (1993), 447--463, as one member of a family interpolating between the two models. The page's form, the product of the projections $1 - e(\eta_j)\,i(X_j)$ onto the joint kernel of contractions acting on both factors, was not found in this form in the sources consulted; the computation is elementary and the edition does not present it as new.}

\section[1983, Lecture Notes 15 (Patterson): the ball and the hyperboloid]{1983, Lecture Notes 15 (Patterson): the ball and the hyperboloid\protect\verifie{Notebook1983_LectureNotes15}}\label{sec:quillen-1983-ln15}

\emph{Context.} The \dossier{1983-Lecture\%20Notes\%2015\%20Patterson}{file \q{Lecture Notes 15 Patterson}} contains Quillen's notes of a lecture of S.~J.~Patterson on Kleinian groups. The notes begin with the ball model of hyperbolic space $\mathbf H^{N+1}$: the unit ball $B^{N+1}$ with the metric $(1 - \|x\|^2)^{-2}\|dx\|^2$. The mathematics is that of the lecture.

\emph{The page.} Page~1 introduces the function
\[
  L(w, w') = \frac{\|w - w'\|^2}{(1 - \|w\|^2)(1 - \|w'\|^2)},
\]
a map $p$ from the ball to $\mathbf R \times \mathbf R^{N+1}$ and the quadratic form $q(y, x) = y^2 - \|x\|^2$. It states that $q(p(x)) = 1$ and that $1 + 2L(x, x') = q(p(x), p(x'))$, the value of the bilinear form of $q$. Since $L$ determines the distance (the notes give $\sinh^2 d = L$), the isometries of the ball come from the orthogonal group $O(1, N + 1)$ of $q$. In the transcription, the map is $p(x) = \bigl(\tfrac12\,(1 + \|x\|^2)/(1 - \|x\|^2),\ x/(1 - \|x\|^2)\bigr)$.

\emph{Intuition.} With $r = \|x\|$, the point $\bigl((1 + r^2)/(1 - r^2),\ 2x/(1 - r^2)\bigr)$ lies on the hyperboloid $y^2 - \|u\|^2 = 1$, because $(1 + r^2)^2 - 4r^2 = (1 - r^2)^2$: this is the stereographic projection of the ball onto the upper sheet. The map as written is half of it, and for that map the two identities are off by a factor $4$.

\begin{encadre}
\begin{theoreme}[After Quillen's notes of a lecture of Patterson, 1983, p.~1]\label{thm:quillen-1983-ln15}
Let $E$ be a real inner product space. For $x \in E$ with $\|x\| \ne 1$ put $P(x) = \bigl((1 + \|x\|^2)/(1 - \|x\|^2),\ 2x/(1 - \|x\|^2)\bigr) \in \mathbf R \times E$, and let $q\bigl((y, u), (y', u')\bigr) = yy' - \langle u, u'\rangle$. Then for $x$, $x'$ of norm different from $1$,
\[
  q(P(x), P(x)) = 1 \qquad\text{and}\qquad 1 + 2L(x, x') = q(P(x), P(x')).
\]
For $p = \tfrac12 P$, the map as written in the notes, $q(p(x), p(x)) = \tfrac14$ and $1 + 2L(x, x') = 4\,q(p(x), p(x'))$.
\end{theoreme}
\end{encadre}

\begin{proof}
Put $s = 1 - \|x\|^2$ and $s' = 1 - \|x'\|^2$. Then $q(P(x), P(x')) = \bigl((2 - s)(2 - s') - 4\langle x, x'\rangle\bigr)/(ss')$. The numerator is $ss' + 2\bigl(2 - s - s' - 2\langle x, x'\rangle\bigr) = ss' + 2\bigl(\|x\|^2 + \|x'\|^2 - 2\langle x, x'\rangle\bigr) = ss' + 2\|x - x'\|^2$. For $x' = x$ this gives $1$. The statements on $p$ follow, since $q$ is bilinear.
\end{proof}

\emph{Verification in Lean.} The file \href{\leanblobmain/Notebook1983_LectureNotes15.lean}{\texttt{Notebook1983\_\allowbreak{}LectureNotes15.\allowbreak{}lean}} defines \texttt{L}, and the two components \texttt{pT} and \texttt{pS} of $P$. The identities of Theorem~\ref{thm:quillen-1983-ln15} are \texttt{q\_p} and \texttt{one\_add\_two\_L}, and those for the map as written are \texttt{q\_half\_p} and \texttt{one\_add\_two\_L\_half}. The hypothesis is $\|x\|^2 \ne 1$, so that the identities hold outside the unit sphere and not only in the ball. The hyperbolic metric, the formula $\sinh^2 d = L$ and the statement on the isometry group are not formalised. The axioms used are the three usual axioms of mathlib.

\litterature{The passage from the ball model to the hyperboloid model is standard (see for instance Ratcliffe, \emph{Foundations of hyperbolic manifolds}, as recalled). The factor is a matter of normalisation in notes taken during a lecture; the handwriting has not been re-read against the facsimile on this point. The edition does not present the statement as new.}

\section[1983, Review Getzler material: composition of Cayley transforms]{1983, Review Getzler material: composition of Cayley transforms\protect\verifie{Notebook1983_ReviewGetzler}}\label{sec:quillen-1983-getzler}

\emph{Context.} The \dossier{1983-Review\%20Getzler\%20material}{file \q{Review Getzler material}} of 1983 contains seven leaves in which Quillen reviews, for himself, the supersymmetric path-integral approach to the index theorem of Friedan--Windey and Getzler. The fermionic part leads to orthogonal transformations of a real vector space, which he writes as Cayley transforms $(1 + X)/(1 - X)$ of skew forms $X$.

\emph{The pages.} On pages~6 and~7 he asks for the composition law: which $Z$ has for Cayley transform the product of those of $X$ and $Y$? The matrices $X$ and $Y$ do not commute, and the first attempts are left unfinished; the margin of page~6 says \q{subtle}. The last line of page~7 gives
\[
  Z = \frac{1}{1 - X}\,(X + Y)\,\frac{1}{1 + XY}\,(1 - X).
\]

\emph{Intuition.} When $X$ and $Y$ commute this is $Z = (X + Y)/(1 + XY)$, the addition formula of the hyperbolic tangent, since the Cayley transform of $\tanh t$ is $e^{2t}$. In general the two identities
\[
  (1 - X)(1 - Y) = (1 + XY) - (X + Y), \qquad (1 + X)(1 + Y) = (1 + XY) + (X + Y)
\]
give $1 - Z$ and $1 + Z$ as products, and the factor $(1 + XY)^{-1}(1 - X)$ common to both cancels in the Cayley transform.

\begin{encadre}
\begin{theoreme}[Quillen, 1983, Review Getzler material, p.~7]\label{thm:quillen-1983-getzler}
Let $X$ and $Y$ be elements of a ring such that $1 - X$, $1 - Y$ and $1 + XY$ are invertible, and put $Z = (1 - X)^{-1}(X + Y)(1 + XY)^{-1}(1 - X)$. Then $1 - Z$ is invertible and
\[
  (1 + X)(1 - X)^{-1}\,(1 + Y)(1 - Y)^{-1} = (1 + Z)(1 - Z)^{-1}.
\]
\end{theoreme}
\end{encadre}

\begin{proof}
Put $u = 1 - X$, $v = 1 - Y$, $w = 1 + XY$, so that $X + Y = w - uv$ and $(1 + X)(1 + Y) = w + (X + Y)$. Then $Z = u^{-1}(w - uv)w^{-1}u = 1 - vw^{-1}u$, so $1 - Z = vw^{-1}u$ is invertible with inverse $u^{-1}wv^{-1}$. Also $1 + Z = u^{-1}\bigl(w + (X + Y)\bigr)w^{-1}u = u^{-1}(1 + X)(1 + Y)w^{-1}u$. Hence $(1 + Z)(1 - Z)^{-1} = u^{-1}(1 + X)(1 + Y)v^{-1}$. Finally $1 + X = 2 - u$ commutes with $u^{-1}$, which gives the left-hand side.
\end{proof}

\noindent The page does not state the hypothesis on $1 + XY$ nor the invertibility of $1 - Z$; for skew-symmetric real matrices $X$, $Y$ the first two hypotheses always hold, and the third is the condition that the product of the two rotations has no eigenvalue $-1$.

\emph{Verification in Lean.} The file \href{\leanblobmain/Notebook1983_ReviewGetzler.lean}{\texttt{Notebook1983\_\allowbreak{}ReviewGetzler.\allowbreak{}lean}} defines \texttt{cayley X} and \texttt{cayleySum X Y} (the element $Z$) with \texttt{Ring.inverse}, in any ring. The formulas for $1 - Z$ and $1 + Z$ are \texttt{one\_sub\_cayleySum} and \texttt{one\_add\_cayleySum}, stated with units \texttt{u}, \texttt{v}, \texttt{w} of the ring. Theorem~\ref{thm:quillen-1983-getzler} is \texttt{isUnit\_one\_sub\_cayleySum} and \texttt{cayley\_mul\_cayley}. The link with orthogonal transformations and skew forms is not formalised. The axioms used are the three usual axioms of mathlib.

\litterature{In the commuting case the formula is the addition law of $\tanh$. We did not search the literature for the noncommutative form beyond a recollection that composition laws of this kind appear for Cayley parametrisations of rotations; the identity is elementary, and the edition does not present it as new.}

\section[Notebook 1984-6: fermions and Clifford generators; a symplectic block lemma]{Notebook 1984-6: Clifford generators from a fermion; a symplectic block lemma\protect\verifie{Notebook1984_6}}\label{sec:quillen-1984-6}

\emph{Context.} On page~2 of \carnet{1984}{6} Quillen passes from a fermionic creation--annihilation pair to the two Clifford generators it defines; on pages~48--49 he analyses a symplectic matrix $\begin{psmallmatrix} A & B\\ C & D\end{psmallmatrix}$ with $B$ invertible (a Bogoliubov transformation, as in Berezin's treatment).

\begin{encadre}
\begin{theoreme}[Quillen, notebook 1984-6, pp.~2, 48--49]\label{thm:quillen-1984-6}
\begin{enumerate}
\item In a ring, let $a^2 = 0$, $(a^*)^2 = 0$, $aa^* + a^*a = 1$, and let $i$ be central with $i^2 = -1$. Then $\gamma_1 = a^* + a$ and $\gamma_2 = i(a^* - a)$ satisfy $\gamma_1^2 = \gamma_2^2 = 1$, $\gamma_1\gamma_2 + \gamma_2\gamma_1 = 0$ and $i\gamma_1\gamma_2 = a^*a - aa^* = 2a^*a - 1$.
\item Over a commutative ring, let $A, B, C, D$ be square matrices with $B$ invertible, $AB^t + BA^t = 0$, $D^tB + B^tD = 0$ and $CB^t + DA^t = 1$. Then $B^{-1}A$ and $DB^{-1}$ are skew-symmetric and $(B^{-1})^t = C - DB^{-1}A$.
\end{enumerate}
\end{theoreme}
\end{encadre}

\begin{proof}
(1) Expand, using $i$ central. (2) $A^t = -B^{-1}AB^t$, so $(B^{-1}A)^t = A^t(B^t)^{-1} = -B^{-1}A$; likewise $D^t = -B^tDB^{-1}$ gives $(DB^{-1})^t = -DB^{-1}$; multiplying $CB^t + DA^t = 1$ on the right by $(B^t)^{-1}$ gives $C - DB^{-1}A = (B^t)^{-1}$.
\end{proof}

\emph{Verification in Lean.} The file \href{\leanblobmain/Notebook1984_6.lean}{\texttt{Notebook1984\_6.\allowbreak{}lean}} proves (1) as \texttt{gamma₁\_sq}, \texttt{gamma₂\_sq}, \texttt{gamma\_anticomm} and \texttt{i\_gamma\_gamma}, and (2) as \texttt{inv\_mul\_skew}, \texttt{mul\_inv\_skew} and \texttt{transpose\_inv\_eq}. Only three of the block relations of the page are needed for (2). The Gaussian kernels built from these matrices are not formalised. The axioms used are the three usual axioms of mathlib.

\litterature{Standard: (1) is the Clifford form of one fermionic mode; (2) is in substance in Berezin, \emph{The method of second quantization} (1966). The edition does not present them as new.}

\section[Notebook 1984-15: pair symmetry of a curvature tensor]{Notebook 1984-15: pair symmetry of a curvature tensor, and the factor $2$\protect\verifie{Notebook1984_15}}\label{sec:quillen-1984-15}

\emph{Context.} On page~55 of \carnet{1984}{15} Quillen rederives the symmetry $R_{ijkl} = R_{klij}$ of a tensor skew in $(i, j)$ and in $(k, l)$ that satisfies the cyclic identity, by adding four cyclic identities.

\begin{encadre}
\begin{theoreme}[Quillen, notebook 1984-15, p.~55]\label{thm:quillen-1984-15}
Let $R_{ijkl}$ take values in an abelian group, with $R_{ijkl} = -R_{jikl} = -R_{ijlk}$ and $R_{ijkl} + R_{jkil} + R_{kijl} = 0$. Then $2R_{ijkl} = 2R_{klij}$; hence $R_{ijkl} = R_{klij}$ when the group has no $2$-torsion. Over $\mathbf F_2$, on four indices, the tensor equal to $1$ exactly when $i, j, k, l$ are distinct and $0 \in \{i, j\}$ satisfies the hypotheses, alternating included, but $R_{0123} = 1 \neq 0 = R_{2301}$.
\end{theoreme}
\end{encadre}

\begin{proof}
Write the cyclic identity at $(i, j, k, l)$, $(j, k, l, i)$, $(k, l, i, j)$, $(l, i, j, k)$, normalise each term by the two skew symmetries, and add the first two and subtract the last two: everything cancels but $2R_{ijkl} - 2R_{klij}$. The counterexample is checked by enumeration; on three indices, a computation over $\mathbf F_2$ shows that pair symmetry still holds.
\end{proof}

\emph{Verification in Lean.} The file \href{\leanblobmain/Notebook1984_15.lean}{\texttt{Notebook1984\_15.\allowbreak{}lean}} proves the theorem as \texttt{two\_smul\_pair\_symm}, \texttt{pair\_symm} and, for the counterexample, \texttt{isCurvature\_R₂} and \texttt{not\_pair\_symm\_zmod\_two} (by \texttt{decide}). The page works over the real numbers, where the factor $2$ is harmless. The axioms used are the three usual axioms of mathlib.

\litterature{Classical (Kobayashi and Nomizu, \emph{Foundations of differential geometry}, I, ch.~V, \S 2). That the factor $2$ cannot be dropped in characteristic~$2$ is the edition's remark; the counterexample was not looked for in print. The edition does not present the statement as new.}


\part{Analysis, integrals and special functions}
\section[Notebook 1972-4: from the microcanonical to the Gibbs distribution]{Notebook 1972-4: harmonic oscillators, from the microcanonical to the Gibbs distribution\protect\verifie{Notebook1972_4}}\label{sec:quillen-1972-4}

\emph{Context.} In April 1972, between a run on the resolution theorem for higher $K$-theory (2 April) and a first candidate for an exact-sequence construction of $K$-theory (19 April), \carnet{1972}{4} contains two entries headed as a self-education in statistical mechanics (15 and 16 April). The first derives the ideal gas law, the adiabatics and the entropy of a monatomic ideal gas; the second asks how the Maxwell--Boltzmann--Gibbs measure $e^{-H/kT}\,dV/Z$ arises from a large number of almost independent copies of a Hamiltonian system, and works it out for harmonic oscillators.

\emph{The page.} Take $N$ independent oscillators, with phase space $\mathbf R^{2N}$, Liouville measure $\omega = \prod dq_i\,dp_i$ and Hamiltonian $H = \sum_i (p_i^2 + q_i^2)/2$. Writing $\omega = dH \wedge \eta$, Quillen computes the Laplace transform $\int e^{-sH}\omega = (2\pi/s)^N$ and reads off the volume $(2\pi)^N E^{N-1}/(N-1)!$ of the shell $H = E$ (pp.~19--20). The same computation with a point mass in the first factor gives the image of the shell measure in the phase plane of the first oscillator; dividing by the volume of the shell, he obtains (p.~23) the probability density
\[
\frac{N-1}{2\pi E}\Bigl(1 - \frac{r^2}{2E}\Bigr)^{N-2}\,\mathrm{Heav}\Bigl(1 - \frac{r^2}{2E}\Bigr)\, r\,dr\,d\theta
\]
on the disc $r^2 \le 2E$. He notes that for fixed $E$ it concentrates at the origin as $N \to \infty$, since each oscillator then has average energy $E/N \to 0$, and passes instead to the limit $E = \lambda N$, where it becomes the Gibbs density $e^{-r^2/2\lambda}/(2\pi\lambda)$, marked with an exclamation sign (p.~24). On p.~25 the same ratio of shell volumes gives the law of the energy $E_1$ of the first oscillator, with limit $e^{-E_1/\lambda}/\lambda$ (the page carries an extra factor $1/2\pi$ coming from its normalisation of $\eta$). The entry continues with the general problem (the law of $X_1$ given $X_1 + \dots + X_N = E$ for independent identically distributed variables) by characteristic functions and Fourier inversion, which is left open.

\begin{encadre}
\begin{theoreme}[Quillen, notebook 1972-4, pp.~23--25]\label{thm:quillen-1972-4}
Let $\lambda > 0$.
\begin{enumerate}
\item For every real $x$, $\dfrac{N-1}{N}\Bigl(1 - \dfrac xN\Bigr)^{N-2} \to e^{-x}$ as $N \to \infty$.
\item For every real $r$,
\[
\frac{N-1}{2\pi\lambda N}\Bigl(1 - \frac{r^2}{2\lambda N}\Bigr)^{N-2}\,\mathbf 1_{\{r^2 \le 2\lambda N\}} \longrightarrow \frac{e^{-r^2/2\lambda}}{2\pi\lambda}.
\]
\item For every real $E_1$, $\dfrac{N-1}{\lambda N}\Bigl(1 - \dfrac{E_1}{\lambda N}\Bigr)^{N-2} \to \dfrac{e^{-E_1/\lambda}}{\lambda}$.
\item For $E > 0$ and $N \ge 2$, $\displaystyle\int_0^E \frac{N-1}{E}\Bigl(1 - \frac uE\Bigr)^{N-2} du = 1$.
\end{enumerate}
\end{theoreme}
\end{encadre}

\noindent Part~(2) is the limit of page~24 with the Heaviside factor written as an indicator; part~(3) is that of page~25 without the factor $1/2\pi$; part~(4) says that the finite-$N$ law of the energy of one oscillator is a probability law. The shell volume itself, which needs the coarea formula, is not formalised.

\begin{proof}
(1) Put $a_N = 1 - x/N$. Then $a_N \to 1$ and $a_N^N \to e^{-x}$; for $N \ge 2$ and $a_N \ne 0$, which holds for $N$ large, $a_N^{N-2} = a_N^N / a_N^2$, and $(N-1)/N \to 1$. (2) and (3) are (1) with $x = r^2/2\lambda$ and $x = E_1/\lambda$, after dividing by $2\pi\lambda$, resp.\ $\lambda$; in (2) the indicator equals $1$ as soon as $2\lambda N \ge r^2$. (4) The function $u \mapsto -(1 - u/E)^{N-1}$ is a primitive of the integrand, and takes the values $-1$ at $0$ and $0$ at $E$.
\end{proof}

\emph{Verification in Lean.} The file \href{\leanblobmain/Notebook1972_4.lean}{\texttt{Notebook1972\_4.\allowbreak{}lean}} proves part~(1) under the name \texttt{tendsto\_pred\_div\_mul\_one\_sub\_pow}, from mathlib's \texttt{Real.tendsto\_one\_add\_div\_pow\_exp}; parts~(2) and~(3) are \texttt{tendsto\_oscillator\_marginal} and \texttt{tendsto\_energy\_marginal}, and part~(4), with $N = m + 2$, is \texttt{integral\_energy\_marginal}.

\litterature{the computation is classical. The equivalence of the microcanonical and canonical ensembles, of which this is the simplest instance, is the subject of A.~I. Khinchin, \emph{Mathematical foundations of statistical mechanics} (Dover, 1949), whose method, as we recall it, goes like the end of the entry through characteristic functions and a limit theorem (the book was not re-read for this section); limit theorems for the first coordinates of a point drawn uniformly from a high-dimensional symmetric set, of which the disc and the simplex here are instances, are of the kind collected by P.~Diaconis and D.~Freedman, \q{A dozen de Finetti-style results in search of a theory}, \emph{Ann.\ Inst.\ H.~Poincaré Probab.\ Statist.} 23 (1987), 397--423. The entry is an exercise of self-instruction, and the edition does not present it as new.}

\section[Notebook 1973-3: Gaussian sums over a translate of a lattice]{Notebook 1973-3: the Gaussian sum over a translate of a lattice, and the invariant of the cross-ratio\protect\verifie{Notebook1973_3}}\label{sec:quillen-1973-3}

\emph{Context.} In April 1973, \carnet{1973}{3} follows several threads around the $K$-theory of rings of integers and of curves: buildings and stability, vector bundles on curves and their slopes, lattice zeta functions. An entry of 26 April (pages~75--86) proposes to treat a lattice $M$ equipped with a positive definite quadratic form $q$ on $M_{\mathbf R}$ as a vector bundle over $\mathbf Z \cup \infty$, and to measure its \q{sections} by the theta series $\theta_E = \sum_{m \in M} e^{-\pi q(m)}$. The Poisson summation formula then plays the part of the Riemann--Roch theorem. An earlier entry, of 14 April (pages~20--27), computes the low-dimensional homology of $PGL_2(k)$ from its action on configurations of points of $\mathbf P^1_k$, which leads to the cross-ratio.

\emph{The page: the theta inequality.} Quillen wants to compare $\theta_E$ with $\theta_{E'}\theta_{E''}$ for an exact sequence $0 \to E' \to E \to E'' \to 0$ of such bundles. Writing the sum over $M$ as a sum over $M''$ of sums over the fibres, he reduces the question to an inequality for a single lattice: the Gaussian sum over a translate $x + M$ should not exceed the sum over $M$. On page~79 he checks it on the line, for $M = \mathbf Z$ and $q(x) = \alpha x^2$: the Fourier expansion of $x \mapsto \sum_m e^{-\pi\alpha(x+m)^2}$ shows that its second derivative at $0$ is negative. On pages~80--81 he proves the general lemma: the Poisson formula writes the shifted sum as a cosine series with positive coefficients, $\cos \le 1$, with equality for every frequency exactly when $x \in M$. He concludes (page~81) that $\theta_E \le \theta_{E'}\theta_{E''}$, with equality if and only if the sequence splits orthogonally, and remarks (page~82) that over a finite field, where $\theta_E = q^{h^0(E)}$, equality is common without splitting.

\emph{Intuition.} The function $f(x) = \sum_{m \in \mathbf Z} e^{-\pi\alpha(x+m)^2}$ is periodic. Its Fourier coefficients are those of the Gaussian, $\alpha^{-1/2} e^{-\pi n^2/\alpha}$, all positive. A function whose Fourier coefficients are all positive reaches its maximum where all the characters equal $1$ at once, at the integers, and only there.

\begin{encadre}
\begin{theoreme}[Quillen, notebook 1973-3, pp.~79--81, in dimension one]\label{thm:quillen-1973-3}
Let $\alpha > 0$. For every real $x$,
\[
\sum_{m \in \mathbf Z} e^{-\pi\alpha(x+m)^2} \;\le\; \sum_{m \in \mathbf Z} e^{-\pi\alpha m^2},
\]
with equality if and only if $x \in \mathbf Z$.
\end{theoreme}
\end{encadre}

\begin{proof}
By the Poisson summation formula, applied to the Gaussian $t \mapsto e^{-\pi\alpha(x+t)^2}$,
\[
\sum_{m \in \mathbf Z} e^{-\pi\alpha(x+m)^2} = \alpha^{-1/2} \sum_{n \in \mathbf Z} e^{-\pi n^2/\alpha} \cos(2\pi n x).
\]
Each term on the right is at most $e^{-\pi n^2/\alpha}$, and the right-hand side at $x = 0$ is $\alpha^{-1/2} \sum_n e^{-\pi n^2/\alpha}$, which the same formula identifies with $\sum_m e^{-\pi\alpha m^2}$. This gives the inequality. If $x \notin \mathbf Z$, the term $n = 1$ is strictly smaller, since $\cos(2\pi x) < 1$, so the inequality is strict. If $x \in \mathbf Z$, the left-hand sum is a reindexing of the right-hand one.
\end{proof}

\noindent The statement is the lemma of page~80 for $M = \mathbf Z$ and $q(x) = \alpha x^2$, which is also the case Quillen tests on page~79. The proof is his. The general lemma, for a lattice in $\mathbf R^n$, needs the Poisson formula over a lattice of $\mathbf R^n$, and the inequality $\theta_E \le \theta_{E'}\theta_{E''}$ follows from it as on page~81; neither is proved here.

\emph{The page: the invariant of the cross-ratio.} On pages~22--25, Quillen lists the six values $z$, $1/z$, $1 - z$, $1/(1-z)$, $1 - 1/z$, $z/(z-1)$ that the cross-ratio of four points takes under the permutations of the points, observes that the Klein four-group acts trivially (a remark he attributes to Atiyah), and proposes the rational function
\[
w(z) = \frac{(z^2 - z + 1)^3}{z^2 (z - 1)^2},
\]
invariant under $z \mapsto 1 - z$ and $z \mapsto 1/z$, hence under $\Sigma_3$. It has triple zeros at the primitive sixth roots of unity and double poles at $0$, $1$, $\infty$, and takes the value $27/4$ on the orbit $\{-1, 1/2, 2\}$, whose stabilizer has order~$2$. He says that $z \mapsto w$ puts the $\Sigma_3$-orbits of $k - \{0, 1\}$ in bijection with $k$.

\begin{proposition}[Quillen, notebook 1973-3, p.~25]\label{prop:quillen-1973-3}
Let $k$ be a field and $z, z' \in k - \{0, 1\}$. Then $w(z') = w(z)$ if and only if $z'$ is one of $z$, $1/z$, $1 - z$, $1/(1-z)$, $z/(z-1)$, $(z-1)/z$.
\end{proposition}

\begin{proof}
Clearing denominators, $w(z') - w(z)$ has numerator
\[
(z'^2 - z' + 1)^3 z^2 (z-1)^2 - (z^2 - z + 1)^3 z'^2 (z'-1)^2 = (z' - z)(z'z - 1)(z' + z - 1)(z'z - z' + 1)(z'z - z' - z)(z'z - z + 1),
\]
a polynomial identity. Each linear factor vanishes exactly at one of the six values.
\end{proof}

\noindent The proposition gives the injectivity on orbits; the surjectivity onto $k$, for $k$ algebraically closed, is not proved here.

\emph{Verification in Lean.} The file \href{\leanblobmain/Notebook1973_3.lean}{\texttt{Notebook1973\_3.\allowbreak{}lean}} proves Theorem~\ref{thm:quillen-1973-3} under the names \texttt{gaussian\_shift\_le} (the inequality), \texttt{gaussian\_shift\_lt} (strict for $x \notin \mathbf Z$) and \texttt{gaussian\_shift\_eq\_iff} (the equality case). The Fourier side is \texttt{gaussian\_shift\_fourier}, derived from \texttt{Complex.tsum\_exp\_neg\_quadratic}, the form of the Poisson summation formula for Gaussians in mathlib; summability is \texttt{summable\_gaussian}. The function $w$ is \texttt{crossRatioInv}; its invariance is \texttt{crossRatioInv\_one\_sub} and \texttt{crossRatioInv\_inv}, the values $27/4$ are \texttt{crossRatioInv\_two}, \texttt{crossRatioInv\_neg\_one}, \texttt{crossRatioInv\_half}, the zeros \texttt{crossRatioInv\_eq\_zero\_iff}, the factorization \texttt{crossRatioInv\_sub\_factor}, and Proposition~\ref{prop:quillen-1973-3} is \texttt{crossRatioInv\_eq\_iff}. The axioms used are the three usual axioms of mathlib.

\litterature{both statements are in the literature. The function $w$ is $j/256$ in terms of the Legendre parameter, $j(\lambda) = 2^8(\lambda^2 - \lambda + 1)^3/(\lambda^2(\lambda - 1)^2)$, a classical formula of the theory of elliptic curves; the value $27/4$ corresponds to $j = 1728$. The view of a Euclidean lattice as a vector bundle over $\operatorname{Spec}\mathbf Z \cup \infty$, with $\log\theta$ as $h^0$ and the Poisson formula as Riemann--Roch, is developed by G.~van der Geer and R.~Schoof, \emph{Selecta Math.} 6 (2000), 377--398, and by J.-B.~Bost, \emph{Theta invariants of Euclidean lattices and infinite-dimensional Hermitian vector bundles over arithmetic curves}, Progress in Math.~334, Birkhäuser, 2020, where the subadditivity of $h^0_\theta$ in short exact sequences, the inequality of page~81, is established; the edition has not checked the exact place. The coset inequality of the theorem is a standard tool of the Gaussian method of W.~Banaszczyk, \emph{Math. Ann.} 296 (1993), 625--635. In dimension one it also follows from the Jacobi triple product. The edition does not present either statement as new.}

\section[Notebook 1982-3: products $\prod(1 + a_n)$]{Notebook 1982-3: $\prod(1 + a_n)$ converges if and only if $\sum a_n$ does\protect\verifie{Notebook1982_3}}\label{sec:quillen-1982-3}

\emph{Context.} On page~18 of \carnet{1982}{3}, and again on page~35, Quillen uses the criterion for infinite products with non-negative terms: $\prod(1 + a_n)$ converges if and only if $\sum a_n < \infty$, to control the convergence of the products defining his determinants.

\emph{Intuition.} Expanding the product, $1 + \sum_{i \in s} a_i \le \prod_{i \in s}(1 + a_i) \le \exp\sum_{i \in s} a_i$.

\begin{encadre}
\begin{theoreme}[Notebook 1982-3, p.~18]\label{thm:quillen-1982-3}
For real $a_i \ge 0$, the product $\prod(1 + a_i)$ converges (unconditionally) if and only if $\sum a_i$ converges; for a sequence, the partial products $\prod_{i<N}(1 + a_i)$ converge if and only if $\sum a_i$ converges. The limit is then at least $1$.
\end{theoreme}
\end{encadre}

\begin{proof}
If $\sum a_i$ converges, the product does, by the upper bound (mathlib). Conversely, the finite products of factors $\ge 1$ grow with the set of indices and are bounded by the limit (eventually, for the unconditional product), so the partial sums are bounded by the lower bound, and a series of non-negative terms with bounded partial sums converges.
\end{proof}

\emph{Verification in Lean.} The file \href{\leanblobmain/Notebook1982_3.lean}{\texttt{Notebook1982\_3.\allowbreak{}lean}} proves the lower bound as \texttt{one\_add\_sum\_le\_prod}, the criterion as \texttt{multipliable\_one\_add\_iff} and the sequential form as \texttt{tendsto\_prod\_range\_iff}; the direct implication is mathlib's \texttt{Real.multipliable\_one\_add\_of\_summable}. The axioms used are the three usual axioms of mathlib.

\litterature{Classical (Knopp, \emph{Theory and application of infinite series}, ch.~VII). Mathlib has the direct implication only. The edition does not present it as new.}

\section[Notebook 1983-4: the cohomological equation on the torus]{Notebook 1983-4: the cohomological equation on the torus and Diophantine numbers\protect\verifie{Notebook1983_4}}\label{sec:quillen-1983-4}

\emph{Context.} In spring 1983, \carnet{1983}{4} studies the Connes algebra of the Kronecker foliation of irrational slope $\theta$ on the torus $T^2$ and its projective modules. On pages~71--72 (its pages~800--801), Quillen computes the Picard group of the topos of leaves of this foliation: the line bundles there are line bundles on $T^2$ equipped with a connection that is flat along the leaves, and the group sought is $H^1(T^2, \mathcal O^*)$, where $\mathcal O$ is the sheaf of smooth functions constant along the leaves.

\emph{The page.} The first step is $H^1(T^2, \mathcal O)$, computed by the fine resolution given by the operator $D = \theta\partial_x - \partial_y$ on $C^\infty(T^2)$. On the exponential $e^{2\pi i(mx + ny)}$, $D$ acts by the scalar $2\pi i(\theta m - n)$, which vanishes only at $(m, n) = (0, 0)$ since $\theta$ is irrational. Quillen deduces that $D$ is invertible off the constants, hence that $H^0$ and $H^1$ are both equal to $\mathbf C$, then, by the exponential sequence, that the Picard group is $\mathbf C^*/\{e^{2\pi i n\theta}\} \times \mathbf Z$.

\emph{Intuition.} Inverting $D$ on the exponentials is not enough: the series \[\sum b_{m,n}\,(2\pi i(\theta m - n))^{-1} e^{2\pi i(mx+ny)}\] must still define a smooth function. The coefficients of a smooth function decay faster than any power of $1 + |m| + |n|$; if the small divisors $|\theta m - n|$ are bounded below only by a power of $|m|^{-1}$, the division costs a fixed power and rapid decay survives. If on the contrary $\theta$ is very well approximated by rationals, one can place tiny coefficients $b_{m,n}$ exactly at the points where $\theta m - n$ is smaller still, and the quotient no longer decays. This is the phenomenon of small divisors, and it is the hypothesis that the page does not state.

\begin{enonce}
The sequences considered are indexed by $\mathbf Z^2$ and complex-valued. A sequence $a$ is \emph{rapidly decreasing} if, for every integer $k$, the quantity $|a_{m,n}|\,(1 + |m| + |n|)^k$ is bounded. A real number $\theta$ is \emph{Diophantine} if there exist $C > 0$ and an integer $N$ such that $|\theta m - n| \ge C\,|m|^{-N}$ for all integers $m \ne 0$ and $n$. The factor $2\pi$ is omitted: $D$ acts by $a_{m,n} \mapsto i(\theta m - n)\,a_{m,n}$.
\end{enonce}

\begin{encadre}
\begin{theoreme}[the edition's statement; missing hypothesis of notebook 1983-4, pp.~71--72]\label{thm:quillen-1983-4}
Let $\theta$ be an irrational real number. The operator $a \mapsto (i(\theta m - n)\,a_{m,n})_{m,n}$ sends rapidly decreasing sequences to rapidly decreasing sequences whose $(0,0)$ term is zero, and it is surjective onto the latter if and only if $\theta$ is Diophantine. More precisely: (1)~if $\theta$ is Diophantine (without any irrationality hypothesis) and $b$ is rapidly decreasing with $b_{0,0} = 0$, the sequence $a_{m,n} = b_{m,n}/(i(\theta m - n))$, $a_{0,0} = 0$, is rapidly decreasing and satisfies $i(\theta m - n)\,a_{m,n} = b_{m,n}$; (2)~if $\theta$ is irrational and not Diophantine, there exists a rapidly decreasing sequence $b$, with $b_{0,0} = 0$, that has no rapidly decreasing preimage.
In particular, the conclusion of the page holds for $\theta$ Diophantine, and fails for a Liouville number.
\end{theoreme}
\end{encadre}

\begin{proof}
Write $w_{m,n} = 1 + |m| + |n|$ and $\sigma_{m,n} = \theta m - n$. We have $|\sigma_{m,n}| \le (|\theta| + 1)\,w_{m,n}$, whence the first assertion.

(1) There exist $c > 0$ and $N$ such that $|\sigma_{m,n}|\,w_{m,n}^N \ge c$ for $(m, n) \ne (0,0)$: if $m \ne 0$, this is the hypothesis, since $|m| \le w_{m,n}$; if $m = 0$, we have $|\sigma_{0,n}| = |n| \ge 1$. Then $|a_{m,n}|\,w^k \le |b_{m,n}|\,w^{k+N}/c$, which is bounded. The lower bound involves only $|m|$, and $|m| \le w$: the dependence on $n$ is no obstacle in this direction.

(2) The negation of the hypothesis provides, for each $j$, a pair $q_j = (m_j, n_j)$ with $m_j \ne 0$ and $|\sigma_{q_j}| < (j+1)^{-1}|m_j|^{-2j}$; irrationality gives $\sigma_{q_j} \ne 0$. Put $b_{q} = \sigma_q / \sqrt{|\sigma_q|}$ if $q$ is one of the $q_j$, and $b_q = 0$ otherwise; thus $b_{0,0} = 0$ and $|b_{q_j}| = \sqrt{|\sigma_{q_j}|} < |m_j|^{-j}$. Here $n$ comes in: since $|\sigma_{q_j}| < 1$, we have $|n_j| \le |\theta|\,|m_j| + 1$, whence $w_{q_j} \le (3 + |\theta|)\,|m_j|$. For $j \ge k$, we get $|b_{q_j}|\,w_{q_j}^k \le (3 + |\theta|)^k |m_j|^{k-j} \le (3 + |\theta|)^k$; the indices $j < k$ are finitely many. Hence $b$ is rapidly decreasing. If $a$ were a preimage of it, we would have $|a_{q_j}| = |\sigma_{q_j}|^{-1/2} > \sqrt{j+1}$, and $a$ would not even be bounded.
\end{proof}

\begin{proposition}\label{prop:quillen-1983-4-liouville}
A Liouville number is not Diophantine. For such a $\theta$, there is therefore a rapidly decreasing sequence with zero $(0,0)$ term that is not in the image of $D$.
\end{proposition}

\begin{proof}
Let $C > 0$ and $N$. For infinitely many denominators $b$, there exists $a$ with $|\theta - a/b| < b^{-(N+2)}$; taking $b \ge 1/C$, we obtain $|\theta b - a| < b^{-(N+1)} \le C\,b^{-N}$, which contradicts the lower bound. A Liouville number is irrational, and Theorem~\ref{thm:quillen-1983-4} applies.
\end{proof}

\begin{remarque}
For $\theta$ non-Diophantine, the image of $D$ in the smooth functions of mean zero is dense but not closed: the quotient $H^1(T^2, \mathcal O)$ computed by the resolution is infinite-dimensional and non-Hausdorff. The computation of the Picard group that follows in the notebook therefore rests on the Diophantine hypothesis. The set of Diophantine reals has full measure, so that the conclusion of the page holds for almost every $\theta$.
\end{remarque}

\begin{sloppypar}
\noindent\emph{Verification in Lean.} The file \href{\leanblobmain/Notebook1983_4.lean}{\texttt{Notebook1983\_4.\allowbreak{}lean}} works at the level of sequences on $\mathbf Z \times \mathbf Z$, without Fourier analysis. It defines the weight \texttt{weight}, rapid decay \texttt{IsRapidDecay}, the predicate \texttt{IsDiophantine} and the symbol \texttt{symbol}. The forward direction is \texttt{IsDiophantine.isRapidDecay\_div} and \texttt{IsDiophantine.symbol\_mul\_div}, from the lower bound \texttt{IsDiophantine.lower\_bound}; the stability of rapid decay is \texttt{IsRapidDecay.symbol\_mul}. The converse is \texttt{exists\_not\_mem\_range\_of\_not\_isDiophantine}, the equivalence is \texttt{surjective\_iff\_isDiophantine}, and the case of Liouville numbers is \texttt{not\_isDiophantine\_of\_liouville} and \texttt{exists\_not\_mem\_range\_of\_liouville}, from mathlib's \texttt{Liouville} predicate. mathlib has neither rapidly decreasing sequences on $\mathbf Z^2$ nor Diophantine numbers in this sense; the file defines them.
\end{sloppypar}

Not formalized: the identification of smooth functions on $T^2$ with rapidly decreasing sequences of their Fourier coefficients, the resolution of the sheaf $\mathcal O$, the non-Hausdorff character of $H^1$ in the Liouville case, the full measure of the Diophantine numbers, and the computation of the Picard group.

\litterature{found in the literature: the phenomenon of small divisors and the cohomological equation on the torus are classical (V.~I.~Arnold, \emph{Geometrical Methods in the Theory of Ordinary Differential Equations}, supplementary chapters; M.~Herman, Publ. Math. IHÉS 49 (1979); A.~Katok, \emph{Combinatorial constructions in ergodic theory and dynamics}, AMS (2003)). For the foliated cohomology of the Kronecker foliation, which is $\mathbf C$ in degree~$1$ if and only if $\theta$ is Diophantine and infinite-dimensional non-Hausdorff otherwise: A.~El Kacimi Alaoui, Compositio Math. 49 (1983); A.~Haefliger, J. Differential Geom. 15 (1980). The statement is classical; the edition does not present it as new and only points out the hypothesis that the page does not state.}

\section[1983, Lecture Notes 3 (Kirwan): from increasing ratios to $v_0^{n-1}v_n \le v_1^n$]{1983, Lecture Notes 3 (Kirwan): from increasing ratios to $v_0^{\,n-1}v_n \le v_1^{\,n}$\protect\verifie{Notebook1983_LectureNotes3}}\label{sec:quillen-1983-ln3}

\emph{Context.} Among the files of 1983 are notes that Quillen took at lectures of other mathematicians. The \dossier{1983-Lecture\%20Notes\%203\%20Kirwan}{file \q{Lecture Notes 3 Kirwan}} records a talk of Frances Kirwan, of February 1983, on the proof of the Alexandrov--Fenchel inequalities through algebraic geometry. The mathematics is that of the lecture, not Quillen's own.

\emph{The page.} For two convex bodies $K_1$, $K_2$ of $\mathbf R^n$, the volume of $\lambda_1K_1 + \lambda_2K_2$ is a polynomial in $\lambda_1$, $\lambda_2$ whose coefficients, up to binomial factors, are the mixed volumes $v_0, \dots, v_n$. Page~2 states the Alexandrov--Fenchel inequalities in the form $v_0/v_1 \le v_1/v_2 \le \dots \le v_{n-1}/v_n$, and notes that they imply $(v_0/v_1)^n \le v_0/v_n$, that is $v_0^{\,n-1}v_n \le v_1^{\,n}$. With $K_2$ a ball this is the isoperimetric inequality, which is the motivation given on the page.

\emph{Intuition.} The quotient $v_0/v_n$ is the product of the $n$ ratios $v_0/v_1, \dots, v_{n-1}/v_n$, and the first is the smallest.

\begin{encadre}
\begin{theoreme}[After Quillen's notes of a lecture of Kirwan, 1983, p.~2]\label{thm:quillen-1983-ln3}
Let $n \ge 1$ and let $v_0, \dots, v_n$ be positive real numbers such that $v_j\,v_{j+2} \le v_{j+1}^2$ for $0 \le j \le n - 2$, that is, $v_0/v_1 \le v_1/v_2 \le \dots \le v_{n-1}/v_n$. Then $(v_0/v_1)^n \le v_0/v_n$, and $v_0^{\,n-1}v_n \le v_1^{\,n}$.
\end{theoreme}
\end{encadre}

\begin{proof}
First $v_0v_{k+1} \le v_1v_k$ for $k \le n - 1$, by induction on $k$: this is an equality for $k = 0$, and $v_0v_{k+2}v_k \le v_0v_{k+1}^2 \le v_1v_kv_{k+1}$, where the second inequality is the induction hypothesis; divide by $v_k$. Then $v_0^{\,k}v_k \le v_1^{\,k}v_0$ for $k \le n$, by induction on $k$: $v_0^{\,k+1}v_{k+1} = v_0^{\,k}(v_0v_{k+1}) \le v_0^{\,k}v_1v_k \le v_1^{\,k+1}v_0$. For $k = n$ this is the first inequality of the statement, and the second after division by $v_0$.
\end{proof}

\emph{Verification in Lean.} The file \href{\leanblobmain/Notebook1983_LectureNotes3.lean}{\texttt{Notebook1983\_\allowbreak{}LectureNotes3.\allowbreak{}lean}} takes a sequence \texttt{v} of real numbers indexed by the natural numbers, positive up to index $n$, with the hypothesis written without division. The two inductions are \texttt{mul\_le\_mul\_of\_ratios} and \texttt{pow\_mul\_le\_pow\_mul}, and the two conclusions of Theorem~\ref{thm:quillen-1983-ln3} are \texttt{div\_pow\_le\_div} and \texttt{pow\_pred\_mul\_le\_pow}. The Alexandrov--Fenchel inequalities themselves, mixed volumes and the isoperimetric inequality are not formalised; mathlib has no mixed volumes. The axioms used are the three usual axioms of mathlib.

\litterature{The inequalities are those of Alexandrov and Fenchel (1930s), and the proof presented in the lecture is that of Khovanskii and Teissier, as recalled. The implication proved here is the standard passage from a log-concave sequence to its extreme terms. The edition does not present it as new.}

\section[Notebook 1986-5: a four-fold embedding and a circle integral]{Notebook 1986-5: a four-fold embedding and $\int_0^{2\pi}(2\sin\frac\theta2)^{2k}\,d\theta$\protect\verifie{Notebook1986_5}}\label{sec:quillen-1986-5}

\emph{Context.} On page~17 of \carnet{1986}{5} Quillen embeds an operator $A$ with $A^2 + B^2 = 1$ as the compression of a symmetry $\varepsilon$ by an isometry $j$; on page~31 he needs the moments of $\lvert 1 - e^{i\theta}\rvert$ on the circle.

\begin{encadre}
\begin{theoreme}[Quillen, notebook 1986-5, pp.~17, 31]\label{thm:quillen-1986-5}
\begin{enumerate}
\item Over a commutative ring in which $2$ is invertible, if $A^2 + B^2 = 1$, the column $j = \bigl(\frac{A+1}2, \frac B2, \frac{A-1}2, \frac B2\bigr)$ and $\varepsilon = \operatorname{diag}(1, 1, -1, -1)$ satisfy $j^tj = 1$ and $j^t\varepsilon j = A$.
\item For $k \in \mathbf N$, $\int_0^{2\pi} (2\sin\frac\theta2)^{2k}\,d\theta = 2\pi\binom{2k}k$.
\end{enumerate}
\end{theoreme}
\end{encadre}

\begin{proof}
(1) Expand. (2) Substitute $\theta = 2x$: the integral is $2\cdot 4^k\int_0^\pi \sin^{2k}x\,dx = 2\pi\, 4^k\prod_{i<k}\frac{2i+1}{2i+2}$ (Wallis), and $4^k\prod_{i<k}\frac{2i+1}{2i+2} = \binom{2k}k$ by induction, from $(k+1)\binom{2k+2}{k+1} = 2(2k+1)\binom{2k}k$.
\end{proof}

\emph{Verification in Lean.} The file \href{\leanblobmain/Notebook1986_5.lean}{\texttt{Notebook1986\_5.\allowbreak{}lean}} proves (1) as \texttt{dot\_self} and \texttt{dot\_eps} (with $A$, $B$ commuting, the transpose standing for the adjoint) and (2) as \texttt{integral\_two\_sin\_half\_pow}, with \texttt{four\_pow\_mul\_prod}, from mathlib's \texttt{integral\_sin\_pow\_even}. The axioms used are the three usual axioms of mathlib.

\litterature{Elementary; (2) is the classical value $\int_0^{2\pi}\lvert 1 - e^{i\theta}\rvert^{2k}\,d\theta = 2\pi\binom{2k}k$. The edition does not present them as new.}

\section[Notebook 1986-12: the geometric-sum bound; the Cayley resolvent bound]{Notebook 1986-12: $\|g^n - 1\| \le n\|g - 1\|$; $|t + 1| - |t - 1| \ge 2\operatorname{Re}t/(1 + |t|)$\protect\verifie{Notebook1986_12}}\label{sec:quillen-1986-12}

\emph{Context.} In \carnet{1986}{12} Quillen estimates functions of unitaries in Schatten ideals through $\|(g^n - 1)/(g - 1)\| \le n$ (p.~23), and bounds the resolvent of the Cayley picture $t(g + 1) + (g - 1)$ of a skew-adjoint operator for $\operatorname{Re}t > 0$ (p.~53).

\begin{encadre}
\begin{theoreme}[Quillen, notebook 1986-12, pp.~23, 53]\label{thm:quillen-1986-12}
\begin{enumerate}
\item In a normed ring, if $\|g^k\| \le 1$ for all $k$, then $\|g^n - 1\| \le n\|g - 1\|$.
\item For $t \in \mathbf C$, $|t + 1|^2 - |t - 1|^2 = 4\operatorname{Re}t$; if $\operatorname{Re}t > 0$, $|t + 1| - |t - 1| \ge 2\operatorname{Re}t/(1 + |t|)$.
\item For a linear isometry $g$ of a complex normed space, $\|((t + 1)g + (t - 1))v\| \ge (|t + 1| - |t - 1|)\|v\|$.
\end{enumerate}
\end{theoreme}
\end{encadre}

\begin{proof}
(1) $g^n - 1 = (\sum_{k<n}g^k)(g - 1)$. (2) Expand; then $(|t + 1| - |t - 1|)(|t + 1| + |t - 1|) = 4\operatorname{Re}t$ and $|t \pm 1| \le 1 + |t|$. (3) The triangle inequality and $\|gv\| = \|v\|$.
\end{proof}

With (2) and (3), $\|(t(g + 1) + (g - 1))^{-1}\| \le (1 + |t|)/(2\operatorname{Re}t)$ wherever the inverse exists, the page's bound.

\emph{Verification in Lean.} The file \href{\leanblobmain/Notebook1986_12.lean}{\texttt{Notebook1986\_12.\allowbreak{}lean}} proves \texttt{norm\_pow\_sub\_one\_le}, \texttt{norm\_add\_one\_sq\_sub}, \texttt{norm\_add\_one\_sub\_ge} and \texttt{norm\_cayley\_ge}. The Schatten-ideal smoothness and the unbounded-operator setting are not formalised. The axioms used are the three usual axioms of mathlib.

\litterature{Elementary; the bound in terms of $g$ was not located as stated (Kato, \emph{Perturbation theory}, ch.~V, for the setting). The edition does not present it as new.}

\section[Notebook 1988-3: the operator $(x^2-1)\partial + 4x$]{Notebook 1988-3: the operator $(x^2 - 1)\partial_x + 4x$ on polynomials\protect\verifie{Notebook1988_3}}\label{sec:quillen-1988-3}

\emph{Context.} On pages~46--47 of \carnet{1988}{3} Quillen studies the operator $L = (x^2 - 1)\partial_x + 4x$ on polynomials, computes it on monomials and asks for its image.

\begin{encadre}
\begin{theoreme}[Quillen, notebook 1988-3, pp.~46--47]\label{thm:quillen-1988-3}
On $\mathbf R[x]$, the operator $L = (x^2 - 1)\partial_x + 4x$ sends $x^n$ to $(n + 4)x^{n+1} - nx^{n-1}$, is injective, and its image is the kernel of the linear form $\varphi(q) = \int_{-1}^1 q(x)(x^2 - 1)\,dx$, a hyperplane since $\varphi(1) = -4/3$.
\end{theoreme}
\end{encadre}

\begin{proof}
The coefficient of $x^{N+1}$ in $Lp$ is $(N + 4)p_N - (N + 2)p_{N+2}$. For $N$ the degree of $p \ne 0$ it is $(N + 4)$ times the leading coefficient, so $Lp \ne 0$. Next, $(x^2 - 1)Lp = ((x^2 - 1)^2p)'$, and $(x^2 - 1)^2$ vanishes at $\pm1$, so $\varphi(Lp) = 0$. Conversely, let $\varphi(q) = 0$, of degree $d + 1$: subtracting a multiple of $L(x^d)$, whose coefficient of $x^{d+1}$ is $d + 4$, lowers the degree without leaving the kernel; one reaches a constant $c$ with $\varphi(c) = -4c/3 = 0$, so $c = 0$.
\end{proof}

\emph{Verification in Lean.} The file \href{\leanblobmain/Notebook1988_3.lean}{\texttt{Notebook1988\_3.\allowbreak{}lean}} proves \texttt{L\_X\_pow}, \texttt{L\_injective}, \texttt{phi\_L}, \texttt{phi\_one} and \texttt{mem\_range\_iff}, by induction on the degree. The page's use of $L$ (inverting it on the hyperplane) is not formalised. The axioms used are the three usual axioms of mathlib.

\litterature{Elementary: $(x^2 - 1)L = \partial_x\circ(x^2 - 1)^2$; the edition does not present it as new.}

\section[Notebooks 1984-13 to 1989-4: Beta integrals]{Notebooks 1984-13, 1988-5, 1988-7, 1988-9, 1989-4: Beta integrals with integer exponents\protect\verifie{Notebook1988_9}}\label{sec:quillen-1988-9}

\emph{Context.} The normalising constants of Quillen's Chern--Simons and cyclic-cocycle formulas are Beta integrals $\int_0^1(t^2 - t)^kt^\ell\,dt$; they recur on page~3 of \carnet{1984}{13}, pages~53 and~85 of \carnet{1988}{5}, page~13 of \carnet{1988}{7}, page~72 of \carnet{1988}{9}, and, over a $2$-simplex, page~56 of \carnet{1989}{4}.

\begin{encadre}
\begin{theoreme}[Quillen, notebooks 1984-13, 1988-5, 1988-7, 1988-9, 1989-4]\label{thm:quillen-1988-9}
For natural numbers $a$, $b$, $i$, $j$, $k$, $\ell$, $m$, $q$:
\begin{enumerate}
\item $\int_0^1 t^a(1 - t)^b\,dt = a!\,b!/(a + b + 1)!$;
\item (1988-9) $\int_0^1(t^2 - t)^kt^\ell\,dt = (-1)^kk!\,(k + \ell)!/(2k + \ell + 1)!$;
\item (1988-5) $\frac1{k!}\int_0^1(t^2 - t)^k\,dt = (-1)^kk!/(2k + 1)!$ and $\frac1{(q+k)!}\int_0^1t^q(t^2 - t)^k\,dt = (-1)^kk!/(q + 2k + 1)!$;
\item (1984-13) $\int_0^1(t^2 - t)^m\,t\,dt/m! = (-1)^m(m + 1)!/(2m + 2)!$ and $\int_0^1(t^2 - t)^m\,dt/m! = (-1)^mm!/(2m + 1)!$;
\item (1988-7) $\int_0^1(m + 1)(t^2 - t)^m\,dt = (-1)^m(m + 1)!\,m!/(2m + 1)!$ and $\int_0^\infty e^{-t^2}t^{2m+1}\,dt = m!/2$;
\item (1989-4) $\int_0^1\int_0^{1-t_1}(1 - t_1 - t_2)^it_1^jt_2^k\,dt_2\,dt_1 = i!\,j!\,k!/(i + j + k + 2)!$.
\end{enumerate}
\end{theoreme}
\end{encadre}

\begin{proof}
(1) By induction on $b$: $t^a(1 - t)^{b+1} = t^a(1 - t)^b - t^{a+1}(1 - t)^b$, and $\frac{a!\,b!}{(a+b+1)!} - \frac{(a+1)!\,b!}{(a+b+2)!} = \frac{a!\,(b+1)!}{(a+b+2)!}$; for $b = 0$ the integral is $1/(a + 1)$. (2)--(5) follow from $(t^2 - t)^k = (-1)^kt^k(1 - t)^k$, and the Gaussian integral from mathlib's $\int_0^\infty x^qe^{-x^p}dx = \frac1p\Gamma(\frac{q+1}p)$. (6) The substitution $t_2 = (1 - t_1)u$ gives $\int_0^c(c - x)^ix^k\,dx = c^{i+k+1}\,i!\,k!/(i + k + 1)!$, and (1) again in $t_1$.
\end{proof}

\emph{Verification in Lean.} The file \href{\leanblobmain/Notebook1988_9.lean}{\texttt{Notebook1988\_9.\allowbreak{}lean}} proves (1) as \texttt{integral\_pow\_mul\_one\_sub\_pow}, without the Gamma function, then \texttt{integral\_sq\_sub\_pow\_mul\_pow}, \texttt{notebook1988\_5\_first} and \texttt{\_second}, \texttt{notebook1984\_13\_first} and \texttt{\_second} (written with $k = m + 1$), \texttt{notebook1988\_7\_beta}, \texttt{notebook1988\_7\_gauss}, \texttt{integral\_scaled} and \texttt{integral\_simplex}. The simplex integral is stated as an iterated integral, not as an integral over the simplex with its measure. The axioms used are the three usual axioms of mathlib.

\litterature{Classical (Euler's Beta integral); mathlib has the complex Beta function and its Gamma formula, but not this real form with natural exponents. The edition does not present them as new.}

\section[Notebook 1988-9: the Chern--Simons values of an idempotent]{Notebook 1988-9: $(2n)!/n!$ grows like $n!\,4^n$\protect\verifie{Notebook1988_9Growth}}\label{sec:quillen-1988-9-growth}

\emph{Context.} On pages~94--97 of \carnet{1988}{9}, in the discussion of entire cyclic cohomology, Quillen computes the values of a Chern--Simons cocycle on the idempotent of $\mathbf C + \mathbf Ce$ and observes that the chain of the idempotent has coefficients $(2n)!/n!$, which he estimates as growing roughly like $n!\,2^n$.

\begin{encadre}
\begin{theoreme}[Quillen, notebook 1988-9, pp.~95--97, corrected]\label{thm:quillen-1988-9-growth}
For every $n$:
\begin{enumerate}
\item $(2n + 1)\,\frac{1}{n!}\int_0^1(t^2 - t)^n\,dt = (-1)^n\,n!/(2n)!$ and $(n + 1)!/(2n + 2)! = n!/(2\,(2n + 1)!)$;
\item $(2n)!/n! = n!\binom{2n}{n}$, and $n!\,2^n \le (2n)!/n! \le n!\,4^n \le (2n + 1)\,(2n)!/n!$;
\item $\bigl((2n)!/n!\bigr)/(n!\,2^n) \to \infty$.
\end{enumerate}
\end{theoreme}
\end{encadre}

\begin{proof}
(1) The Beta integral gives $\int_0^1(t^2 - t)^n\,dt = (-1)^n(n!)^2/(2n + 1)!$ (section~\ref{sec:quillen-1988-9}). (2) $2^n \le \binom{2n}{n}$ by induction from $(n + 1)\binom{2n+2}{n+1} = 2(2n + 1)\binom{2n}{n}$; $\binom{2n}{n} \le 4^n \le (2n + 1)\binom{2n}{n}$ are in mathlib. (3) The ratio is at least $2^n/(2n + 1) \ge n/3$.
\end{proof}

\emph{Verification in Lean.} The file \href{\leanblobmain/Notebook1988_9Growth.lean}{\texttt{Notebook1988\_9Growth.\allowbreak{}lean}} proves \texttt{chernSimons\_value}, \texttt{factorial\_ratio}, \texttt{factorial\_div\_eq}, \texttt{two\_pow\_le\_centralBinom}, \texttt{factorial\_div\_bounds} and \texttt{tendsto\_ratio}. \emph{Correction}: the page's estimate $n!\,2^n$ is only a lower bound; the growth is that of $n!\,4^n$ up to a factor $2n + 1$ (the slip may be the transcription's); the conclusion that the idempotent pairs with slowly growing cochains is unaffected. The cocycle itself and entire cyclic cohomology are not formalised. The axioms used are the three usual axioms of mathlib.

\litterature{Classical (central binomial coefficients; Connes's entire cyclic cohomology for the context); the edition does not present it as new.}


\clearpage\addtocontents{toc}{\protect\addvspace{2.25em}}% after the last part
\section*{Statement on the use of artificial intelligence}\addcontentsline{toc}{section}{Statement on the use of artificial intelligence}

The statements proved here are read from transcriptions of the notebooks produced by Anthropic language models (Claude) under editorial rules written by the author; these transcriptions have not been fully reviewed by a person. The formal proofs were written in Lean~4 by Claude Opus~5.5, under the author's direction, and Lean mechanically checks their correctness against the mathlib library. The same model drew from them the written-out proofs of this document. The author reviewed the choice of statements and the text, and takes full responsibility for them. Lean guarantees that a statement holds under the hypotheses as written. It guarantees neither that the statement faithfully renders the page, nor that it is new.

\begin{thebibliography}{99}\small
\bibitem{ass1994} J.~Avron, R.~Seiler and B.~Simon, \q{The index of a pair of projections}, \emph{Journal of Functional Analysis} 120 (1994), no.\,1, 220--237; doi:10.1006/jfan.1994.1031.
\bibitem{lean4} L.~de Moura and S.~Ullrich, \q{The Lean~4 theorem prover and programming language}, in \emph{Automated Deduction -- CADE~28}, Lecture Notes in Computer Science 12699, Springer, 2021, 625--635; doi:10.1007/978-3-030-79876-5\_37.
\bibitem{mathlib2020} The mathlib Community, \q{The Lean mathematical library}, in \emph{Proceedings of the 9th ACM SIGPLAN International Conference on Certified Programs and Proofs (CPP~2020)}, ACM, 2020, 367--381; doi:10.1145/3372885.3373824.
\bibitem{beer2008} G.~Beer, \q{Embeddings of bornological universes}, \emph{Set-Valued Analysis} 16 (2008), no.\,4, 477--488; Zbl~1171.54012.
\bibitem{beervipera2012} G.~Beer and M.~C.~Vipera, \q{The Alexandroff one-point compactification as a prototype for extensions}, \emph{Advances in Mathematics} 231 (2012), nos.\,3--4, 1598--1618; doi:10.1016/j.aim.2012.07.008.
\bibitem{hu1949} S.-T.~Hu, \q{Boundedness in a topological space}, \emph{Journal de Mathématiques Pures et Appliquées} (9) 28 (1949), 287--320.
\end{thebibliography}

\end{document}
